\chapter{C2: the directed weak source}
\label{ch:cert-7}
Formula tags are local to their certificate section. All source and testing fields remain separately labelled.

\section{\texorpdfstring{A freshly derived full-state sign and projective join}{A freshly derived full-state sign and projective join}}
\label{component:pf77de02d80c4}
% Scoped mathematical certificate; see local_identity_audit.closure.json for exact changes.

\label{pf77de02d80c4:sec:rederived-state-join}

This section is a new proof derived from the complete weak-C2 parent retained
in the earlier R21 checkpoint.  It does not claim byte identity with the
missing late-R21 state-join file.  Its conclusion is a reduction of the full
parent to one positive state-ratio chart.  The certificate in that remaining
chart is a separate hypothesis; no global finite-lift conclusion is drawn
here.

Let $K_c,K_t\supset\mathbb R$ be real closed fields, occupying separately
labelled tensor factors.  All quantities $u,v,W,X,Y,Z$ below belong to $K_t$;
all six coefficients $n,m,A,C,D,\alpha$ belong to $K_c$ and are nonnegative.
Write $D_u=1+u^2$, $D_v=1+v^2$, and define the complete parent
\begin{align}
 F={}&nD_uD_vX^2+m\{v(W+uY)-X-uZ\}^2
       +AD_uD_vY^2+CD_v(uW-Y)^2\\
    &+D(uW+uvX-Y-vZ)^2+\alpha D_u(vY-Z)^2+K,\tag{SJ1}\\
 K={}&\tfrac12(uW+uvX-Y-vZ)^2
       +\tfrac12D_u(vY-Z)^2
       +\tfrac12D_v(X+uZ)^2-\tfrac12D_uD_vZ^2.\tag{SJ2}
\end{align}
Products involving source and testing quantities mean tensor products.
For example, $nD_uD_vX^2$ means
$n\otimes(D_uD_vX^2)$; a testing-only term means $1\otimes K$.
The positive packets in (SJ1) are sums of squares in
$K_c\otimes_{\mathbb R}K_t$, since each coefficient and each testing
multiplier has a square root in its own real closed field.

\paragraph{Sign gauges preserve the entire parent.}
For $\varepsilon,\eta\in\{1,-1\}$, the substitution
\[
 (u,v,W,X,Y,Z)\longmapsto
 (\varepsilon u,\eta v,W,\eta X,\varepsilon Y,
  \varepsilon\eta Z)                                      \tag{SJ3}
\]
preserves every summand of (SJ1), including (SJ2).  Indeed the roots of the
$m,C,D,\alpha$ packets acquire, respectively, the signs
$\eta,\varepsilon,\varepsilon,\varepsilon\eta$, and the other roots also
change only by a sign.  Thus $u,v$ may both be taken nonnegative without
altering the source coefficients or the full expression.

\paragraph{An exact positive hard block.}
Set
\[
 t=\frac{X+uZ}{D_u},\qquad r=\frac{Y-uW}{D_u},\qquad
 w=\frac{W+uY}{D_u},\qquad z=\frac{Z-uX}{D_u}.             \tag{SJ4}
\]
The inverse transformation is
$X=t-uz$, $Z=z+ut$, $Y=r+uw$, $W=w-ur$.
Let $H=1+u^2+v^2$ and $d=H^2-u^2v^2$.  Both
$H-uv=1+(u-v)^2+uv$ and $H+uv$ are positive for $u,v\geq0$;
in fact $H>|uv|$ for arbitrary $u,v$.  Consequently $d>0$.
Define
\[
 \mathsf A=\frac12\begin{pmatrix}H&-uv\\-uv&H\end{pmatrix},\qquad
 \mathsf B=\frac12\begin{pmatrix}-u^2v&-uv^2\\uv^2&u^2v\end{pmatrix}.
\]
The matrix $\mathsf A$ is positive definite.  Direct expansion and square
completion give the exact identity
\begin{align}
 \frac K{D_u}
 ={}&\left(\binom tr+\mathsf A^{-1}\mathsf B\binom wz\right)^{\!T}
       \mathsf A
       \left(\binom tr+\mathsf A^{-1}\mathsf B\binom wz\right)\\
    &+\frac{D_uD_v}{2d}
       \{u^2v^2(w^2+z^2)-2uvH wz\}.                    \tag{SJ5}
\end{align}
The rational matrix used in the shift is explicitly
\[
 \mathsf A^{-1}\mathsf B=
 \frac1d\begin{pmatrix}
 -u^2v(1+u^2)&-uv^2(1+v^2)\\
 uv^2(1+v^2)&u^2v(1+u^2)
 \end{pmatrix}.                                         \tag{SJ6}
\]
For any $a,b\in K_t$, the hard quadratic form is
\[
 \binom{a}{b}^{\!T}\mathsf A\binom{a}{b}
   =\frac{H-uv}{4}(a+b)^2+\frac{H+uv}{4}(a-b)^2.         \tag{SJ7}
\]
This also gives an explicit square decomposition of the first term of
(SJ5), with square roots taken only in $K_t$.

If $u,v\geq0$ and $wz\leq0$, the second term of (SJ5) is nonnegative.
It too is a square in $K_t$ (or can be split into its three nonnegative
terms).  Hence $K\geq0$, and (SJ1) is a tensor sum of squares immediately.
In particular the axes $w=0$ and $z=0$ are included, without division by
either state coordinate.  The same identity, without sign gauges, proves
the sharper implication $K<0\Longrightarrow uvwz>0$.

\paragraph{Reduction of the remaining states.}
The only remaining case after (SJ3) is $K<0$, $u,v>0$, $wz>0$.
Then $w\ne0$.  Divide all four state coordinates by $w$.  Equation (SJ4)
now has $w'=1$ and $z'=z/w>0$.  The ratio $z'$ is arbitrary: it may be
infinitesimal, bounded, or unbounded in $K_t$.  No estimate in this
reduction replaces $z'$ by a bounded coefficient.

The parent is homogeneous of degree two in the state coordinates, so
\[
 F(u,v;W,X,Y,Z)=w^2F(u,v;W/w,X/w,Y/w,Z/w).               \tag{SJ8}
\]
If the parent on the right has a certificate $\sum_i h_i^2$ in
$K_c\otimes_{\mathbb R}K_t$, then the parent on the left has certificate
$\sum_i((1\otimes w)h_i)^2$.  This proves that certificates for the full
parent with $u,v>0$, $w=1$, and arbitrary positive $z$ suffice for all
finite affine testing parameters and all real states.

\paragraph{Literal projective endpoints.}
For clarity, the projective endpoints can be checked without a limiting
argument.  Homogenize $u=U_1/U_0$, $v=V_1/V_0$, and put
$D_U=U_0^2+U_1^2$, $D_V=V_0^2+V_1^2$, with neither pair identically zero.
Multiplication by $U_0^2V_0^2$ turns (SJ2) into
\begin{align}
 K_h={}&\tfrac12(U_1V_0W+U_1V_1X-U_0V_0Y-U_0V_1Z)^2\\
       &+\tfrac12D_U(V_1Y-V_0Z)^2
        +\tfrac12D_V(U_0X+U_1Z)^2-\tfrac12D_UD_VZ^2.     \tag{SJ9}
\end{align}
At the four axes this is, exactly,
\begin{align}
 U_0=0:\quad &K_h=\tfrac12U_1^2\{(V_0W+V_1X)^2+(V_1Y-V_0Z)^2\},\\
 V_0=0:\quad &K_h=\tfrac12V_1^2D_U(X^2+Y^2),\\
 U_1=0:\quad &K_h=\tfrac12U_0^2D_V(X^2+Y^2),\\
 V_1=0:\quad &K_h=\tfrac12V_0^2\{(U_1W-U_0Y)^2+(U_0X+U_1Z)^2\}.
                                                               \tag{SJ10}
\end{align}
All other homogenized packets in (SJ1) retain nonnegative coefficients
and are sums of squares.  Thus the endpoints are admitted directly,
including the intersections of these axes.  Equations (SJ3)--(SJ10)
establish the stated full-state, sign, and endpoint reduction.  They do
not establish the still-required positive-ratio chart certificate.

\paragraph{Verification scope.}
The companion exact checker verifies the inverse transformation, the
entire parent under all four gauges, the hard-block and Schur identities,
homogeneity, homogenization, and all four endpoint formulas.  These are
algebraic corroborations of the proof above; they are not a formal proof
of the complete R21 theorem and do not certify other missing components.


\section{\texorpdfstring{Reconstructed R2 source entry and full two-packet extraction}{Reconstructed R2 source entry and full two-packet extraction}}
\label{component:p667905169990}
% Scoped mathematical certificate; see local_identity_audit.closure.json for exact changes.

\label{p667905169990:sec:r2-reconstructed-source-entry}

\paragraph{Actual source and conventions.}
This is a fresh mathematical reconstruction of the retained RE1--RE8
source-entry formulas.  The original late-R21 file is absent.  The
identities are checked again against the unchanged native source
formulas, and the boundedness proof is given in full.  This lemma is
structural: it supplies neither a scalar contact conductor nor an
all-testing R2 certificate.  In particular, restoring this lemma does
not restore the missing thin-source, normal, or outer-chart proofs.

Use the original weak-C2 native source $(a,B,t,r,h)$, with
$0<r,h<1$, $a,B,t>0$, the original strict source gaps, and $n\geq0$.
In R2, $t$ is positive unbounded while
\[
 B,\qquad z=aB,\qquad y=tB
\]
are positive infinitesimals.  There is no bound on $a$ and no lower
bound on $r$, $h$, $1-h$, or $e=1-r$.  Put $l=1-h$ and $S=1/t$.
Define the actual native arrays
\[
 c=1+zr,\quad s_c=1+B(l+ar),\quad
 R=1+hz,\quad g=R-ae>0,\quad j_D=c-l/t>0,
\]
\begin{align*}
 \Delta={}&a+(1+z)+(1+z)B+2z(1+z)r
          +a^3B^2r^2+2a^2Br\\
          &+tB\{1+a+z+a^2Br(2-r)\},\\
 M={}&(1+z)(1+ar)(h+t)-a(1+y)e\,l.
\end{align*}
The six weights used throughout this section are exactly
\[
\begin{aligned}
 n&=\frac{(1+ar)c s_c-r\Delta}{2r\Delta},&
 m&=\frac{a^2(1+y)e c s_c}{2g\Delta},\\
 C&=\frac{cM}{2l\Delta},&
 A&=\frac{tc(1+ar)s_c}{2h\Delta},\\
 D&=\frac{(1+z)(1+ar)(1-atBr)s_c}{2t j_D\Delta},&
 \alpha&=\frac{B^2(1+ar)M}{2(1-hB)\Delta}.
\end{aligned}                                                     \tag{RER1}
\]
All denominators are pure source.  The original positivity assumptions
are retained without adding a rate comparison.

\paragraph{Physicality controls the apparently unbounded source scale.}
The exact numerator identity is
\[
 (1+ar)(1+zr)\{1+B(l+ar)\}-r\Delta=eQ-BK,                \tag{RER2}
\]
where
\[
 K=(a+1)(h+t)(1+z),
\]
\[
 Q=1+B+y+z(h+t+2r)
       +z^2\{t(1+er)+h(1+r)\}+Bz(h+t).                 \tag{RER3}
\]
Every term of $Q$ is nonnegative.  Since $r\Delta>0$, original
physicality is exactly $eQ\geq BK$, including equality when $n=0$.
Put $\kappa=ay=tz$.  Before assuming that $\kappa$ is bounded, direct
termwise inspection of (RER3), using $B,y,z$ infinitesimal, gives
$Q\leq C_0(1+\kappa)$ for a fixed positive real $C_0$.
The strict gap $g>0$ gives $ae<R\sim1$, so multiplication of the
physicality inequality by $a$ yields
\[
 \kappa(a+1)(1+hS)(1+z)\leq aeQ<RQ
                                      \leq C_1(1+\kappa)           \tag{RER4}
\]
for a fixed positive real $C_1$.

If $a$ is bounded, $\kappa=ay$ is infinitesimal.  If $a$ is unbounded,
move $C_1\kappa$ to the left in (RER4).  Since
$a+1-C_1\geq(a+1)/2$, one obtains
$\kappa\leq2C_1/(a+1)$, again infinitesimal.  In either case,
(RER4) also bounds $(a+1)\kappa$.  Returning to (RER3) and then to
$eQ\geq BK$ proves
\[
 ay\in\mathrm I,\qquad a(a+1)y=O(1),\qquad Q\sim1,
 \qquad \frac{(a+1)y}{e}=O(1).                         \tag{RER5}
\]
No division by $n$ occurred; the argument includes the face $n=0$.

\paragraph{Uniform source normalizations.}
Put
\[
 k=\frac{1+ar}{1+a},\qquad p=\frac{ae}{R}\in(0,1),
 \qquad \mu=m(1-p)=\frac{mg}{R},\qquad\delta_n=nr.
\]
The quantity $k$ is a positive bounded unit.  Indeed $k\leq1$; if
$a\leq1$, then $k\geq1/(1+a)\geq1/2$.  If $a>1$, use
$ae<R=1+hz$ to get $k>(a-hz)/(1+a)$, which has a positive real lower
bound.  Direct termwise division of the displayed polynomial
$\Delta$ by $1+a$ gives $\Delta/(1+a)\sim1$.  Also
\[
 \frac{M}{t(1+ar)}\sim1:
\]
the first term has ratio $(1+z)(1+hS)\sim1$, while the absolute
ratio of the second is at most $R(1+y)S/(1+ar)$ and hence is
infinitesimal.  Since $c,s_c,j_D,1-hB\sim1$, (RER1) now gives
\[
 \frac{SC\,l}{k}\sim\frac12,\quad
 \frac{SA\,h}{k}\sim\frac12,\quad
 \frac{D}{Sk}\sim\frac12,\quad
 \frac{\alpha}{Sy^2(1+a)k^2}\sim\frac12.               \tag{RER6}
\]
For the $D$ relation use $atB=ay=\kappa\in\mathrm I$ from
(RER5).  Thus all original $h$ and $l$ openings are present.

Two exact bounded products remove the remaining nonuniform source
poles:
\[
 mg=\frac{a^2e(1+y)c s_c}{2\Delta},\qquad
 \delta_n=\frac{eQ-BK}{2\Delta}.                       \tag{RER7}
\]
The quantities $mg$, $mg/a$, $mg/a^2$, $\mu$, $\delta_n$, and
$nr^2$ are bounded.  For example
$a^2e/(1+a)\leq R a/(1+a)=O(1)$, while
$0\leq\delta_n\leq eQ/(2\Delta)=O(1)$.
These estimates do not imply that $m$ or $n$ themselves are bounded.

\paragraph{Exact extraction from the entire parent.}
Let testing $u,v,W,X,Y,Z$ be arbitrary, independently of the source;
put $D_u=1+u^2$, $D_v=1+v^2$.  The complete original quarter-parent is
\begin{align*}
 \mathcal P={}&m(vW-X+uvY-uZ)^2+nD_uD_vX^2
     +CD_v(uW-Y)^2+AD_uD_vY^2\\
 &+(D+\tfrac12)(uW+uvX-Y-vZ)^2
     +(\alpha+\tfrac12)D_u(vY-Z)^2\\
 &+\tfrac12D_v(X+uZ)^2-\tfrac12D_uD_vZ^2.
\end{align*}
Define the mixed rows
\[
 q_a=v(W+uY),\qquad N=X+uZ-pq_a,\qquad \xi=X+ruZ.
\]
Then the following identity retains both actual positive packets,
without setting either normal to zero:
\[
 \mathcal P=mpN^2+nD_uD_v\xi^2+G,                       \tag{RER8}
\]
where
\begin{align*}
 G={}&\mu\{N^2-2q_aN+(1-p)q_a^2\}
        +\delta_nD_uD_v(-2uZX-ru^2Z^2)\\
 &+CD_v(uW-Y)^2+AD_uD_vY^2\\
 &+(D+\tfrac12)(uW+uvX-Y-vZ)^2
        +(\alpha+\tfrac12)D_u(vY-Z)^2\\
 &+\tfrac12D_v(X+uZ)^2-\tfrac12D_uD_vZ^2.
\end{align*}
It follows simply by expanding
$m\{(1-p)q_a-N\}^2-mpN^2$ and
$n(X^2-\xi^2)=nr(-2uZX-ru^2Z^2)$.
Both retained packets are finite separated SOS: their source weights
are nonnegative, $D_uD_v$ is pure testing, and $N,\xi$ are allowed
mixed rows.  When $n=0$ the second packet vanishes literally.
The exact identity and all estimates above are the reconstructed
structural conclusion; no positivity conclusion about $G$ is asserted.


\section{\texorpdfstring{A five-root local certificate on an A-unbounded weak-C2 subchart}{A five-root local certificate on an A-unbounded weak-C2 subchart}}
\label{component:p9ec3d6bbd717}
% Scoped mathematical certificate; see local_identity_audit.closure.json for exact changes.

\label{p9ec3d6bbd717:sec:r9-weak-A-five-root}

\begin{theorem}
Retain the original physical C2 source, its exact stationary contact,
all strict source gaps, and $n\ge0$.
Normalize $d=1$. Suppose $B,z,y$ are positive infinitesimals,
$r,G/B,j,y/B,L/B$ are positive bounded units, and $H/y$ is
infinitesimal. Impose in addition
\[
 z/B\text{ and }1-r\text{ are positive bounded units}.       \tag{WA1}
\]
Then the entire original parent is a finite tensor sum of squares
for arbitrary independent testing half-angles $u,v$ that are
infinitesimal or zero, including either sign and complex states.
There is no comparison of a testing scale with $B,H,z,y$.
The total-scale parameter $\delta_{\rm tot}$ is positive infinitesimal.
This is a local testing certificate on the stated source subchart;
it does not assert that its intersection with the final residual
complement is nonempty, or supply an unproved remote-testing join.
\end{theorem}

No membership in the residual complement is needed in this proof.
When applied to the original residual problem, its existing complement
condition is still intersected with the source chart; no nonemptiness
of that intersection is inferred.

The original omitted event coefficients remain $\lambda=\beta=0$.
Here $B=H+L$, and the actual contact in the positive gauge is
\[
 (u_c,v_c;W_c,X_c,Y_c,Z_c)
 =\left(\sqrt z,\sqrt y;\frac B{\sqrt{zy}},
                  -r\sqrt z,\frac H{\sqrt y},1\right).
\]
The source ratio $\ell=B/z$ below is not an event coefficient.

\paragraph{The source pivots survive the small H opening.}
Write $h=H/B$, $\kappa=B/y$, $\ell=B/z$. Both $\kappa$ and
$\ell$ are positive bounded units, and $h$ is infinitesimal.
The original formulas imply that $n,m,C,D,\alpha$ are bounded,
that $m,C,D$ are positive bounded units, and that $Ah$ is a
positive bounded unit. Thus $A$ is unbounded and
$Ah^2,\sqrt A\,h,A^{-1/2}$ are infinitesimal. No bound on the
relative rates of $h$ and $B$ is used.

Here is an explicit check of the radial pivot. With $T_c=B+zr$,
the original formula is
\[
 \frac\alpha{B^2}
   =\frac{T_c}{\delta}\frac{M/B^2}{2(1-H)},\qquad
 M-yG=H(1+z)T_c+z\{yj+rL(1+y)\}>0.                  \tag{WA2}
\]
Consequently $M/B^2\ge(y/B)(G/B)$ is bounded below by a
positive real, and it is bounded above by the original formula
for $M$ and the unit ratios in (WA1). Since $T_c/\delta$ is a
positive bounded unit, $\alpha/B^2$ is a positive bounded unit
uniformly as $H/y\to0$. Likewise
\[
 m=\frac{z^2(1+y)(1-r)(1+zr)(1+zr+L)}{2G\delta}
\]
is a positive bounded unit under (WA1); the local residue argument
uses only $n\ge0$. Finally $A$ is the only unbounded retained coefficient, so
$\delta_{\rm tot}=1/(1+n+m+C+A+D+\alpha)$ is positive
infinitesimal. These arguments preserve all original coefficient
and source-gap hypotheses.

\paragraph{Five homogeneous contact roots.}
Put $U=u^2$, $V=v^2$, and define
\[
 \xi=X+ruZ,\qquad q=\sqrt A\,(Y-huW),\qquad
 R_1=uW-\kappa vZ,\quad R_2=vW-\ell uZ,\quad
 T=uvW-BZ.                                             \tag{WA3}
\]
Every root vanishes at the actual source contact. Inverses and
radicals in their coefficients are pure source. The large $A$
packet is retained in $q$; no positive packet is subtracted from
the parent and no physicality hypothesis is transferred to a
modified remainder.

Use the following source coefficients for the exact scalar graph:
\begin{align*}
 a&=(C+D+\tfrac12)(1-h)^2+Ah^2,&a_4&=Ah^2,\\
 b_2&=2mh+C(1-h)^2+(A+\alpha+\tfrac12)h^2,
 &b_4&=h^2(m+A+\alpha+\tfrac12),\\
 c&=\alpha,&c_2&=nr^2+(m+\tfrac12)(1-r)^2+\alpha,\\
 c_4&=nr^2,&d_2&=(n+\tfrac12)r^2+2Dr,\\
 d_4&=(n+D+\tfrac12)r^2.
\end{align*}
Thus the graph on $\xi=q=0$ is
\[
 \begin{split}
 F_g/4={}&[aU+a_4U^2+(m+b_2U+b_4U^2)V]W^2\\
 &+[c+c_2U+c_4U^2+(D+d_2U+d_4U^2)V]Z^2
       -2uv(e_0+e_1U)WZ,
 \end{split}
\]
where $e_0,e_1$ are exactly the original coefficients in (WS3).
Set
\[
 \begin{gathered}
 \beta_s=\frac{Bd_4+\kappa c_4}{\kappa\ell^2},\qquad
 \gamma_s=b_4+\frac{d_4}{\kappa\ell},\\
 j_0=\frac{Bd_4-a_4\kappa^2\ell+d_2\ell}
                   {2\kappa\ell^2},\qquad
 j_1=\frac{d_4}{2\kappa\ell},\qquad J_s(U)=j_0+j_1U,\\
 f_s=\frac12\left(b_2-\frac{c_4}{\ell^2}
                -\frac c{B^2}-\frac{Bd_4}{\kappa\ell^2}\right).
 \end{gathered}                                        \tag{WA4}
\]
The four original scalar contact equations (value and the three
logarithmic derivatives in $u,v,W/Z$) give the exact identity
\[
 F_g/4=(R_1,R_2,T)\,\mathsf S\,(R_1,R_2,T)^T,\qquad
 \mathsf S=
 \begin{pmatrix}
 a+a_4U&uvJ_s&-vJ_s\\
 uvJ_s&m+\beta_sU&u(f_s-j_1U)\\
 -vJ_s&u(f_s-j_1U)&c/B^2+\gamma_sU
 \end{pmatrix}.                                       \tag{WA5}
\]
This is a different scalar identity from the excluded polynomial
two-block splitter (JG7). In particular, the last diagonal retains
the full radial mass $\alpha/B^2$ and the Gram depends on the
independent testing variables. All source coefficients in (WA4)
are bounded; no uncancelled $1/B$ occurs.

\paragraph{Both full first normals and the complete hard block.}
For the $\xi$ normal define
\[
 \begin{gathered}
 p_0=\frac{D+r(n+1/2)}\kappa,\qquad
 p_1=\frac{r(n+D+1/2)}\kappa,\qquad
 q_1=\frac{nr+Bp_1}\ell,\\
 f_x=-mh+(D+\tfrac12)(1-h)-p_0-q_1,\\
 \mathsf L_x=\bigl(uv(p_0+p_1U),\ -m+q_1U,\
                                u(f_x-p_1U)\bigr).
 \end{gathered}                                        \tag{WA6}
\]
For the unnormalized $\eta=Y-huW$ normal set
\[
 \begin{gathered}
 p_y=-(C+D+\tfrac12)(1-h)+Ah,\qquad p_y'=Ah,\\
 k_v=-C(1-h)+Ah+m+(\alpha+\tfrac12)h,\qquad
 k_{uv}=h(A+m+\alpha+\tfrac12),\\
 k_z=-m(1-r)+r(D+\tfrac12)-\alpha-\tfrac12,\\
 q_y=\frac{-k_z-\kappa p_y'-Bk_{uv}}\ell,\qquad
 f_y=k_v-q_y,\\
 \mathsf L_y=\bigl(p_y+p_y'U,\ uvq_y,\
                                      v(f_y+k_{uv}U)\bigr).
 \end{gathered}                                        \tag{WA7}
\]
The two remaining original source stationarity equations, in $X,Y$,
give these exact representations of the complete first normals.
No first-normal term is discarded. Define also
\[
 \begin{split}
 h_{xx}&=n(1+U)(1+V)+m+(D+\tfrac12)UV+\tfrac12(1+V),\\
 h_{xy}&=-uv(m+D+\tfrac12),\\
 h_{yy}&=A(1+U)(1+V)+C(1+V)+D+\tfrac12
                   +mUV+(\alpha+\tfrac12)(1+U)V.
 \end{split}
\]
Then the \emph{entire original four-state parent} is exactly
\[
 F/4=\mathcal R^T\mathsf G(u,v)\mathcal R,\qquad
 \mathcal R=(\xi,q,R_1,R_2,T)^T,
\]
\[
 \mathsf G=
 \begin{pmatrix}
 h_{xx}&h_{xy}/\sqrt A&\mathsf L_x\\
 h_{xy}/\sqrt A&h_{yy}/A&\mathsf L_y/\sqrt A\\
 \mathsf L_x^T&\mathsf L_y^T/\sqrt A&\mathsf S
 \end{pmatrix}.                                       \tag{WA8}
\]
Block notation in (WA8) has sizes $1,1,3$. Every coefficient of
this finite polynomial matrix in $u,v$ is bounded and pure source.
The checker first expands the original signed core and all six
original event packets to verify the complete hard block and both
normal rows. It then verifies (WA5)--(WA7) from the corresponding
four plus two contact equations. Thus (WA8) is not merely a graph
or Hessian certificate.

\paragraph{Strict residue and a legal separated Gram lift.}
For infinitesimal independent $u,v$, the residue of (WA8), in the
printed root order, is
\[
 \overline{\mathsf G}=
 \begin{pmatrix}
 \bar n+\bar m+1/2&0&0&-\bar m&0\\
 0&1&0&0&0\\
 0&0&\bar C+\bar D+1/2&0&0\\
 -\bar m&0&0&\bar m&0\\
 0&0&0&0&\overline{\alpha/B^2}
 \end{pmatrix}\succ0.                                \tag{WA9}
\]
Indeed the nontrivial $2\times2$ determinant is
$\bar m(\bar n+1/2)>0$, and (WA2) supplies the last positive
pivot. The $q$ couplings disappear because $A^{-1/2}$ is
infinitesimal while $Ah$ is bounded. Every other omitted entry
has a factor $u,v$, or $uv$.

For completeness, strict bounded residue here licenses a finite
tensor SOS without a mixed Cholesky or inverse. Let
$\mathsf G_0=\mathsf G(0,0)$, a positive definite pure-source
matrix whose inverse is bounded. Whiten by its pure-source square
root and write the resulting matrix as
\[
 I+\sum_{j=1}^N H_j t_j,
\]
where each $H_j$ is a bounded pure-source symmetric matrix and
each $t_j$ is a nonconstant pure-testing monomial in $u,v$, hence
infinitesimal. Split $I$ into $N$ equal copies and orthogonally
diagonalize each $H_j$ purely in the source field. A resulting
scalar multiplier is $1/N+\ell t_j$, with $\ell$ bounded.
Choose a positive real $R>N|\ell|$ and put
$a=N\ell/R$, $b=Rt_j$. Then $|a|<1$, $|b|<1$, and
\[
 \frac1N+\ell t_j
 =\frac1{2N}\{(1+a)(1+b)+(1-a)(1-b)\}.               \tag{WA10}
\]
Each product is the square of a pure-source square root times
a pure-testing square root. This supplies finitely many legal
tensor roots for (WA8). Only finitely many source eigenvalues
occur, so one real $R$ can be chosen for all of them. If $N=0$
the source Gram itself is already sufficient.

This construction works for either sign and for zero $u$ or $v$,
and permits all relative source/testing rates. Applying the real
identity separately to the real and imaginary state parts proves
the complex-state assertion. It uses pure-testing square roots of
$1\pm Rt_j$, and consequently makes no claim to a global
polynomial SOS in the testing variables. There is therefore no
conflict with the obstruction proved on the physical SO1 family.



\section{\texorpdfstring{R9 source data and the finite coefficient calculation}{R9 source data and the finite coefficient calculation}}
\label{component:pd709027c7805}
% Scoped mathematical certificate; see local_identity_audit.closure.json for exact changes.

This appendix fixes all coefficient arrays used in the local theorem.
It can be read with the original physical-source entry in the supplied
R13 frontier audit. Every displayed formula is source-only.
Put $r=1-e$, $z=aB$, $y=tB$, and
\[
 c_a=1+aBr,\quad S_c=1+B(1-h+ar),\quad
 g=1+a(hB+r-1),\quad j=c_a-(1-h)/t,
\]
\begin{align*}
\Delta={}&a+(1+aB)+(1+aB)B+2aB(1+aB)r+a^3B^2r^2+2a^2Br\\
 &+tB\{1+a+aB+a^2Br(2-r)\},\\
M={}&(1+aB)(1+ar)(h+t)-a(1+tB)e(1-h).
\end{align*}
The physical stationary event coefficients are
\[
\begin{aligned}
n&=\frac{(1+ar)c_aS_c-r\Delta}{2r\Delta},&
m&=\frac{a^2(1+tB)e c_aS_c}{2g\Delta},\\
C&=\frac{c_aM}{2(1-h)\Delta},&
Ah&=\frac{t c_a(1+ar)S_c}{2\Delta},\\
D&=\frac{(1+aB)(1+ar)(1-atBr)S_c}{2tj\Delta},&
\alpha&=\frac{B^2(1+ar)M}{2(1-hB)\Delta}.
\end{aligned}
\]
The following scalar and normal arrays are exact rational functions of
the native stationary source. Each terminal certificate proves the source
bounds it needs at its own independent openings; the arrays themselves
have no bounded-$a$ or old R9-domain restriction.

\paragraph{Scalar arrays.}
The notation $Ah$ in these formulas means the product of $A$ and $h$.
\[
\begin{aligned}
q_s&=nr^2+(m+\tfrac12)e^2,\qquad q=q_s+\alpha,\\
a_w&=(C+D+\tfrac12)(1-h)^2+Ah\,h,\\
b_2&=2mh+C(1-h)^2+Ah\,h+(\alpha+\tfrac12)h^2,\\
b_4&=h^2(m+\alpha+\tfrac12)+Ah\,h,\\
d_2&=(n+\tfrac12)r^2+2Dr,\qquad
d_4=(n+D+\tfrac12)r^2,\\
e_0&=me+D(1-h)+\alpha h+\tfrac12,\\
e_1&=mhe+r(D+\tfrac12)(1-h)+(\alpha+\tfrac12)h,\\
c_F&=(b_2-a^2nr^2)/t^2,\qquad
\lambda=(m-a^2q)/(ta^2B),\\
f_{A0}&=Ah\,h/t^2+d_2-2e_1/t+2anr^2/t,\\
f_{C0}&=(\alpha-c_Ft^2B^2)/(taB^2),\\
f_G'&=-\{a_w/t^2+D-2e_0/t+2aq/t+2c_FtB\}/(aB).
\end{aligned}
\]
The potentially singular $a,B$ factors cancel in the combinations
used in RC1. One direct proof is to clear $\mathcal D$ and the printed
monomials and expand; all three resulting numerator polynomials are
zero. Equivalently, the original value and two first derivatives at
contact give precisely RC1. Thus RC1 is an exact identity, including
arbitrarily small $a$ and literal $n=0$.

\paragraph{All first-normal arrays.}
For the angular weight $q=q_s+\alpha$ define
\[
\begin{aligned}
K_X'&=-nr+(m+\tfrac12)e-aq,\\
K_E'&=a_w/t^2-e_0/t+aq/t,\\
K_Y&=\{-(C+D+\tfrac12)(1-h)+Ah\}/t+D-\alpha,\\
x_3&=2\{(D+\tfrac12)[(1-h)/t-1]+e/2-nr(1+a/t)-mh/t\},\\
x_5&=-2r(n+D+\tfrac12),\\
a_E&=2Ah\,h/t^2-2e_1/t,\qquad b_E=2b_2/t^2,\qquad c_E=2b_4/t^2,\\
a_Y&=2\{Ah/t-me+r(D+\tfrac12)-\alpha-\tfrac12\},\\
b_Y&=2\{-C(1-h)+Ah+m+(\alpha+\tfrac12)h\}/t,\\
c_Y&=2\{Ah+h(m+\alpha+\tfrac12)\}/t.
\end{aligned}
\]
\paragraph{The six exact contact balances.}
With $x=aB$, $y=tB$, the scalar balances are
\[
 f_{A0}=\lambda-2d_4x-b_4y/t^2,\qquad
 f_{C0}=\lambda+b_4y/t^2,\qquad
 f_G'=3\lambda-d_4x+b_4y/t^2.\tag{RC1}
\]
The complete first-normal balances are
\[
 2K_E'+a_Ex+b_Ey+c_Exy+2\lambda x=0,
\]
\[
 2K_X'+x_3y+x_5xy-2t\lambda x=0,\qquad
 2K_Y+a_Yx+b_Yy+c_Yxy=0.\tag{RC6}
\]
Substitution of the native weights and multiplication by the displayed
source denominators gives six zero polynomial numerators. Equivalently
these are the original contact value and its five first derivatives.
They underlie the generic whole-parent identity BA7 (also printed as
R10.9). BA6 explicitly gives its entire hard quadratic. Those identities
hold before any source size assumption; the bounded-$a$, R8 and R10
certificates separately establish their coefficient bounds and pivots.


\section{\texorpdfstring{The shared C1--C2 boundary and a positive source partition}{The shared C1--C2 boundary and a positive source partition}}
\label{component:pc024f5300388}
% Scoped mathematical certificate; see local_identity_audit.closure.json for exact changes.

\label{pc024f5300388:sec:r13-source-overlap}

This section records an exact source-only identity. It does not by itself
admit any new independent testing domain. Its purpose is to identify the
same physical boundary in the two remaining source charts and to retain
the nonnegative slack when that boundary opens.

Use the original weak-C2 variables $B=H+L$, $z,y>0$ and $0<r<1$,
with the unchanged source gaps and nonnegative $n$. Write $e=1-r$,
$K=1+z$, $T_c=B+zr$, and let $\delta$ be the original positive
C2 denominator. Define the positive source polynomial
\[
 P_n=B+B^2+2Bzr+BHz+yB(1+z)+Hrz^2+yz^2re.
 \tag{OP1}
\]
Then the following identity holds on the entire original source:
\[
 \boxed{\quad eP_n=(1+z)(H+y)(B+zr)+2r\delta n.\quad}
 \tag{OP2}
\]
Every factor in the first summand is positive, and the second summand
is nonnegative, including literal $n=0$. In particular the pure-source
weights
\[
 q={ (1+z)(H+y)(B+zr)\over eP_n}>0,
 \qquad s_n={2r\delta n\over eP_n}\ge0
 \quad\hbox{satisfy}\quad q+s_n=1.                   \tag{OP3}
\]
Both square roots are legitimate source-field elements. Neither weight
is assumed to have nonzero residue.

\begin{proof}
The unchanged formula for $n$ is
\[
 2r\delta n=(B+zr)(1+zr)(1+zr+B-H)-r\delta.
\]
Substitution of the printed polynomial for $\delta$, followed by
collecting terms, gives OP2. This is a polynomial identity, so no
source boundary is obtained by continuity and no opening has been
inverted except the displayed positive pure-source factors.
\end{proof}

\paragraph{Exact correspondence of the source contacts.}
Use a subscript $1$ for native C1 notation in this paragraph. The map
\[
 x_1=z,\qquad \eta_1={y\over B},\qquad
 \zeta_1={B\over z},\qquad r_1={H\over B},\qquad
 \theta_1=e,
 \tag{OP4}
\]
with $k_1=\sqrt z$ and the native parameter
$\Delta_1=B/(z\sqrt y)$, gives
\[
 (u_c,v_c;W_c,X_c,Y_c,Z_c)
 =\left(\sqrt z,\sqrt y;{B\over\sqrt{zy}},
                       -r\sqrt z,{H\over\sqrt y},1\right)
\]
from the original native C1 contact formula. Thus it is the same
contact, including all state coordinates. Direct substitution in the
native positive polynomial $P_d$ gives $BP_d=P_n$.
The native C1 beta coefficient at this formal contact therefore obeys
\[
 \beta_1=-{r\delta n\over eP_n}=-{s_n\over2},
 \qquad \beta_1+\tfrac12={q\over2}>0.                 \tag{OP5}
\]
For $n=0$, the original C2 parent has
$\lambda=n=\beta=0$. It is literally an original C1 parent,
up to an optional positive normalization of $C$. The native source
entry and uniqueness at stationary contact identify its coefficients
with OP4. For $n>0$, OP5 makes the formal C1 beta negative. Hence
OP4 does \emph{not} map that source to the nonnegative C1 event cone,
and no physical C1 all-tests theorem may be imported there.



\section{\texorpdfstring{Reconstructed R1 unbounded-$q$ source and scalar certificate}{Reconstructed R1 unbounded-q source and scalar certificate}}
\label{component:p9e725a065583}
% Scoped mathematical certificate; see local_identity_audit.closure.json for exact changes.

\label{p9e725a065583:sec:reconstructed-r1-qo-source-scalar}

\paragraph{Scope of the dictionary and budgets.}
The identities below use the actual native weak-C2 stationary coefficients,
with $d=1$, $\lambda_{\rm event}=\beta=0$, all strict native gaps and
$n\ge0$. Source quantities lie in $K_c$ and testing quantities in a
separately labelled real closed extension $K_t$; all inverses and radicals
are pure source or pure testing. The rational identities have no source-size
restriction. Their R1 estimates assume $B,T,x,y$ infinitesimal; the later
RL and EN certificates establish their own replacement estimates.

For the native source write $a_s,B_s,t_s,r,h$, and put
$c_s=1+a_sB_sr$, $R_s=1+ha_sB_s$, $K_s=1+a_sB_s$.
The formerly used auxiliary ratio is simply
\[
 q_{\rm old}=R_s^2/(a_st_sK_s^2).\tag{RSR1}
\]
No size restriction on $q_{\rm old}$ or $q_z$ is imposed in this
refined dictionary. The actual full-R1 estimates, including arbitrary
positive $q_z=jT$, are justified by R1N1--R1N11.

\subsubsection*{The exact compact source and its uniform bounds}

Define pure-source quantities
\[
\begin{gathered}
 j=\frac{a_sc_s}{R_s},\quad B=\frac{B_sR_s}{c_s},\quad
 T=\frac{t_sc_s}{R_s},\\
 p=\frac{a_s(1-r)}{R_s},\quad H=\frac{1-h}{t_sc_s},\\
 x=jB,\quad y=TB,\quad q_z=jT .
\end{gathered}
 \tag{RSR2}
\]
The native gaps give $0<p<\min(1,j)$ and $0<H<1$.
In the R1 estimates $B,T,x,y$ are positive infinitesimals, while
$j$ and $q_z=jT$ are arbitrary positive. The exact dictionary itself
also applies at the different source scales used in RL and EN. Direct inversion gives
\[
 r=\frac{j-p}{j+px},\quad h=\frac{1-HT}{1+HTx},\quad
 a_s=\frac{j+px}{1+HTx},\quad
 B_s=\frac{B(1+HTx)}{1+pB},\quad
 t_s=\frac{T(1+pB)}{1+HTx}.
 \tag{RSR3}
\]
In particular the compact variables describe the actual source,
including its $r\to0$, $p\to1$, $H\to0$, and $H\to1$ openings.
Set
\[
 K=j+1-p,\qquad E=1+B(HT+p-1),\qquad
 W_0=(1-H)j+1-p.
 \tag{RSR4}
\]
Then $\tfrac12(j+1)<K<j+1$, $E\sim1$, and $W_0>0$.
The following polynomials contain the entire stationary dictionary:
\begin{align*}
\mathcal P={}&j+1+B\{(j+1)^2-p\}\\
 &+TB\{Hj+j+1+B[Hj(2j+2-p)+p(2j+1-p)]\}\\
 &+HjT^2B^2(1+Bp+BHj),\\
\mathcal L={}&K+T(W_0-H)+BT[pK+Hj(1-p)]+BHT^2(W_0-1),\\
\mathcal N={}&p+B(-j^2+3jp-j-p^2+p)\\
 &+BT(Hj^2+Hj-j^2+jp-j+p)\\
 &+B^2T(2Hj^2p-Hj^2-Hjp^2+2Hjp
             -2j^2p+3jp^2-jp-p^3+p^2)\\
 &+B^2HT^2j(Hj-j+p)(1+Bp).
 \tag{RSR5}
\end{align*}
With $\mu=m(1-p)$, the exact six original stationary weights are
\[
\begin{aligned}
 \mu&=\frac{jp(1+y)(1+x)(1+Hy)}{2\mathcal P},&
 n&=\frac{\mathcal N}{2(j-p)\mathcal P},\\
 C&=\frac{\mathcal L}{2HT\mathcal P},&
 A&=\frac{T(1+x)(1+Hy)(K+Hjy)}{2(1-HT)\mathcal P},\\
 D&=\frac{(1+Hy)(K+Hjy)\{1-y[(1-H)j-p]\}}
             {2T(1-H)\mathcal P},&
 \alpha&=\frac{B^2(K+Hjy)\mathcal L}{2E\mathcal P}.
\end{aligned}
 \tag{RSR6}
\]
For example, the original stationary denominator and its source mass
satisfy
\[
 \Delta=\frac{(1+x)\mathcal P}{(1+Bp)(1+HTx)},\qquad
 M=\frac{(1+x)\mathcal L}{(1+HTx)^2}.
\]
The two native strict gaps become
\[
 g=\frac{(1-p)(1+x)}{1+HTx},\qquad
 j_{\rm gap}=\frac{(1-H)(1+x)}{1+Bp}.
\]
These equations verify the source origin of RSR6 directly; no new
stationary solution has been substituted for the original one.

Every term after $j+1$ in $\mathcal P$, divided by $j+1$, is
infinitesimal: use $B(j+1)=B+x\in\mathrm I$, $T\in\mathrm I$,
$p<1$, and $H<1$. The same termwise estimate, using
$K>(j+1)/2$, applies to $\mathcal L-K$. Thus
\[
 \mathcal P\sim j+1,\qquad \mathcal L\sim K.
 \tag{RSR7}
\]
These are relative equivalences, valid also when $j$ is infinitesimal
or unbounded. They show, in particular, that $A/T$, $THC$, and
$T(1-H)D$ are positive bounded units, $\mu/p$ is bounded, and
$h_0=\mu+1/2$ is a positive bounded unit. Moreover $\alpha/B$ is
infinitesimal. Only the displayed source openings can make $m,n,C,D$
unbounded.

The original physicality condition gives a useful uniform restriction.
Regroup the exact polynomial $\mathcal N$ as
\[
\mathcal N=p\{1+B(3j+1-p)+BT(j+1)\}
 -Bj(j+1)\{1+(1-H)T\}+\mathcal R.
 \tag{RSR8}
\]
Here $\mathcal R$ is precisely the last two lines of $\mathcal N$
in RSR5. Since $0<p<\min(1,j)$, each coefficient in its $B^2T$
line is bounded in absolute value by a fixed multiple of $j(j+1)$.
The $B^2HT^2$ line has the same bound, multiplied by the bounded
factor $T(1+Bp)$. Consequently
\[
 |\mathcal R|\leq C_0BT\,Bj(j+1)
\]
for a fixed real $C_0$. All corrections multiplying $p$ and
$Bj(j+1)$ in RSR8 are infinitesimal. As $n\geq0$ is exactly
$\mathcal N\geq0$, it follows that
\[
 \eta_*:=\frac{Bj(j+1)}p\ \hbox{is bounded},\qquad
 \eta_*\in\mathrm I\ \Longrightarrow\
 n\sim\frac{p}{2(j-p)(j+1)}>0.
 \tag{RSR9}
\]
Indeed $\mathcal N=p(1+\varepsilon_1)-Bj(j+1)(1+\varepsilon_2)$
with $\varepsilon_i\in\mathrm I$. This proves both assertions without
dividing by $n$ and includes literal $n=0$ in the bounded-$\eta_*$
alternative. It also gives $nr$ bounded.

\subsubsection*{The complete graph and its actual contact circuit}

The quarter-parent, before any packet is removed, is
\begin{align*}
\mathscr P={}&m\{vW+uvY-X-uZ\}^2
 +CD_v(uW-Y)^2+AD_uD_vY^2\\
&+(D+\tfrac12)(uW+uvX-Y-vZ)^2
 +(\alpha+\tfrac12)D_u(vY-Z)^2+nD_uD_vX^2\\
&+\tfrac12D_v(X+uZ)^2-\tfrac12D_uD_vZ^2,
\qquad D_u=1+u^2,\ D_v=1+v^2.
 \tag{RSR10}
\end{align*}
The source-free rotated state coordinate is $w=(W+uY)/D_u$.
On its nonzero chart set $w=1$ and $z=(Z-uX)/D_u$. Define
the two complete state normals
\[
 N=X+uZ-pD_uv,\qquad M=uW-Y-HD_uvz.
\]
Their exact inverse is
\[
\begin{aligned}
 X&=N/D_u+pv-uz,& Z&=z+uN/D_u+puv,\\
 Y&=u-M/D_u-Hvz,& W&=1+uM/D_u+uHvz.
\end{aligned}
 \tag{RSR11}
\]
The scalar studied here is exactly
$F_{\rm full}=(\mathscr P/D_u)|_{N=M=0}$.
No $n$ packet, $m$ packet, or state-normal term is silently removed.
Substitution of RSR11 into RSR10 gives the complete full-parent
normal expression RFN7 in Section~\ref{component:p261f71af685d}, with
$N_0=N/D_u$, $M_0=M/D_u$ and the retained packet
$mpD_u^2N_0^2$ restored.

Use the separated rows
\[
 U=u/\sqrt x,\quad V=v/\sqrt y,\quad Z_*=z/\sqrt{q_z},\qquad
 Q_*=U-VZ_*,\quad B_*=V-UZ_*,\quad R_*=Z_*-UV,
\]
\[
 \Psi=U^2+V^2+Z_*^2+U^2V^2Z_*^2-4UVZ_*.
 \tag{RSR12}
\]
These are rows; $U,V,Z_*$ are not assumed bounded coefficients.
The actual contact is $(U,V,Z_*;N,M)=(1,1,1;0,0)$.
The exact native formulas make the complete parent value and all five
first derivatives vanish there. On the positive core,
\[
 \Psi=(U-V)^2+Z_*^2(1-UV)^2+2UV(1-Z_*)^2
 \tag{RSR13}
\]
is a legal separated sum of squares: the last root is
$\sqrt2\,(uv)^{1/2}(xy)^{-1/4}(1-Z_*)$, with its radicals in
the two separate real closed fields.

Let the ten positive monomial coefficients and the seven circuit
masses be exactly RB6--RB8. In particular
\[
 M_a=Ax,\quad M_b=\{\mu(1-p)+(n+\tfrac12)p^2\}y,\quad
 M_c=\alpha q_z,\quad M_{ac}=nj^2y,
\]
\[
 f_c=M_b-M_{ac},\quad q_c=M_a-f_c,\quad
 e_c=(n+\tfrac12)p^2y^2-g_c,\quad d_c=M_{b^2c}-k_c.
\]
The original contact balances give the exact scalar identity
\begin{align*}
 F_{\rm full}={}&q_cQ_*^2+M_{ac}B_*^2+r_cR_*^2
  +d_cV^2Q_*^2+e_cV^2B_*^2\\
 &+f_c\Psi+g_cV^2\Psi+k_cZ_*^2(V^2-1)^2.
 \tag{RSR14}
\end{align*}
The identities RB7 and the two odd coefficient balances are independently
checked on the actual source. Thus RSR14 is a finite polynomial identity,
not a Taylor approximation to the contact.

Here is an explicit source-only formula for its radial mass. Put
$\mathcal F=K+T\mathcal F_1+T^2\mathcal F_2$, where
\begin{align*}
\mathcal F_1={}&-2BHjp+BHj+2Bjp-2Bp^2+2Bp\\
 &-Hj-Hp-H+j+1-p^2/j+p/j,\\
\mathcal F_2={}&-B^2H^2jp+B^2Hjp-B^2Hp^2+B^2Hp\\
 &-BH^2j-BH^2p+BHj-BHp^2/j+BHp/j.
 \tag{RSR15}
\end{align*}
Every displayed term divided by $K$ is bounded, using $p/j<1$;
terms containing $Bj$ are even infinitesimal. Hence $\mathcal F\sim K$.
Direct source substitution yields
\[
\begin{aligned}
 f_c&=\frac{Tx^2\mathcal F}{2\mathcal P}>0,&
 q_c&=Ax-f_c\sim\frac{TxK}{2(j+1)}>0,\\
 r_c&=\frac{B^2jT(1-p)(1+x)(1+Hy)\mathcal L}
                {2E\mathcal P}>0,&
 g_c&=\frac{HT^2j^2B^3\mathcal L}{2\mathcal P}>0.
\end{aligned}
 \tag{RSR16}
\]
In particular
\[
 f_c\sim\frac{Tx^2K}{2(j+1)},\qquad
 \frac{r_c}{f_c}\sim\frac{1-p}{j},\qquad
 \frac{g_c}{f_c}=Hy\frac{\mathcal L}{\mathcal F}.
 \tag{RSR17}
\]

For the remaining signed mass define
\begin{align*}
\mathcal K={}&pW_0+
 B(1-p)\{j[H(p+1)-p]-(1-p)(p-H)\}\\
 &+TBH\{(p+1)(1-H)j+(1-p)(BHj+p+1-H)\}\\
 &+T^2B^2H^2j(1-H).
 \tag{RSR18}
\end{align*}
Then
\[
 k_c=-\frac{T^2B^2j(1+x)\mathcal K}{2E\mathcal P}.
\]
The identity
\begin{align*}
\mathcal K-pW_0+Bp(1-p)(j+1)
={}&B(1-p)\{Hj(p+1)+H(1-p)+p^2\}\\
 &+TBH\{(p+1)(1-H)j+(1-p)(BHj+p+1-H)\}\\
 &+T^2B^2H^2j(1-H)
 \tag{RSR19}
\end{align*}
has a nonnegative right side. Since $W_0\geq1-p$ and
$B(j+1)\in\mathrm I$, it proves $\mathcal K>0$ and $k_c<0$.
The displayed finite arrays also give
\[
 \mathcal K/K=O(p+B),\qquad
 |k_c|/(q_cy)=O(p+B).
 \tag{RSR20}
\]
All constants in these estimates are real and uniform over the source
openings; neither $j-p$ nor $1-H$ has been bounded below.

\subsubsection*{Paying the signed masses, with the full current reserved}

Put
\[
 h_0=\mu+\tfrac12,\qquad
 M_e=\frac{nh_0j^2}{n+h_0},\qquad J_c=M_ey,
 \qquad e_0=\tfrac12p^2y^2-g_c.
 \tag{RSR21}
\]
These are pure-source quantities and $0\leq J_c\leq M_{ac}$.
Writing $e_c=np^2y^2+e_0$ preserves the entire second correlated
current packet. The conclusion is the following precise scalar
statement: after subtracting the reserved positive packets
\[
 (M_{ac}-J_c)B_*^2+np^2y^2V^2B_*^2,
 \tag{RSR22}
\]
the remainder is a separated SOS and retains, for example, one half
of $q_cQ_*^2+r_cR_*^2+f_c\Psi+J_cB_*^2$.
When $n=0$ the two reserved packets and $J_c$ vanish literally.
This assertion supplies precisely the effective current mass needed
when the full $N$ normal is later completed against the original $n$
packet. It is not a completion of that normal itself.

Here and below ``infinitesimal cost'' means that the budget multiplies
the original positive packet by a bounded separated Gram coefficient
with zero residue. Finite sums of such costs can be paid while retaining
any fixed fraction smaller than one. Multiplication of the explicit
SOS identities by pure-source nonnegative scalars and pure-testing
squares is legal. The remaining diagonal coefficients are of the form
$1-\varepsilon$, with $\varepsilon$ infinitesimal in the bounded tensor
algebra, so the inherited bounded-Gram completion applies. Its actual
text is the paragraph ``Why a strict bounded Gram suffices'' in
Section~\ref{component:p472f434581f3}; only that field-independent
lemma is imported, not the C14.2 source assumptions in the same file.
For clarity, the elementary step is, for a pure-source infinitesimal
$s$ and bounded pure-testing $t$, to choose real $a,\varepsilon>0$
with $|t|<a$ and $|s|<\varepsilon/a$ and use
\[
 \varepsilon\pm st=\tfrac12\{(\varepsilon/a+s)(a\pm t)
                         +(\varepsilon/a-s)(a\mp t)\}.
\]
Every summand has separated nonnegative factors and hence a separated
square root. Interchange the fields for the other type of infinitesimal
entry and use a finite real diagonal budget. This proves the exact
SOS lifting step used here. No square root or inverse of a mixed sum
is used.

First $e_0^-\leq g_c$. The exact budget
\[
 (1+V^2)\Psi-V^2B_*^2
 =Q_*^2+R_*^2+V^2(1-UVZ_*)^2
 \tag{RSR23}
\]
pays its negative part from the circuit at cost
\[
 \frac{e_0^-}{f_c}+\frac{e_0^-}{f_cy}v^2
 \ \leq\ Hy\frac{\mathcal L}{\mathcal F}
       +H\frac{\mathcal L}{\mathcal F}v^2\in\mathrm I.
\]
Its positive part is retained as a square. Also $d_c>0$ because
$M_{b^2c}>0$ and $k_c<0$; the $d_c$ and $g_c$ terms cause no loss.

For $k_c$, use the exact identity
\[
 Z_*(V^2-1)=-VQ_*-R_*.
 \tag{RSR24}
\]
The $VQ_*$ part is paid from $q_cQ_*^2$ with coefficient
$2|k_c|v^2/(q_cy)$, which is infinitesimal by RSR20. The remaining
$R_*$ part is paid in the following exhaustive three cases. Fixed
real multipliers such as the factor two do not change any cost claim.

\paragraph{Case 1: $1-p$ has a positive real lower bound.}
The exact ratio is
\[
 \frac{|k_c|}{r_c}
 =\frac{T\mathcal K}{(1-p)(1+Hy)\mathcal L}\in\mathrm I.
 \tag{RSR25}
\]
This uses only $\mathcal K/\mathcal L=O(p+B)$ and $T\in\mathrm I$.
Neither current mass nor a bound on $j$ is required.

\paragraph{Case 2: $1-p\in\mathrm I$ and $\eta_*$ has a positive real
lower bound.}
Here $p\sim1$, so $Bj(j+1)$ is a bounded positive unit. Since $B$ is
infinitesimal, $j$ is unbounded and $j^2B$ is a positive bounded unit.
Equations RSR16--RSR20 give
\[
 \frac{|k_c|}{f_c}\in\mathrm I,\qquad
 \frac{|k_c|}{f_cx}
 =O\!\left(\frac{T(p+B)}{j^2B}\right)\in\mathrm I.
 \tag{RSR26}
\]
The exact circuit budget
\[
 (1+U^2)\Psi=R_*^2+B_*^2+U^2Q_*^2+U^2(1-UVZ_*)^2
 \tag{RSR27}
\]
therefore pays $|k_c|R_*^2$ at cost
$|k_c|/f_c+(|k_c|/(f_cx))u^2\in\mathrm I$.
This case includes any intersection with literal $n=0$.

\paragraph{Case 3: $1-p,\eta_*\in\mathrm I$.}
Now $p\sim1$, $j>p$, and RSR9 gives $n>0$ together with its exact
relative estimate. The source mass $M_e$ is a positive bounded unit.
For the lower bound use
\[
 \frac{p^2}{M_e}
 =\frac{p^2}{h_0j^2}+\frac{p^2}{nj^2}=O(1).
\]
For the upper bound use $M_e\leq h_0j^2$ when $j$ is bounded,
and $M_e\leq nj^2\sim p/2$ when $j$ is unbounded. In particular,
\[
 \frac{|k_c|}{f_c}\in\mathrm I,\qquad
 \frac{|k_c|}{q_cy}=O(1),\qquad
 \frac{|k_c|}{J_cx}
 =O\!\left(\frac{T(p+B)(1+1/n)}{j^2}\right)=O(1).
 \tag{RSR28}
\]
The last bound follows directly from
$1/n\sim2(j-p)(j+1)/p$, with $p\sim1$ and $j>p$.
There is an exact separated SOS budget
\[
 2\{\Psi+V^2Q_*^2+U^2B_*^2\}-R_*^2\quad\hbox{is SOS}.
 \tag{RSR29}
\]
For completeness, set $t_*=UV$, $\delta_*=U-V$,
$E_*=Z_*(1-t_*)$, $T_*=t_*(1-Z_*)$, and
$J_*=T_*-\delta_*^2Z_*$. Its SOS decomposition is
\[
 (E_*+T_*)^2+2\delta_*^2+4t_*(1-Z_*)^2
                       +2J_*^2+4t_*\delta_*^2Z_*^2.
\]
All roots are separated as in RSR13. Applying RSR29 costs a multiple
of $|k_c|/f_c$ on the circuit, $|k_c|v^2/(q_cy)$ on $Q_*$,
and $|k_c|u^2/(J_cx)$ on $B_*$. Each is infinitesimal by RSR28.

The three cases exhaust the source because $0<1-p<1$ and $\eta_*$
is bounded. A finite number of the displayed allocations therefore
proves the scalar assertion following RSR22 with a strict retained
fraction. All coefficient bounds depend only on the source and
$u,v$; no bound on the independent state coordinate $z$ was used.
The proof allows arbitrary relative rates of all source and testing
infinitesimals. Literal $n=0$ is covered by Cases 1 and 2; Case 3
excludes it by the original physicality equation itself.



\section{\texorpdfstring{Reconstructed full R1 conductor in the unbounded source-ratio sector}{Reconstructed full R1 conductor in the unbounded source-ratio sector}}
\label{component:p261f71af685d}
% Scoped mathematical certificate; see local_identity_audit.closure.json for exact changes.

\label{p261f71af685d:sec:fn-reconstructed}

\paragraph{Use of the complete normal calculation.}
The parent identities below are exact for the compact stationary dictionary.
The source estimates and square budgets give the full local R1 certificate
under the full-R1 source hypotheses verified in R1N1--R1N11, with arbitrary
positive $jT$. The later RL and EN bodies reuse only the printed exact
identities and supply their own source estimates.

\paragraph{Fields, source scope, and notation.}
Let $K_c,K_t\supset\mathbb R$ be real closed fields, with bounded rings
$B_c,B_t$ and infinitesimal ideals $\mathfrak m_c,\mathfrak m_t$.
The ring of certificate expressions is
$\mathcal T=K_c\otimes_{\mathbb R}K_t$, and its bounded coefficient ring
is $\mathcal B=B_c\otimes_{\mathbb R}B_t$.  A bounded coefficient with
zero real residue belongs to
$\mathfrak n=\mathfrak m_c\otimes B_t+B_c\otimes\mathfrak m_t$.
Every radical or inverse used below is pure source, pure testing, or
an explicitly separated product of such factors.  No order on
$\mathcal T$ and no inverse of a mixed expression is used.

Use the positive compact source coordinates
\[
 x=jB,\qquad y=TB,\qquad q_z=jT,\qquad
 0<p<\min\{1,j\},\qquad 0<H<1,                         \tag{RFN1}
\]
where $B,T,x,y\in\mathfrak m_c$ and $q_z>0$ is arbitrary. These are
the compact coordinates of the complete original R1 sector; the absence
of any $q_z$ restriction is verified in R1N1--R1N11.  The original
strict source gaps remain in force and $n\geq0$, including $n=0$.
For reference, their native substitution is
\[
 r=\frac{j-p}{j+px},\quad
 h=\frac{1-HT}{1+HTx},\quad
 a_s=\frac{j+px}{1+HTx},\quad
 B_s=\frac{B(1+HTx)}{1+pB},\quad
 t_s=\frac{T(1+pB)}{1+HTx}.                              \tag{RFN2}
\]
The six weights $n,m,C,A,D,\alpha$ are always the original source
weights at (RFN2), rather than weights of an auxiliary stationary
source.  Put $\mu=m(1-p)$, $h_0=\mu+\tfrac12$,
$k=(j+1-p)/(j+1)$, and $\eta_*=Bj(j+1)/p$.
In this source sector RS proves
\begin{gather}
 \tfrac12<k<1,\quad \eta_*=O(1),\quad
 \mu=O(p),\quad h_0\asymp1,\quad A=O(T),\quad
 \alpha\asymp B^2(j+1),                                  \tag{RFN3a}\\
 TCH\sim k/2,\qquad T(D+\tfrac12)(1-H)\sim k/2,\qquad
 \frac{n(j-p)}j=O(1),                                   \tag{RFN3b}\\
 f_c\asymp Tx^2,\qquad q_c\asymp Tx,\qquad
 r_c\asymp TjB^2(1-p),                                  \tag{RFN3c}\\
 \eta_*\in\mathfrak m_c\quad\Longrightarrow\quad
 n\sim\frac{p}{2(j-p)(j+1)}>0.                           \tag{RFN3d}
\end{gather}
Here $\asymp$ means that the ratio is a positive bounded unit, and
$\sim$ means that the ratio has residue one.  All estimates are uniform
at the independent openings $H$, $1-H$, $j-p$, and $1-p$.
In particular $B(j+1)=B+x$ and $\alpha$ are infinitesimal.
No comparison between source and testing infinitesimals is allowed.

\paragraph{The local conclusion.}
For all positive independent $u,v\in\mathfrak m_t$, normalize the pure
testing coordinate $w=(W+uY)/(1+u^2)$ to one and allow the other pure
testing coordinate $z=(Z-uX)/(1+u^2)>0$ to be arbitrary.  The conclusion
is a finite tensor sum of squares for the complete original parent,
with both state normals unrestricted.  In particular $z$ need not be
bounded.  The separate full-state sign reduction and outer-testing join
are needed to pass from this local conclusion to the whole R1 row.

\subsubsection*{The complete parent and the correlated source mass}

Write $D_u=1+u^2$, $D_v=1+v^2$.  The original quarter-parent is
\begin{align}
 \mathcal P={}&m(vW-X+uvY-uZ)^2+nD_uD_vX^2
  +CD_v(uW-Y)^2+AD_uD_vY^2\\
 &+(D+\tfrac12)(uW+uvX-Y-vZ)^2
  +(\alpha+\tfrac12)D_u(vY-Z)^2\\
 &+\tfrac12D_v(X+uZ)^2-\tfrac12D_uD_vZ^2.               \tag{RFN4}
\end{align}
Use the two mixed normal rows
\[
 N_0=\frac{X+uZ-pD_uv}{D_u},\qquad
 M_0=\frac{uW-Y-HD_uvz}{D_u}.                            \tag{RFN5}
\]
Their exact inverse, with $w=1$, is
\[
 \begin{aligned}
 X&=N_0+pv-uz,&Z&=z+uN_0+puv,\\
 W&=1+u(M_0+Hvz),&Y&=u-M_0-Hvz.
 \end{aligned}                                           \tag{RFN6}
\]
Only the pure testing unit $D_u$ is inverted.  Retain the nonnegative
packet $mpD_u^2N_0^2$ and define
$G=(\mathcal P-mpD_u^2N_0^2)/D_u$.  Direct substitution gives the
\emph{entire} expression
\begin{align}
 G={}&\mu D_u\{N_0^2-2vN_0+(1-p)v^2\}
       +nD_v(N_0+pv-uz)^2\\
 &+CD_uD_v(M_0+Hvz)^2+A D_v(u-M_0-Hvz)^2\\
 &+(D+\tfrac12)D_u(M_0-(1-H)vz)^2\\
 &+(\alpha+\tfrac12)
   \{u(1-p)v-(1+Hv^2)z-uN_0-vM_0\}^2\\
 &+\tfrac12D_uD_v(N_0+pv)^2
       -\tfrac12D_v(z+uN_0+puv)^2.                      \tag{RFN7}
\end{align}
Thus $G$ is exactly quadratic in $(N_0,M_0)$:
\[
 G=F+2L_NN_0+2L_MM_0+H_NN_0^2+H_MM_0^2
                  +2(\alpha+\tfrac12)uvN_0M_0,          \tag{RFN8}
\]
where $F=G|_{N_0=M_0=0}$ and
\[
 \begin{split}
 H_N&=h_0D_u+\tfrac12v^2+\alpha u^2+nD_v,\\
 H_M&=D_u(CD_v+D+\tfrac12)+AD_v+(\alpha+\tfrac12)v^2.
 \end{split}                                           \tag{RFN9}
\]
There are no omitted higher state-normal terms.

\paragraph{Exact scalar interface, before spending the $n$-mass.}
Set
\[
 U=\frac u{\sqrt x},\quad V=\frac v{\sqrt y},\quad
 Z_* =\frac z{\sqrt{q_z}},\qquad
 Q_*=U-VZ_*,\quad B_*=V-UZ_*,\quad R_*=Z_*-UV,
\]
\[
 \Psi=U^2+V^2+Z_*^2+U^2V^2Z_*^2-4UVZ_*.                 \tag{RFN10}
\]
All three normalized testing variables are positive, but none is
assumed bounded.  A legal separated certificate is
\[
 \Psi=(U-V)^2+Z_*^2(1-UV)^2+2UV(1-Z_*)^2.               \tag{RFN11}
\]
For example, $\sqrt{UV}$ is the product of the source factor
$(xy)^{-1/4}$ and the testing factor $\sqrt{uv}$.

The exact RS/RB scalar identity for the actual source is
\begin{align}
 F={}&q_cQ_*^2+nj^2y B_*^2+r_cR_*^2+d_cV^2Q_*^2
   +(np^2y^2+e_0)V^2B_*^2\\
   &+f_c\Psi+g_cV^2\Psi+k_cZ_*^2(V^2-1)^2,
 \qquad e_0=\tfrac12p^2y^2-g_c.                        \tag{RFN12}
\end{align}
For an entirely algebraic definition of the masses, put
$K_0=CH^2+(D+\tfrac12)(1-H)^2$ and
\[
\begin{array}{ll}
 M_a=Ax,&M_b=\{\mu(1-p)+(n+\tfrac12)p^2\}y,\\
 M_c=\alpha q_z,&
 M_{ab}=\{A+\alpha(1-p)^2+\mu(1-p)+\tfrac12(1-p)^2\}xy,\\
 M_{ac}=nxq_z=nj^2y,&
 M_{bc}=\{K_0+AH^2+2H\alpha+H-\tfrac12\}yq_z,\\
 M_{b^2}=(n+\tfrac12)p^2y^2,&M_{abc}=(K_0+n)xyq_z,\\
 M_{b^2c}=H^2(A+C+\alpha+\tfrac12)y^2q_z,&
 M_{ab^2c}=CH^2xy^2q_z.
\end{array}
\]
Then
\[
\begin{split}
 k_c&=(M_a-M_b+M_{ac}+M_{ab^2c}-M_{bc})/2,\\
 r_c&=M_c-M_b+M_{ac}-k_c,\quad f_c=M_b-M_{ac},\quad q_c=M_a-f_c,\\
 d_c&=M_{b^2c}-k_c,\qquad g_c=M_{ab^2c}.
\end{split}
\]
The source contact identities make (RFN12) exact, including both odd
monomials of $F$.  It is checked coefficientwise against (RFN7), not
only at the contact point.

Define the unnormalized current $b=v-uz/j=\sqrt y B_*$ and the
pure source quantities
\[
 c_b=\frac{nj}{n+h_0},\qquad M_e=\frac{nh_0j^2}{n+h_0}.
                                                               \tag{RFN13}
\]
The whole scalar pieces $nj^2b^2+np^2v^2b^2$ in (RFN12) must be kept
until their correlations with $N_0$ have been completed.  The full
$N_0$-linear term below contains $njb+npv^2b$, so these pieces and the
corresponding hard terms give
\[
 n(N_0+jb)^2+n v^2(N_0+pb)^2.
\]
Complete only the first square, using a pure source inverse:
\[
 h_0N_0^2+n(N_0+jb)^2
   =(n+h_0)(N_0+c_bb)^2+M_e b^2.                       \tag{RFN14}
\]
Keep $nv^2(N_0+pb)^2$ whole.  Put
$\widehat N=\sqrt{n+h_0}(N_0+c_bb)$.
The other hard terms in (RFN9) are the nonnegative whole squares
\[
 h_0(uN_0)^2+\tfrac12(vN_0)^2+\alpha(uN_0)^2.             \tag{RFN15}
\]
They are still kept whole after the shift; no term $nu^2$ or $nv^2$
is treated as a bounded perturbation.

For $n>0$ one has the exact bounded-shift identity
\[
 \frac{c_b^2}{M_e}=\frac{n}{h_0(n+h_0)}=O(1).
                                                               \tag{RFN16}
\]
At $n=0$, set $c_b=M_e=0$ and omit the effective current row.  Formula
(RFN14) remains literal, and no division by $n$ or $\sqrt{M_e}$ is
made in that case.

RS pays the signed $e_0$ and $k_c$ terms in (RFN12) at infinitesimal
cost after (RFN14), with the effective current mass $M_e y$, and leaves
fixed positive real fractions of
\[
 q_cQ_*^2,\qquad M_e b^2,\qquad r_cR_*^2,\qquad f_c\Psi. \tag{RFN17}
\]
The second term is omitted when $n=0$.  This is the precise scalar
allocation needed here; an allocation that spends the raw $nj^2y$
mass before (RFN14) would not suffice.

\subsubsection*{Finite exact budgets for the full normal rows}

In the rest of the proof, a row is called free if any prescribed fixed
real multiple of its square can be purchased at infinitesimal cost
from (RFN17).  This has an exact meaning: a finite SOS identity writes
its square below a finite sum of those retained squares, with bounded
separated coefficients in $\mathfrak n$.  Subtracting any fixed real
multiple changes their strict Gram diagonals only infinitesimally.

\paragraph{Three elementary circuit budgets.}
Direct expansion gives
\begin{align}
 (1+U^2)\Psi-R_*^2
     &=B_*^2+U^2Q_*^2+U^2(1-UVZ_*)^2,                  \tag{RFN18a}\\
 (1+V^2)\Psi-V^2B_*^2
     &=Q_*^2+R_*^2+V^2(1-UVZ_*)^2.                   \tag{RFN18b}
\end{align}
There is also a useful opposite-current budget.  Put
$t_*=UV$, $\delta_*=U-V$, $E_*=Z_*(1-t_*)$,
$T_*=t_*(1-Z_*)$, and $J_*=T_*-\delta_*^2Z_*$.  Then
\begin{align}
 2(\Psi+V^2Q_*^2+U^2B_*^2)-R_*^2
 ={}&(E_*+T_*)^2+2\delta_*^2+4t_*(1-Z_*)^2\\
   &+2J_*^2+4t_*\delta_*^2Z_*^2.                      \tag{RFN18c}
\end{align}
Every term on the right is a legal separated square.  Thus a
nonnegative multiplier $a$ buys $aR_*^2$ at infinitesimal cost whenever
\[
 a/f_c\in\mathfrak n,\qquad a/(q_cy)=O(1),\qquad
 a/(M_e yx)=O(1),                                      \tag{RFN18d}
\]
since the last two diagonal costs have the additional testing factors
$v^2$ and $u^2$.  This criterion is used only when $M_e>0$.

\paragraph{An exact radial budget.}
The row $R_u=v(u^2-x)$ is free in all source subcases, including the
weak opening $1-p\in\mathfrak m_c$.  Here is an explicit finite
certificate, so that this assertion does not rest on pointwise
nonnegativity.  Write $D_*=V(U^2-1)$ and
\[
\begin{split}
 t_7={}&(Z_*-U^2V,\ V-U^2V,\ U-U^2V,\ UZ_*-U^2V,\\
       &\hspace{25mm}UV-U^2V,\ UVZ_*-U^2V,\ U^2-U^2V)^T,\\
 t_3={}&(1-UV,\ V-UV,\ U-UV)^T,
\end{split}
\]
\[
 G_7=\begin{pmatrix}
 4&-2&-3/2&0&-1/2&0&1/2\\
 -2&3&0&-3/2&0&-1&2\\
 -3/2&0&4&3/2&-2&-3&0\\
 0&-3/2&3/2&4&-3/2&0&-3/2\\
 -1/2&0&-2&-3/2&3&2&-1\\
 0&-1&-3&0&2&8&-4\\
 1/2&2&0&-3/2&-1&-4&4
 \end{pmatrix},\qquad
 G_3=\begin{pmatrix}3&-2&-2\\-2&2&2\\-2&2&3\end{pmatrix}.
\]
The exact identity is
\[
 4\Psi+4U^2Q_*^2+4U^2R_*^2-D_*^2
   =t_7^TG_7t_7+4VZ_*(1-U)^2+UZ_*t_3^TG_3t_3.         \tag{RFN19}
\]
The leading principal minors of $G_7$ are
$4,8,101/4,1049/16,907/16,3279/16,743/16$; those of $G_3$ are
$3,2,2$.  Both are positive definite real matrices.  Their real
Cholesky factors and the separated square roots of $VZ_*$ and $UZ_*$
therefore make the right side a finite legal SOS.  Denote it by
$\mathcal W(U,V,Z_*)$.

Since $R_u=x\sqrt y D_*$, the exhaustive source cases are as follows.

\emph{Case 1: $1-p$ is a positive unit.}
Multiply (RFN19) by $x^2y$.  The costs relative to
$f_c\Psi,q_cQ_*^2,r_cR_*^2$ are, respectively,
\[
 O(B),\qquad O(Bu^2),\qquad O(u^2/(1-p)).
\]
All belong to $\mathfrak n$.

\emph{Case 2: $p\sim1$ and $\eta_*$ is a positive unit.}
Then $j$ has a positive real lower bound, and $j^2B$ is a positive
unit.  Adding $4U^2$ times (RFN18a) to (RFN19) gives
\[
 D_*^2\ \preceq_{\rm sos}\
 4(1+U^2+U^4)\Psi+4U^2Q_*^2.
\]
After multiplication by $x^2y$, the costs are
\[
 O\!\left(B+u^2/j+u^4/(j^2B)\right),\qquad O(Bu^2),
\]
again infinitesimal.  The notation $a\preceq_{\rm sos}b$ means
$b-a\in\Sigma\mathcal T^2$ and is only shorthand for the displayed
exact identities.

\emph{Case 3: $p\sim1$ and $\eta_*\in\mathfrak m_c$.}
Formula (RFN3d) gives $n>0$ and $1/M_e=O(1)$; it also gives
$p^2/M_e=O(1)$.  Put $E^{\rm rad}=Z_*(U^2-1)$.
Permuting $V$ and $Z_*$ in (RFN19) yields
\[
 \mathcal W(U,Z_*,V)
   =4\Psi+4U^2Q_*^2+4U^2B_*^2-(E^{\rm rad})^2.
\]
The identity $D_*=UE^{\rm rad}+(U^2-1)B_*$ then gives, exactly,
\begin{align}
 &8U^2\Psi+8U^4Q_*^2+
   \{8U^4+2(U^2-1)^2\}B_*^2-D_*^2\\
 &\hspace{10mm}=2U^2\mathcal W(U,Z_*,V)
                 +\{UE^{\rm rad}-(U^2-1)B_*\}^2.        \tag{RFN20}
\end{align}
After multiplication by $x^2y$, the three costs are
\[
 O(u^2/j),\qquad O(u^4/j),\qquad
 \frac{8u^4+2(u^2-x)^2}{M_e}.
\]
They are infinitesimal.  These three cases prove the asserted free
radial row without assuming a rate for any testing variable.

\paragraph{The smaller current helpers.}
If $\eta_*\in\mathfrak m_c$, (RFN3d) implies
\[
 \frac{p^2}{M_e}
 =\frac{p^2}{h_0j^2}+\frac{p^2}{nj^2}=O(1),
\]
because $p/j<1$ and
$p(j-p)(j+1)/j^2\leq p+p/j$.
Thus $(pv^2b)^2$ costs only $(p^2/M_e)v^4$ of the effective current.
If $\eta_*$ is a positive unit instead, then
\[
 \frac{p^2y}{f_c}=O\!\left(\frac{p+p/j}{\eta_*}\right)=O(1).
\]
Multiply (RFN18b) by $p^2y^2v^2$.  This buys $(pv^2b)^2$ with circuit
cost
\[
 \frac{p^2y}{f_c}(yv^2+v^4)\in\mathfrak n.
\]
The same identity multiplied by $Tp^2u^2y^2$ shows that
$\sqrt T\,puvb$ is free: its cost is either
$(p^2/M_e)Tu^2v^2$ or
$(p^2y/f_c)Tu^2(y+v^2)$.

\subsubsection*{The entire two-normal conductor}

\paragraph{First normal: exact source cancellation and allocation.}
Put
\[
 q=vz-Tu,\qquad R=z-uv/B,\qquad R_u=v(u^2-x),
\]
\[
 K_c=\frac p2+j\{H(\alpha+\tfrac12)-\tfrac12\}+n(p-j),
 \qquad
 K_R=-\mu-(\alpha+\tfrac12)(1-p)+\frac TjK_c.
\]
Both $K_c/j$ and $K_R$ are bounded by (RFN3).  The complete first normal
is exactly
\[
 L_N=\{j(n-\alpha)+p(n+\tfrac12)v^2\}b
                  +K_RR_u+\frac{K_c}{j}uvq.            \tag{RFN21}
\]
In particular the sign of the $\mu$ term in $K_R$ is negative.
For a direct source check, the difference between this formula and
the derivative of (RFN7) is $vS_N$, where
\[
 S_N=np-\mu+\tfrac p2+
       \{-\mu-(\alpha+\tfrac12)(1-p)\}x
       +j(\alpha-n)+yK_c=0.
\]
The last equality is an exact rational identity of the actual native
source under (RFN2), not an asymptotic contact equation.

After the two correlated $n$-squares are formed, the remaining linear
coefficient is
\[
 L_N^{\rm rem}=-j\alpha b+\tfrac12pv^2b
                       +K_RR_u+\frac{K_c}{j}uvq.
\]
The $q$ row is strict because
$q_cQ_*^2=(q_c/(Tx))(q^2/T)$, with $q_c/(Tx)$ a positive unit.
Its coefficient in $L_N^{\rm rem}$ relative to $q/\sqrt T$ is
$(K_c/j)uv\sqrt T\in\mathfrak n$.
The $pv^2b$ helper is free by the preceding paragraph.
For the remaining small row, the exact identity
\[
 j\alpha b=-\frac\alpha B R_u-\alpha uR
\]
has $\alpha/B=O(B+x)\in\mathfrak m_c$.
When $1-p$ is a unit, the square of $\alpha uR$ costs
$O(B^2(j+1)^2u^2/(1-p))\in\mathfrak n$ against $r_cR_*^2$.
When $p\sim1$, use $B_*^2\preceq_{\rm sos}(1+U^2)\Psi$ from
(RFN18a).  The circuit cost of $(j\alpha b)^2$ is
\[
 O\!\left(B^3(j+1)^2+B^2(j+1)^2u^2/j\right)\in\mathfrak n,
\]
since $j$ is bounded below and $B(j+1)=B+x$ is infinitesimal.
Finally, $K_RR_u$ has bounded coefficient; buy an arbitrarily large
fixed real $R_u^2$ pivot using (RFN19)--(RFN20).

The shift $N_0=\widehat N/\sqrt{n+h_0}-c_bb$ introduces no new
unbounded Gram coefficient: $1/\sqrt{n+h_0}=O(1)$, and for $n>0$
$c_b/\sqrt{M_e}=O(1)$ by (RFN16).  The only bounded, potentially
non-infinitesimal couplings are therefore those of $R_u$ to
$\widehat N$ and $\sqrt{M_e}b$.  A sufficiently large fixed real radial
pivot absorbs both while preserving a strict positive real residue.
For $n=0$ the latter coupling is exactly absent.

\paragraph{Second normal: cancellation of both source openings.}
Define
\[
 K_d=CH-(D+\tfrac12)(1-H),\qquad
 C_q=K_d+AH+\tfrac12
       +v^2\{CH(1+u^2)+AH+(\alpha+\tfrac12)H\}.
\]
The complete second normal is
\[
 L_M=C_q q+\mu uvb+v(\alpha-BTCHu^2)R.                  \tag{RFN22}
\]
The actual source satisfies the three exact identities
\[
\begin{split}
 S_0&=T(K_d+AH+\tfrac12)-A=0,\\
 S_1&=jK_d+\mu+xTCH=0,\\
 S_2&=jK_d+xTCH+T\{CH+AH+(\alpha+\tfrac12)H\}
          -A-(\alpha+\tfrac12)(1-p)+\alpha/B=0.
\end{split}
\]
Before imposing these identities, the complete derivative minus
(RFN22) is exactly
$uS_0+u^2vzS_1/j+uv^2(S_2-S_1)$.
This records all state terms and rules out a hidden $j$-weighted
radial error.  In particular $S_1$ gives
$K_d=-\mu/j-BTCH=O(1)$ even when $C$ or $D$ is unbounded.
Consequently $TC_q=O(T+v^2)\in\mathfrak n$.

Normalize
\[
 M_0=\sqrt{TH(1-H)}\,\widehat M.
\]
The full diagonal (RFN9), not only its zero-angle value, obeys
\[
 TH(1-H)H_M\sim k/2>0.                                 \tag{RFN23}
\]
Indeed the $C$ and $D$ terms contribute
$(k/2)D_u\{(1-H)D_v+H\}$ up to infinitesimal bounded coefficients,
and all remaining terms are infinitesimal after this normalization.
This is uniform at both independent openings.
The normalized $C_qq$ coupling has coefficient
$T\sqrt{H(1-H)}C_q\in\mathfrak n$ on the strict row $q/\sqrt T$.
The normalized $\mu uvb$ row is free because $\mu=O(p)$ and
$\sqrt T\,puvb$ has already been purchased.

It remains to allocate the entire opposite row, including its
higher testing terms.  Put
\[
 J=B(j+1)+u^2\in\mathfrak n,\qquad
 a_R=Tq_zv^2(\alpha-BTCHu^2)^2.
\]
The factor $H(1-H)$ may be discarded from the cost bound, since it is
bounded.  Since $\alpha-BTCHu^2=O(BJ)$, the exact squared multiplier
$a_R$ satisfies
\[
 \frac{a_R}{f_c}=O((T/j)v^2J^2),\qquad
 \frac{a_R}{q_cy}=O(v^2J^2).                            \tag{RFN24a}
\]
If $1-p$ is a unit, its direct cost against $r_cR_*^2$ is
$O(Tv^2J^2/(1-p))\in\mathfrak n$.
If $p\sim1$ and $\eta_*$ is a unit, apply (RFN18a); the cost is
$(a_R/f_c)(1+u^2/x)\in\mathfrak n$, using that $j^2B$ is a unit.
If $p\sim1$ and $\eta_*\in\mathfrak m_c$, (RFN3d) gives
\[
 \frac{a_R}{M_e yx}
     =O\!\left(\frac{T(1+1/n)}{j^2}v^2J^2\right)=O(1),  \tag{RFN24b}
\]
because $(1+1/n)/j^2=O(1)$ in this source case.  The other two
conditions in (RFN18d) follow from (RFN24a).  The exact budget
(RFN18c) therefore buys this row at infinitesimal cost as well.
These estimates treat the whole square
$(\alpha-BTCHu^2)^2$; there is no truncation at a testing degree.

\paragraph{Full Gram, legal roots, and restoration.}
After the normalizations, the entire remaining mixed hard term is
\[
 2(\alpha+\tfrac12)uv\sqrt{TH(1-H)}
    \left(\frac{\widehat N}{\sqrt{n+h_0}}-c_bb\right)\widehat M.
\]
Its coefficients relative to $\widehat N$, $\widehat M$, and, when
present, $\sqrt{M_e}b$ are infinitesimal; (RFN16) controls the shift.
All other remaining normal couplings were allocated above.  Choose a
fixed sufficiently large radial pivot, and then purchase all the
finitely many free helpers.  Their scalar costs lie in $\mathfrak n$,
so fixed positive fractions of (RFN17) survive.  The complete bounded
Gram has strict positive real residue: the $\widehat N$ pivot is one,
the $\widehat M$ pivot has residue $\overline{k}/2$, the retained scalar pivots
are positive, and the only bounded radial off-diagonals have been
absorbed by the chosen real radial pivot.

The separated strict bounded-Gram lemma now gives a finite sum of
squares in $\mathcal T$.  Its row entries need not be bounded.  All
budget remainders are explicit finite separated sums of squares, all
source-only completions use positive pure source denominators, and
all state-coordinate inverses are pure testing.  Restore the whole
squares (RFN15), the retained $nv^2(N_0+pb)^2$ square, the scalar
budget remainders, the $mpD_u^2N_0^2$ packet, and the pure testing
factor $D_u$.  This certifies the complete original parent (RFN4).
Every bounded coefficient used in the argument depends only on the
source and on $u,v$.  The arbitrary state ratio $z$ remains inside
retained rows throughout.  Hence the proof has no bounded-state
restriction, and it includes literal $n=0$ without any limiting
argument.



\section{\texorpdfstring{A new complete R1 theorem with no source-ratio split}{A new complete R1 theorem with no source-ratio split}}
\label{component:pbaab204a8ab1}
% Scoped mathematical certificate; see local_identity_audit.closure.json for exact changes.

\label{pbaab204a8ab1:sec:new-r1-all-tests}

\paragraph{Statement and exact source.}
Use the original common-II source, with $d=1$, $\lambda=\beta=0$,
the six native coefficients of the recovered
Section~\ref{component:pd709027c7805}, every original strict gap,
and $n\geq0$.  In the original R1 sector,
\[
 t_s,\quad B_s,\quad a_sB_s,\quad t_sB_s>0
 \quad\hbox{are infinitesimal}.                         \tag{R1N1}
\]
There is no bound on $a_s$, on any original positive opening, or on
the source ratio previously called $q_{\rm old}$.  The conclusion is:
the entire original parent is a finite separated tensor sum of
squares at every independent real projective measurement pair and
every real or complex state.  This proof treats R1 only.

Source and testing quantities belong to separate real closed fields
$K_c,K_t\supset\mathbb R$.  All SOS assertions are in
$K_c\otimes_{\mathbb R}K_t$.  We use the actual recovered strict
bounded-Gram lemma: a bounded finite Gram with positive definite
real residue gives tensor SOS, with no boundedness requirement on
its row entries.  No order or mixed inverse is imposed on the
tensor ring.

Put $c_s=1+a_sB_sr$, $R_s=1+ha_sB_s$, and define
\[
\begin{aligned}
 j&=\frac{a_sc_s}{R_s},\quad B=\frac{B_sR_s}{c_s},\quad
 T=\frac{t_sc_s}{R_s},\quad p=\frac{a_s(1-r)}{R_s},\\
 H&=\frac{1-h}{t_sc_s},\qquad x=jB,\quad y=TB,\quad q_z=jT.
\end{aligned}                                          \tag{R1N2}
\]
The native strict gaps give $0<p<\min(1,j)$ and $0<H<1$.
They also give $1-h<t_sc_s$, so $h\sim1$.
Consequently $B,T,x,y$ are infinitesimal, while
\emph{$q_z$ is an arbitrary positive source quantity}.
In fact $x=a_sB_s$ and $y=t_sB_s$ exactly.
The inverse source substitution is
\[
 r=\frac{j-p}{j+px},\quad h=\frac{1-HT}{1+HTx},\quad
 a_s=\frac{j+px}{1+HTx},\quad
 B_s=\frac{B(1+HTx)}{1+pB},\quad
 t_s=\frac{T(1+pB)}{1+HTx}.                             \tag{R1N3}
\]
All arguments below therefore concern the actual original source.

\subsubsection*{Source bounds without a bound on \(jT\)}

The exact dictionary RSR4--RSR6 in the freshly reconstructed
actual-source and scalar proof remains valid under
(R1N2)--(R1N3), because it is an exact rational substitution.
For completeness the two polynomials controlling its denominators
and current masses are
\begin{align*}
 \mathcal P={}&j+1+B\{(j+1)^2-p\}\\
 &+TB\{Hj+j+1+B[Hj(2j+2-p)+p(2j+1-p)]\}\\
 &+HjT^2B^2(1+Bp+BHj),\\
 \mathcal L={}&K+T\{(1-H)j+1-p-H\}\\
 &+BT[pK+Hj(1-p)]+BHT^2\{(1-H)j-p\},
 \qquad K=j+1-p .
\end{align*}
The estimates
\[
 \mathcal P\sim j+1,\qquad \mathcal L\sim K,\qquad
 \frac12<k:=\frac K{j+1}<1                               \tag{R1N4}
\]
use only $T,B,Bj=x$ infinitesimal and $p,H<1$.
Indeed dividing each correction in $\mathcal P$ by $j+1$ bounds it
by a finite sum of $B+x$, $TB(1+B+x)$ and
$T^2B^2(1+B+x)$; dividing $\mathcal L-K$ by $K$ bounds it
by a finite sum of $T$, $TB$ and $BT^2$.
In particular no product $jT$ is bounded in this argument.

The exact nonnegative numerator of $n$ has the regrouping
\[
 \mathcal N=p\{1+B(3j+1-p)+BT(j+1)\}
       -Bj(j+1)\{1+(1-H)T\}+\mathcal R ,
 \quad |\mathcal R|\le C_0BT\,Bj(j+1).                   \tag{R1N5}
\]
This is the finite polynomial regrouping RSR8; its coefficient bound
uses $p<\min(1,j)$ only.  Since
$n=\mathcal N/[2(j-p)\mathcal P]\ge0$, it follows that
\[
 \eta_*:=\frac{Bj(j+1)}p=O(1),\qquad
 \eta_*\in\mathrm I
 \Longrightarrow n\sim\frac{p}{2(j-p)(j+1)}>0.           \tag{R1N6}
\]
Thus literal $n=0$ is retained.  The same dictionary gives, with
$\mu=m(1-p)$,
\[
\begin{aligned}
 \mu&=O(p),\quad \mu+\tfrac12\asymp1,\quad
 \alpha\asymp B^2(j+1),\quad \frac{n(j-p)}j=O(1),\\
 \frac AT&\sim\frac k2,\quad
 TCH\sim\frac k2,\quad TD(1-H)\sim\frac k2.
\end{aligned}                                          \tag{R1N7}
\]
All statements are uniform at the independent openings $H$, $1-H$,
$j-p$, and $1-p$.

Put
\[
 g_*=\frac{CD}{C+D},\qquad h_*=\frac D{C+D}.
\]
Writing $U_c=TCH$ and $U_D=TD(1-H)$ gives the exact source identity
\[
 Tg_*=\frac{U_cU_D}{(1-H)U_c+HU_D}\sim\frac k2.          \tag{R1N8}
\]
The denominator is a bounded positive unit even when either opening
is infinitesimal.  In particular
$1/C,1/D,1/g_*$ are infinitesimal and
$Ag_*\sim k^2/4$.  Also $(C+D)^{-1/2}=O(\sqrt T)$.

\subsubsection*{The complete small-angle theorem extends to all R1 sources}

The reconstructed proof Section~\ref{component:p261f71af685d}, RFN4--RFN24,
is a full-parent proof, not just a restriction to a scalar graph.
We now justify its application under the larger source domain
(R1N1), rather than importing its former domain label.

First, all of its exact identities hold for the rational dictionary
(R1N3), without any size assumption on $q_z$.  The parent, its two
normals, their quadratic block, the scalar circuit, the radial Gram
and the three source cancellation identities were all proved as
identities of this dictionary.  Second, every uniform source bound
used in that proof is exactly (R1N4)--(R1N7) or the following scalar
bound:
\[
 f_c\asymp Tx^2,\quad q_c\asymp Tx,\quad
 r_c\asymp TjB^2(1-p),\quad
 g_c/f_c=O(Hy),\quad |k_c|/(q_cy)=O(p+B).                \tag{R1N9}
\]
To check the latter bounds without $jT\in\mathrm I$, use the explicit
finite polynomials RSR15 and RSR18: each term of
$\mathcal F_1/K,\mathcal F_2/K$ is bounded using $p/j<1$
and $Bj=x$, and $\mathcal K/K=O(p+B)$ by the same inequalities.
Thus $\mathcal F\sim K$ and all statements (R1N9) follow from
RSR16--RSR20.  No $jT$ factor occurs in these estimates.

The signed scalar allocation retains the effective mass
$M_e=nh_0j^2/(n+h_0)$, $h_0=\mu+1/2$, after the complete
correlated $n$-square completion.  Its three source cases use
$1-p$ a unit, or $p\sim1$ with $\eta_*$ a unit, or
$p\sim1$ with $\eta_*$ infinitesimal.  The proofs of these cases use
only (R1N4)--(R1N9).  In the last case $M_e$ is a unit; in the second,
$j^2B$ is a unit.  Literal $n=0$ occurs only in the first two cases.

For clarity, the following are all potentially nonuniform costs in
the complete normal lift, after cancellation of the source factors.
The symbols $u,v$ denote independent infinitesimal tests.
\[
\begin{array}{c|c}
 \text{purchased row or source case}&\text{bounded diagonal costs}\\ \hline
 v(u^2-x),\ 1-p\text{ unit}
       &O(B),\ O(Bu^2),\ O(u^2/(1-p))\\
 v(u^2-x),\ p\sim1,\ \eta_*\text{ unit}
       &O(B+u^2/j+u^4/(j^2B)),\ O(Bu^2)\\
 v(u^2-x),\ p\sim1,\ \eta_*\in\mathrm I
       &O(u^2/j),\ O(u^4/j),\
            [8u^4+2(u^2-x)^2]/M_e\\
 pv^2(v-uz/j),\ \eta_*\in\mathrm I
       &(p^2/M_e)v^4\\
 pv^2(v-uz/j),\ \eta_*\text{ unit}
       &(p^2y/f_c)(yv^2+v^4)
\end{array}                                                   \tag{R1N10}
\]
Every cost is infinitesimal for its stated case.  The source quotients
$p^2/M_e$ or $p^2y/f_c$ used in the last two lines are bounded by
(R1N6) and $p<\min(1,j)$.

The full opposite normal row has squared multiplier
\[
 a_R=Tq_zv^2(\alpha-BTCHu^2)^2,\qquad
 J=B(j+1)+u^2 .
\]
The occurrence of the possibly unbounded $q_z$ cancels exactly in
its three cost ratios:
\[
 a_R/f_c=O((T/j)v^2J^2),\qquad
 a_R/(q_cy)=O(v^2J^2),\qquad
 a_R/(M_eyx)=O\!\left(T(1+1/n)v^2J^2/j^2\right).        \tag{R1N11}
\]
The last ratio is used only for $p\sim1,\eta_*\in\mathrm I$,
where $(1+1/n)/j^2$ is bounded.  In the other two cases the
direct opposite-current or the $(1+U^2)\Psi$ budget applies exactly
as in RFN24.  The remaining normal coefficients are bounded
multiples of $uv\sqrt T$, $T+v^2$ or the rows in (R1N10);
their estimates require only $B,T,x,y$ infinitesimal.
This checks all complete first-normal and hard-normal coefficients.
The variable $z$ remains inside retained rows, with no bound on it.

Consequently RFN4--RFN24 proves the following enlarged local statement:
for every original R1 source (R1N1), the entire parent is tensor SOS
when $u,v>0$ are infinitesimal, $w=(W+uY)/(1+u^2)=1$, and
$z=(Z-uX)/(1+u^2)>0$ is arbitrary.  The complete source-only
$n$-completion includes $n=0$.  The freshly proved SJ1--SJ10
state/sign reduction therefore gives every state for these small
angles, as well as every literal zero-angle and projective endpoint.

\subsubsection*{Outer testing angles: two exact current identities}

We include the outer mechanism and its precise positivity condition.
Let
\[
 R=uW-Y,\quad C_t=v(uX-Z),\quad B_t=X+uZ,\quad
 N=R+h_*C_t,\quad D_u=1+u^2,\quad D_v=1+v^2,
\]
and let $K$ be the signed quarter-core
\[
 K=\tfrac12(R+C_t)^2+\tfrac12D_u(vY-Z)^2
             +\tfrac12D_v B_t^2-\tfrac12D_uD_vZ^2.       \tag{R1N12}
\]
All unselected original nonnegative events are kept as whole squares.

If $1/v$ is bounded, use
\[
 K_0=2\sqrt{g_*}C_t,\quad
 B_0=\sqrt{2D_v/D_u}\,B_t,\quad
 c_u=D_u^{-1/2},\quad t_u=uD_u^{-1/2},\quad
 \eta_v=\sqrt{D_v}/v .
\]
The exact full selected-parent identity is
\begin{align*}
 4CD_vR^2+4D(R+C_t)^2+4K
 ={}&4(C+D)N^2+4Cv^2R^2+2(R+C_t)^2\\
 &+2D_u(vY-Z)^2+(K_0,B_0)Q(K_0,B_0)^T,\\
 Q={}&
 \begin{pmatrix}
 1-c_u^2\eta_v^2/(2g_*)&t_u\eta_v/\sqrt{2g_*}\\
 t_u\eta_v/\sqrt{2g_*}&1
 \end{pmatrix}.
\tag{R1N13}
\end{align*}
Since $1/g_*$ is infinitesimal and all the testing factors are bounded,
$\overline Q=I$.  This treats every other angle and every state,
including Bob infinity by homogeneous clearing.

For $v$ infinitesimal put $\epsilon=1/g_*$, $q=\sqrt\epsilon$,
$\gamma=Ag_*$, and
\[
 K_1=2\sqrt{g_*}C_t+\sqrt\epsilon Y,\quad
 B_1=B_t-uvY,\quad V_0=vY,\quad Y_1=\sqrt{\epsilon D_u}\,Y,
\]
\[
 L_0=\tfrac12qc_u(K_1-c_uY_1)-t_uvB_1+c_uvV_0.
\]
Then $L_0=\sqrt{D_u}\,v(vY-Z)$ exactly, and
\begin{align*}
 &4CD_vR^2+4D(R+C_t)^2+4AD_uD_vY^2+4K\\
 &\quad=4(C+D)N^2+4Cv^2R^2+2(R+C_t)^2+\mathcal G,\\
 \mathcal G={}&K_1^2+2D_v(B_1^2+V_0^2)
 +(4\gamma D_v-c_u^2)Y_1^2\\
 &+2qv(K_1-c_uY_1)V_0-2L_0^2 .
\tag{R1N14}
\end{align*}
The complete four-root Gram has residue
\[
 \operatorname{diag}(1,2,2,\overline{k}^{\,2}
                                      -\overline{c_u}^{\,2}).    \tag{R1N15}
\]
Thus it is strict whenever the last entry is positive.  In particular
it covers every unbounded $u$, and covers every noninfinitesimal
$u$ whenever $\overline k=1$.  All coefficients in the identity
are bounded functions of the displayed normalized homogeneous Alice
coordinates, so it includes the literal Alice-infinite endpoint.

\subsubsection*{The remaining bounded nonzero Alice chart}

The only outer chart left by (R1N13)--(R1N15) has $v$ infinitesimal,
$u$ a bounded positive unit, and $\overline k<1$.  Since
$1-k=p/(j+1)$, this implies that $p,j$ are positive units and $j$
is bounded.  Therefore $\eta_*\in\mathrm I$ and (R1N6) gives the
exact relative critical limits
\[
 m\sim\frac{jp}{2(1-p)(j+1)},\qquad
 n\sim\frac{p}{2(j-p)(j+1)}.                              \tag{R1N16}
\]
These statements allow either or both coefficients to be unbounded.

Choose bounded positive real caps $m_c\le m$, $n_c\le n$ as specified
below and keep their nonnegative excesses.  Retain all of $\alpha$
as its original square.  Define six rows
\[
 \widehat N=\sqrt{C+D}(R+h_*C_t),\quad
 P=C_t/\sqrt T,\quad Y_0=\sqrt T\,Y,\quad
 X,\quad Z,\quad V=vY.                                  \tag{R1N17}
\]
They need not be independent or bounded.  Put
$\delta=(C+D)^{-1/2}=O(\sqrt T)$, so
$R=\delta\widehat N-h_*\sqrt T\,P$.
The selected \emph{entire} parent is exactly
\begin{align*}
 \mathcal P_{\rm sel}={}&
 \widehat N^2+Cv^2R^2+\tfrac12(R+C_t)^2
       +(Tg_*)P^2+(A/T)D_uD_vY_0^2\\
 &+m_c\{(v/u)R+(D_u/u)V-X-uZ\}^2+n_cD_uD_vX^2\\
 &+\tfrac12D_u(V-Z)^2+\tfrac12D_v(X+uZ)^2
                                      -\tfrac12D_uD_vZ^2 .
\tag{R1N18}
\end{align*}
This is obtained by the exact completion
$CD_vR^2+D(R+C_t)^2=(C+D)N^2+g_*C_t^2+Cv^2R^2$.
It keeps every state-normal term.

The row relation
\[
 PY_0=V(uX-Z)                                            \tag{R1N19}
\]
allows an arbitrary real zero shift
$2b(PY_0-V(uX-Z))$.  Keep the two explicitly nonnegative squares
$Cv^2R^2$ and $(R+C_t)^2/2$ whole.  The remaining complete Gram
on (R1N17) is bounded: $u,1/u$ are bounded, $Tg_*,A/T$ are
units, and the $R$ contribution in the $m_c$ square has coefficients
$\delta v/u$ and $h_*\sqrt T v/u$, both infinitesimal.
Its residue is the direct sum of a unit pivot and the two blocks
\[
 \begin{pmatrix}k_0/2&b\\b&k_0(1+u_0^2)/2\end{pmatrix},\qquad
 Q_3(m_c,n_c;b),                                        \tag{R1N20}
\]
where subscripts denote real residues and $Q_3$ is the quadratic
form in $(X,Z,V)$
\begin{align*}
 Q_3={}&m_c\{(D_u/u)V-X-uZ\}^2+n_cD_uX^2\\
       &+\tfrac12D_uV^2-D_uVZ+\tfrac12(X+uZ)^2
                                            -2bV(uX-Z).
\tag{R1N21}
\end{align*}
In (R1N21), $u=u_0>0$ is real.  The choices of $b,m_c,n_c$
that make both blocks strict are now explicit.

\paragraph{Interior openings.}
Suppose $p_0<1$ and $j_0>p_0$.  Use initially
$m_0=j_0p_0/[2(1-p_0)(j_0+1)]$ and
$n_0=p_0/[2(j_0-p_0)(j_0+1)]$, and take $b=k_0/2$.
The first leading principal minor of $Q_3$ is
$m_0+n_0D_u+1/2>0$; the other two are
\[
 \frac{p_0u_0^2D_uK_0}
      {4(j_0-p_0)(j_0+1)^2(1-p_0)},\qquad
 \frac{p_0u_0^2D_uK_0^2}{8(j_0-p_0)(j_0+1)^3},          \tag{R1N22}
\]
where $K_0=j_0+1-p_0$.  They are positive.
The first block in (R1N20) has determinant
$k_0^2u_0^2/4>0$.  Positive definiteness is open, so fixed
real caps slightly below $m_0,n_0$ preserve strictness and are
smaller than the actual source coefficients.

\paragraph{Only the \(n\) opening is infinite.}
If $j_0=p_0<1$, then $n$ is unbounded and
$m_0=j_0^2/[2(1-j_0)(j_0+1)]$ is finite positive.
At $b=k_0/2=1/[2(j_0+1)]$, restrict (R1N21) to the kernel
$X=0$ of the $n$ packet.  Its first diagonal on $(Z,V)$ is
$u_0^2(m_0+1/2)>0$, and its determinant is
$u_0^2k_0^2/4>0$.  Choose $m_c<m_0$ sufficiently close and
then a sufficiently large fixed real $n_c$.  The whole block is
strict, with both caps below the original weights.

\paragraph{Only the \(m\) opening is infinite.}
If $p_0=1<j_0$, then $m$ is unbounded and
$n_0=1/[2(j_0^2-1)]$ is finite positive.
Take a real $b$ just above $k_0/2=j_0/[2(j_0+1)]$, with
\[
 b<\min\{k_0\sqrt{D_u}/2,\ j_0/[2(j_0-1)]\}.
\]
This interval is nonempty.  The first block in (R1N20) is strict.
On the kernel $(D_u/u)V-X-uZ=0$ of the $m$ packet, use $(X,V)$
coordinates.  The first diagonal of (R1N21) is
$D_u/[2(j_0^2-1)]>0$, and its determinant is
\[
 -\frac{D_u^2[\,2b(j_0-1)-j_0\,][\,2b(j_0+1)-j_0\,]}
        {4u_0^2(j_0^2-1)}>0.                            \tag{R1N23}
\]
Choose $n_c<n_0$ sufficiently close and then a sufficiently large
real $m_c$.  This gives the required full strict block.

\paragraph{Both openings are infinite.}
If $p_0=j_0=1$, both $m,n$ are unbounded and $k_0=1/2$.
Choose
$1/4<b<\sqrt{D_u}/4$.
The common kernel of the two event packets is
$X=0$, $Z=D_uV/u_0^2$.  On it the rest of (R1N21) is exactly
\[
 \frac{D_u(4b-1)}{2u_0^2}V^2>0.                        \tag{R1N24}
\]
A sufficiently large real cap on both packets therefore gives
a strict block.

The elementary finite-cap fact used in these cases is as follows:
if a real positive semidefinite event Gram is positive definite on
the complement of its kernel, and the remaining real quadratic
form is positive definite on that kernel, a sufficiently large
real multiple of the event Gram makes their sum positive definite.
It follows by a block Schur complement, and works for the sum of
two event Grams as well.  Thus the choices above prove strictness
of the \emph{complete} bounded Gram (R1N18), without restricting
any normal or state coordinate.

\subsubsection*{Exhaustive testing and state join}

The sign gauges SJ3 preserve every original source packet and the
whole core, so affine $u,v$ may be taken nonnegative.
If $v$ is not infinitesimal, (R1N13) applies.
If $v$ is infinitesimal and $u$ is not infinitesimal, either
(R1N15) applies or the complete six-root argument
(R1N16)--(R1N24) applies.  These cases exhaust the outer angles.
If $u,v>0$ are both infinitesimal, the enlarged local FN theorem
proved above applies after SJ's full-state reduction to $w=1,z>0$;
SJ gives the other state signs and the axes directly.
Literal homogeneous measurement endpoints are given by SJ10 or
the homogeneous current identities, without a limiting argument.
The zero state is immediate.  For a complex state, the real
parent matrix splits its value into its real and imaginary
quadratic evaluations, to which the same proof applies.

Restore every capped excess, the original $n,m,\alpha$ packets
not used in a particular chart, and all cleared positive pure
testing homogeneous factors.  All normalizations, radicals and
inverses are pure source or pure testing.  This proves the
complete original R1 source row at all independent tests and
states, without a source-ratio split or a comparison of source
and testing rates.

\paragraph{Verification scope.}
The companion checker verifies the exact source-coordinate map,
the harmonic source identity, both full outer-current identities,
the new full six-root parent representation, the zero relation,
all positive-minor and kernel formulas, and their original
source coefficients.  The source bounds and exhaustive chart
argument above are analytical.  This theorem uses actual complete
source and FN proof texts, not a historical status receipt.
It makes no assertion about the still-separate R2--R4 rows.


\section{\texorpdfstring{A fresh complete outer-testing certificate for the original R2 source}{A fresh complete outer-testing certificate for the original R2 source}}
\label{component:p8f8780350b13}
% Scoped mathematical certificate; see local_identity_audit.closure.json for exact changes.

\label{p8f8780350b13:sec:r2-rederived-outer-join}

This proof treats every independent testing pair outside the positive
double-infinitesimal angle chart of the original R2 source.  It uses an
explicit full-parent Bell reflection and the two-current and six-root
identities that were freshly proved for R1.  Only those universal algebraic
identities and their real matrix arguments are reused.  No assertion that
the reflected source is an original R1 source is made.

\subsubsection*{Actual source and exact reflected parent}

The source is the original weak-C2 tuple, with $d=1$, $\lambda=\beta=0$,
the six coefficients in (RER1), all original strict gaps, $n\geq0$,
$a,B,t_s>0$, and $0<r,h<1$.  R2 means
\[
t_s\text{ unbounded},\qquad B, z=aB, y=t_sB\text{ infinitesimal}.
\tag{R2O1}
\]
Put $S=1/t_s$, $l=1-h$, $e=1-r$, $R_s=1+hz$,
$g_s=R_s-ae>0$, and $k=(1+ar)/(1+a)$.  The rederived exact source
lemma (RER1)--(RER7) gives
\[
k\text{ a positive bounded unit},\qquad k<1,\qquad
SC l\sim SAh\sim D/S\sim k/2,
\qquad \Delta/(1+a)\sim1.
\tag{R2O2}
\]
It also gives $ay$ infinitesimal, $Q\sim1$, and the exact expressions
\[
n=\frac{eQ-BK_s}{2r\Delta},\qquad
m=\frac{a^2e(1+y)c s_c}{2g_s\Delta},\qquad
BK_s=(1+a)y(1+hS)(1+z),
\tag{R2O3}
\]
where $c=1+zr\sim1$ and $s_c=1+B(l+ar)\sim1$.
There is no positive lower bound assumed on $r,h,l$, or $g_s$.

For arbitrary finite testing angles define the unnormalized testing
state transformation
\[
\begin{split}
G\psi=(&uvW-uX-vY+Z,\ -vW+X-uvY+uZ,\\
       &-uW-uvX+Y+vZ,\ W+vX+uY+uvZ).
\end{split}
\tag{R2O4}
\]
It has $G^TG=D_uD_v I$.  The Bell transformation on observables is
\[
(A'_0,A'_1,B'_0,B'_1)=(-B_1,-B_0,-A_1,-A_0).
\]
After the displayed real local basis gauges it sends the angles to
$(u',v')=(v,u)$ and the source coefficients to
\[
(\lambda',n',m',C',A',D',\alpha',\beta')
 =(\alpha,m,n,C,D,A,\lambda,\beta).
\tag{R2O5}
\]
Direct substitution in the whole parent gives
\[
F_{c'}(v,u;G\psi)=D_uD_v F_c(u,v;\psi).
\tag{R2O6}
\]
The checker verifies the coefficient of every original packet in this
single full polynomial identity, as well as $G^TG$.  The transformed
$\lambda'=\alpha$ event is kept as its whole positive square.
Because $D_uD_v>0$ is pure testing, a certificate on the left restores
the right by dividing each root by $\sqrt{D_uD_v}$ within $K_t$.
No source or testing contact is identified in this operation.

In the next two subsections write the transformed testing angles again
as $u,v$, and denote the reflected coefficients by
\[
\mathsf C=C,\quad\mathsf D=A,\quad\mathsf A=D,
\quad\mathsf m=n,\quad\mathsf n=m.
\]
Thus the symbols on the right refer to the original R2 source.
Define the pure-source quantities
\[
g_* =\frac{\mathsf C\mathsf D}{\mathsf C+\mathsf D},\qquad
h_* =\frac{\mathsf D}{\mathsf C+\mathsf D},\qquad
\gamma=\mathsf A g_*.
\tag{R2O7}
\]
If $U_C=SC l$ and $U_A=SAh$, then exactly
\[
Sg_* =\frac{U_CU_A}{hU_C+lU_A}\sim k/2,
\qquad \gamma\sim k^2/4.
\tag{R2O8}
\]
The denominator is a positive bounded unit even when $h$ or $l$ is
infinitesimal.  In particular $1/\mathsf C$, $1/\mathsf D$,
$1/g_*$ are infinitesimal,
$(\mathsf C+\mathsf D)^{-1/2}=O(\sqrt S)$, and $0<h_*<1$.
These are the complete source bounds needed for the first two outer
current identities.

\subsubsection*{Universal outer current identities}

Set
\[
R_t=uW-Y,\quad C_t=v(uX-Z),\quad B_t=X+uZ,\quad
J=R_t+C_t,\quad N=R_t+h_*C_t,
\]
and let $K$ be the exact signed quarter-core
\[
K=\tfrac12J^2+\tfrac12D_u(vY-Z)^2
 +\tfrac12D_vB_t^2-\tfrac12D_uD_vZ^2.
\tag{R2O9}
\]
All unselected original events, including $\lambda'$, are retained
as whole nonnegative packets.

If $1/v$ is bounded, put
\[
K_0=2\sqrt{g_*}C_t,\quad B_0=\sqrt{2D_v/D_u}\,B_t,
\quad c_u=D_u^{-1/2},\quad t_u=uD_u^{-1/2},\quad
\eta_v=\sqrt{D_v}/v.
\]
The entire selected-parent identity is
\begin{align*}
4\mathsf C D_vR_t^2+4\mathsf D J^2+4K
={}&4(\mathsf C+\mathsf D)N^2+4\mathsf C v^2R_t^2
 +2J^2+2D_u(vY-Z)^2\\
&+(K_0,B_0)
\begin{pmatrix}
1-c_u^2\eta_v^2/(2g_*)&t_u\eta_v/\sqrt{2g_*}\\
t_u\eta_v/\sqrt{2g_*}&1
\end{pmatrix}
\binom{K_0}{B_0}.
\tag{R2O10}
\end{align*}
The finite Gram is bounded and has real residue $I$, because $c_u,t_u$
and $\eta_v$ are bounded while $1/g_*$ is infinitesimal.  The actual
separated bounded-Gram lemma makes it tensor SOS.  The roots themselves
need not be bounded.  This gives every state and every other finite angle
when transformed Bob is noninfinitesimal, equivalently when the original
Alice angle is noninfinitesimal.

For positive infinitesimal transformed $v$, put
$q=g_*^{-1/2}$ and define
\[
K_1=2\sqrt{g_*}C_t+qY,\quad
B_1=B_t-uvY,\quad V_0=vY,\quad Y_1=q\sqrt{D_u}Y,
\]
\[
L_0=\tfrac12q c_u(K_1-c_uY_1)-t_uvB_1+c_uvV_0.
\]
Then $L_0=\sqrt{D_u}\,v(vY-Z)$ exactly, and
\begin{align*}
&4\mathsf C D_vR_t^2+4\mathsf D J^2
       +4\mathsf A D_uD_vY^2+4K\\
&\quad=4(\mathsf C+\mathsf D)N^2+4\mathsf C v^2R_t^2
 +2J^2+\mathcal G,\\
\mathcal G={}&K_1^2+2D_v(B_1^2+V_0^2)
 +(4\gamma D_v-c_u^2)Y_1^2\\
&+2qv(K_1-c_uY_1)V_0-2L_0^2.
\tag{R2O11}
\end{align*}
This is the complete four-root Gram, including all normal terms.  Its
residue in the ordered frame $(K_1,B_1,V_0,Y_1)$ is
\[
\operatorname{diag}(1,2,2,k_0^2-c_{u,0}^2).
\tag{R2O12}
\]
Subscript $0$ denotes real residue.  Hence it is strict whenever its
last entry is positive.  This includes every unbounded transformed
$u$ and every bounded positive-unit transformed $u$ when $k_0=1$.
All inverse and square-root factors in these identities belong to one
field at a time.

\subsubsection*{The remaining unit angle and both coefficient openings}

It remains to treat transformed $v$ infinitesimal, transformed $u$ a
bounded positive unit, and $k_0<1$.  The source identity
\[
1-k=\frac{ae}{1+a},\qquad 0<ae<R_s\sim1
\tag{R2O13}
\]
shows that $a,e$ are positive bounded units: if $a$ were infinitesimal
or unbounded, or if $e$ were infinitesimal while $a$ is bounded, then
$1-k$ would be infinitesimal.  Thus in the current chart $a$ is a
positive unit, $e$ is a positive unit, and both $1+a$ and $\Delta$
are units.

Use the opening-safe source coordinates
\[
j=R_s/a,\qquad p=e,\qquad j-p=g_s/a>0.
\tag{R2O14}
\]
They satisfy $0<p<\min(1,j)$, $j,p$ are positive bounded units, and
the residue of $(j+1-p)/(j+1)$ equals $k_0$.  The exact relative
critical coefficient identities are
\[
\frac{\mathsf m}{jp/[2(1-p)(j+1)]}
 =\left(Q-\frac{BK_s}{e}\right)\frac{R_s+a}{R_s\Delta},
\qquad
\frac{\mathsf n}{p/[2(j-p)(j+1)]}
 =(1+y)c s_c\frac{R_s+a}{\Delta}.
\tag{R2O15}
\]
Their right sides have residue one.  Indeed $BK_s/e$ is infinitesimal
here by (R2O3), and all other normalized factors have the limits stated
in (RER).  No division by $r$, $g_s$, $n$, or an infinitesimal gap
was used to assert a bounded error.  The displayed quotients cancel
those openings exactly before the limit is taken.  In particular the
critical chart itself forces original $n>0$, even though the theorem
as a whole retains $n=0$.

Choose positive real caps $m_c\leq\mathsf m$ and
$n_c\leq\mathsf n$, with the choices specified below.  Keep both
nonnegative excesses as the original whole packets.  Define six rows
\[
\widehat N=\sqrt{\mathsf C+\mathsf D}\,N,\quad
P=C_t/\sqrt S,\quad Y_0=\sqrt S\,Y,
\quad X,\quad Z,\quad V=vY.
\tag{R2O16}
\]
They can be unbounded or linearly dependent.  If
$\delta=(\mathsf C+\mathsf D)^{-1/2}$, then
$R_t=\delta\widehat N-h_*\sqrt S P$.  The selected entire
quarter-parent is exactly
\begin{align*}
\mathcal P_{\rm sel}={}&\widehat N^2+\mathsf C v^2R_t^2
 +\tfrac12J^2+(Sg_*)P^2
 +(\mathsf A/S)D_uD_vY_0^2\\
&+m_c\{(v/u)R_t+(D_u/u)V-X-uZ\}^2
 +n_cD_uD_vX^2\\
&+\tfrac12D_u(V-Z)^2+\tfrac12D_v(X+uZ)^2
 -\tfrac12D_uD_vZ^2.
\tag{R2O17}
\end{align*}
The full-state row relation is
\[
PY_0=V(uX-Z).
\tag{R2O18}
\]
Add $2b(PY_0-V(uX-Z))=0$ for a real $b$.  Keep the two displayed
positive squares $\mathsf C v^2R_t^2$ and $J^2/2$ whole.  The
remaining Gram is bounded: $u,1/u$ are bounded, $Sg_*$ and
$\mathsf A/S$ are units, and the $R_t$ contributions in the
$m_c$ square have coefficients $\delta v/u$ and
$h_*\sqrt S v/u$, both infinitesimal.  Its residue is the direct sum
of a unit pivot, the matrix
\[
\begin{pmatrix}k_0/2&b\\b&k_0(1+u_0^2)/2\end{pmatrix},
\tag{R2O19}
\]
and the following real quadratic form on $(X,Z,V)$:
\[
Q_3=m_c\{(D_u/u)V-X-uZ\}^2+n_cD_uX^2
 +\tfrac12D_uV^2-D_uVZ+\tfrac12(X+uZ)^2
 -2bV(uX-Z).
\tag{R2O20}
\]
In this formula $u=u_0>0$ is real.  We now give all finite-cap choices;
no opening of either source coefficient is suppressed.

If $p_0<1$ and $j_0>p_0$, use initially
\[
m_0=\frac{j_0p_0}{2(1-p_0)(j_0+1)},\quad
n_0=\frac{p_0}{2(j_0-p_0)(j_0+1)},\quad b=k_0/2.
\]
The three leading principal minors of $Q_3$ are
\[
m_0+n_0D_u+\tfrac12,\quad
\frac{p_0u_0^2D_u(j_0+1-p_0)}
 {4(j_0-p_0)(j_0+1)^2(1-p_0)},\quad
\frac{p_0u_0^2D_u(j_0+1-p_0)^2}
 {8(j_0-p_0)(j_0+1)^3},
\tag{R2O21}
\]
all strictly positive.  Matrix (R2O19) has determinant
$k_0^2u_0^2/4>0$.  Choose real caps slightly below $m_0,n_0$;
openness of real positive definiteness and (R2O15) ensure they remain
strict and are below the actual coefficients.

If $j_0=p_0<1$, the coefficient $\mathsf n$ is unbounded and
$m_0=j_0^2/[2(1-j_0)(j_0+1)]$ is finite positive.  Take
$b=1/[2(j_0+1)]$.  On the kernel $X=0$ of the $\mathsf n$ packet,
the first diagonal and determinant of $Q_3$ on $(Z,V)$ are
\[
u_0^2(m_0+1/2)>0,\qquad
u_0^2/[4(j_0+1)^2]>0.
\tag{R2O22}
\]
Choose $m_c<m_0$ sufficiently close and then a sufficiently large
real $n_c$.  A real block Schur complement makes the full matrix strict.

If $p_0=1<j_0$, the coefficient $\mathsf m$ is unbounded and
$n_0=1/[2(j_0^2-1)]$ is finite positive.  Choose a real $b$ with
\[
\frac{j_0}{2(j_0+1)}<b<
\min\left\{\frac{k_0\sqrt{D_u}}2,
                 \frac{j_0}{2(j_0-1)}\right\}.
\]
The interval is nonempty because $u_0>0$.  On the kernel
$(D_u/u)V-X-uZ=0$ of the $\mathsf m$ packet, the first diagonal
and determinant in $(X,V)$ are
\[
\frac{D_u}{2(j_0^2-1)}>0,\qquad
-\frac{D_u^2[2b(j_0-1)-j_0][2b(j_0+1)-j_0]}
       {4u_0^2(j_0^2-1)}>0.
\tag{R2O23}
\]
Choose $n_c<n_0$ sufficiently close, then a sufficiently large real
$m_c$.  Matrix (R2O19) and the full $Q_3$ are strict.

If $p_0=j_0=1$, both coefficients are unbounded and $k_0=1/2$.
Choose $1/4<b<\sqrt{D_u}/4$.  On the common event kernel
$X=0$, $Z=D_uV/u_0^2$, the remaining form is
\[
\frac{D_u(4b-1)}{2u_0^2}V^2>0.
\tag{R2O24}
\]
Sufficiently large real caps on both events make the full form strict.
The finite-cap step in each case is the elementary real fact that a
quadratic form positive definite on the kernel of a nonnegative event
Gram becomes positive definite after a sufficiently large multiple of
that event Gram is added; a block Schur complement proves it.

These four cases exhaust $0<p_0\leq\min(1,j_0)$ with positive $j_0,p_0$.
The separated bounded-Gram lemma now proves the entire selected parent
(R2O17), with every state normal retained.  Restore the original excesses
and all unused packets.  The real caps are used only on this testing
chart: no globally physical capped source is asserted, so this argument
does not conflict with the finite-contact cap exclusion (CAP15).

\subsubsection*{Exact outer and projective join}

Return to the original testing angles.  The full-parent sign gauges
(SJ3) permit $u,v\geq0$ without changing source coefficients.  Literal
zeros and projective infinities are the exact signed-core squares
(SJ10), plus the untouched event squares.  For finite positive angles:
\begin{itemize}
\item if the original $u$ is noninfinitesimal, use (R2O10);
\item if the original $u$ is infinitesimal and $v$ is unbounded,
      use (R2O11)--(R2O12);
\item if the original $u$ is infinitesimal and $v$ is a positive
      bounded unit, use (R2O11) when $k_0=1$, and the complete
      six-root argument (R2O13)--(R2O24) when $k_0<1$.
\end{itemize}
These are exactly all finite positive angle pairs outside the case in
which both are infinitesimal.  Every argument is a whole-parent state
Gram, so no additional state chart or state-ratio restriction is present.
Equation (R2O6) restores the original frame and positive pure-testing
factor.  The zero state and the real/imaginary split for complex states
are immediate.  No source-to-testing rate comparison was made.

\paragraph{Assembly interface and remaining obligation.}
The original R2 outer testing region, every state, and all literal
projective endpoints are proved here.  A full-parent certificate on the
remaining original positive $u,v$ double-infinitesimal chart, for the
same unchanged source and all states, combines with this result to give
the whole original R2 row.  That local certificate is a separate
hypothesis until its actual proof is independently accepted.  This
outer theorem does not itself close any other source row or the global
finite semidefinite-lift criterion.


\section{\texorpdfstring{R2 with an unbounded source ratio and a bounded relative current}{R2 with an unbounded source ratio and a bounded relative current}}
\label{component:pb88ccd68a4bf}
% Scoped mathematical certificate; see local_identity_audit.closure.json for exact changes.

\label{pb88ccd68a4bf:sec:r2-thin-new}

\paragraph{Statement and domain.}
Use the actual original weak-C2 source RER1--RER8, normalized by
$d=1$, with $\lambda=\beta=0$ and every original strict gap.
Suppose $t,a$ are positive unbounded, while $B,aB,tB$ are
positive infinitesimals, and put $e=1-r$.  Assume
\[
             \lambda_n=n/(tB)=O(1),\qquad n\geq0.       \tag{TU1}
\]
The entire original parent is tensor SOS for every independent
projective test and every state.  Literal $n=0$ is included.
No comparison of a testing quantity with a source quantity is made.
This theorem treats precisely this source sector; it is not a
completion of the other R2 sectors or the global finite-lift proof.

The bounded-Gram and finite toric allocation lemmas used below are
the actual retained SC16--SC20 and QI1--QI10 proofs.  We use their
printed finite algebra and prove its source bounds on the present
domain.  We do not apply the physical C1 theorem to a formal source
whose event coefficient could be negative.

\subsubsection*{An exact common contact and bounded source comparison}

Set
\[
 S=1/t,\quad z=1/a,\quad q=a(tB),\quad w=ae,\qquad
 b=\sqrt S,\quad g=\sqrt z,\quad c=\sqrt q.
                                                               \tag{TU2}
\]
Thus $B=Szq$, $e=zw$, the actual squared angles are
$u_c^2=Sq$ and $v_c^2=zq$, and $S,z,q$ are positive
infinitesimals.  Keep $h$ as the original C2 source parameter.
The strict $m$ gap gives
\[
 R=1+Sqh,\quad M=R-w>0,
 \qquad 0<w<R\sim1.                                      \tag{TU3}
\]
Both original openings $h,1-h$ remain arbitrary positive quantities.

The exact RER2 identity, divided by $e>0$, gives
\[
 Q_{\rm RER}
 =\tau\{(1+z)(1+Sh)(1+Sq)+2r(z\Delta)\lambda_n\},
 \qquad \tau=q/(zw),\quad \lambda_n=n/(qz).
\]
Here $Q_{\rm RER}\sim1$, $z\Delta\sim1$ and $r\sim1$.
The bracket is a positive bounded unit, since $\lambda_n\geq0$
is bounded and its first term is a positive unit.  Consequently
\[
 \tau\text{ is a positive bounded unit},\qquad
 \overline\tau=\frac1{1+2\overline\lambda_n},\qquad
 q=zw\tau.
                                                               \tag{TU4}
\]
No division by $n$ has occurred.  This sector includes the earlier
thin condition $\rho_n=n/e^2\in\mathrm I$: in that case the same
RER identity gives $\tau\sim1$ and
$\lambda_n=\rho_nw/\tau\in\mathrm I$ directly, without first
assuming (TU1).

For the comparison only, substitute into the native C1 rational
formulas NC1--NC4
\[
 k^2=Sq,\quad \eta=1/S,\quad \zeta=z,\quad
 r_{\rm C1}=h,\quad t_{\rm C1}=w,\quad
 a_{\rm C1}=S(1-h).                                     \tag{TU5}
\]
Their $\theta$ is $e=zw$ and their $R,M$ are exactly (TU3).
Write a superscript $0$ for these formal C1 coefficients.
All their source denominators have the requisite positive signs:
\[
\begin{gathered}
 K=1+Sq,\quad J=1+Sq(1-e)-S(1-h),\quad
 E=1-q(1-e),\quad Q=1-Sqzh,
\end{gathered}
\]
and $J,E,Q$ are positive bounded units.  Their common polynomial is
\[
 P=1+qz+Sq(2+z-2e)+Sq^2(z+e-e^2)+S^2hq^2(1+z-e).
                                                               \tag{TU6}
\]
In particular $P\sim1$.  We make no sign claim for $\beta^0$.

The two sources have exactly the same state contact.  With $Z_c=1$,
\[
 u_c=bc,\quad v_c=cg,\quad W_c=bg,\quad
 Y_c=hu_cW_c,\quad X_c=-(1-e)u_c.                       \tag{TU7}
\]
This follows either by substitution in the two native dictionaries,
or from $Y_c=hu_cW_c$, $X_c=-ru_cZ_c$ and
$Z_c/W_c=\sqrt{at}$ in the actual C2 source.  In particular
the comparison does not move the contact.

Use the actual, not formal, normalized coefficients
\[
 c_b=SC,\ a_b=SA,\ D_b=D/S,\ \alpha_b=\alpha/(Sz),\quad
 h_2=mzM,\quad h_b=zw/2.
                                                               \tag{TU8}
\]
They satisfy $m=h_2/(zM)$ and $\beta+1/2=h_b/(zw)$.
Use the same normalizations with superscript $0$, and let $\delta$
mean actual minus formal.  The following explicit polynomials make
all needed cancellations finite and opening-safe.  Put $G=1+z-e$,
and define
\begin{align*}
 P_C={}&1+qz+S\{1+q(2+2z-e)+q^2z\}\\
 &+S^2\{q[h(1+z)+1-e]
          +q^2[hez-e(1+z)+2z+1]\}+S^3hq^2G,\\
 P_A={}&1+qz+S\{q(2+z-2e)+q^2z(1-e)\}
                 +S^2q^2G(G-hz),\\
 P_D={}&1+q(z+e)+S\{q[h(1+z)+1-e]+q^2(z+e-e^2)\}
                 +S^2hq^2G,\\
 P_m={}&1+qz+S\{q(2+z-2e)+q^2z\}
                 +S^2q^2\{(1-e)^2+z(1-e)+hez\},\\
 P_b={}&(1+z)(1+qz)+Sq(1+z)(2+z-2e)\\
       &+Sq^2z(1+z-e^2)+S^2q^2G^2.
\end{align*}
Then the exact rational identities are
\[
\begin{aligned}
 (1-h)\delta c_b&=-\frac{n(1-e)P_C}{zP},&
 h\delta a_b&=-\frac{n(1-e)P_A}{zP},\\
 \delta D_b&=-\frac{n(1-e)(1+qz)P_D}{zJP},&
 \delta h_2&=\frac{n(1-e)(1+qz)P_m}{zP},\\
 \delta h_b&=\frac{n(1-e)P_b}{zP},&
 \delta\alpha_b&=
 \frac{n(1-e)q^2(1+qz)[(1+S)e-S(1-h)(1+z)]}{zQP}.
\end{aligned}                                                   \tag{TU9}
\]
These identities are proved simply by substituting the six printed
RER1 weights and the six NC4 weights, clearing the indicated pure
source denominators, and multiplying.  They include $n=0$; the
right sides define no quotient by $n$.

After substituting $e=zw$, the polynomial
\[
       P_\Delta=\frac{P_b-(1+qz)P_m}{z}                 \tag{TU10}
\]
is a genuine finite polynomial in $S,z,q,w,h$.  Indeed the numerator
specialized at $z=0$ is zero identically, so ordinary monic division
by $z$ has zero remainder.  Consequently
\[
 \delta(h_b-h_2)=n(1-e)P_\Delta/P.
                                                               \tag{TU11}
\]
Every polynomial in (TU6), (TU9), (TU10) has bounded coefficients
and bounded arguments.  From (TU4) and (TU9)--(TU11) we obtain
\[
\begin{gathered}
 (1-h)\delta c_b,\ h\delta a_b,\ \delta D_b
                     =q\lambda_n\,O(1),\qquad
 \delta h_2,\delta h_b=q\lambda_n\,O(1),\\
 \delta\alpha_b/q^2=\lambda_n q\,O(S+e),\qquad
 \delta(h_b-h_2)/(qz)=\lambda_n O(1).
\end{gathered}                                                   \tag{TU12}
\]
Here $O(1)$ denotes an explicit bounded rational source array with
denominators among the units $P,J,Q,R$; it does not hide a state
or testing factor.  These conclusions require no lower bound on
$h$, $1-h$, or $M$.

In particular $c_b\asymp1/(1-h)$, $a_b\asymp1/h$,
$D_b\sim1/2$, $\alpha_b/q^2\sim1/2$, and
$h_2/q,h_b/q$ positive bounded units with common residue
$1/2+\overline\lambda_n$.  We also have
\[
 h_b/w=z/2,\qquad h_b-h_2=O(qz),\qquad
 H_0:=4h_2/R+4h_b/w\geq2z,
 \quad H_0\asymp q+z.                                  \tag{TU13}
\]
The last statements use the formal exact SC2 difference together
with (TU11), and hold also when $M$ is arbitrarily small.

\subsubsection*{The complete additional current, including its normal}

Let the testing $u,v$ be positive infinitesimals and first take
$Z=1$ on its pure-testing state chart.  Use the exact SC coordinates
\[
 Q_b=(W+uY)/b,\quad Y_b=Y/S,\quad p=w/R,\quad
 N=((X+u)/b-pvQ_b)/g,
\]
\[
 H=uQ_b/b,\quad L=vQ_b/g,\quad
 \sigma=c_b+a_b,\quad \rho_0=Kh/R,\quad
 N_y=\sqrt\sigma\{D_uY_b-\rho_0H\}.
                                                               \tag{TU14}
\]
Both $N,N_y$ vanish at (TU7).  The exact SC3 full-parent identity
is an algebraic identity for arbitrary $c_b,a_b,D_b,\alpha_b,h_2,h_b$;
no C1 coefficient relation is used in deriving it.  With (TU8),
it represents all the actual C2 parent except its additional term
$4nD_uD_vX^2/S$.  Retain its original nonnegative $m$ packet
$4h_2wN^2/(MR)$, incorporate the second paired packet into
$H_0N^2$, and retain $4h_b(vN)^2/w$ exactly as in the uniform
QI normal lemma.  Define $\widehat N=\sqrt{H_0}N$.

Put $h_t=H/R$, and introduce the same four positive separated
monomials and their geometric mean as in QI:
\[
 A_*=ch_t,\quad B_*=cL/R,\quad C_*=uv/b,\quad
 D_*=qg,\qquad \Gamma=(A_*B_*C_*D_*)^{1/4}.
                                                               \tag{TU15}
\]
Their exact relations include
$\Gamma^2=qh_tv$ and $B_*C_*=(c/g)v^2h_t$.
The extra monomial
\[
                   U_*=cgu/b
\]
has the two exact identities
\[
       U_*B_*=\Gamma^2,\qquad
       cgX/b=czN+eB_*-U_*.
                                                               \tag{TU16}
\]
The entire additional current therefore has the following square
completion, with no state normal set to zero:
\begin{align*}
 \frac{4n}{S}D_uD_vX^2
 ={}&4\lambda_nD_uD_v(U_*-B_*-czN)^2\\
 &+4\lambda_nD_uD_v
   \{2(1-e)\Gamma^2-(1-e^2)B_*^2
                         -2(1-e)czB_*N\}.
                                                               \tag{TU17}
\end{align*}
The first line is a nonnegative separated SOS and vanishes at the
actual contact, including its unrestricted state-normal coordinate.
Its coefficient is literally zero if $n=0$.  The remaining scalar
changes have bounded coefficients on $B_*^2$ and $\Gamma^2$.
Their possibly nonzero real residues will cancel exactly with
the source comparison below.  The full additional first-normal coefficient, on
$B_*\widehat N$, is infinitesimal because
\[
       \lambda_n cz/\sqrt{H_0}=O(\lambda_n c\sqrt z)
                                            \in\mathrm I.
                                                               \tag{TU18}
\]
In particular (TU17) does not estimate or discard $u/b$, which can
be arbitrarily large.  The new positive square contains that entire
row, while the exact product relation in (TU16) places its cross
term on the old angular circuit.

\subsubsection*{Why the full QI allocation is stable here}

We verify the actual finite hypotheses of QI1--QI10 rather than
an abstract continuity assertion.  Its coefficient arrays have
only the pure-testing denominators $D_u$ and positive pure-source
units after cancellation of the two $h$ openings.  The latter
cancellation is preserved exactly because
\[
 1-\rho_0=(1-h)/R,\qquad \rho_0=Kh/R.
                                                               \tag{TU19}
\]
For example all perturbations of
$c_b(1-\rho_0)^2$, $a_b\rho_0^2$,
$c_b(1-\rho_0)$ and $a_b\rho_0$ are $q\lambda_n$
times bounded rational arrays, by (TU9).  Moreover
\[
 \delta(\sigma\rho_0-c_b)
       =\rho_0\delta a_b-(1-\rho_0)\delta c_b
                             =q\lambda_n O(1).
                                                               \tag{TU20}
\]
This is the complete potentially singular harmonic first-normal
coefficient after multiplication by $\sqrt\sigma$.
Both $\sigma$ and its formal comparison are comparable to
$[h(1-h)]^{-1}$, and their ratio has residue one.
Their square-root ratios and inverses are consequently bounded
pure-source units.  Thus this change of hard normal loses no
opening bound.

In the exact generic scalar QI2, the changes of $a,b_*,c_*$ are
$q\lambda_n$ times bounded arrays in $u,v$.  This follows term by
term from (TU12), (TU19): its only $C,A$ occurrences are the first
two expressions in (TU19); its $D_b$ occurrences have bounded
factors; its $\alpha_b$ occurrences have bounded factors and the
additional small factors printed in QI2.  Consequently the
formal specialization $a_0,b_0,c_0$ and its exact $U,V$ divisions
in QI3--QI4 can be used unchanged, while the remaining $q$
coquotients change by bounded arrays.  Write
$\lambda_0=\overline\lambda_n$.  Their complete leading changes
are respectively $0,-8\lambda_0,4\lambda_0,0$: substituting
$S=z=q=u=v=0$ in the bounded coquotients (TU9) gives
$\delta a/q=-4\lambda_0$, $\delta b_*/q=0$,
$\delta c_*/q=-4\lambda_0$, and
$\delta(R^2P_L)/q=4\lambda_0$.  Thus before the extra current
the four residues are
\[
 \overline{a_1-c_1}=2,\qquad
 \overline{b_1+2c_1}=-8-8\lambda_0,\qquad
 \overline{R^2P_L/q}=2+4\lambda_0,\qquad
 \overline{4\alpha_bD_u/q^2}=2.
                                                               \tag{TU21}
\]
Equation (TU17) adds $-4\lambda_0$ to the $B_*^2$ residue
and $+8\lambda_0$ to the $\Gamma^2$ residue.  Consequently
the \emph{combined} full scalar has exactly the original four
residues $(2,-8,2,2)$ on $(A_*^2,\Gamma^2,B_*^2,D_*^2)$.
All differences from those residues are infinitesimal bounded
arrays, because $\lambda_n$ itself is bounded.  This cancellation
is why a thinness assumption is unnecessary on (TU1).
The one additional bound QI uses from physical $\beta\geq0$ is
also available directly, without any sign for $\beta^0$:
\[
  w/(c/g)=\sqrt{w/\tau}=O(1),\qquad w=O(1).
                                                               \tag{TU22}
\]
Thus the exact dangerous cross term $-4(wv/(c/g))D_uB_*C_*$
still has infinitesimal coefficient.

For the complete first $N$ normal, (TU13) gives exactly the two
bounds used in the uniform QI normal lemma:
$g/\sqrt{H_0}$ and $c/\sqrt{H_0}$ are bounded, and
\[
 \frac{h_bD_v-h_2}{\sqrt{qH_0}}=O(z+v^2).
                                                               \tag{TU23}
\]
The $H$ term and all other terms in SC14 therefore have the same
infinitesimal coefficients on the QI rows; (TU18) handles the whole
extra term.  For the harmonic normal, (TU20) adds a
$q\lambda_n$ bounded array to the left side of QI10.  It can be
absorbed into its existing $qK_q$ coquotient.  All other terms are
the same bounded rational arrays multiplied by the opening-safe
quantities in (TU12), (TU19).  The exact $D_b$ root and its two
normal terms in QI9 remain unchanged as algebraic identities.
The complete hard block therefore retains the real residue
$\operatorname{diag}(1,4)$ in $(\widehat N,N_y)$.

For clarity, here is the finite allocation being invoked.  Use
the QI roots
\[
 E_*=h_t-v,\quad F_*=vh_t-gD_*,\quad
 R_F=\frac{\sqrt{g\Gamma}}c(\Gamma-D_*),\quad
 R_E=\sqrt\Gamma(\sqrt{A_*}-\sqrt{cv}),
\]
together with the complete finite toric root list SC16.
The exact identities QI7--QI8 give strict real square pivots on
$F_*,R_F,R_E$; QI6 gives the strict $E_*$ pivot and twice the
four-monomial circuit.  Every remaining coefficient is an
infinitesimal bounded rational array in the testing $u,v$, with
the same finite denominators and divisions just verified.
All removed packets, including (TU17), vanish at (TU7).
The actual source contact value, its three scalar first jets,
and both full first-normal contact values consequently vanish
exactly for this remaining expression.  Applying the two finite
monic divisions SC17--SC20 and the displayed QI7 first-jet
identities therefore puts the entire scalar correction and
both complete first normals on this strict root list with
infinitesimal coefficients.  The coefficient functions are
independent of the state ratio $Q_b/Z$; no boundedness of that
ratio is assumed.  This is precisely the constructive QI
allocation, with its source inequalities checked in
(TU12)--(TU23), not a conclusion from a scalar restriction alone.

The resulting finite Gram is bounded and has positive definite
real residue.  The strict bounded-Gram lemma supplies tensor SOS.
Restore the original positive paired packets and the new full
square in (TU17).  This proves the local assertion for the entire
parent on the positive small-angle, negative-core chart.

\subsubsection*{States, signs, and all remaining tests}

Every inverse and radical used above is pure source, pure testing,
or the explicitly separated positive monomial radical in (TU15).
For $Z\ne0$ the normalization $Z=1$ is pure testing; multiplying
back by $Z^2$ restores homogeneity.  At $Z=0$ the original
pure-testing core restriction is nonnegative.  The exact core
Schur identity and sign gauges used in the retained QI proof
force $W+uY>0$ on the remaining negative-core chart, and hence
all four monomials are positive.  Otherwise the core and all
original event packets already give the direct exit.  Literal
zero-angle charts have that same original exit.  Complex states
split into their real and imaginary parts, since the real
measurement matrices and full quadratic parent are unchanged.

The new R2 outer theorem R2O1--R2O24 treats every projective
testing pair outside the double-infinitesimal original chart,
with no condition on $a$ or $n/e^2$.  Its hypotheses are exactly
the original RER source conditions already retained here.
Combining that result with this local proof gives the claimed
all-tests, all-states theorem on (TU1).  No other source branch
is admitted by this conclusion.

\paragraph{Verification scope.}
The associated checker verifies the exact native contact map,
all six comparison formulas (TU9), the finite division (TU10),
the four complete residue cancellations in (TU21),
the whole generic SC reduction, the entire additional-current
completion (TU17), and its actual contact zero.  Ordered-field
bounds, the complete finite QI allocation and the outer-chart
join are analytical parts of this proof, not consequences of
the number of checker assertions.


\section{\texorpdfstring{R2 with both the native source ratio and relative current unbounded}{R2 with both the native source ratio and relative current unbounded}}
\label{component:pef021eabacde}
% Scoped mathematical certificate; see local_identity_audit.closure.json for exact changes.

\label{pef021eabacde:sec:r2-unbounded-current-new}

\paragraph{Precise assertion.}
Retain the original RER source, all its strict gaps, and the notation
TU2--TU16 of the preceding bounded-relative-current proof.
Assume the native $a,t$ are positive unbounded and
\[
             \lambda=\frac{n}{tB}=\frac{n}{qz}
                   \text{ is positive unbounded}.                 \tag{UL1}
\]
Then the whole original weak-C2 parent is tensor SOS for every
independent projective test and every state. The argument below
complements TU1; it never assigns a sign to the formal C1 coefficient
$\beta^0$, and never replaces the original physical source by that
formal source.

\subsubsection*{Source bounds without a thinness assumption}

The exact RER physicality equation already used in TU4 is
\[
 Q_{\rm RER}=\tau\{A+2r(z\Delta)\lambda\},\qquad
 \tau=\frac q{zw},\qquad
 A=(1+z)(1+Sh)(1+Sq).
\]
Here $A,Q_{\rm RER},r,z\Delta$ are positive units with residue one.
Since $\lambda$ is unbounded, this proves, with no cancellation of
mixed factors,
\[
 \tau\in\mathrm I,\qquad \lambda\tau\sim\frac12,\qquad
 \lambda q\sim\frac{zw}{2},\qquad
 \frac n{z^2w}\sim\frac12.                                     \tag{UL2}
\]
All four divisions are pure source. Recall $0<w<R\sim1$.
The exact TU9 formulas therefore imply
\[
\begin{gathered}
 (1-h)\delta c_b,\ h\delta a_b,\ \delta D_b,\
            \delta h_2,\delta h_b=\lambda q\,O(1),\\
 \delta\alpha_b/q^2=\lambda q\,O(S+zw),\qquad
 h_2\asymp zw,\quad h_b=zw/2,\\
 H_0=4h_2/R+2z\asymp z,\qquad H_0\geq2z,\\
 \frac g{\sqrt{H_0}}=O(1),\quad
 \frac c{\sqrt{H_0}}=O(1),\quad
 \frac{\lambda q}{H_0}=O(1),\quad \frac n{z^2}=O(1).
\end{gathered}                                                   \tag{UL3}
\]
In particular $\lambda q$ is infinitesimal. The displayed coefficient
arrays have denominators only among the pure-source units
$P,J,Q,R$ after the original $h,1-h$ openings cancel.
The actual and formal $\sigma=c_b+a_b$ differ by a relative
infinitesimal. Their square-root ratio is a bounded unit.
No lower bound on $w,h,1-h$ or the strict gap $M$ is assumed.

Use the complete generic SC3 identity with the actual coefficients.
Retain its positive $m$ packet and use $\widehat N=\sqrt{H_0}N$.
The formal comparison is used only for finite scalar and normal
identities. Its $q$-coquotients and the QI coefficient arrays are
bounded here: their derivation requires $S,z,q$ infinitesimal and
$w$ bounded, but not $\beta^0\geq0$. The sole QI bound that is no
longer available, $w/(c/g)=O(1)$, is replaced explicitly below.

\subsubsection*{The entire extra current supplies a new angular pivot}

Use the exact positive-chart rows TU15--TU16, denoted here by
$A_*,B_*,C_*,D_*,U_*$ and $\Gamma$. Define
\[
 J_3=\sqrt\lambda(U_*-B_*),\qquad
 E_*=h_t-v,\qquad F_*=vh_t-gD_*.
                                                               \tag{UL4}
\]
The first line of the exact whole-current identity TU17 is
\[
 4D_uD_v\left(J_3-d_N\widehat N\right)^2,\qquad
 d_N=\frac{\sqrt\lambda\,cz}{\sqrt{H_0}}
       =\sqrt{\frac{nz}{H_0}}\in\mathrm I.                       \tag{UL5}
\]
Thus it supplies a strict real pivot on $J_3$, with only
infinitesimal changes to the hard normal Gram. Its scalar
remainder is still exactly
$8\lambda D_uD_v(1-e)\Gamma^2
-4\lambda D_uD_v(1-e^2)B_*^2$.
Its complete first-normal remainder, retained below, is
$-8\lambda D_uD_v(1-e)czB_*N$.

Write $\widetilde c,\widetilde a,\widetilde d,\widetilde\alpha,
\widetilde m,\widetilde b$ for the six differences
$\delta c_b,\delta a_b,\delta D_b,\delta\alpha_b,
\delta h_2,\delta h_b$ divided by $\lambda q$.
They are defined explicitly by the six right sides in TU9, using
$n=\lambda qz$; their apparent opening poles cancel in every
combination below. Put $U=u^2,V=v^2$, and retain
$\rho_0=Kh/R$, $a_t=1-\rho_0+(w/R)V$.
Define the exact rational arrays
\[
\begin{split}
\mathsf a={}&4D_vR^2\{\widetilde c(1-\rho_0)^2
                      +\widetilde a\rho_0^2/D_u\}
 +4\widetilde\alpha zS^2VR^2\rho_0^2/D_u
 +4\widetilde d S^2R^2a_t^2,\\
\mathsf b={}&-8\widetilde\alpha zSR\rho_0
                         -8\widetilde d SRa_tD_u,\\
\mathsf c={}&4\widetilde dD_u^2,\\
\mathsf l={}&4R\{\widetilde m(1-w/R)
                              +\widetilde b(w/R)D_v\}.
\end{split}                                                     \tag{UL6}
\]
For example $\delta a/q=\lambda\mathsf a$ in QI2,
and $\delta(R^2P_L/q)=\lambda\mathsf l$.
Direct specialization of TU9 gives, keeping $S,z,w,h$ exact,
\[
\begin{split}
 (\mathsf a-\mathsf c)|_{q=U=V=0}&=0,\\
 (\mathsf b+2\mathsf c+8D_uD_v(1-e))|_{q=U=V=0}&=0,\\
 (\mathsf l-4D_uD_v(1-e^2))|_{q=U=V=0}&=0.
\end{split}                                                     \tag{UL7}
\]
These are stronger than residue equalities. Clear the unit
denominators and divide successively by the monic variables
$q,U,V$. Each of the three expressions is consequently
$qF_q+UF_U+VF_V$, with finite bounded rational coefficient arrays.
The division does not invert $q,U,V$ in the tensor ring.

The combined full scalar correction to QI6 therefore consists of
bounded-array multiples of $\lambda q,\lambda U,\lambda V$
on $A_*^2,\Gamma^2,B_*^2$, an infinitesimal coefficient on
$E_*^2$, and the infinitesimal $\alpha$ correction on $D_*^2$.
The following exact identities allocate every apparently unbounded
testing term:
\[
\begin{aligned}
\lambda U A_*^2&=\lambda q(uE_*+bC_*)^2,\\
\lambda V A_*^2&=\lambda q(F_*+gD_*)^2,\\
\lambda U\Gamma^2&=\lambda q\,bC_*(uE_*+bC_*),\\
\lambda V\Gamma^2&=\lambda q(F_*+gD_*)
                                    (F_*+gD_*-vE_*),\\
\sqrt\lambda\,uB_*&=\sqrt{nS/z^2}\,(F_*+gD_*),\\
\sqrt\lambda\,vB_*&=\sqrt n\,C_*-vJ_3 .
\end{aligned}                                                   \tag{UL8}
\]
Every coefficient introduced on squaring the last two identities
is infinitesimal by UL2--UL3 and $S,u,v\in\mathrm I$.
Finally the formerly dangerous QI6 term has the exact replacement
\[
 \frac{wg}{c}\,vB_*C_*
       =wz\,C_*^2-\frac{wg}{c\sqrt\lambda}\,vC_*J_3,
 \qquad
 \frac{wg}{c\sqrt\lambda}
          =w\sqrt{\frac{z^2}{n}}=O(1).                          \tag{UL9}
\]
Both resulting coefficients are infinitesimal. This retains the
entire term for arbitrary independent source and testing rates.

\subsubsection*{Both unrestricted state-normal corrections}

The new complete first $N$ coefficient, including the remainder
of TU17, differs from the formal generic SC coefficient by
\[
\begin{split}
\delta\mathcal L_N={}&
 8\lambda c(1-e)\mathsf F_N B_*
 +8\lambda q\widetilde d\,ubvg(Sa_tH-D_uv),\\
\mathsf F_N={}&
 \frac{P_bD_v-(1+qz)P_m}{P}-zD_uD_v
                 =V-zUD_v+qF_{Nq}.
\end{split}                                                     \tag{UL10}
\]
The last equality is a finite monic $q$ division; its coquotient
has bounded arrays because the numerator at $q=0$ is exactly
$V-zUD_v$. After division by $\sqrt{H_0}$, its $q$ part is
infinitesimal on $B_*$ by UL3. The other parts use
\[
\begin{split}
\lambda cV B_*&=\lambda qg\,vC_*
                            -\sqrt{\lambda q}\,v^2J_3,\\
\lambda czU B_*&=\lambda qbg\,v(uE_*+bC_*).
\end{split}                                                     \tag{UL11}
\]
Their coefficients on $\widehat N$ are infinitesimal since
$g/\sqrt{H_0}$ and $\sqrt{\lambda q/H_0}$ are bounded.
The $\widetilde d$ term in UL10 uses
$uH=R(uE_*+bC_*)$ and $uv=bC_*$; it has the same
infinitesimal coefficient bounds. Thus this is a lift of the
complete first normal, with no state component put to zero.

For the second normal it is convenient to compare the coefficient
after multiplication by $\sqrt\sigma$. Differentiating the full
generic SC3 identity, including the restored $D$ packet, gives
\[
 \sqrt\sigma\,\delta\mathcal L_y
       =8\lambda q\{\mathsf F_HH+\mathsf F_vv\},                 \tag{UL12}
\]
where
\[
\begin{split}
\mathsf F_H={}&D_v\{\rho_0\widetilde a
                    -(1-\rho_0)\widetilde c\}
 -D_v\widetilde aU\rho_0/D_u
 +\widetilde\alpha zS^2V\rho_0/D_u
 -\widetilde dS^2a_t,\\
\mathsf F_v={}&\widetilde dSD_u-\widetilde\alpha zS.
\end{split}
\]
These arrays are bounded through both source openings. Exactly,
\[
 \mathsf F_H|_{q=U=V=0}
       =\frac{(1-e)S}{1-S(1-h)},\qquad
 \mathsf F_v|_{q=U=V=0}
       =-\frac{(1-e)S}{1-S(1-h)}.                              \tag{UL13}
\]
As $R=1$ there, finite monic division gives
$R\mathsf F_H+\mathsf F_v=qF_q+UF_U+VF_V$ with bounded
arrays. Since $H=R(E_*+v)$, the whole UL12 coefficient is
therefore an infinitesimal combination of $E_*,F_*,A_*,C_*,D_*$:
use only
\[
 qv=cA_*-qE_*,\qquad Uv=buC_*,\qquad
 Vv=v(F_*+gD_*-vE_*).                                         \tag{UL14}
\]
The factor $1/\sqrt\sigma$ is bounded.
Changing the formal harmonic-normal normalization to the actual
one multiplies its previously bounded arrays by a pure-source
unit with residue one.

For completeness, the unchanged formal first $N$ coefficient is
bounded after the actual $\sqrt{H_0}$ normalization:
$h_2^0,h_b^0=O(q)$ and $h_b^0-h_2^0=O(qz)$, while
$g/\sqrt{H_0}$ and $c/\sqrt{H_0}$ are bounded. These are exactly
the two estimates used in the QI normal allocation.
All remaining hard coefficients from the original $C,A,D,\alpha$
lines have the same residues as before; the relative $C,A$
changes and the absolute $D,\alpha$ changes are infinitesimal.
Together with UL5 the complete hard block on
$\widehat N,N_y,J_3$ has a positive definite real residue.

\subsubsection*{A finite contact lift for the additional row}

It remains to allocate the full scalar correction on roots that
vanish at contact. Merely bounding a coefficient on $A_*,B_*,C_*,D_*$
would not suffice. Keep the complete strict toric root list SC16,
the QI rows $E_*,F_*,R_F,R_E$, and add
\[
 V_3=\sqrt{\lambda\Gamma}(\sqrt{U_*}-\sqrt{B_*}).
\]
Since $U_*B_*=\Gamma^2$, the exact identities are
\[
\begin{split}
 J_3^2&=4V_3^2+\lambda(\sqrt{U_*}-\sqrt{B_*})^4,\\
 \Gamma J_3&=\frac{V_3^2}{\sqrt\lambda}
                      -2\sqrt\lambda\,\Gamma(B_*-\Gamma).
\end{split}                                                     \tag{UL15}
\]
Reserve half the $J_3^2$ coefficient in UL5 for the first
identity. It supplies strict real pivots on both $J_3,V_3$,
with the quartic square retained nonnegative. The normal
coupling in UL5 remains an infinitesimal hard Gram correction.

Here are the additional first-jet coefficient checks needed for
the second identity, so that its factor $\sqrt\lambda$ causes
no hidden loss. The only new scalar terms linear in $J_3$ in
UL8--UL9 have coefficients, up to bounded arrays, of the forms
\[
 \epsilon_B=-2\sqrt n\,v,\qquad
 \epsilon_T=\frac{wg}{c\sqrt\lambda}\,v .
                                                               \tag{UL16}
\]
Both have bounded finite testing difference arrays and are
infinitesimal. At the actual pure-source contact $v_c=cg$,
\[
 \sqrt\lambda\,\epsilon_B(u_c,v_c)=-2n\in\mathrm I,\qquad
 \sqrt\lambda\,\epsilon_T(u_c,v_c)=wz=e\in\mathrm I.             \tag{UL17}
\]
Multiplication by any of the bounded arrays produced in UL7
preserves these conclusions. No other scalar term introduces
$J_3$ linearly: the other five identities in UL8, and the
$\lambda q$ scalar corrections, use only the old QI rows.

Write a term $\epsilon J_3A_i$ as
$\epsilon J_3(A_i-\Gamma)+\epsilon\Gamma J_3$.
The first summand is already an infinitesimal quadratic on the
strict list. Divide $\epsilon-\epsilon(u_c,v_c)$ finitely by
$u-u_c,v-v_c$ and apply the exact SC18 conductors to
$\Gamma(u-u_c)$ and $\Gamma(v-v_c)$. This allocates the
nonconstant part as infinitesimal quadratic corrections involving
$J_3$. For the constant part use UL15. Its $V_3^2$
coefficient is infinitesimal, and its remaining coefficient on
$\Gamma(B_*-\Gamma)$ is infinitesimal by UL17.

There are also new bare-$\Gamma$ $E_*,F_*$ terms when the
first five identities in UL8 are expanded and their $C_*,D_*$
factors are written as $(A_i-\Gamma)+\Gamma$. Put
$L=\lambda q$ and $M_1=nS/z^2$, both infinitesimal.
The additional $E_*$ coefficients are bounded-array multiples
of $Lbu$ and $Lgv$; the additional $F_*$ coefficients are
bounded-array multiples of $Lg$ and $M_1g$.
At $u_c=bc,v_c=cg$ their required contact ratios are
\[
 \frac{Lbu_c}{c}=LS,\qquad
 \frac{Lgv_c}{c}=Lz,\qquad
 \frac{Lg}{g}=L,\qquad
 \frac{M_1g}{g}=M_1 .
                                                               \tag{UL18}
\]
All are infinitesimal, with bounded finite difference arrays
before these contact substitutions. The exact QI7 identities
therefore allocate these additional terms by the same finite
division as the old $E_*,F_*$ terms. This check is necessary:
the two new normal rows cannot be ignored when applying the
contact-jet argument.

After these explicit additional steps, the remainder has exactly
the QI form: a bounded infinitesimal quadratic on the listed
roots, the now-verified bare-$\Gamma$ $E_*,F_*$ couplings, and
a toric expression
$\Gamma\sum_i\mu_i(A_i-\Gamma)+\Gamma^2\sigma$.
The old $E_*,F_*$ contact coefficient ratios remain the bounded
ones in QI7--QI8: their unchanged $q=0$ arrays depend on bounded
$w$ and do not require $w/(c/g)$ bounded.
Every retained packet, including UL5 and UL15, vanishes at the
actual contact. The complete actual scalar consequently has zero
contact value and all three zero first jets. Apply the two
explicit finite divisions SC17--SC20. Their three-jet balance
makes the constant toric first-jet coefficients equal; SC19 and
the variance roots then allocate this common term. This puts
the entire scalar on a bounded Gram with strictly positive
real residue.

The full normal lifts UL10--UL14 have exact contact zero as well.
Their $J_3,E_*,F_*$ terms already vanish there. The sum of their
remaining four monomial coefficients is zero at contact, so one
finite SC18 division allocates them on the same strict list.
Every coefficient function in these steps depends only on the
testing $u,v$. No bound on the state ratio $Q_b/Z$ is used, and
the third contact jet is used only for its exact toric balance.

The strict bounded-Gram lemma now proves tensor SOS for the whole
original parent on the small-angle positive negative-core chart.
All its radicals are radicals of separated positive monomials;
every inverse is pure source or pure testing.

\subsubsection*{Projective completion and scope}

The state, sign, zero-axis and real/imaginary-state arguments are
the same exact original-core arguments at the end of TU.
In particular $Z=0$ is an original positive restriction, and the
negative-core $Z=1$ chart forces $W+uY>0$, permitting the
separated radicals above. The theorem R2O covers every testing
angle chart outside $u,v$ both infinitesimal without any
condition on $\lambda$ or the native $a$. This proves UL1 for
all tests and states.

Combining this theorem with TU covers the entire original R2
sector with native $a$ unbounded: its nonnegative
$\lambda=n/(tB)$ is either bounded, including zero, or positive
unbounded. The complementary bounded-$a$ theorem remains a
separate input to the final R2 join. This local result alone is
not the universal tensor-SOS theorem or finite-lift existence.

\paragraph{Verification boundary.}
The focused exact checker verifies UL7, UL10, UL13, the full
generic first-normal derivatives before any graph restriction,
all six UL8 conductors, UL9, UL11, UL14--UL17 and the actual
contact of $J_3$. The native source comparisons are the separately
verified TU9 dependency. The source inequalities, bounded finite
arrays, finite contact allocation and projective join are the
analytical proof above, not consequences of assertion counts.


\section{\texorpdfstring{The entire original R2 source, with all independent tests and states}{The entire original R2 source, with all independent tests and states}}
\label{component:p913a09bcdcda}
% Scoped mathematical certificate; see local_identity_audit.closure.json for exact changes.

\label{p913a09bcdcda:sec:r2-complete-rederived}

\paragraph{Theorem and unchanged source.}
Let \(K_c,K_t\) be independent real closed extensions of
\(\mathbb R\).  Use the original common-II stationary source
with \(d=1\), \(\lambda=\beta=0\), the six actual rational
coefficients \(n,m,C,A,D,\alpha\) in RER1, and all original
source gaps and coefficient nonnegativity conditions.
In particular \(n\geq0\); the strict \(m,D\) denominator
gaps are
\[
 g_s=1+a(hB+r-1)>0,\qquad
 j_D=1+aBr-(1-h)/t>0 .
\]
Assume the original R2 size conditions
\[
 a,B,t>0,\quad 0<r,h<1,\qquad
 t\text{ unbounded},\qquad B,aB,tB\in\mathrm I_c .
                                                               \tag{R2A1}
\]
Then the complete original parent is a finite sum of squares
in \(K_c\otimes_{\mathbb R}K_t\) at every independent projective
testing pair and every real or complex state.
All allowed source openings and literal \(n=0\) are included.

For specificity, in finite testing half-angle coordinates the
quarter-parent being certified is exactly
\[
\begin{split}
 \mathcal P={}&m(vW-X+uvY-uZ)^2+nD_uD_vX^2
       +CD_v(uW-Y)^2+AD_uD_vY^2\\
 &+(D+\tfrac12)(uW+uvX-Y-vZ)^2
       +(\alpha+\tfrac12)D_u(vY-Z)^2\\
 &+\tfrac12D_v(X+uZ)^2-\tfrac12D_uD_vZ^2,
 \qquad D_u=1+u^2,\quad D_v=1+v^2 .
\end{split}                                                     \tag{R2A2}
\]
The statement is about this entire quadratic in the four state
coordinates.  It does not specialize either state normal to zero
and does not use a scalar-only positivity assertion.

\subsubsection*{Exhaustion of the original source}

The exact source-entry lemma RER1--RER8 derives its bounds from
the original inequality
\[
 (1-r)Q\geq BK,
\]
including equality when \(n=0\).  In particular it proves
\(a(tB)\in\mathrm I_c\) before any comparison source is used.
No earlier-source exclusion is added to (R2A1).

Every positive source \(a\) lies in precisely one of the following
two size classes.

\emph{Bounded \(a\).}
Apply the whole-parent theorem RL1--RL24, with the exact circuit
and complete-normal identities SH1--SH8.
Its first source alternative is exactly bounded native \(a\),
with no restriction on \(e=1-r\), \(n/e^2\), or their ratios.
Its correlated \(n\)-completion is literal at \(n=0\).
The proof retains both complete normals and every positive
rotated state ratio.

\emph{Unbounded \(a\).}
Set
\[
 S=1/t,\quad z=1/a,\quad q=a(tB),\quad
 w=a(1-r),\quad \Lambda=\frac{n}{tB}=\frac{n}{qz}.
                                                               \tag{R2A3}
\]
These are the exact original-source coordinates TU2.
The additional nonnegative source quantity \(\Lambda\) is either
bounded, including zero, or positive unbounded.
If it is bounded, apply TU1--TU23.
If it is unbounded, apply UL1--UL18.
The formal C1 coefficients in those two proofs are only a finite
algebraic comparison.  Neither proof assumes their physicality
or imports a C1 theorem at a negative event coefficient.
The entire actual \(n\)-packet and both complete normal
corrections are retained.

These alternatives exhaust all sources in (R2A1).
There is no intermediate size class between bounded and
positive unbounded in a real closed extension of \(\mathbb R\).
In particular \(n=0\) belongs to the bounded-\(\Lambda\)
alternative whenever \(a\) is unbounded; it is not removed by
a division by \(n\).

\subsubsection*{Exhaustion of tests and states}

The fresh exact sign and projective lemma SJ1--SJ10 permits
nonnegative finite testing half-angles by real state gauges
that preserve every event packet.  Literal zero angles and
homogeneous infinite endpoints have direct core SOS identities.
On the remaining positive finite angles, R2O1--R2O24 proves
the entire original parent whenever at least one angle is
not infinitesimal.  Its physical reflection, all four
finite-cap opening cases, and the original positive excesses
are part of that proof; no reflected R1 source hypothesis
is imported.

It remains only to consider \(u,v>0\) both infinitesimal.
The source alternatives above supply a complete local certificate
for every such pair and every state.  RL uses the nonzero
pure-testing coordinate \(w=(W+uY)/D_u\) to normalize
\(w=1\), and permits the resulting \(z>0\) to be arbitrary.
TU and UL use the pure-testing \(Z=1\) chart, with no bound on
the other state ratio.  Their respective original-core sign
identities justify these normalizations and give direct SOS
on the complementary state charts.  Thus no state chart is
lost when the two exact frames are used in different source
sectors.

Every normalization is reversed by multiplying the roots by
a positive pure-source or pure-testing factor.  The zero
state is immediate.  Since the measurement matrices are real,
the quadratic parent on a complex state is the sum of its
values on the real and imaginary parts.  This proves the
stated real and complex all-state conclusion.

\subsubsection*{Exact proof dependencies and scope}

The substantive new mathematical leaves are:
\[
\begin{array}{c|l}
 \text{source or testing region}&\text{complete proof}\\ \hline
 a=O(1)&\mathrm{RL1\!-\!RL24,\ SH1\!-\!SH8}\\
 a\text{ unbounded},\ \Lambda=O(1)&\mathrm{TU1\!-\!TU23}\\
 a,\Lambda\text{ unbounded}&\mathrm{UL1\!-\!UL18}\\
 \text{outer testing angles and all endpoints}&
                   \mathrm{R2O1\!-\!R2O24,\ SJ1\!-\!SJ10}.
\end{array}                                                     \tag{R2A4}
\]
They use the actual source-entry identities RER1--RER8,
the reconstructed exact compact source and scalar identities
RSR, the complete parent RFN4--RFN16, and the printed finite
QI/SC circuit and contact-division lemmas where explicitly
invoked.  These are proof bodies and exact finite identities;
historical status files and assertion counts are not premises.

This theorem completes the whole original weak-C2 row R2.
It does not by itself assert the global universal tensor-SOS
theorem, produce an explicit finite LMI, bound an LMI size,
or assert external proof verification.  Those are separate
assembly and deliverable questions.


\section{\texorpdfstring{R2 with bounded \(a\), or an infinitesimal paired-event parameter}{R2 with bounded a, or an infinitesimal paired-event parameter}}
\label{component:p393d3faa269a}
% Scoped mathematical certificate; see local_identity_audit.closure.json for exact changes.

\label{p393d3faa269a:sec:r2-shifted-local}

\paragraph{Exact statement.}
Use the actual original weak-C2 source and all its original strict
gaps, with \(d=1\), \(\lambda=\beta=0\), and \(n\geq0\).
Suppose \(t_s\) is positive unbounded and
\(B_s,a_sB_s,t_sB_s\) are positive infinitesimals.
Put \(e=1-r\), \(R_s=1+ha_sB_s\).  On the source subset
\[
       a_s=O(1)\quad\hbox{or}\quad
       p=a_s e/R_s\in\mathrm I_c,                     \tag{RL1}
\]
the entire parent is separated tensor SOS for every independent
projective testing pair and every real or complex state.
Literal \(n=0\) and every opening of \(r,h,1-h,g\) are retained.
For the all-testing conclusion, the local proof below is joined
explicitly to R2O1--R2O24.

The two fields and the bounded-Gram convention are those of EN1--EN4.
In particular no mixed inverse is used.  Every coefficient bound
below refers to a finite array in \(B_c\otimes B_t\); the rows need
not be bounded.

\subsubsection*{Actual source entry and bounds}

Use the exact compact source dictionary RSR2--RSR8, with
\[
 j=a_sc_s/R_s,\quad B=B_sR_s/c_s,\quad T=t_sc_s/R_s,\quad
 p=a_se/R_s,\quad H=(1-h)/(t_sc_s),\quad c_s=1+a_sB_sr.
\]
Set
\[
 S=1/T,\quad \ell=HT,\quad \bar h=1-\ell,\quad
 x=jB,\quad y=TB,\quad B=Sy,\quad q_z=j/S .
                                                               \tag{RL2}
\]
Thus \(S,x,y\) are positive infinitesimals,
\(0<p<\min(1,j)\), \(0<\ell<1\), and \(j\asymp a_s\).
Both \(\ell\) and \(\bar h\) may be arbitrarily small.
The original physicality inequality RER2--RER7 gives \(jy\in\mathrm I_c\).
In terms of the actual finite polynomials \(\mathcal P,\mathcal L,
\mathcal N,\mathcal F,\mathcal K\) of RSR, put
\[
 K=j+1-p,\quad k=K/(j+1),\quad
 F_b=j+1+p(1-p)/j,\quad E=1+B(\ell+p-1).
                                                               \tag{RL3}
\]
Here \(1/2<k<1\) and \(1\le F_b/(j+1)<2\).  Direct substitution
\(B=Sy,T=1/S,H=S\ell\) in those finite polynomials gives
\[
 \mathcal P/(j+1)\sim1,\qquad
 S\mathcal L/K\sim1,\qquad S\mathcal F/F_b\sim1.         \tag{RL4}
\]
For clarity, the terms after division contain only bounded
\(p/j,j/(j+1),K/(j+1)\), multiplied by \(S,y\), or \(jy\).
No inverse of \(\ell\), \(\bar h\), or \(j-p\) is involved.

The same substitution in the exact numerator of \(n\) yields
\[
 \mathcal N=p(1+\epsilon_p)
                -yj(j+1)(1+\epsilon_y),\qquad
 \epsilon_p,\epsilon_y\in\mathrm I_c.                  \tag{RL5}
\]
For an explicit finite division certificate, the zero-\(p\) term is
\[
 \mathcal N|_{p=0}
 =-yj(j+1)\{1+S(1-\ell)\}
             \{1+S\ell yj/(j+1)\}.
\]
The quotient \((\mathcal N-\mathcal N|_{p=0})/p\) is a
finite polynomial.  This grouping can be read termwise from RSR5:
the first
coefficient is \(1+y(j+1)+O(Sy(j+1)+Sy^2(j+1))\);
the second coefficient is \(1+O(S+y+Syj)\).
Using \(\mathcal N\geq0\), therefore,
\[
 \eta_2=\frac{yj(j+1)}p=O(1),\qquad
 \eta_2\in\mathrm I_c\ \Longrightarrow\
 n\sim\frac{p}{2(j-p)(j+1)}>0 .                        \tag{RL6}
\]
The estimate includes the face \(\mathcal N=0\), which can occur
only in the positive-unit \(\eta_2\) case.

Let \(\mu=m(1-p)\) and \(h_0=\mu+1/2\).
The complete source formulas now imply
\[
\begin{gathered}
 \mu\sim\frac{jp}{2(j+1)},\quad h_0\asymp1,\quad
 CH\sim k/2,\quad SA\bar h\sim k/2,\quad D/S\sim k/2,\\
 \alpha\sim\frac{Sy^2K^2}{2(j+1)},\quad
 f_c\sim\frac{j^2y^2F_b}{2(j+1)},\quad
 r_c\sim\frac{jy^2(1-p)K}{2(j+1)} .
\end{gathered}                                                   \tag{RL7}
\]
These use the exact masses in RFN12 and the closed RSR formulas.
All displayed positive masses are strictly positive.
Write \(\kappa=-k_c>0\).  The exact positive-surplus formula
RSR19 proves this sign.  Its substituted finite polynomial gives
\[
 \frac{\kappa}{r_c}=O(p+Sy)\quad(p\in\mathrm I_c),\qquad
 \frac{\kappa}{jy^2}=O(1).                              \tag{RL8}
\]
For the remaining case in (RL1), \(p\) and \(j\) are positive
bounded units, and
\[
 \frac{\kappa}{f_c}\sim\frac{pK}{jF_b}\leq\frac12,
 \qquad jF_b-2pK=(j-p)K\geq0.                          \tag{RL9}
\]
The inequality is a bound on real residues when combined with
the preceding asymptotic equality.  On this case \(\eta_2\) is
infinitesimal, and the effective current defined below is a
positive bounded unit.

\subsubsection*{The entire parent and its harmonic shift}

Retain the complete original quarter-parent RFN4.  Its exact inverse
state change at \(w=(W+uY)/(1+u^2)=1\) is
\[
 X=N_0+pv-uz,\quad Z=z+uN_0+puv,\quad
 W=1+u(M_0+Hvz),\quad Y=u-M_0-Hvz.                    \tag{RL10}
\]
Here \(z>0\) is arbitrary.  Retain \(mpD_u^2N_0^2\) and divide
the remaining whole parent by \(D_u=1+u^2\).
The exact expression is RFN7, with both complete state normals,
hard entries RFN9, and scalar RFN12.

Use the additional exact change
\[
 q=vz-Tu,\qquad L=M_0+Hq,\qquad
 M_0=L-Hq,\qquad L=\sqrt{S\ell\bar h}\,\widehat L .
                                                               \tag{RL11}
\]
All divisions are pure source or pure testing.  In particular
this change does not invert the harmonic mixed unit
\(1+\rho u^2\).
The complete diagonal of \(\widehat L\) is
\[
 S\ell\bar h\{D_u(CD_v+D+\tfrac12)+AD_v+
                   (\alpha+\tfrac12)v^2\},
                                                               \tag{RL12}
\]
with positive real residue \(\overline k/2\) and bounded
coefficient array through both openings.

Use the circuit coordinates \(U,V,Z_*,Q_*,B_*,R_*,\Psi\)
of SH1--SH2, and \(b=v-uz/j=\sqrt y B_*\).
Before spending any scalar \(n\)-mass, complete the exact correlated
block, as in RFN13--RFN16:
\[
 h_0N_0^2+n(N_0+jb)^2
 =(n+h_0)(N_0+c_bb)^2+M_e b^2,\quad
 c_b=\frac{nj}{n+h_0},\quad M_e=\frac{nh_0j^2}{n+h_0}.  \tag{RL13}
\]
Keep \(nv^2(N_0+pb)^2\) and all three additional positive hard
squares RFN15 whole, including after the change
\(\widehat N=\sqrt{n+h_0}(N_0+c_bb)\).
At \(n=0\), put \(c_b=M_e=0\) and omit the current row.
For \(n>0\), \(c_b^2/M_e=n/[h_0(n+h_0)]\) is bounded.

\paragraph{The exact shifted scalar.}
Let \(F\) be RFN12.  The shifted scalar is
\(F-2HqL_M+H^2q^2H_M\).
Use the whole \(d_cV^2Q_*^2\) term in its \(q^2\) diagonal.
Using the exact actual-source identity \(S_0=0\) in RFN22,
the resulting coefficient on \((\sqrt S q)^2\) is exactly
\[
 D_q=S\{A\bar h^2+\ell^2(C+D+\tfrac12)D_u\}
       -\frac{f_c}{Tx}-TH^2C\,u^2v^2
       -\frac{k_cv^2}{yTx}.                            \tag{RL14}
\]
It has bounded array and positive real residue \(\overline k/2\).
Indeed \(f_c/(Tx)=O(jy)\) is infinitesimal, \(TH^2C=\ell CH\)
is bounded, and \(\kappa/(yTx)=\kappa/(jy^2)\) is bounded.
Its equivalent \(Q_*^2\) mass has order \(jy\).
The only extra scalar cross terms are
\[
 -2Hq\{\mu uvb+v(\alpha-yCHu^2)R\},\qquad
 R=z-uv/(Sy).                                        \tag{RL15}
\]
The budgets below allocate their entire expressions.

The remaining signed scalar terms are exactly
\(e_0V^2B_*^2+k_cZ_*^2(V^2-1)^2\), where
\(e_0=p^2y^2/2-g_c\).
The negative part of \(e_0\) is at most \(g_c\).
SH2 purchases it at relative circuit cost
\[
 \frac{g_c}{f_c}(1+V^2)
   =O\bigl(S(y+v^2)\bigr)\in\mathrm I.                \tag{RL16}
\]
Use \(Z_*(V^2-1)=-VQ_*-R_*\) for the other signed term.
If \(p\) is infinitesimal, (RL8) pays its opposite part
directly from \(r_cR_*^2\), and its \(VQ_*\) cost is
\(O(v^2)\).
Otherwise \(p,j,M_e\) are positive units.  Young's exact
square identity with parameter \(1/6\), followed by SH1 at
\(a=3/2\), gives the legal SOS domination
\[
 \kappa Z_*^2(V^2-1)^2
 \preceq_{\rm sos}
 \frac74\kappa\Psi+
 \frac{21}{2}\kappa V^2Q_*^2+
 \frac72\kappa V^2B_*^2 .                             \tag{RL17}
\]
The circuit cost is at most \(7/8+o(1)\), by (RL9).
The two other relative costs are \(O(v^2)\).
Thus fixed positive real fractions of \(D_q(\sqrt S q)^2\),
\(r_cR_*^2\), \(f_c\Psi\), and \(M_eb^2\) remain;
the last is absent only at \(n=0\).
The opposite mass \(r_c\) was not spent in the second case.

\subsubsection*{All helper rows and both complete normals}

As usual, a free row permits any fixed real square pivot at
infinitesimal cost from the preceding retained scalar forms.
SH2 proves that both \(E_u=u(uvz-jy)\) and \(E_v=v(uvz-jy)\)
are free.  The row \(R_u=v(u^2-x)\) is free as well.
For \(p\) infinitesimal, use the actual RFN19 rational Gram:
its three costs are \(O(S^2y)\), \(O(Syu^2)\),
and \(O(Su^2/(1-p))\).
When \(p,j,M_e\) are units, apply SH1 to the \(R_*^2\)
term in RFN19.  The resulting costs are
\[
 O(S^2y+Su^2/j),\quad O(Syu^2+Su^2v^2),\quad
 O(jSu^2v^2/M_e),                                    \tag{RL18}
\]
all infinitesimal.

The rows \(pv^2b\) and \(\sqrt S\,\mu uvb\) are free.
If \(\eta_2\) is infinitesimal, (RL6) and (RL13) imply
\[
 \frac{p^2}{M_e}
  =\frac{p^2}{h_0j^2}+\frac{p^2}{nj^2}=O(1).
                                                               \tag{RL19}
\]
They then cost \(O(v^4)\) and \(O(Su^2v^2)\).
If \(\eta_2\) is a positive unit, the actual source gives
\(p^2y/f_c=O(1)\) and \(\mu=O(p)\).
The second identity in SH2 buys the same rows with costs
\[
 O\bigl(v^2(y+v^2)\bigr),\qquad
 O\bigl(Su^2(y+v^2)\bigr),                             \tag{RL20}
\]
including \(n=0\).

The full opposite term \(\sqrt S\,v(\alpha-yCHu^2)R\)
is free by SH4--SH8.  Also \(j\alpha b\) is free:
\[
 j\alpha b=-(\alpha/B)R_u-\alpha uR,\qquad
 \alpha/B=O(y(j+1))\in\mathrm I_c.
                                                               \tag{RL21}
\]
For \(p\) infinitesimal, the direct \(r_c\) cost for
\(\alpha uR\) is \(O(Sy^2(j+1)^2u^2)\).
In the unit \(p,j,M_e\) case use SH1; the three costs
are \(O(Sy^2u^2)\), \(O(Sy^2u^2v^2)\), and
\(O(Sjy^2u^2v^2/M_e)\).  All are infinitesimal.
These statements account for the whole source \(\alpha\), not
only its leading coefficient.

\paragraph{The first normal.}
The exact RFN21 source cancellations give, after (RL11),
\[
 L_N'=[j(n-\alpha)+p(n+\tfrac12)v^2]b
             +K_sR_u+\kappa_NE_v,                     \tag{RL22}
\]
\[
 K_s=-\mu+(\alpha+\tfrac12)(\ell+p-1),\qquad
 \kappa_N=-(n+\tfrac12)(1-p/j).
\]
Both coefficients are bounded.  For \(\kappa_N\), use the
exact numerator of \(n\) and
\(n(j-p)/j=O(p/[j(j+1)])\); this cancels the \(r\) opening.
After the entire correlated completion (RL13), the remaining
normal is
\(-j\alpha b+\tfrac12pv^2b+K_sR_u+\kappa_NE_v\).
Every row is free by (RL18)--(RL21).
Division by \(\sqrt{n+h_0}\) is bounded.
The shift generates couplings to \(\sqrt{M_e}b\) with bounded
coefficients \(c_b/\sqrt{M_e}\); they are absent at \(n=0\).
Large fixed real helper pivots absorb these couplings by exact
Young identities.

\paragraph{The second normal and the entire hard block.}
The complete RFN22 normal after (RL11) is
\[
 L_L=(C_q-HH_M)q+\mu uvb+v(\alpha-yCHu^2)R,\qquad
 C_q-HH_M=-D-H(C+D+\tfrac12)u^2 .                      \tag{RL23}
\]
The coefficient of \(\sqrt S q\) after the normalization
of \(L\) is infinitesimal; \(D\) is infinitesimal and
\(H(C+D+1/2)\) is bounded.
The other two terms are free by (RL20) and SH6--SH8.
Exactly the same rows pay (RL15), whose first factor is
\(\sqrt S q\) with a bounded source multiplier.

The remaining hard off-diagonal is
\((\alpha+1/2)uv\sqrt{S\ell\bar h}/\sqrt{n+h_0}\),
which is infinitesimal.  Its \(c_b b\) correction has the
same bound after division by \(\sqrt{M_e}\).
The complete finite Gram therefore has a positive definite
real residue, with principal normal pivots
\[
                   1,\quad\overline k/2,\quad\overline k/2
                                                               \tag{RL24}
\]
on \(\widehat N,\widehat L,\sqrt S q\), together with the
retained strict scalar and purchased helper pivots.
All extra positive hard squares and both original retained
current packets are restored whole.  The strict bounded-Gram
lemma proves tensor SOS for the entire parent.

\subsubsection*{State and outer-chart exhaustion}

The fresh full-state sign lemma SJ1--SJ10 permits positive
\(u,v\), and on the remaining negative-core chart gives
\(wz>0\).  Divide by the nonzero pure-testing \(w\) to obtain
the normalization used above with arbitrary \(z>0\).
No estimate in this proof uses its size.  Literal axes and
homogeneous endpoints have the direct original-core SOS exit.
Real and imaginary state parts separate.
The exact outer theorem R2O1--R2O24 covers every other testing
pair on all original R2 sources.  This proves the stated
all-tests conclusion on (RL1).

The complementary source \(a_s\) unbounded and \(p\) a positive
unit is not inferred from this proof.  It requires the separate
TU source comparison and additional-current conductor.


\section{\texorpdfstring{Two exact circuit budgets and the complete R2 opposite normal}{Two exact circuit budgets and the complete R2 opposite normal}}
\label{component:p09b2107cc05e}
% Scoped mathematical certificate; see local_identity_audit.closure.json for exact changes.


Use the compact actual-source dictionary of RSR1--RSR8.  In the R2
source put \(S=1/T\), \(x=jSy\), \(B=Sy\), and \(H=S\ell\), where
\(0<\ell<1\).  Both testing angles \(u,v\) are positive
infinitesimals.  Normalize the rotated testing coordinate \(w=1\);
the other rotated coordinate \(z>0\) is arbitrary, including
unbounded values.  No coefficient estimate below involves \(z\).

Set
\[
 U=\frac{u}{\sqrt x},\quad V=\frac v{\sqrt y},\quad
 Z_*=\frac z{\sqrt{j/S}},\quad
 Q_*=U-VZ_*,\quad B_*=V-UZ_*,\quad R_*=Z_*-UV,
\]
\[
 \Psi=U^2+V^2+Z_*^2+U^2V^2Z_*^2-4UVZ_*.
\]
All monomials and radicals separate between the two fields.

\paragraph{A parameterized opposite-square identity.}
For every positive real \(a>1\), let \(K=a/(a-1)\),
\(d=U-V\), \(t=UV\), \(E=Z_*(1-t)\), \(T_*=t(Z_*-1)\).
Exact multiplication gives
\[
\begin{split}
 a\Psi+K V^2(Q_*^2+B_*^2)-R_*^2
 ={}&a d^2+(a-1)(E-T_*/(a-1))^2\\
 &+(2at+KV^4)(1-Z_*)^2+2KZ_*(t-V^2)^2 .
\end{split}                                                     \tag{SH1}
\]
The right side is a finite legal separated SOS.  For example
\(a=3/2\) gives \(K=3\).  This is a different budget from the
older one with \(U^2B_*^2\); neither budget dominates the other
uniformly at independent source and testing rates.

The two symmetric current budgets are
\[
\begin{split}
 (1+U^2)\Psi&=R_*^2+B_*^2+U^2Q_*^2
                         +U^2(1-UVZ_*)^2,\\
 (1+V^2)\Psi&=Q_*^2+R_*^2+V^2B_*^2
                         +V^2(1-UVZ_*)^2 .
\end{split}                                                     \tag{SH2}
\]
In particular the two actual rows
\[
 E_u=u(uvz-jy),\qquad E_v=v(uvz-jy)                       \tag{SH3}
\]
are free from a fixed positive fraction of \(f_c\Psi\) whenever
\(f_c\asymp j^2y^2\).  Their respective relative circuit costs
are \(O(x+u^2)\) and \(O(y+v^2)\), both infinitesimal.  Here
\emph{free} means that an arbitrarily large fixed real multiple of
the row square can be inserted at infinitesimal cost in the
strict circuit Gram.  It is not a pointwise bound without an
SOS witness: multiply (SH2) by \(j^2y^2x\) or \(j^2y^3\).

\paragraph{The full opposite row after the harmonic shift.}
Let
\[
 \begin{gathered}
 q=vz-u/S,\quad R=z-\frac{uv}{Sy},\quad R_v=u(v^2-y),\\
 d_s=\frac{BCH(1-p)(1+x)}{E},\qquad E=1+B(HT+p-1).
 \end{gathered}
\]
The exact actual source identity is
\[
                 \frac{\alpha}{y}-CHx=d_s.             \tag{SH4}
\]
Indeed it is equivalent to
\((K+Hjy)-jE=(1-p)(1+x)\), after canceling the displayed
positive source factors in the RSR dictionary.
The two exact testing identities
\[
 u(u^2-x)=S E_u-Su^2q,\qquad
 v(u^2-x)=S E_v-Suvq                                  \tag{SH5}
\]
now give the complete identity
\[
\begin{split}
 \sqrt S\,v(\alpha-yCHu^2)R
 ={}&[\alpha-yCHu^2-CHu^2(v^2-y)]\,(\sqrt S q)\\
 &+\sqrt S\,CH(v^2-y)E_u-\frac{d_s}{\sqrt S}R_v .
\end{split}                                                     \tag{SH6}
\]
There is no restriction on \(z\), and no normal has been set to
zero.  The first coefficient in (SH6) is infinitesimal, because
\(CH=O(1)\), \(y,u,v\) are infinitesimal, and
\(\alpha=O(Sy^2(j+1))\) is infinitesimal by \(jy\in\mathrm I\).
The second coefficient is infinitesimal and (SH3) is free.

For the last term write
\[
 R_v=-\sqrt x\,y(Q_*+VR_*),\qquad
 \frac{d_s}{\sqrt S}
   =\sqrt S\,y\,CH(1-p)(1+x)/E.                         \tag{SH7}
\]
Suppose the scalar has retained strict pivots of respective orders
\(jy\) and \(jy^2(1-p)\) on \(Q_*^2,R_*^2\), as in the actual
R2 source.  Twice Young's identity and (SH7) buy this row at
relative costs
\[
 O(S^2y^4(1-p)^2),\qquad
 O(S^2y^2(1-p)v^2),                                   \tag{SH8}
\]
respectively.  Both are infinitesimal even when \(1-p\) has
arbitrarily small positive order.  No inverse of that opening
remains.  Multiplication by the additional bounded source
factor \(\sqrt{\ell(1-\ell)}\), used to normalize the complete
second normal, preserves the estimate.

Thus the entire opposite term in that normal is controlled
uniformly through all R2 source ratios, once the stated scalar
pivots have been retained.  This lemma resolves that normal
allocation only.  It does not assert the signed scalar allocation
or all-source R2 positivity.



\section{\texorpdfstring{The positive-unit $e$ subdomain of weak-C2 R4}{The positive-unit e subdomain of weak-C2 R4}}
\label{component:p78dfeedee062}
% Scoped mathematical certificate; see local_identity_audit.closure.json for exact changes.

\label{p78dfeedee062:sec:r17-r4-eunit}

\begin{theorem}
Keep the original physical stationary C2 source, $d=1$, literal
$\lambda_{\rm event}=\beta=0$, $n\ge0$, all strict source gaps, the
positive-infinitesimal total-scale intersection, and every inherited
prior-admission exclusion. Suppose the ordered source is in R4 and
$e=1-r$ is a positive bounded unit. Then the entire original four-state
parent is a finite tensor SOS for every independent projective testing
pair and every real or complex state. This admits precisely the
$e$-unit subdomain of R4; the $e$-infinitesimal part is not claimed.
\end{theorem}

\paragraph{Source entry.}
Use the exact WE1 source parameters
\[
 a=z/B,\quad t=y/B,\quad h=H/B,\quad l=1-h,
 \quad g=1-ae+haB,\quad j=1+aBr-l/t.
\]
On R4, $B,z,y,j$ are positive infinitesimals and $t,h,l$ positive
bounded units. Since $e$ is a unit, $g>0$ forces $a$ bounded.
The exact $j$ identity gives $t_0=l_0$. In the following residue
formulas the letters on the right denote their real residues, with
$h+l=1$, $e=1-r$, and $g=1-ae\ge0$. Put
$M_*=1+ar-ael=g+a(r+eh)>0$. Joint regular reduction of WE1 gives
\[
 (nr)_0={e\over2(1+a)},\quad
 (mg)_0={a^2e\over2(1+a)},\quad
 (Ah)_0={l(1+ar)\over2(1+a)},\quad
 (Dj)_0={1+ar\over2l(1+a)},                         \tag{R4E1}
\]
\[
 C_0={M_*\over2l(1+a)},\qquad
 \alpha_{*,0}:=(\alpha/B^2)_0={(1+ar)M_*\over2(1+a)}.\tag{R4E2}
\]
Thus $A,C,\alpha_*,Dj$ are positive bounded units, while $D$ is
unbounded. The coefficient $n$ is either a positive unit or unbounded,
according as $r$ is a unit or infinitesimal. In particular literal
$n=0$ is impossible here by R4E1. No coefficient boundary was discarded.
If $g$ is a unit then $\mu=m/a^2$ is a positive bounded unit with
\[
                 \mu_0={e\over2g(1+a)}.              \tag{R4E3}
\]
If $g$ is infinitesimal, $a$ is a unit, $m$ unbounded, and $mg$ a
positive unit. These are joint regular formulas: after clearing $r,g,j$
by the displayed products, the remaining WE1 denominators are products
of $t,l,1-hB,\Delta$, with $\Delta_0=1+a>0$. No relative gap rates
are imposed. Write $N_*=n+1/2$, so $N_*^{-1}$ and $N_*^{-1/2}$ are
bounded and $nr$ is bounded in all source cases.

\paragraph{Exact removal of the divergent D channel.}
Use the complete WA3--WA8 identity in the scaled roots
\[
 {\cal R}=(\xi,\eta,R_1,R_2',T)^T,
 \quad \xi=X+ruZ,\quad\eta=Y-huW,
\]
\[
 R_1=uW-vZ/t,\quad R_2'=avW-uZ,\quad T=uvW-BZ,
 \qquad F/4={\cal R}^TG{\cal R}.
\]
The exact full-matrix identity RD3--RD6, proved independently of any
small-$h$ hypothesis in the R17 small-$r$ R8 lemma, gives
\[
 L_D=uv\xi-\eta+lR_1+ruvR_2'-arvT
     =uW+uvX-Y-vZ+jvZ,
\]
\[
 F/4=D L_D^2+{\cal R}^T(G|_{D=0}+DjH_D){\cal R}.       \tag{R4E4}
\]
Here the only nonzero symmetric entries of $H_D$ are
$(13)=tuv$, $(15)=-tu$, $(34)=rtuv$, $(35)=-artv$.
The underlying root syzygy is exact, and this identity contains the
full hard block and both complete first normals. Formal $D=0$ does
not assert a physical modified source. Put $\widehat D=\sqrt D L_D$
and solve the source-unit coefficient $l$:
\[
 R_1={1\over l}\left(D^{-1/2}\widehat D+\eta
       -uvN_*^{-1/2}\widehat\xi-ruvR_2'+arvT\right),
 \qquad\widehat\xi=\sqrt{N_*}\xi.                  \tag{R4E5}
\]
Every coefficient is bounded and separated, with no mixed inverse.

\paragraph{All local tests when g is a unit.}
The $D$-free matrix after the $\widehat\xi$ normalization has bounded
pure-source coefficients. Indeed all scalar occurrences of $n$ are
through $nr^2$, and all normal occurrences through $nr$; the hard
coefficient $n/N_*$ is bounded. $A,C,m,\alpha_*,\mu$ are bounded,
and $t,l,g$ have unit inverse. The only possible $a$ pole after using
$R_2'$ is
\[
 Q_*={b_2-\alpha_*\over a},\quad
 b_2=2a^2\mu h+Cl^2+(Ah)h+(B^2\alpha_*+1/2)h^2.     \tag{R4E6}
\]
The exact WE3 numerator vanishes at $a=0$ for arbitrary $B,h,r,t$.
After cancellation its denominator has only unit factors on this
$g$-unit chart; no $D$ or $j$ occurs. All other scaled WA coefficients
have the same regular rewrites as RD7. Thus R4E5 and bounded $Dj$
give a bounded full five-root Gram on
$(\widehat\xi,\eta,\widehat D,R_2',T)$.

At the local testing residue $u=v=0$, $D^{-1/2}$ vanishes and
$R_1=\eta/l$. The complete residue has independent pivots
$A_0/l_0^2$, $1$, $\alpha_{*,0}$ on $\eta,\widehat D,T$.
Its remaining $\widehat\xi,R_2'$ block is
\[
 \begin{pmatrix}
 1+a_0^2\mu_0\nu&-a_0\mu_0\sqrt\nu\\
 -a_0\mu_0\sqrt\nu&\mu_0
 \end{pmatrix},\qquad
 \nu=(N_*^{-1})_0\ge0.                              \tag{R4E7}
\]
The determinant is exactly $\mu_0>0$. This includes $r$ infinitesimal
($\nu=0$), $a$ infinitesimal, and every independent relative rate.
For reference, the $\eta$ pivot is the exact cancellation
\[
 (A+C+1/2)+{2[-(C+1/2)l+Ah]\over l}
       +{(C+1/2)l^2+Ah^2\over l^2}=A/l^2.
\]
No scalar or normal term is omitted in this computation.

\paragraph{All local tests when g is infinitesimal.}
Here $a$ is a unit. Use the unscaled $R_2=vW-uZ/a$, substituting
$R_2'=aR_2$ in R4E4--R4E5, and the exact
centered-$m$ matrix identity R5.2--R5.4 before R4E4. Namely, retain
\[
 (m-K)R_{mc}^2,\qquad R_{mc}=-\xi+uv\eta+(1-g)R_2+huT,
\]
and replace the remaining matrix by
$G|_{m=K}+\tau H_m$, where $\tau=(m-K)g$ is bounded with positive
residue and
\[
 (H_m)_{14}=-1,\quad(H_m)_{25}=v,\quad
 (H_m)_{44}=2-g,\quad(H_m)_{45}=hu.
\]
Choose a sufficiently large positive real $K$; $m-K>0$ for every
such $K$ since $m$ is unbounded. This identity and R4E4 commute,
as $H_m$ has no $D$ dependence. All remaining coefficients are
bounded: $a$ is a unit, $m$ has been replaced by $K$, $\tau$ is
bounded, and the only unbounded $n,D$ are treated by R4E5.
The same three residue pivots remain. The other block is now
\[
 \begin{pmatrix}
 1+K\nu&-(K+\tau_0)\sqrt\nu\\
 -(K+\tau_0)\sqrt\nu&K+2\tau_0
 \end{pmatrix},\qquad
 \det=K+2\tau_0-\tau_0^2\nu.                         \tag{R4E8}
\]
Since $0\le\nu\le2$, any real $K>2\tau_0^2+1$ makes it strict.
No source opening is compared with a testing rate. Restore the exact
retained square. Equations R4E7--R4E8 and WA10 therefore prove the
entire local parent for all independent infinitesimal or zero $u,v$,
with arbitrary states and either sign. The finite separated lift
whitens the strict pure-source matrix and handles its finitely many
infinitesimal testing monomials; it introduces no mixed radicals.

\paragraph{Finite remote tests.}
If both finite angles have nonzero residue, use the inherited
all-noncommuting-residues theorem, at arbitrary coefficient scales
(the precise entry is the inherited Section~\ref{component:p2c3fc22497a5}).
Its hypotheses apply to this unchanged physical parent. It remains
to give the two one-unit-angle charts. In these remote charts,
positive original event coefficients may be clipped at positive real
lower bounds and their exact positive excess packets retained; no
contact assertion is made for the clipped polynomial.

If $u$ is a bounded unit and $v$ infinitesimal or zero, clip $n,A,C$
at positive real lower bounds $n_*,A_*,C_*$, and retain all their
excesses, the entire $m$ and $\alpha$ packets. Normalize the original
$D$ event by $\widehat J=\sqrt D(uW+uvX-Y-vZ)$ and solve
\[
 W={D^{-1/2}\widehat J+Y+vZ-uvX\over u}.
\]
The whole remainder has bounded coefficients on
$(\widehat J,X,Y,Z)$. Its residue is
\[
 \widehat J^2+n_*D_uX^2+\tfrac12(X+uZ)^2+A_*D_uY^2\succ0,
 \qquad D_u=1+u_0^2.                                \tag{R4E9}
\]
Here and in the remaining residue displays the tests denote their real
residues. Positivity is strict because $u_0\ne0$. All source cases,
including unbounded $n,m$, have been covered by exact excess packets.

If $u$ is infinitesimal or zero, $v$ a bounded unit, and $a$ a unit,
$m$ has a positive real lower bound even when it is unbounded.
Clip $n,m,A,C$ at positive real bounds and retain their excesses and
$\alpha$. Solve instead
$Z=(uW+uvX-Y-D^{-1/2}\widehat J)/v$. The bounded remainder residue is
\[
 \widehat J^2+(n_*+1/2)D_vX^2+m_*(vW-X)^2
                  +(A_*+C_*+1/2)D_vY^2\succ0.       \tag{R4E10}
\]
This is strict since $m_*>0$ and $v_0\ne0$.

Finally let $a$ be infinitesimal in this chart. Then $g_0=1$,
$\mu=m/a^2$ is a unit, and
\[
       A_0={l\over2h},\qquad C_0={1\over2l}.          \tag{R4E11}
\]
Keep $S=aW$ as an extra legal root and write $R=uW$ everywhere
else. Clip only $n$ at a positive real $n_*$, retaining its excess
and the $\alpha$ packet. Retain the actual $A,C,m,D$ and the whole
signed core. Use
\[
 R=Y+vZ-uvX+D^{-1/2}\widehat J,\qquad
 mM^2=\mu(vS-aX+auvY-auZ)^2.                         \tag{R4E12}
\]
There is no inverse of $a$ or $u$ here. The complete remainder is
a bounded matrix on $(S,X,Y,\widehat J,Z)$. Its separate residue
pivots are $\mu_0V$, $(n_*+1/2)D_v$, and $1$ on
$S,X,\widehat J$. The remaining $Y,Z$ matrix is
\[
 \begin{pmatrix}
 (l+V)/(2h)&-v/2\\
 -v/2&V(h+V)/(2l)
 \end{pmatrix},\qquad
 \det={V^2(1+V)\over4hl}>0,
 \quad V=v_0^2>0.                                   \tag{R4E13}
\]
This exact computation retains the entire signed core; it is not a
state-slice argument. The possibly dependent extra root $S$ causes
no difficulty for a positive Gram. WA10 at the real angle residues
now proves each complete finite remote chart.

\paragraph{Projective endpoints and the join.}
The source bounds give positive real lower bounds for $n,C,A,D$.
Clip them at $n_*,C_*,A_*,D_*>0$, retain their whole positive
excess packets, and retain the entire original $m,\alpha$ packets.
All remaining source coefficients are real constants. The following
three exact homogeneous charts give the complete endpoint join,
without any further source hypothesis.

For $u$ unbounded put $p=1/u$ and use $p^2F/4$. If $v$ is bounded,
the selected residue at $p=0$ is
\[
 n_*D_vX^2+C_*D_vW^2+A_*D_vY^2
       +(D_*+1/2)(W+vX)^2+\tfrac12(vY-Z)^2\succ0.     \tag{R4E14}
\]
It is strict on $(X,W,Y,Z)$ for every real $v$, including zero.
The full homogeneous selected matrix has bounded polynomial entries.

For $v$ unbounded put $q=1/v$ and use $q^2F/4$. If $u$ is bounded,
write $R=uW$; every remaining occurrence of $W$ is through $R$.
The selected residue at $q=0$ is
\[
 (n_*+1/2)D_uX^2+C_*(R-Y)^2
       +(A_*+1/2)D_uY^2+D_*(uX-Z)^2\succ0.           \tag{R4E15}
\]
It is strict on $(X,R,Y,Z)$ for every real $u$, including zero,
and the exact selected polynomial uses no inverse of $u$.

If both tests are unbounded use $p=1/u,q=1/v$ and $p^2q^2F/4$,
with the legal roots $P=pZ,Q=qZ$. The exact half-core becomes
\[
 \tfrac12\{(qW+X-pqY-P)^2+(1+p^2)(Y-Q)^2
                       +(1+q^2)(p^2X^2+2XP-P^2)\}.
\]
Together with the selected event packets its residue is
\[
 (n_*+1/2)X^2+C_*W^2+A_*Y^2+D_*(X-P)^2
                                  +\tfrac12(Y-Q)^2\succ0,       \tag{R4E16}
\]
strict on $(X,W,Y,P,Q)$. Dependent roots do not affect positivity.
All finite reciprocal deviations are pure testing infinitesimals,
so the separated strict-Gram lift applies in all three charts.
Restore the original positive excesses and $m,\alpha$ packets.
Homogeneous restoration uses pure-testing squares off the endpoints;
literal $p=0$ or $q=0$ uses the polynomial endpoint value itself.

The finite local, finite noncommuting, both finite one-unit-angle
charts, and the homogeneous infinity charts exhaust all independent
tests. All four original state coordinates and both normal directions
were retained. Polynomial axis values, sign gauges, and application to
real and imaginary state parts supply the literal axes, both signs,
and complex states. The result is an all-tests admission of the
positive-unit-$e$ source subdomain of R4, with its original physicality,
strict gaps, coefficient signs, and residual intersections unchanged.


\section{\texorpdfstring{Weak-C2 R3 with a positive-unit $1-r$ is admitted}{Weak-C2 R3 with a positive-unit 1-r is admitted}}
\label{component:p1c9902e9a6f2}
% Scoped mathematical certificate; see local_identity_audit.closure.json for exact changes.

\label{p1c9902e9a6f2:sec:r18-r3-eunit}

\begin{theorem}
Keep the entire original physical stationary weak-C2 source, $d=1$,
$\lambda_{\rm event}=\beta=0$, every strict gap and nonnegative
coefficient, and every original residual intersection. Suppose
$B,z,y,l=1-h$ are positive infinitesimals, $t=y/B$ a positive bounded
unit, and $e=1-r$ a positive bounded unit. Then the full original
parent is a finite tensor SOS for every independent projective testing
pair and every real or complex state. These are precisely the
positive-unit-$e$ sources of R3. The remaining $e$-infinitesimal
part of R3 is not asserted here.
\end{theorem}

\paragraph{Source estimates including all simultaneous openings.}
Put $a=z/B$, $g=1-ae+haB$ and $j=1+aBr-l/t$. The positivity
of $g$ and the unit lower bound for $e$ imply that $a$ is bounded.
Moreover $h\sim1$ and $j\sim1$. In the following formulas variables
on the right denote their real residues, $e=1-r>0$ and
$g=1-ae\ge0$. Direct joint regular reduction of the unchanged WE1
formulas gives
\[
 (nr)_0={e\over2(1+a)},\qquad
 (mg)_0={a^2e\over2(1+a)},\qquad
 A_0={t(1+ar)\over2(1+a)},\qquad
 D_0={1+ar\over2t(1+a)},                         \tag{R3E1}
\]
\[
 (Cl)_0={(1+ar)(1+t)\over2(1+a)},\qquad
 \alpha_{*,0}:=(\alpha/B^2)_0
       ={(1+ar)^2(1+t)\over2(1+a)}.              \tag{R3E2}
\]
Thus $A,D,Cl,\alpha_*$ are positive bounded units, $C$ is unbounded,
and $n$ has a positive real lower bound, even when $r$ is
infinitesimal and $n$ is unbounded. In particular literal $n=0$
is impossible on this source domain, by the positive first formula
in R3E1; no coefficient boundary was discarded. Write
$N=n+1/2$, so $N^{-1}$ is bounded. If $g$ is a unit, then
$\mu=m/a^2$ is a positive unit with
\[
              \mu_0={e\over2g(1+a)}.             \tag{R3E3}
\]
If $g$ is infinitesimal, $a$ is a unit, $m$ is unbounded and $mg$
is a positive unit. These joint statements use only the displayed
products: the WE1 denominators after clearing $r,g,l$ have unit
factors $t,h,1-hB,\Delta$ with $\Delta_0=1+a$.
There is no comparison among the $B,l,r,g$ openings. The original
positive-infinitesimal total-scale condition follows already from
$C$ unbounded.

\paragraph{The full local matrix with its divergent C pivot retained.}
Use the exact complete WA3--WA8 parent identity in the roots
\[
 \xi=X+ruZ,\quad \eta=Y-huW,\quad
 R_1=uW-vZ/t,\quad R_2'=avW-uZ,\quad T=uvW-BZ.
\]
Normalize its first two rows by
$\widehat\xi=\sqrt N\xi$ and $\widehat\eta=\sqrt C\eta$.
First suppose $g$ is a unit. The complete resulting five-root
matrix has bounded pure-source polynomial coefficients in $u,v$.
Here is the coefficient check, including the potentially singular
source openings. The hard entries use $n/N$, $C/C$, and the bounded
coefficients $A,D,m,\alpha$, together with bounded $N^{-1/2}$ and
infinitesimal $C^{-1/2}$. All scalar occurrences of $n$ are through
$nr^2$, and all first-normal occurrences are through $nr$, both
bounded. All non-hard occurrences of $C$ are through $Cl$ or
$Cl^2$, both bounded. Thus the entire $\eta$ normal row becomes
infinitesimal after division by $\sqrt C$. This includes the full
WA7 terms, not merely their zero-angle values.

The WE5 scaled-$R_2'$ formulas bound every coefficient through
$a\to0$ except the apparently singular
\[
 Q_*={b_2-\alpha_*\over a},\qquad
 b_2=2a^2\mu h+Cl^2+Ah^2+(B^2\alpha_*+1/2)h^2.  \tag{R3E4}
\]
The inherited exact WE3 numerator vanishes at $a=0$ for arbitrary
$B,h,r,t$. Its quotient has a unit denominator here too: the $l$
denominator cancels in $Cl^2$ before division, and the remaining
factors $g,h,t,j,1-hB,\Delta$ are units on this $g$-unit chart.
Consequently $Q_*$ is bounded jointly through $a,l,r\to0$.
This bounds the entire scaled scalar block and both normal rows,
without any $a/l$ or $l/a$ hypothesis.

At $u=v=0$ the full residue has separate pivots $1,A_0,\alpha_{*,0}$
on $\widehat\eta,R_1,T$. Its $\widehat\xi,R_2'$ block is
\[
 \begin{pmatrix}
 1+a_0^2\mu_0\nu&-a_0\mu_0\sqrt\nu\\
 -a_0\mu_0\sqrt\nu&\mu_0
 \end{pmatrix},\qquad
 \nu=(N^{-1})_0\ge0,\qquad \det=\mu_0>0.          \tag{R3E5}
\]
This includes $a$ infinitesimal, $r$ infinitesimal, and their
independent simultaneous opening. The $R_1$ pivot is $A_0$ because
$Cl^2\sim0$ and $h\sim1$; it does not require a new soft root.

If $g$ is infinitesimal, use the unscaled $R_2=R_2'/a$, as $a$
is a unit. Apply the exact R5.2--R5.4 centered-$m$ identity, keeping
\[
 (m-K)R_{mc}^2,\qquad
 R_{mc}=-\xi+uv\eta+(1-g)R_2+huT,
\]
for a sufficiently large positive real $K$. In the remainder
replace only the formal matrix coefficient $m$ by $K$ and add
$\tau H_m$, where $\tau=(m-K)g$ is bounded and $H_m$ is the exact
R5.4 matrix. The remaining source coefficients, especially
$\alpha_*$, stay at their actual values. This is a whole-matrix
identity, not a physical-source alteration. The same coefficient
bounds hold; all normal entries divided by $\sqrt C$ vanish.
The three separate residue pivots remain $1,A_0,\alpha_{*,0}$;
the last block becomes
\[
 \begin{pmatrix}
 1+K\nu&-(K+\tau_0)\sqrt\nu\\
 -(K+\tau_0)\sqrt\nu&K+2\tau_0
 \end{pmatrix},\qquad
 \det=K+2\tau_0-\tau_0^2\nu.                       \tag{R3E6}
\]
Since $0\le\nu\le2$, any real $K>2\tau_0^2+1$ makes it strict.
Thus WA10 applies to the entire local parent in both source cases,
for all independent infinitesimal or zero $u,v$, either sign and
all states. Restore the exact retained square in the $g$-small case.

\paragraph{Every finite remote testing chart.}
If both finite testing residues are nonzero, use the inherited
arbitrary-coefficient-scale noncommuting-residue theorem. It remains
to treat one nonzero finite angle residue. In these charts select
positive real lower bounds for original positive event coefficients
when indicated and retain their exact positive excess packets;
no stationary-contact assertion is made for the selected remainder.

Use the actual original $C$ event with
\[
       J_C=\sqrt C\,(uW-Y),\qquad d_C=C^{-1/2}\in\Inf.
                                                               \tag{R3E7}
\]
If $u$ is a bounded unit and $v$ infinitesimal or zero, clip $n,A,D$
at positive real lower bounds and keep their positive excesses and
the whole original $m,\alpha$ packets. Solve
$W=(Y+d_CJ_C)/u$. The selected complete remainder has bounded
coefficients in $(J_C,X,Y,Z)$ and residue
\[
 J_C^2+n_*(1+U)X^2+A_*(1+U)Y^2
                         +\tfrac12(X+uZ)^2\succ0,               \tag{R3E8}
\]
where $u$ on the right is its nonzero real residue and $U=u^2$.
The $D$ event has zero residue and causes no difficulty.

If $u$ is infinitesimal or zero, $v$ a bounded unit and $a$ a unit,
then $m$ has a positive real lower bound (also when $g$ is small).
Clip $n,m,A,D$ at positive bounds, retain all excesses and $\alpha$,
and solve $Y=uW-d_CJ_C$. The complete selected residue is
\[
 (1+V)J_C^2+(n_*+1/2)(1+V)X^2
       +m_*(vW-X)^2+D_*VZ^2\succ0,\quad V=v_0^2>0.              \tag{R3E9}
\]
This retains every state direction and the whole signed core.

Finally suppose $a$ is infinitesimal in the latter testing chart.
Then $g\sim1$, $\mu$ is a positive unit, and R3E1 gives
$A_0=t_0/2$, $D_0=1/(2t_0)$. Keep the additional legal roots
$S=aW$, $R=uW$ and use $Y=R-d_CJ_C$. Clip only $n$, retain its
excess and $\alpha$, and retain the actual bounded $A,D,\mu$.
The $m$ event is exactly
\[
 m(vW-X+uvY-uZ)^2
       =\mu(vS-aX+auvY-auZ)^2.                    \tag{R3E10}
\]
The entire bounded matrix on $(J_C,S,X,R,Z)$ has separate
positive pivots $1+V$, $\mu_0V$ and $(n_*+1/2)(1+V)$.
Its remaining $R,Z$ block is
\[
 \begin{pmatrix}
 [t_0(1+V)+V]/2&-v_0/2\\
 -v_0/2&V/(2t_0)
 \end{pmatrix},\qquad
 \det={V^2(t_0+1)\over4t_0}>0.                    \tag{R3E11}
\]
The dependent roots $R,S$ are allowed in a positive Gram; no inverse
of $u$ or $a$ is introduced. WA10 at the real testing residues
establishes every finite remote chart above.

\paragraph{Projective join and exact scope.}
The source estimates give positive real lower bounds for $n,A,C,D$.
Therefore the complete selected-event endpoint polynomials R4E14,
R4E15 and R4E16 apply verbatim: clip these four original event
coefficients, retain all their excesses and the entire $m,\alpha$
packets, and use the original degree-$(2,2)$ homogeneous parent.
Those exact identities cover respectively one infinite Alice angle,
one infinite Bob angle, and both infinite angles, including the
zero values in the other projective chart. Their strict Grams do
not depend on any R4-specific source relation. Thus every finite
and infinite testing pair is covered. Polynomial axes, both signs
and the real/imaginary state split give literal zeros and complex
states. All inverses and radicals used are pure-source or pure-test.

This admits exactly R3 with $1-r$ a positive bounded unit. It neither
claims the remaining $1-r$-infinitesimal sources empty nor changes
the historical weak-source partition. No numerical table is a
premise. The focused checker verifies the new source residues, full
local residue blocks and complete finite remote residues; unchanged
WA, centered-$m$, and homogeneous endpoint identities are imported.

\paragraph{Limits of the discrete symmetry.}
In the proved cyclic order $(\lambda,D,C,A,\alpha,m,\beta,n)$,
the stabilizer of the omitted pair $\{\lambda,\beta\}=\{0,6\}$
inside the CHSH dihedral action consists of the identity and
$i\mapsto6-i$. The latter is exactly
\[
 (u,v;W,X,Y,Z;n,m,C,A,D,\alpha)
 \longmapsto(1/v,1/u;Z,X,Y,W;n,D,\alpha,A,m,C).
\]
It exchanges strong $C$/weak $\alpha$ with the opposite directed
orientation. Returning to the chosen weak orientation restores the
original source, so this known discrete action supplies no additional
within-orientation routing of R1--R3. This statement concerns only
the established dihedral action, not all conceivable analytical
changes of variables.


\section{\texorpdfstring{R3 and R4 at infinitesimal \(e\), with all source ratios and tests}{R3 and R4 at infinitesimal e, with all source ratios and tests}}
\label{component:p1de8b5ffab69}
% Scoped mathematical certificate; see local_identity_audit.closure.json for exact changes.

\label{p1de8b5ffab69:sec:new-r34-all-tests}

\paragraph{Statement and exact domain.}
Use the original common-II stationary source, its six native weights
$n,m,C,A,D,\alpha$, all its strict source gaps, $d=1$,
$\lambda=\beta=0$, and $n\geq0$. Suppose $B_s,a_sB_s,t_sB_s$
are positive infinitesimals, $t_s,h$ are positive bounded units,
and $e=1-r$ is positive infinitesimal. In addition impose either
\[
 \begin{array}{ll}
 \mathrm{R3}:&1-h\text{ is infinitesimal},\\
 \mathrm{R4}:&1-h\text{ is a positive unit and }
       1+a_sB_sr-(1-h)/t_s\text{ is infinitesimal}.
 \end{array}                                                    \tag{EN1}
\]
There is no condition on the size of $a_s$, $a_s^2B_s$,
$a_s e$, or their ratios. The conclusion is finite separated tensor
SOS for the entire original parent at every independent projective
testing pair and every state, including signs, axes, and literal $n=0$.
The already recovered positive-unit-$e$ proofs then complete the
whole original R3 and R4 rows. No other weak-C2 row is inferred.

The source and testing fields are independent real closed extensions
$K_c,K_t\supset\mathbb R$. A bounded coefficient means an element
of $B_c\otimes B_t$; its residue is the homomorphism to $\mathbb R$.
The strict bounded-Gram lemma is used only for matrices with positive
definite real residue. Row entries may be arbitrarily large.
Every inverse and radical below is pure source, pure testing, or a
product of such factors.

\subsubsection*{One exact compact source for both ends}

Use the exact compact variables of the reconstructed actual-source
dictionary RSR2--RSR6:
\[
\begin{aligned}
c_s&=1+a_sB_sr,\quad R_s=1+ha_sB_s,\quad
j=a_sc_s/R_s,\quad B=B_sR_s/c_s,\quad T=t_sc_s/R_s,\\
p&=a_s e/R_s,\quad H=(1-h)/(t_sc_s),\qquad
x=jB=a_sB_s,\quad y=TB=t_sB_s,\quad q_z=jT .
\end{aligned}                                                    \tag{EN2}
\]
Thus $0<p<\min(1,j)$, $0<H<1$, and
\[
B,x,y\in\mathrm I_c,\quad T\in\mathrm U_c,\quad
1-HT\in\mathrm U_c,\quad p/j\in\mathrm I_c.              \tag{EN3}
\]
The last assertion is exact source entry:
$e=p(1+x)/(j+px)\sim p/j$.
In R3, $H$ is infinitesimal. In R4, $1-H$ is infinitesimal
and $0<\overline T<1$. In particular
$H(1-H)\in\mathrm I_c$ in both rows.
The native inverse is
\[
r=\frac{j-p}{j+px},\quad h=\frac{1-HT}{1+HTx},\quad
a_s=\frac{j+px}{1+HTx},\quad
B_s=\frac{B(1+HTx)}{1+pB},\quad
t_s=\frac{T(1+pB)}{1+HTx}.                              \tag{EN4}
\]
The exact six weights are RSR6 at this substitution.
Those rational identities, including every contact first jet, were
checked against the native source without size restrictions. Their
former R1 size hypothesis is not imported.

Put $\sigma=1+(1-H)T$, $K=j+1-p$, $E=1+B(HT+p-1)$, and
use the explicit polynomials $\mathcal P,\mathcal L,\mathcal N,
\mathcal F,\mathcal K$ in RSR5, RSR15, RSR18. Then
\[
\frac{\mathcal P}{j+1}\sim1,\qquad
\frac{\mathcal L}{j+1}\sim\sigma,\qquad
\frac{\mathcal F}{j+1}\sim\sigma.                       \tag{EN5}
\]
Here $T,\sigma,(1-HT)^{-1}$ are bounded positive units.
To prove these uniform statements, divide the printed finite
polynomials by $j+1$. The corrections in $\mathcal P$ contain
$B+x$ or $TB$; those in $\mathcal L$ after
$(j+1)\sigma$ contain $p/(j+1)$ or $B+x$.
In $\mathcal F$ the only additional quotient is
$p(1-p)/(j(j+1))$, infinitesimal by (EN3).
Every other term contains $B$ or $Bj/(j+1)$.
Thus no restriction on $j$ or $jT$ is needed.

The exact physicality numerator has the grouping
\[
\mathcal N=p\{1+B(3j+1-p)+BT(j+1)\}
-Bj(j+1)\{1+(1-H)T\}+\mathcal R,\quad
|\mathcal R|\le C_0BT\,Bj(j+1).                         \tag{EN6}
\]
The bound follows termwise from RSR5 using $p<\min(1,j)$ and
bounded $T$. Since $n=\mathcal N/[2(j-p)\mathcal P]\ge0$,
\[
\begin{aligned}
&\eta_*=\frac{Bj(j+1)}p=O(1),\qquad
0\le n=O\!\left(\frac{p}{j(j+1)}\right)\in\mathrm I_c,\\
&\eta_*\in\mathrm I_c\Longrightarrow
n\sim\frac{p}{2(j-p)(j+1)}>0 .
\end{aligned}                                                    \tag{EN7}
\]
This includes $\mathcal N=0$ literally.

Write $\mu=m(1-p)$, $h_0=\mu+1/2$. The full source formulas give
\[
\begin{aligned}
&\mu=O(p),\quad h_0\asymp1,\quad
\alpha/[B^2(j+1)]\sim\sigma/2,\quad \alpha/B\in\mathrm I_c,\\
&A\sim\frac{T}{2(1-HT)},\qquad
CH\sim\frac{\sigma}{2T},\qquad
D(1-H)\sim\frac1{2T}.
\end{aligned}                                                    \tag{EN8}
\]
Consequently $A$ and the finite one of $C,D$ have positive real
lower bounds. The other coefficient is unbounded. The exact
normalized source limits are
\[
\begin{array}{c|c|c|c}
 & C&A&D\\ \hline
R3&+\infty&T/2&1/(2T)\\
R4&1/(2T)&T/[2(1-T)]&+\infty .
\end{array}                                                    \tag{EN9}
\]
Only real residues of the bounded entries are meant in this table.
The coefficient $m$ may be infinitesimal, a unit, or unbounded.

\subsubsection*{The scalar allocation in the complete two-normal frame}

Use the entire parent RFN4 and the exact full-state change RFN5--RFN9,
with $w=(W+uY)/(1+u^2)=1$ and arbitrary $z=(Z-uX)/(1+u^2)>0$.
The identities are rational identities of (EN4), so they retain
both complete normals and every state term in the present source.
The tests in this subsection satisfy $u,v>0$ infinitesimal.
For clarity, the inverse is
\[
X=N_0+pv-uz,\quad Z=z+uN_0+puv,\quad
W=1+u(M_0+Hvz),\quad Y=u-M_0-Hvz.                       \tag{EN10}
\]
Retain the actual positive packet $mpD_u^2N_0^2$, and divide the
remaining entire parent by the pure testing factor $D_u=1+u^2$.
Its exact expansion is RFN7--RFN9. In particular the hard entries are
\[
\begin{aligned}
&H_N=h_0D_u+\tfrac12v^2+\alpha u^2+nD_v,\\
&H_M=D_u(CD_v+D+\tfrac12)+AD_v+(\alpha+\tfrac12)v^2,\qquad
H_{NM}=(\alpha+\tfrac12)uv .
\end{aligned}                                                    \tag{EN11}
\]

Keep the separated rows
$U=u/\sqrt{x}$, $V=v/\sqrt{y}$, $Z_*=z/\sqrt{q_z}$,
$Q_*=U-VZ_*$, $B_*=V-UZ_*$, $R_*=Z_*-UV$, and
\[
\Psi=U^2+V^2+Z_*^2+U^2V^2Z_*^2-4UVZ_*
=(U-V)^2+Z_*^2(1-UV)^2+2UV(1-Z_*)^2.                 \tag{EN12}
\]
The last radical is separated, for example
$\sqrt{UV}=(xy)^{-1/4}\sqrt{uv}$.
The exact scalar identity is RFN12 with the same actual circuit
masses. From (EN5), the closed formulas RSR16--RSR20 now give
\[
\begin{aligned}
&f_c\sim Tx^2\sigma/2,\quad
q_c\sim\frac{Tx}{2(1-HT)},\quad
r_c\sim TjB^2(1-p)\sigma/2,\\
&g_c/f_c=Hy\,\mathcal L/\mathcal F\in\mathrm I_c,\quad
k_c<0,\quad |k_c|/(q_cy)=O(p+B).
\end{aligned}                                                    \tag{EN13}
\]
In particular all three displayed positive masses are positive.
The sign of $k_c$ follows from the same exact positive-surplus
identity RSR19: its error bound uses $B(j+1)\in\mathrm I_c$,
not $T$ infinitesimal.

Before spending any current mass, perform the exact source-only
completion
\[
\begin{aligned}
h_0N_0^2+n(N_0+jb)^2
&=(n+h_0)(N_0+c_bb)^2+M_eb^2,\\
b&=v-uz/j,\quad c_b=\frac{nj}{n+h_0},\quad
M_e=\frac{nh_0j^2}{n+h_0}.
\end{aligned}                                                    \tag{EN14}
\]
Keep $nv^2(N_0+pb)^2$ and the positive extra hard terms RFN15
as whole squares. This reserves exactly the correlated $n$ packets
of the complete parent. If $n=0$, set $c_b=M_e=0$ and omit that
row. For $n>0$, $c_b^2/M_e=n/[h_0(n+h_0)]$ is bounded.

The signed scalar allocation differs from R1 in one source split.
The $e_0$ cost is unchanged: its negative part is at most $g_c$,
and RFN18b buys it at cost
$H(\mathcal L/\mathcal F)(y+v^2)\in\mathrm I$.
Also $d_c>0$. Use $Z_*(V^2-1)=-VQ_*-R_*$ for $k_c$.
The $VQ_*$ cost is $O((p+B)v^2)$, infinitesimal.
For its opposite-row part the exhaustive cases are:
\[
\begin{array}{c|l}
p\in\mathrm I_c&
 |k_c|/r_c=
 T\mathcal K/[(1-p)(1+Hy)\mathcal L]=O(p+B)\in\mathrm I_c\\
p\in\mathrm U_c,\ \eta_*\in\mathrm U_c&
 j\text{ unbounded},\ j^2B\in\mathrm U_c,\ 
 |k_c|/f_c\in\mathrm I_c,\ |k_c|/(f_cx)=O(1)\\
p\in\mathrm U_c,\ \eta_*\in\mathrm I_c&
 j\text{ unbounded},\ M_e\in\mathrm U_c,\ 
 |k_c|/f_c\in\mathrm I_c,\\
& |k_c|/(q_cy)=O(1),\ |k_c|/(M_eyx)=O(1).
\end{array}                                                    \tag{EN15}
\]
In the second case $j$ is unbounded because $p/j$ is infinitesimal;
$j^2B\asymp p\eta_*$. In the third,
$nj^2\sim p/2$ and $n$ is infinitesimal, so $M_e\sim p/2$.
The closed formulas show
$|k_c|/f_c=O((p+B)/j)$ in the last two cases and
$|k_c|/(M_eyx)=O((p+B)(1+1/n)/j^2)$ in the third.
These bounds are uniform through $p\to1$.

Case two uses RFN18a:
$R_*^2\preceq_{\rm sos}(1+U^2)\Psi$.
Case three uses RFN18c:
$R_*^2\preceq_{\rm sos}2(\Psi+V^2Q_*^2+U^2B_*^2)$.
The non-infinitesimal source ratios in (EN15) acquire the testing
factors $u^2$ or $v^2$. Thus all costs are infinitesimal.
Literal $n=0$ lies in the first or second case; it is not discarded.
After this finite allocation, fixed positive fractions of
\[
q_cQ_*^2,\quad r_cR_*^2,\quad f_c\Psi,\quad M_eb^2       \tag{EN16}
\]
remain, with the last row absent exactly at $n=0$.

\subsubsection*{Both complete normal arrays, with the new endpoint weights}

The exact derivatives RFN21--RFN22, and their exact source cancellations
$S_N=S_0=S_1=S_2=0$, hold unchanged. Let
$q=vz-Tu$, $R=z-uv/B$, $R_u=v(u^2-x)$. In particular
\[
\begin{aligned}
L_N^{\rm rem}
&=-j\alpha b+\tfrac12pv^2b+K_RR_u+(K_c/j)uvq,\\
L_M&=C_qq+\mu uvb+v(\alpha-BTCHu^2)R,\\
K_c&=p/2+j\{H(\alpha+1/2)-1/2\}+n(p-j),\\
K_R&=-\mu-(\alpha+1/2)(1-p)+(T/j)K_c,\\
K_d&=CH-(D+1/2)(1-H)=-\mu/j-BTCH,\\
C_q&=K_d+AH+1/2+
v^2\{CH(1+u^2)+AH+(\alpha+1/2)H\}.
\end{aligned}                                                    \tag{EN17}
\]
Every displayed coefficient other than a retained row is bounded:
$p/j$ and $n$ are infinitesimal, $\mu/j=O(p/(j+1))$ is bounded,
and $A,CH,D(1-H),T$ are bounded by (EN8).
In particular $K_R$, $K_c/j$, and $C_q$ are bounded.

Here is an explicit audit of all costs, to avoid importing the former
$T\in\mathrm I_c$ hypothesis of RFN.
The exact radial Gram RFN19--RFN20 purchases $R_u^2$ at infinitesimal
cost from (EN16). Its three cases are still $1-p$ a unit;
$p\sim1,\eta_*$ a unit; and $p\sim1,\eta_*$ infinitesimal.
Their costs are respectively
\[
\begin{array}{c|l}
1-p\in\mathrm U_c&O(B),\ O(Bu^2),\ O(u^2/(1-p))\\
p\sim1,\eta_*\in\mathrm U_c&
 O(B+u^2/j+u^4/(j^2B)),\ O(Bu^2)\\
p\sim1,\eta_*\in\mathrm I_c&
 O(u^2/j),\ O(u^4/j),\
 [8u^4+2(u^2-x)^2]/M_e .
\end{array}                                                    \tag{EN18}
\]
They require only $T$ bounded above and below, because $T$ cancels
between $y=TB$ and the masses (EN13).
The small-current budgets NSR3--NSR5 likewise apply unchanged:
both $pv^2b$ and $\sqrt T\,\mu uvb$ are free.
For $\eta_*$ infinitesimal their coefficients are
$O(v^4)$ and $O(Tu^2v^2)$, using $p^2/M_e=O(1)$.
For $\eta_*$ a unit the coefficients are
$O(v^2(y+v^2))$ and $O(Tu^2(y+v^2))$, using
$p^2y/f_c=O(1)$. Each is infinitesimal.

The row $q/\sqrt T$ is strict by (EN13).
Its $N$ coupling has coefficient $(K_c/j)uv\sqrt T$ infinitesimal.
The identity $j\alpha b=-(\alpha/B)R_u-\alpha uR$ handles
the remaining $N$ term. For $1-p$ a unit the latter square costs
$O(B^2(j+1)^2u^2/(1-p))$ against $r_cR_*^2$.
For $p\sim1$, RFN18a bounds the entire $j\alpha b$ square by
$O(B^3(j+1)^2+B^2(j+1)^2u^2/j)$ of the circuit.
These are infinitesimal. A sufficiently large fixed real radial
pivot pays the bounded $K_R$ coupling, including the bounded shift
$c_b/\sqrt{M_e}$ from (EN14).

Normalize the complete second normal by
$M_0=\sqrt{TH(1-H)}\,\widehat M$. The full hard diagonal (EN11)
then has residue
\[
\overline{TH(1-H)H_M}=
\begin{cases}
(1+\overline T)/2,&R3,\\
1/2,&R4 .
\end{cases}                                                    \tag{EN19}
\]
Every term of the complete diagonal is bounded: the only divergent
factors occur in $CH$ and $D(1-H)$.
The normalized $q$ coupling is
$T\sqrt{H(1-H)}C_q$, now infinitesimal because
$H(1-H)$ is infinitesimal. This is the replacement for the
$T$-infinitesimal argument in R1.
The normalized current coupling is a bounded multiple of the
already purchased $\sqrt T\,\mu uvb$.

For the entire remaining opposite row put $J=B(j+1)+u^2$ and
$a_R=Tq_zv^2(\alpha-BTCHu^2)^2$.
Since $\alpha-BTCHu^2=O(BJ)$, its costs are
\[
\begin{aligned}
a_R/f_c&=O((T/j)v^2J^2),&
a_R/(q_cy)&=O(v^2J^2),\\
a_R/r_c&=O(Tv^2J^2/(1-p)),&
a_R/(M_eyx)&=O(T(1+1/n)v^2J^2/j^2).
\end{aligned}                                                    \tag{EN20}
\]
The direct $r_c$ cost applies for $1-p$ a unit.
For $p\sim1,\eta_*$ a unit, RFN18a applies using $j^2B$ a unit.
For $p\sim1,\eta_*$ infinitesimal, RFN18c applies using
$(1+1/n)/j^2=O(1)$. Thus the entire square, including every
higher testing degree, is free. No state ratio occurs in a cost.

Finally the complete mixed hard term is
$2(\alpha+1/2)uv\sqrt{TH(1-H)}
(\widehat N/\sqrt{n+h_0}-c_bb)\widehat M$.
Its bounded coefficients have zero residue.
The strict bounded-Gram argument therefore applies exactly as in
the final paragraph of RFN: choose a large fixed radial pivot, buy
the finitely many free rows, and retain positive fractions of
(EN16). The two hard pivots are $1$ and the positive number
(EN19). Restore all whole packets in (EN14), RFN15, and
$mpD_u^2N_0^2$, and the pure testing factor $D_u$.
This proves the full local parent for all infinitesimal $u,v>0$
and arbitrary $z>0$, with both unrestricted state normals.
The exact SJ state/sign reduction supplies every state and the
literal axes in this local chart.

\subsubsection*{Every other testing angle, with weak \(m\) retained}

This outer join uses only (EN9), positivity of the six native
weights, and exact full-parent identities. Set
$\widetilde m=m/(1+m)$, a bounded positive source coefficient below
$m$, and retain the difference as the whole original $m$ square.
Keep the entire $n$ and $\alpha$ events throughout.
For the single unbounded coefficient $C$ or $D$, retain its excess
over a sufficiently large positive real cap. Other selected
coefficients have the actual bounded positive residues (EN9).

\paragraph{Finite Bob with nonzero residue and bounded Alice.}
Let $v$ have nonzero finite residue, $u$ be bounded, and use the
five full-state rows
\[
P=uW,\quad Y,\quad X,\quad Z,\quad S=\sqrt{\widetilde m}\,W.
                                                               \tag{EN21}
\]
The selected $m$ event is the entire square
$[vS-\sqrt{\widetilde m}(X-uvY+uZ)]^2$.
All other terms use $P,Y,X,Z$ with bounded coefficients.
In R3 the strong row is $P-Y$; on its kernel the selected remaining
form on $(Y,X,Z)$ is
\[
\frac T2D_uD_vY^2+\frac{1+T}{2T}v^2(uX-Z)^2+K_0 .
\]
In R4 the strong row is $P+uvX-Y-vZ$; on its kernel it is
\[
\frac{D_vv^2}{2T}(Z-uX)^2+
\frac{T}{2(1-T)}D_uD_vY^2+K_0 ,
\quad
K_0=\frac{D_u(vY-Z)^2+D_v(X+uZ)^2-D_uD_vZ^2}{2}.
\]
The leading principal minors in the order $(Y,X,Z)$ are
\[
\begin{array}{c|c|c}
 &R3&R4\\ \hline
1&
\dfrac{D_u(TD_v+v^2)}2&
\dfrac{D_u(T+v^2)}{2(1-T)}\\[4pt]
2&
\dfrac{D_u(TD_v+v^2)(Tu^2v^2+Tv^2+T+u^2v^2)}{4T}&
\dfrac{D_uD_v(T+v^2)(T+u^2v^2)}{4T(1-T)}\\[4pt]
3&
\dfrac{v^2(1+T)D_u^2D_v(Tu^2+v^2)}{8T}&
\dfrac{v^2D_u^2D_v^2(Tu^2+v^2)}{8T(1-T)} .
\end{array}                                                    \tag{EN22}
\]
All are strictly positive at the real residues, even $u=0$.
The $S$ direction has the positive coefficient $v^2$ in the
selected $m$ event. A sufficiently large real cap on the strong
row makes the entire five-row residue positive definite.
This is the real finite-cap Schur lemma: if a bounded real
quadratic form is positive on the kernel of a linear map, adding
a sufficiently large real multiple of the squared map makes it
positive definite. The unbounded actual coefficient permits
every such cap. This argument includes $m$ infinitesimal.

\paragraph{Infinitesimal Bob and Alice a finite positive unit.}
Use $R=uW-Y$, $B_t=X+uZ$, $V_Z=vZ$, and $Y$; the testing inverse
$1/u$ is bounded. The full selected parent is a bounded quadratic
on these four rows because
$v(W+uY)=(v/u)(R+D_uY)$ and
$J=R+uvB_t-D_uV_Z$.
Its residue splits off $(\widetilde m+1/2)B_t^2$ and the form
\[
CR^2+AD_uY^2+(D+1/2)(R-D_uV_Z)^2
-D_uYV_Z-\tfrac12D_uV_Z^2.                            \tag{EN23}
\]
In R3 restrict to the strong-row kernel $R=0$.
After removing a positive factor $D_u$, the matrix on $(Y,V_Z)$ is
$\left(\begin{smallmatrix}T/2&-1/2\\-1/2&
(1+T)D_u/(2T)-1/2\end{smallmatrix}\right)$;
its determinant is $(1+T)u^2/4>0$.
In R4 the strong kernel is $R=D_uV_Z$.
On $(R,Y)$ its first diagonal is $1/(2T)-1/(2D_u)>0$,
and its determinant is $u^2/[4(1-T)]>0$.
A large real cap again makes (EN23) strict.

\paragraph{Infinite Bob and bounded Alice.}
Put $w=1/v$, divide the homogeneous parent by $v^2$, and use
the five rows (EN21). Select any fixed positive real lower caps
$c_0,a_0,d_0$ below $C,A,D$, and the actual $\widetilde m$ event.
The full quadratic is bounded; its residue is
\[
[S+\sqrt{\widetilde m}\,uY]^2+
c_0(P-Y)^2+(a_0+\tfrac12)D_uY^2+
d_0(uX-Z)^2+\tfrac12D_uX^2.                           \tag{EN24}
\]
This is positive definite on the five rows, for every bounded
residue of $u$. It includes $u=0$ and $m$ infinitesimal.

\paragraph{Infinite Alice.}
Let $s=1/u$ and keep the entire $m,n,\alpha$ events.
For finite $v$ of nonzero residue, after division by $u^2$
the selected residue on $(W,Y,X,Z)$ is
\[
c_0D_vW^2+a_0D_vY^2+
(d_0+\tfrac12)(W+vX)^2+\tfrac12(vY-Z)^2,
\]
which is positive definite. For $v$ infinitesimal use instead
$(N,W,Y,V)=(sX+Z,W,Y,vX)$.
The complete signed core after division by $u^2$ is
\[
\frac12\{W+(1+s^2)V-sY-svN\}^2
+\frac12(1+s^2)v^2Y^2
-(1+s^2)Y(vN-sV)-\frac12(1+s^2)(vN-sV)^2
+\frac12D_vN^2 .
\]
Adding the selected $C,A,D$ packets gives strict residue
$c_0W^2+a_0Y^2+(d_0+1/2)(W+V)^2+N^2/2$.
Every full coefficient is bounded; no $s/v$ is used.
If also Bob is infinite, divide by $u^2v^2$ and use
$(W,X,Y,S_Z,T_Z)=(W,X,Y,sZ,wZ)$.
The complete selected residue is
\[
c_0W^2+a_0Y^2+d_0(X-S_Z)^2+\tfrac12X^2
+\tfrac12(Y-T_Z)^2,                                   \tag{EN25}
\]
again positive definite. The full quadratic follows directly from
$J/(uv)=wW+X-swY-S_Z$ and the exact signed core;
its coefficients are bounded polynomials in $s,w$.

These charts exhaust all nonnegative projective angles:
both infinitesimal are the complete local theorem, and every other
pair is one of (EN21)--(EN25) or (EN23).
All strict-matrix coefficients use only bounded pure source or
pure testing factors. Restore every selected excess, the untouched
$m,n,\alpha$ packets, and the homogeneous testing squares.
The exact SJ sign gauges, literal projective endpoints and
real/imaginary state split finish the asserted all-tests theorem.

\paragraph{Audit scope and dependencies.}
The source dictionary and exact scalar/normal identities are the
freshly reconstructed RSR, NSR, RFN sources; only their explicitly
rechecked algebraic identities and finite SOS budgets are used.
Equations (EN3)--(EN20) prove all changed source and coefficient
bounds afresh. The entire outer join (EN21)--(EN25) is new.
The focused checker corroborates source coefficient limits,
full-parent expressions, complete hard entries, and every printed
positive kernel matrix. Analytical uniform estimates and the
exhaustive join are the proof above, not an inference from its receipt.


\section{\texorpdfstring{The complete ordered weak-C2 shape R5 is admitted}{The complete ordered weak-C2 shape R5 is admitted}}
\label{component:p8da9ea6ed9e5}
% Scoped mathematical certificate; see local_identity_audit.closure.json for exact changes.

\label{p8da9ea6ed9e5:sec:r17-r5}

\begin{theorem}
Normalize $d=1$. Retain the original physical C2 source, its stationary contact, every strict
source gap, all coefficient signs, the positive-infinitesimal total-scale
intersection and every original prior-admission exclusion. Suppose
$B,z,y,r>0$ are infinitesimal and $t=y/B$, $h=H/B$, $l=1-h=L/B$,
and $j=1+zr-l/t$ are positive bounded units. Then the entire original
four-state parent is a finite tensor SOS at every independently chosen
projective testing pair, with both signs, literal zero/infinite angles,
and real or complex states. These are precisely the unchanged R5 source
hypotheses. There is no comparison of source and testing rates.
\end{theorem}

\paragraph{Source entry and all gaps.}
Put $a=z/B$ and $g=G/B=1-a(1-r-H)$. Since $r,H$ are infinitesimal,
$g>0$ implies $a$ bounded and $a_0\in[0,1]$. Write $\rho=a_0$,
$\chi=h_0\in(0,1)$, $\tau=t_0>1-\chi$ and $\ell=1-\chi$,
$\delta=\tau-\ell>0$. The inequality for $\tau$ follows from $j$
being a positive unit. The exact source formulas WE1 give jointly
\[
\begin{gathered}
 (nr)_0={1\over2(1+\rho)},\quad
 A_0={\tau\over2\chi(1+\rho)},\quad
 C_0={\chi+\tau-\rho\ell\over2\ell(1+\rho)},\\
 D_0={1\over2(1+\rho)\delta},\quad
 (\alpha/B^2)_0={\chi+\tau-\rho\ell\over2(1+\rho)}.
\end{gathered}       \tag{R5.1}
\]
All displayed quantities are positive units; $n$ is unbounded and
$\alpha$ infinitesimal. Thus literal $n=0$ is impossible on this shape,
by the exact identity $r(2n+1)\asymp1$, rather than being discarded.
If $\rho<1$, $\mu=m/a^2$ is bounded with residue
$1/[2(1-\rho^2)]$. If $\rho=1$, $g$ is infinitesimal,
$m$ unbounded and $(mg)_0=1/4$. These statements retain arbitrary
independent $r,g,B$ openings. In particular the source is automatically
in the positive-infinitesimal total-scale regime. No nonemptiness of
the residual intersection is asserted.

\paragraph{An exact centered-excess identity for the entire Gram.}
Use the unnormalized version of the established complete WA3--WA8 identity,
with roots
\[
 \xi=X+ruZ,\quad\eta=Y-huW,\quad
 R_1=uW-t^{-1}vZ,\quad R_2=vW-a^{-1}uZ,\quad T=uvW-BZ.
\]
Let $G$ be its five-by-five matrix, so $F/4={\cal R}^TG{\cal R}$.
Both complete normal rows and the complete hard block are included.
Write the original $m$ event as $M=vW-X+uvY-uZ$. Its centered row is
\[
 R_{mc}=M-gvW=-\xi+uv\eta+(1-g)R_2+huT.             \tag{R5.2}
\]
The equality uses $1-g=a(1-r-H)$ and $H=hB$. It vanishes at the
actual contact. Put $L_m=(-1,uv,0,1-g,hu)^T$.
For any positive real $K<m$, there is the exact matrix identity
\[
 G=G\big|_{m=K}+(m-K)L_mL_m^T+(m-K)g\,H_m,          \tag{R5.3}
\]
where the only nonzero entries of $H_m$ are
\[
 (H_m)_{14}=(H_m)_{41}=-1,\quad
 (H_m)_{25}=(H_m)_{52}=v,\quad
 (H_m)_{44}=2-g,\quad
 (H_m)_{45}=(H_m)_{54}=hu.                           \tag{R5.4}
\]
All other actual source coefficients, especially $\alpha/B^2$, remain
fixed in $G|_{m=K}$. This is a formal matrix substitution, not an assertion
that a modified source is physical or satisfies contact. Equivalently,
at the original event level,
\[
 mM^2-(m-K)R_{mc}^2
 =KM^2+2(m-K)g\,vWM-(m-K)g^2v^2W^2.                \tag{R5.5}
\]
The retained square contains every radial and normal jet of the excess.
Equations R5.3--R5.4 prove boundedness of the remainder directly; no
Hessian-subtraction premise is imported.

\paragraph{All infinitesimal or zero testing angles.}
If $\rho\in(0,1)$, apply the pure-source congruence replacing $\xi$ by
$\widehat\xi=\sqrt n\xi$. Every coefficient is bounded: $nr$ is
bounded, $nr^2$ infinitesimal, $a,t$ bounded units, and $A,C,D,m$ and
$\alpha/B^2$ bounded. At $u=v=0$ the full residue consists of a unit
$\widehat\xi$ pivot, a positive $m_0$ pivot on $R_2$, a positive
$(\alpha/B^2)_0$ pivot on $T$, and the $\eta,R_1$ block
\[
 \begin{pmatrix}
 A_0+C_0+D_0+1/2&-(C_0+D_0+1/2)(1-\chi)+A_0\chi\\
 -(C_0+D_0+1/2)(1-\chi)+A_0\chi&
 (C_0+D_0+1/2)(1-\chi)^2+A_0\chi^2
 \end{pmatrix}.                                    \tag{R5.6}
\]
Its determinant is $A_0(C_0+D_0+1/2)>0$.

If $\rho=0$, replace $R_2$ by $aR_2$ as in WE4--WE5.
The exact divisibility $(b_2-\alpha/B^2)/a$ from WE3 holds for arbitrary
$h$ and $r$; its canceled denominator is a unit at $a=r=B=0$ with
$h,l,t,j$ units. Hence every coefficient remains bounded, including
all full normal entries. The $R_2$ pivot becomes $\mu_0=1/2$;
R5.6 and the other pivots remain positive. No $h/a$ bound is used.

If $\rho=1$, retain $(m-K)R_{mc}^2$ using R5.3, for any fixed
positive real $K$. Since $a$ is a unit, all coefficients in the
remainder are bounded after $\widehat\xi=\sqrt n\xi$.
The constant $\widehat\xi,R_2$ coupling is infinitesimal, while
its $R_2$ pivot has residue $K+2(mg)_0=K+1/2$.
All other zero-angle pivots are R5.6, $1$, and $(\alpha/B^2)_0$.
Thus the complete remainder has strict residue in every source case.
The finite separated Gram lift WA10 applies for independent infinitesimal
$u,v$, including zeros and either sign. Adding the retained exact square
proves the complete local parent without any state restriction.

\paragraph{The finite remote determinant when $\rho<1$.}
For bounded tests, normalize $\widehat X=\sqrt nX$ in the original
parent. Its hard residue is $(1+U)(1+V)\widehat X^2$, and the remaining
three-state form on $(W,Y,Z)$ is
\[
\begin{split}
 P={}&m_0(vW+uvY-uZ)^2+C_0(1+V)(uW-Y)^2
   +A_0(1+U)(1+V)Y^2+D_0(uW-Y-vZ)^2\\
 &+\tfrac12\{(uW-Y-vZ)^2+(1+U)(vY-Z)^2
        +(1+V)U Z^2-(1+U)(1+V)Z^2\},                
\end{split}\tag{R5.7}
\]
where here $u,v$ denote the real testing residues and $U=u^2,V=v^2$.
Set
$\mathcal D=8\chi(1-\rho)(1+\rho)^3\ell\delta>0$.
The exact determinant of its symmetric matrix is
\[
 { (1+U)(1+V)\over\mathcal D}
 \left[(\chi+\tau)(\tau U-\rho V)^2
 +(\chi+\tau-\rho\ell)
 \{\tau\delta U^2V+[\chi+(1-\rho)\tau]UV^2\}\right]. \tag{R5.8}
\]
For $U,V>0$ it is strictly positive. For $U>0,V=0$ it is strictly
positive; for $U=0,V>0$ it is strictly positive when $\rho>0$.
The restriction of R5.7 to $Z=0$ is strictly positive in these cases:
its signed core there is
$\tfrac12[(uW-Y)^2+(1+U)VY^2]$, and the $A,C,m$ packets supply
the stated two-state rank. This restriction and R5.8 prove strict
positive definiteness of the entire three-state residue. All source
coefficients are bounded, so WA10 around the real testing residues
proves these complete remote charts.

\paragraph{Finite remote tests with the unbounded m opening.}
Let $\rho=1$ and extract the exact square in R5.5 before normalizing
$\widehat X=\sqrt nX$. Its remainder has bounded coefficients.
At a real testing pair, its three-state residue is
\[
 B_{\rm fin}+K M_0^2+2\kappa vW M_0,\qquad
 M_0=vW+uvY-uZ,\quad\kappa=1/4,                     \tag{R5.9}
\]
where $B_{\rm fin}$ is R5.7 with the $m_0$ packet omitted and
$A_0,C_0,D_0$ evaluated at $\rho=1$.
For $U,V>0$, restrict to $M_0=0$, so $W=uZ/v-uY$.
Its two-state determinant equals $V^{-1}\lim_{\rho\uparrow1}
\det(P)/m_0$. Formula R5.8 gives
\[
 { (1+U)(1+V)\over16\chi\ell\delta V}
 \left[(\chi+\tau)(\tau U-V)^2
 +(\chi+\tau-\ell)\{\tau\delta U^2V+\chi UV^2\}\right]>0.\tag{R5.10}
\]
At $Z=0$ on that kernel, $W=-uY$ and the positive $A,C$ terms,
together with the nonnegative signed-core restriction, give a positive
$Y^2$ coefficient. Hence $B_{\rm fin}|_{M_0=0}$ is positive definite.
When $u=0,v\ne0$, the kernel is $W=0$ and its form is
$(A_0+C_0+1/2)(1+V)Y^2+D_0(Y+vZ)^2\succ0$.
When $v=0,u\ne0$, the kernel is $Z=0$ and its form is
$(C_0+D_0+1/2)(uW-Y)^2+A_0(1+U)Y^2\succ0$.
For each such real test pair, elementary real completion on the
transverse coordinate $M_0$ now gives a sufficiently large positive
real $K$ for which R5.9 is strictly positive definite. This choice is
legal: $m$ is unbounded, so $m-K>0$ for every real $K$; $K$ never
compares a source rate with a testing rate. Use that $K$ in the exact
full-parent identity R5.5 and apply WA10 to the bounded remainder.

\paragraph{The small-a first-axis chart.}
The only finite remote case left is $\rho=0$, $u_0=0$, $v_0\ne0$.
Use the five roots
\[
 S=aW,\quad R=uW,\quad \widehat X=\sqrt nX,\quad Y,\quad Z
\]. Write the original $m$ packet as
$\mu(vS-aX+auvY-auZ)^2$ and replace $uW$ everywhere else by $R$.
This is exactly the complete identity WE8 with unnormalized $Y$ and
$A(1+u^2)(1+v^2)Y^2$ retained. All coefficients are bounded.
The residue has separate pivots $\mu_0v_0^2$ on $S$ and
$1+v_0^2$ on $\widehat X$. Its remaining $(R,Y,Z)$ form is
\[
 C_0(1+V)(R-Y)^2+A_0(1+V)Y^2+D_0(R-Y-vZ)^2
 +\tfrac12\{(R-Y)^2+VY^2-2vRZ\}.                  \tag{R5.11}
\]
Its restriction to $Z=0$ is positive definite and its determinant is
\[
 {V^2(\chi+\tau)^2(1+V)\over8\chi\ell\delta}>0.    \tag{R5.12}
\]
This supplies the required strict full residue at the literal axis too.

\paragraph{Every reciprocal chart and literal infinity.}
For unbounded $u$ put $p=1/u$ and use the polynomial parent $p^2F$.
If $v$ is bounded, the $\widehat X$ residue is $1+V$ and the other
three-state residue is exactly
\[
 [C_0(1+V)+D_0+1/2]W^2+A_0(1+V)Y^2
                         +(m_*+1/2)(vY-Z)^2,         \tag{R5.13}
\]
which is strict for every real $v$. Here $m_*=m_0$ when bounded,
and $m_*=K$ after centered extraction if $g$ is infinitesimal.
The additional centered correction in R5.5 has a factor $p$ and
vanishes in residue. This covers $p=0$ literally.

For unbounded $v$, write $q=1/v$ and use $q^2F$. If $u$ is bounded,
the $\widehat X$ pivot is $1+U$ and the other residue is
\[
 m_*(W+uY)^2+C_0(uW-Y)^2+(A_0+1/2)(1+U)Y^2+D_0Z^2.\tag{R5.14}
\]
For $\rho<1$ this is strict unless $\rho=u_0=0$. For $\rho=1$,
the bounded correction is $2\kappa W(W+uY)$; its restriction to
$W+uY=0$ is strict, and a sufficiently large real $K$ again makes
the entire residue strict. If $\rho=u_0=0$, use $S=aW,R=uW$,
$\widehat X,Y,Z$ as above in the complete WE10 identity. Its residue is
\[
 \widehat X^2+\mu_0S^2+C_0(R-Y)^2+(A_0+1/2)Y^2+D_0Z^2\succ0.
                                                               \tag{R5.15}
\]
Every coefficient of this entire polynomial chart is bounded.

Finally, if both angles are unbounded, use $p=1/u,q=1/v$ and
$p^2q^2F$, with the five roots $\widehat X,W,Y,P=pZ,Q=qZ$.
The exact signed-core part after homogeneity is one half of
\[
 (qW+X-pqY-P)^2+(1+p^2)(Y-Q)^2
                    +(1+q^2)(p^2X^2+2XP-P^2).       \tag{R5.16}
\]
All other original packets are retained in their homogeneous forms;
in particular the $m$ packet is $(pW-pqX+Y-Q)^2$ and its centered
row subtracts $gpW$. After $\widehat X=\sqrt nX$ and any needed
centered extraction, the complete residue is
\[
 \widehat X^2+C_0W^2+A_0Y^2+D_0P^2
                                  +(m_*+1/2)(Y-Q)^2\succ0.       \tag{R5.17}
\]
Its coefficients are bounded; the centered correction again contains
$p$ and vanishes in residue. Both $p=0$ and $q=0$ are literal
polynomial values. The roots $P,Q$ need not be independent to use
this positive Gram.

\paragraph{Join and tensor legality.}
The local chart, finite remote cases, the one small-a finite axis,
each single reciprocal chart and the double reciprocal chart exhaust
all independently chosen testing types. Every use of WA10 whitens a
strict pure-source matrix at real testing residues and splits finitely
many bounded-source times infinitesimal-testing coefficients into
positive pure-source/pure-testing products. There is no mixed inverse
or radical. The roots themselves use only pure-source radicals such
as $\sqrt n$ and pure-testing monomials. Exact positive excess squares
are restored in every chart. Away from literal endpoints, homogeneous
restoration uses only pure-testing squares; at the endpoints the
polynomial value is used. The real identities apply separately to real
and imaginary state parts. All four original state coordinates, both
normal directions, all hard couplings, arbitrary source openings and
the unchanged physicality/total-scale/prior-exclusion intersection
have therefore been retained. This is an all-tests admission of R5.


\section{\texorpdfstring{R6 with bounded z/B is also an all-tests admission}{R6 with bounded z/B is also an all-tests admission}}
\label{component:pd5fa361b4623}
% Scoped mathematical certificate; see local_identity_audit.closure.json for exact changes.

\label{pd5fa361b4623:sec:r17-r6-bounded-a}
\begin{theorem}
Retain every original R6 hypothesis and normalize $d=1$: $B,z,y>0$
infinitesimal; $t,h,l,j,r$ positive bounded units; $g$ infinitesimal;
the unchanged physicality, contact, coefficient signs, total scale and
prior exclusions. If $a=z/B$ is bounded, then the entire original
parent is finite tensor SOS at all independent projective tests,
including literal axes/infinities, either sign, and complex states.
This admits the whole stated subdomain. The remaining R6 sources have
$a$ unbounded and $1-r$ infinitesimal; no claim about that endpoint is
made here.
\end{theorem}

Write $e=1-r$. Since $1-g=a(e-H)$, the added condition is equivalent
to $e$ being a positive unit; $a_0=1/e_0$. Denote the real residues
by $e\in(0,1)$, $h\in(0,1)$, $l=1-h$, $t>l$, and $\delta=t-l$.
The exact source formulas, with no rate comparison, give
\[
\begin{gathered}
 n_0={e^2\over2(1-e^2)},\quad
 A_0={t\over2h(1+e)},\quad
 C_0={h+t-el\over2l(1+e)},\\
 D_0={1\over2(1+e)\delta},\quad
 (\alpha/B^2)_0={h+t-el\over2e(1+e)},\quad
 \kappa=(mg)_0={1\over2(1+e)}.
\end{gathered}                       \tag{R6.1}
\]
Every listed quantity is a positive unit; $m$ is unbounded. In
particular $n$ cannot be literally zero on this admitted subdomain.
The unbounded-$a$ endpoint, where $n=0$ may occur, is not erased.

Use the exact full-parent excess identity R5.2--R5.5 with a positive
real cap $K$, retaining $(m-K)R_{mc}^2$. The remainder has bounded
coefficients in the unchanged five WA roots, because $a,t$ and all
remaining event weights are bounded and $mg$ bounded. At zero testing
angles its $\eta,R_1$ block is R5.6 and its radial $T$ pivot is
$(\alpha/B^2)_0$. Its $\xi,R_2$ block is
\[
 \begin{pmatrix}n_0+1/2+K&-K-\kappa\\
                 -K-\kappa&K+2\kappa\end{pmatrix},
\quad\det=K(n_0+1/2)+2\kappa(n_0+1/2)-\kappa^2.     \tag{R6.2}
\]
Since $n_0+1/2=\kappa/(1-e)$, its determinant is positive for every
$K>0$. Thus the complete local remainder has strict residue, and WA10
proves all infinitesimal or zero tests, keeping all hard and normal terms.

For a finite remote pair, use the original four states. After centered
extraction the bounded remainder is, in residue,
\[
 B_{\rm fin}+K M^2+2\kappa vWM,\qquad
 M=vW-X+uvY-uZ,                                      \tag{R6.3}
\]
where $B_{\rm fin}$ is the complete original parent at the source
residues R6.1, with $m$ and $\alpha$ omitted. The restriction to
$M=0$ is obtained by the exact substitution $X=vW+uvY-uZ$.
On $(W,Y,Z)$ its determinant is
\[
 { (1+U)^2(1+V)^2\over\mathcal D}
 \left[(h+t)(etU-V)^2+(h+t-el)
          \{t(h+t-e)U^2V+hUV^2\}\right],             \tag{R6.4}
\]
where $\mathcal D=8h(1-e)(1+e)^3l\delta>0$.
Every coefficient outside the displayed square is strictly positive:
$h+t>1>e$ follows from $t>l$. Hence the determinant is positive
whenever $(U,V)\ne(0,0)$, including both axes. The restriction to
$Z=0$ is positive definite: the original signed core there is the sum
of the three squares
$[J^2+(1+U)VY^2+(1+V)X^2]/2$, and the positive $A,C,n$ packets
force $Y=W=0$ on $M=Z=0$ at every such remote pair.
Thus $B_{\rm fin}|_{M=0}$ is positive definite. A sufficiently large
positive real $K$ makes R6.3 strictly positive definite on the full
four-state space. Use this $K$ in the exact source identity and apply
WA10 around the real testing residues. All coefficients remain bounded;
there is no mixed division or source/test comparison.

For unbounded $u$ and bounded $v$, use the full polynomial chart
$p^2F$, $p=1/u$. Its excess correction contains $p$ and disappears in
residue. The centered $m$ row has residue $vY-Z$. On its kernel the
complete finite remainder on $(X,W,Y)$ is
\[
 n_0(1+V)X^2+C_0(1+V)W^2+A_0(1+V)Y^2
                                      +(D_0+1/2)(W+vX)^2\succ0.\tag{R6.5}
\]
A positive cap gives a strict full residue. For unbounded $v$ and bounded
$u$, use $q^2F$, $q=1/v$. The $m$ row is $W+uY$ and the bounded
correction is $2\kappa W(W+uY)$. On its kernel, the complete residue is
\[
 (n_0+1/2)(1+U)X^2+C_0(1+U)^2Y^2
             +(A_0+1/2)(1+U)Y^2+D_0(uX-Z)^2\succ0.  \tag{R6.6}
\]
Again a sufficiently large real $K$ gives a strict full residue.
These assertions include every finite real value of the other angle,
with either sign or zero.

For both angles unbounded use the exact double homogeneous parent and
core R5.16, now in the five roots $(X,W,Y,P=pZ,Q=qZ)$, without normalizing
$X$. Its complete residue after centered extraction is
\[
 (n_0+1/2)X^2+C_0W^2+A_0Y^2+D_0(X-P)^2
                                  +(K+1/2)(Y-Q)^2\succ0.        \tag{R6.7}
\]
The $X,P$ block has determinant $D_0(n_0+1/2)>0$ and the $Y,Q$ block
has determinant $A_0(K+1/2)>0$. The excess correction has a factor $p$
and disappears in residue. All homogeneous coefficients are bounded.
Literal $p=0,q=0$ are the polynomial chart values, not limiting divisions.

These cases exhaust all tests on the stated source subdomain. The
separated finite lift, exact excess restoration, pure-testing homogeneous
restoration and real/imaginary split are exactly those established in
the R5 join. They retain every original normal and hard block, all signs,
all independent rates and the original source/prior-exclusion intersection.


\section{\texorpdfstring{The bounded-z/B subdomain of R7 is admitted for all tests}{The bounded-z/B subdomain of R7 is admitted for all tests}}
\label{component:p9a223c96c696}
% Scoped mathematical certificate; see local_identity_audit.closure.json for exact changes.

\label{p9a223c96c696:sec:r17-r7-bounded-a}
\begin{theorem}
Normalize $d=1$ and retain every original R7 condition: $B,z,y>0$
infinitesimal; $t=y/B$ a positive bounded unit; $h=H/B$ and $g=G/B$
infinitesimal; the exclusion of simultaneously infinitesimal $r,j$;
original physicality, stationary contact, coefficient signs, strict gaps,
total-scale intersection and prior exclusions. If $a=z/B$ is bounded,
the entire original parent is finite tensor SOS at every independent
projective testing pair. Both signs, literal zero/infinite angles and
complex states are included.
\end{theorem}
The complementary unbounded-$a$ R7 source, including its physical
$n=0$ face, is proved in Section~\ref{component:p405d96527133}.
The present theorem and proof retain their bounded-$a$ scope.

\paragraph{Complete source entry.}
Write $e=1-r$ and $l=1-h$. Since $1-g=a(e-H)$, bounded $a$ is equivalent
on this source shape to $e$ being a positive unit; $a_0=1/e_0$.
We use the letters $e\in(0,1]$ and $t\ge1$ for their real residues.
The strict gap $j=1+zr-l/t$ gives $j_0=1-1/t$. The ordered R7 exclusion
says that $e=1$ and $t=1$ cannot hold together. The original source formulas
give jointly
\[
 C_0={t-e\over2(1+e)}>0,\quad
 (Ah)_0={t\over2(1+e)}>0,\quad
 (\alpha/B^2)_0={t-e\over2e(1+e)}>0,\quad
 \kappa=(mg)_0={1\over2(1+e)}.                        \tag{R7.1}
\]
Thus $A,m$ are unbounded, $C$ and $\alpha/B^2$ are positive units,
and $\alpha$ is infinitesimal. If $e<1$, then
\[
 n_0={e^2\over2(1-e^2)}>0;
\]
if $e=1$, $r$ is infinitesimal and $n$ is unbounded, with $(nr)_0=1/4$.
If $t>1$, then $D_0=1/[2(1+e)(t-1)]$ is a positive unit. If $t=1$,
then $j$ is infinitesimal and $D$ unbounded with $(Dj)_0=\kappa$;
in this case $e<1$ by the retained exclusion. These are exhaustive
source alternatives, with no rate restriction. They also show that
literal $n=0$ is impossible on this admitted subdomain and that total
scale is positive infinitesimal. The $n=0$ boundary on the unbounded-$a$
complement is not removed by this statement.

\paragraph{A complete D-event gauge, including its exact zero syzygy.}
We use the full WA parent, now in the roots
\[
\begin{gathered}
 {\cal R}=(\xi,\eta,R_1,R'_2,T)^T,
 \quad\xi=X+ruZ,\quad\eta=Y-huW,\\
 R_1=uW-vZ/t,\quad R'_2=avW-uZ,\quad T=uvW-BZ.
\end{gathered}
\]
Let $G$ be the exact full-parent matrix obtained from WA8 by leaving
$\eta$ unnormalized and scaling its fourth root by $a$.
The following gauge is the exact D identity developed in the parallel
R8 work; it is valid before any shape specialization. Set
\[
\begin{split}
 d_D&=(uv,-1,l,ruv,-arv)^T,\\
 S_D&=(0,0,(aB-U)v,-(B-V/t)u,U-aV/t)^T,\\
 L_D&=(-2rtu,0,0,-r^2tu,ar^2t)^T.
\end{split}
\]
Then $S_D^T{\cal R}=0$ identically, and
\[
 d_D^T{\cal R}=J+jvZ=:R_{Dc},\qquad
 J=uW+uvX-Y-vZ.                                     \tag{R7.2}
\]
In particular this is the actual contact-zero D row, not an uncentered
positive packet. The exact matrix identity is
\[
 \partial_DG=d_Dd_D^T+\tfrac12(S_DL_D^T+L_DS_D^T)+jH_D,\tag{R7.3}
\]
where the only nonzero entries of $H_D$ are
\[
 (H_D)_{13}=tuv,\quad(H_D)_{15}=-tu,\quad
 (H_D)_{34}=rtuv,\quad(H_D)_{35}=-artv
\]
and their symmetric partners. Consequently the full original parent is
\[
 F/4=D R_{Dc}^2+{\cal R}^T(G|_{D=0}+DjH_D){\cal R}. \tag{R7.4}
\]
The notation $G|_{D=0}$ keeps all other actual coefficients, including
$\alpha/B^2$, unchanged; no modified physical source is asserted.
All full first normals and hard entries are included in R7.3.
This gauge has been independently checked here against the entire WA
matrix, its zero syzygy and the original four-state D event.

\paragraph{One strict local Gram for all source alternatives.}
First apply R7.4, and then retain the centered m excess from R5.2--R5.4,
with positive real cap $K$. The m identity is unaffected by the D gauge:
$H_D$ and its syzygy are independent of $m$ and the D row is independent
of $m$. The remaining full Gram has bounded coefficients after
\[
 \widehat\eta=\sqrt A\,\eta,
\]
and, if $r$ is infinitesimal, also $\widehat\xi=\sqrt n\,\xi$.
Indeed $Ah,nr,mg,Dj$ are bounded, $a,t,l^{-1}$ are bounded units,
$Ah^2$ is infinitesimal, and every large m or D coefficient has been
removed in its complete contact-zero square. No ratio between $A,n,m,D$
is used. This bounds the entire scalar block, both full normal rows and
the hard block, not only the graph.

At $u=v=0$, all entries of $H_D$ vanish. The $\widehat\eta$ pivot is
$1$, its couplings are infinitesimal, the $R_1$ pivot is $C_0+1/2>0$,
and the $T$ pivot is $(\alpha/B^2)_0>0$. If $r$ is a unit, the remaining
$\xi,R'_2$ block is the pure-source rescaling by $\operatorname{diag}(1,a_0^{-1})$
of R6.2. Its determinant is positive because
$n_0+1/2=\kappa/(1-e)$. If $r$ is infinitesimal, normalization by
$\sqrt n$ instead gives separate positive pivots $1$ and
$(K+2\kappa)/a_0^2$; the cross coefficient is infinitesimal.
Thus the complete local remainder has strict residue for every source
alternative. WA10 gives a finite separated tensor SOS for every
independent infinitesimal or zero $u,v$. Restore both retained exact
squares. This proves the full local parent, including its complete
normal/radial jets, without transferring physicality to a capped source.

\paragraph{Remote entry and a finite-dimensional cap lemma.}
For remote charts use the original states and normalize
$\widehat Y=\sqrt A Y$. If $r$ is infinitesimal, normalize also
$\widehat X=\sqrt nX$. First retain the m excess. If $D$ is unbounded,
retain a D excess too, using the exact full-event identity
\[
 DJ^2-(D-K_D)(J+jvZ)^2
 =K_DJ^2-2(D-K_D)jvZJ-(D-K_D)j^2VZ^2.               \tag{R7.5}
\]
All coefficients in the remaining full parent are bounded. The hard
normalized state pivots are positive and decouple in residue; all other
coefficients are bounded and divided by $\sqrt A$ or $\sqrt n$ in
couplings to these pivots. Thus the soft residue has $Y=0$, and also
$X=0$ when $r$ is infinitesimal.

We repeatedly use the following elementary real fact. If a real quadratic
form $B$ is positive definite on the common kernel of real linear rows
$L_1,\ldots,L_k$, then $B+K\sum L_i^2$ is positive definite for sufficiently
large positive real $K$, after any fixed bounded cross terms vanishing on
that common kernel have been included in $B$. This follows by choosing
real coordinates on the kernel and its orthogonal complement and taking
the Schur complement. Here the relevant original event coefficients are
unbounded, so every such real cap leaves positive excesses. The choice
uses only real testing residues, never a comparison of source and testing
rates. The positive hard pivots can then be included in the same strict
full-parent Gram.

\paragraph{Finite remote tests when j is a unit.}
Let $U=u^2,V=v^2$ refer now to the real testing residues, not both zero.
After $Y=0$, the centered m row has residue
$M=vW-X-uZ$. Write $B_{\rm fin}$ for the entire finite parent with its
$m$ packet removed; the bounded cap correction is $2\kappa vWM$ and
vanishes on $M=0$.

If $r$ is a unit, retain $X$ and restrict by $X=vW-uZ$. The two-state
form on $(W,Z)$ has determinant
\[
 { (1+U)(1+V)\over4(1-e)(1+e)^2(t-1)}
       \{(etU-V)^2+(t-e)^2U^2V\}>0.                 \tag{R7.6}
\]
Its $Z=0$ restriction is positive: the positive $n,C$ packets and the
nonnegative original signed core at $Z=0$ force $W=0$ at every remote
pair. Hence it is positive definite, and the cap lemma gives a strict
full residue.

If $r$ is infinitesimal, then $e=1$, $t>1$ and $X$ is another hard state.
For $v\ne0$, restrict by $W=uZ/v$. The remaining scalar is
\[
 {Z^2\over4(t-1)V}\{(tU-V)^2+(t-1)^2U^2V\}>0.      \tag{R7.7}
\]
For $v=0,u\ne0$, $M=0$ forces $Z=0$ and the scalar on $W$ is
$(C_0+D_0+1/2)UW^2>0$. These include the literal first axis and both
angle signs. Again the cap lemma supplies strict residue for the complete
bounded remainder.

\paragraph{Finite remote tests with both m and D unbounded.}
Here $t=1$, $e<1$, $r=1-e$ is a unit and $n$ is finite. With $Y=0$ the
two centered rows have residues
\[
 M=vW-X-uZ,\qquad J=uW+uvX-vZ.
\]
They have rank two at every finite remote pair. Their common kernel is
spanned in $(W,X,Z)$ by
\[
 (v(1+U),\ V-U,\ u(1+V)).                            \tag{R7.8}
\]
On this vector, the complete finite parent with both large packets
removed has value
\[
 { (1+U)(1+V)\over2(1-e)(1+e)}
       \{(eU-V)^2+(1-e)^2U^2V\}>0.                  \tag{R7.9}
\]
The bounded residue corrections are $2\kappa vWM$ and
$-2\kappa vZJ$; both vanish on that common kernel. The cap lemma therefore
gives positive real caps for a strict three-state soft residue; the hard
$\widehat Y$ pivot completes the full parent. The higher source terms
proportional to $mg^2$ and $Dj^2$ have infinitesimal coefficients and are
retained in the bounded remainder, not discarded from the exact identity.

\paragraph{Every reciprocal chart.}
For unbounded $u$, use $p=1/u$ and the entire polynomial parent $p^2F$.
The m row on $Y=0$ is $-Z$. If $D$ is bounded, its kernel $Z=0$ has
complete residue
\[
 n_0(1+V)X^2+C_0(1+V)W^2+(D_0+1/2)(W+vX)^2\succ0, \tag{R7.10}
\]
with $X$ omitted and its normalized hard pivot retained if $n$ is
unbounded. If $D$ is unbounded, its additional row is $W+vX$; on the
common kernel the finite form is
$(1+V)(n_0+C_0V)X^2>0$. In this case $n_0>0$.
The m and D bounded cross corrections contain $p$ and vanish in residue.
The cap lemma applies for every bounded real $v$, including zero.

For unbounded $v$, use $q=1/v$ and $q^2F$. The m row on $Y=0$ is $W$.
If $D$ is bounded, the complete residue on $W=0$ is
\[
 (n_0+1/2)(1+U)X^2+D_0(uX-Z)^2\succ0,               \tag{R7.11}
\]
again with $X$ replaced by its hard normalized pivot when $n$ is
unbounded. If $D$ is unbounded, the additional row is $uX-Z$ and the
common-kernel form is $(n_0+1/2)(1+U)X^2>0$.
The bounded corrections are proportional to $W^2$ and to $Z(uX-Z)$,
so vanish on the relevant common kernels. This covers every bounded
real $u$, including zero.

For both angles unbounded, use $p^2q^2F$ and the exact core R5.16 in
$(W,X,Y,P=pZ,Q=qZ)$, with $Y$ normalized and $X$ normalized if needed.
On the soft states $Y=0$ the m row is $-Q$. For bounded $D$, the finite
residue on $Q=0$ is
\[
 (n_0+1/2)X^2+C_0W^2+D_0(X-P)^2\succ0,              \tag{R7.12}
\]
with the same convention for a hard $X$. If $D$ is unbounded, its row
is $X-P$ and the common-kernel form is
$(n_0+1/2)X^2+C_0W^2>0$. The surviving D correction is proportional to
$P(X-P)$, while the m correction contains $p$; they vanish on the common
kernel. Large real caps give a strict full residue. These identities
are polynomial in $p,q$ and include literal $p=0,q=0$.

\paragraph{Exact join and scope.}
The local, finite remote and three reciprocal configurations exhaust every
independent projective testing pair on the stated source domain. Each
bounded remainder has strict residue and thus the finite separated lift
WA10, while every extracted excess is restored as its exact positive
contact-zero square. The proof uses no mixed inverse or radical; all
normalizations and homogeneous restoration factors are purely source or
purely testing. Every full normal, hard coupling, retained equality and
source opening remains in the identities. Applying the real proof to the
real and imaginary state parts completes the complex assertion. This is
an all-tests admission of bounded-$a$ R7, not an admission of its still
unbounded-$a$ endpoint or of the omitted simultaneous small-$r$, small-$j$
intersection.


\section{\texorpdfstring{The remaining R6 boundary: $a^2B$ infinitesimal}{The remaining R6 boundary: a2B infinitesimal}}
\label{component:p75b395956985}
% Scoped mathematical certificate; see local_identity_audit.closure.json for exact changes.

\label{p75b395956985:sec:r6-b-inf}

\paragraph{Source scope and relation to the preceding theorem.}
Retain the complete original weak-C2 source: $d=1$, $\lambda=\beta=0$,
$n\geq0$, the actual stationary contact, every strict gap, total scale,
and the complement of prior admissions. On R6, $t,h,l=1-h,r$ and
the original $j$ gap are positive bounded units, while $g,B,x=aB,y=tB$
are positive infinitesimals. This section treats $a$ unbounded and
$b=a^2B$ positive infinitesimal. The notation $s=1/a$, $w_s=ae$,
$B=s^2b$, $x=sb$, $e=s(1+hsb-g)$ is unchanged from (RB1).
The exact physical identity (RB2), equivalently RT2, remains in force.
It bounds $b$ on all unbounded-$a$ R6 and gives $\rho_0=1/2$ here,
where $\rho=n/e^2$. In particular this subdomain has $n>0$; the
literal $n=0$ source face is covered by the preceding $b$-unit theorem.

\begin{theorem}
On this entire original R6 source intersection with $b=a^2B$
infinitesimal, the complete original parent is a finite separated
source--testing tensor sum of squares at every independent projective
testing pair and every real or complex state. Both angle signs, all
literal projective endpoints, and all relative opening rates of
$s,b,g$ are included. No mixed inverse or mixed square root is used.
\end{theorem}

Together with the $b$-unit theorem and the inherited bounded-$a$
R6 theorem, this proves the entire R6 row. It does not by itself
admit another weak-C2 row or the global finite-lift theorem.

\paragraph{A regular source chart.}
Keep $R=1+hx$, $K=1+x$, $T=1+xr$, $k=\sqrt{x}$ and the contact
parameters $\delta=R/(K\sqrt{bt})$, $E=tK/R$ from (RB8). Define
\[
 \varepsilon=\delta^{-1},\qquad \jmath=\delta k,
 \qquad k=\varepsilon\jmath.
 \tag{RI1}
\]
Both $\varepsilon$ and $\jmath$ are positive infinitesimals. In this
section $\jmath$ denotes the contact scale, not the original source
gap called $j$. The exact source substitution is
\[
 \begin{gathered}
 x=\varepsilon^2\jmath^2,\quad K=1+x,\quad R=1+hx,\quad
 E=tK/R,\\
 b=\frac{\varepsilon^2R^2}{K^2t},\qquad
 s=\frac{t\jmath^2K^2}{R^2}.
 \end{gathered}                                                   \tag{RI2}
\]
All denominators in (RI2) are positive source units. There is no
comparison of a testing variable with either source infinitesimal.

Use precisely the complete contact-centered remainder $G$ in (RB11),
with the full coordinate inverse (RB7). Thus the two retained packets
on the contact chart are $(m-K_m)N^2$ and
$nD_uD_v(X+\gamma v\omega)^2$, where
\[
 p=(1+hx-g)/R,\quad c=l/(tT),\quad
 \gamma=rK/(sR),\quad D_u=1+u^2,\quad D_v=1+v^2.
\]
Its exact contact is now
\[
 (u,\omega,v;N,M_B)
   =(\varepsilon\jmath,\jmath,\varepsilon\jmath^2E;0,0).
 \tag{RI3}
\]
Substitution of (RI3) in the complete contact formulas
(RB7)--(RB9a) makes both normals and the retained $n$ row vanish.

\paragraph{Two exact source cancellations.}
The quantities $mg$, $n\gamma$, $n\gamma^2$ and
$\alpha/(s^3b^2)$ have canceled denominators that are units at
$s=b=g=0$. Introduce
\[
 \mu_1=mg/R,\quad \mu_2=(mg)g/R^2,\quad
 F_0=p^2/2-n\gamma^2-2n\gamma p+\mu_2,\qquad
 A_N=p-2\mu_1-2n\gamma.
 \tag{RI4}
\]
Both $F_0$ and $A_N$ are divisible by $b$ in this rational unit ring.
Here is a short exact verification, independent of a series expansion.
At $b=0$, put $q=1-g$. Direct cancellation in the original source gives
\[
 \begin{gathered}
 p=q,\quad mg=\frac{q}{2(1+s)},\quad
 n\gamma=\frac{sq}{2(1+s)},\quad
 n\gamma^2=\frac{q(1-sq)}{2(1+s)}.
 \end{gathered}                                                    \tag{RI5}
\]
Substitution makes both expressions in (RI4) identically zero.
The denominators are units before this specialization; hence the
polynomial numerators are divisible by $b$, without introducing a
pole. Their needed residues are
\[
 \frac{F_0}{b}\sim\frac{h+t}{2},\qquad
 \left.\frac{A_N}{b}\right|_{s=g=0}=0,\qquad
 \left.\frac{A_N}{bs}\right|_{g=0}\sim h+t.
 \tag{RI6}
\]
The last quotient also has a unit canceled denominator at $s=b=0$.
Every statement in (RI5)--(RI6) is separately checked from the original
source formulas. In particular the source estimate is uniform in
$g/s$, $s/b$ and all other opening ratios.

\paragraph{Complete bounded raw support, including both first normals.}
Use the separated raw monomials
\[
 \mathsf A=\varepsilon\omega v,\quad
 \mathsf B=v,\quad \mathsf C=Eu\omega,\quad
 \mathsf D=\varepsilon^2\jmath^3E,
 \qquad \mathsf x=\mathsf B-\mathsf C.
 \tag{RI7}
\]
Set $H=\alpha+1/2$, $L=AD_v+Hv^2$, and $V_0=1-puv\omega$.
At zero normals, (RB7) gives
\[
 q_B=D_uvV_0,\quad uW-Y=cD_uvV_0,\quad
 Y=u\omega-cvV_0,\quad
 J=-(1-c)D_uvV_0,
\]
and the entire scalar has the following exact allocation:
\[
 G_0=f_A\mathsf A^2+f_{BB}\mathsf B^2+
 f_{BC}\mathsf B\mathsf C+f_{CC}\mathsf C^2+
 \varrho\mathsf D^2,                               \tag{RI8}
\]
where
\[
\begin{aligned}
 f_A={}&D_u^2D_vF_0/\varepsilon^2,\\
 f_{BB}={}&[Cc^2D_u^2D_v+(D+1/2)(1-c)^2D_u^2
                  +D_uLc^2]V_0^2\\
          &+2HD_ucV_0-D_u/2-D_u^2\mu_2v^2\omega^2,\\
 f_{BC}={}&[-2D_uLcV_0-2HD_u+2n\gamma D_uD_v]/E,\\
 f_{CC}={}&[D_uL+D_uD_vn\gamma^2v^2]/E^2,\\
 \varrho={}&\alpha D_u/\mathsf D^2.
\end{aligned}                                                      \tag{RI9}
\]
These coefficients are bounded. In particular
\[
 \frac{F_0}{\varepsilon^2}=
       \frac{R^2}{K^2t}\frac{F_0}{b},\qquad
 \frac{\alpha}{\mathsf D^2}=
       \frac1t\frac{\alpha}{s^3b^2}.
 \tag{RI10}
\]
The term that would give $u^2/\varepsilon^2$ in the first raw
coefficient has instead been placed in $f_{CC}$, using the exact
identity $u^2(v\omega)^2=(v\mathsf C/E)^2$.
The final term of $f_{BB}$ likewise allocates the whole
$-D_u^2\mu_2v^4\omega^2$ contribution without dividing by
$\varepsilon$. These are identities for the full scalar.

Write the exact normal expansion as
$G=G_0+L_NN+L_MM_B+G_{\rm hard}$, with $G_{\rm hard}$ quadratic.
The two complete linear coefficients are
\[
 L_N=R_A\mathsf A+R_B\mathsf B+R_C\mathsf C,
 \qquad L_M=S_B\mathsf B+S_C\mathsf C,               \tag{RI11}
\]
where
\[
\begin{aligned}
 R_A={}&D_uA_N/\varepsilon,\\
 R_B={}&D_u(p-2n\gamma)v^2\omega\\
 &-2uv[Cc^2D_uD_v+(D+1/2)(1-c)^2D_u+Ac^2D_v]V_0\\
 &-2Hcuv(1+cv^2V_0),\\
 R_C={}&2cuvL/E,\\
 S_B={}&2\{D_u[CcD_v-(D+1/2)(1-c)]+cL\}V_0+2H,\\
 S_C={}&-2L/E.
\end{aligned}                                                      \tag{RI12}
\]
Here $A_N/\varepsilon=
\varepsilon R^2(A_N/b)/(K^2t)$ is bounded. Thus no first-normal
pole has been discarded. Every hard coefficient is bounded as well:
the only testing inverse in the full change of state is $D_u^{-1}$,
and all original source coefficients in (RB11) are bounded after the
displayed cancellation. At $u=v=0$ the complete hard matrix is exactly
\[
 \operatorname{diag}(K_m+1/2,A+C+D+1/2).             \tag{RI13}
\]

\paragraph{A bounded conductor for these complete coefficients.}
Define the contact differences
\[
 \mathsf y=\mathsf A-k\mathsf C,\quad
 \mathsf z=u\mathsf B-k\mathsf C,\quad
 \mathsf h=k\mathsf C-\mathsf D.
\]
The scaled conductor identities are
\[
\begin{aligned}
 \mathsf C(u-k)&=\mathsf z-u\mathsf x,\\
 \varepsilon\mathsf C(\omega-\jmath)
       &=\mathsf y-\varepsilon\omega\mathsf x,\\
 \mathsf C(v-\varepsilon\jmath^2E)
       &=E(\omega\mathsf z+\jmath\mathsf y-k\omega\mathsf x),\\
 \mathsf Cv(\omega-\jmath)
       &=E\omega(\omega-\jmath)\mathsf z+
         E\jmath\omega\mathsf y-Ek\omega^2\mathsf x.
\end{aligned}                                                       \tag{RI14}
\]
These are SE5--SE6 with $w=\omega$, also checked directly here.
For a monic contact division
\[
 S=(u-k)T_u+(\omega-\jmath)T_\omega+
           (v-\varepsilon\jmath^2E)T_v,
 \qquad T_\omega=\varepsilon U+vV,
\]
the exact bounded lift is
\[
\begin{aligned}
 \mathsf CS={}&\mathsf x[-uT_u-\varepsilon\omega U
                    -Ek\omega^2V-Ek\omega T_v]\\
 &+\mathsf y[U+E\jmath\omega V+E\jmath T_v]\\
 &+\mathsf z[T_u+E\omega(\omega-\jmath)V+E\omega T_v].
\end{aligned}                                                       \tag{RI15}
\]
No inverse of a small source parameter or a contact difference occurs.

We verify the hypothesis of (RI15) for the full coefficients, rather
than import the different small-$e$ source assumptions of SE8.
Clear the finite source unit denominators and $D_u^2$. They can be
chosen independent of $\omega$ and even in $u,v$. In (RI9), all
$\omega$ dependence is through $uv\omega$ and $v^2\omega^2$;
the other testing dependence is even in $u,v$. The same property
holds after subtracting any model of the form used below.

For the scalar square-lift algorithm CC4, substitute
$\mathsf A=\mathsf y+k\mathsf C$,
$\mathsf B=\mathsf x+\mathsf C$,
$\mathsf D=k\mathsf C-\mathsf h$ into (RI8).
The coefficient requiring its first conductor is
\[
 S=k^2f_A+f_{BB}+f_{BC}+f_{CC}+k^2\varrho.
\]
Its zero contact value follows by domain cancellation, as in CC4.
Divide monically in the order $u$, then $\omega$, then $v$.
After $u=k=\varepsilon\jmath$, the $\omega$ divided difference
belongs to $(\varepsilon v,v^2)$. Accordingly choose
$T_\omega=\varepsilon U+vV$ with $U,V\in(v)$.
Modulo $\varepsilon$, the remaining $v$ polynomial is even, so
$T_v\in(\varepsilon,v)$; also
$T_u|_{\varepsilon=u=v=0}=0$.

After this first conductor the coefficient of $\mathsf h$ remains
$-2k\varrho$. The sole gradient relation is
$(-2k,1,1,-1)$ on
$(\mathsf x,\mathsf y,\mathsf z,\mathsf h)$; the determinant of
the first three contact gradients is
$-\varepsilon^3\jmath^4E^2\ne0$. The adjugate cancellation of
CC5 therefore removes the contact values without putting an inverse
of this determinant into any coefficient. Subtract the contact
multiple $2k\varrho_c(-2k,1,1,-1)$.
Each of the four resulting coefficient remainders $J_i$ vanishes at
contact. Its restriction at $\varepsilon=u=v=0$ is independent of
$\omega$: for the $\mathsf x$ coefficient this follows from
$2f_{BB}+f_{BC}$; the other three restrictions are zero by the
$U,V,T_u,T_v$ properties just established and the factor $k$.
Consequently the next monic $\omega$ divided difference of every
$J_i$ belongs to $(\varepsilon,v)$. Formula (RI15) applies again.

The last scalar repair uses the exact circuit relation
\[
 k\mathsf C(\mathsf y+\mathsf z-\mathsf h-2k\mathsf x)
       =Q_1^2-\mathsf y\mathsf z,
\]
where $Q_1$ is the first of the separated radical pairings below.
This completes the same restricted scalar square lift as CC6.
For (RI11), one conductor suffices: after replacing
$\mathsf A,\mathsf B$ as above, the coefficient multiplying
$\mathsf C$, restricted at $\varepsilon=u=v=0$, is independent of
$\omega$ directly from (RI12). Its divided difference therefore
lies in $(\varepsilon,v)$ and (RI15) applies. The full scalar and
both full normals have thus been lifted over the bounded coefficient
ring. The restricted raw support $\mathsf A^2,\mathsf B^2,
\mathsf B\mathsf C,\mathsf C^2,\mathsf D^2$ is essential;
an unrestricted ideal claim is neither needed nor asserted.

The original stationary parent and the two retained contact-zero
packets give the exact scalar value and full first-jet identities
used here. The contact (RI3) is independent of $g$. Hence taking
the exact coefficient difference divisible by $g$ preserves these
identities and every raw-support condition.

\paragraph{Strict residues of the complete contact lift.}
Suppose now that $u,v,\omega$ are all positive infinitesimals, as
in the remaining contact state chart. Give
\[
 (\varepsilon,\jmath,u,\omega,v)
       \text{ weights }(1,1,2,1,3).
\]
Then $\mathsf x$ has weight three and the five circuit roots have
weight five. Every contact substitution and (RI14)--(RI15) preserves
the weights. Put $\kappa=(h+t)/(2t)$ and
\[
 \Lambda=\frac{1-s+s^2+u^2-[2l+h(h+t)]k^2}{2ht},\qquad
 \mu=\frac1{h(1+s)}.
\]
Direct exact collection of (RI9) through weight ten gives
\[
 G_0=\Lambda\mathsf x^2+\kappa T_4+O_{\rm wt}(12)
                   \quad\pmod g,                 \tag{RI16}
\]
where
\[
 T_4=\mathsf A^2+(u\mathsf B)^2+(k\mathsf C)^2+
          \mathsf D^2-4k^2\mathsf B\mathsf C.
\]
The two complete normal collections are
\[
 L_N=O_{\rm wt}(8),\qquad
 L_M=\mu\mathsf x+O_{\rm wt}(7)\quad\pmod g.
 \tag{RI17}
\]
For transparency, the only source terms needed for (RI16)--(RI17),
besides (RI6), are the following exact Taylor coefficients at $g=0$:
\[
\begin{aligned}
 A&=\frac t{2h}(1-s+s^2)+O_{\deg(s,b)}(3),\\
 C&=\frac{h+t}{2l}-\frac{t+1}{2l}s+
                      \frac{t+1}{2l}s^2+O_{\deg(s,b)}(3),\\
 D&=\frac{1-s+s^2}{2(t-l)}-
        \frac{sbt(h+t)}{2(t-l)^2}+O_{\deg(s,b)}(3),\\
 n\gamma&=\frac s2-\frac{s^2}2-\frac{sb(h+t)}2+
                         O_{\deg(s,b)}(3).
\end{aligned}                                                       \tag{RI18}
\]
Here the error denotes exact degree filtration in a rational ring
with unit denominators. Also $s=t\jmath^2+O_{\rm wt}(6)$,
$c=(l/t)(1-k^2)+O_{\rm wt}(6)$ and
$E=t(1+lk^2)+O_{\rm wt}(8)$. The weight eight bound for $L_N$
follows term by term from (RI6), (RI12).
At $g=0$, after the displayed unit denominators are cleared, all
scalar weighted monomials have even weight, and all $M_B$ linear
monomials have odd weight. Hence the next possible degrees are
twelve and seven, respectively.

The stronger coefficient-wise statement also holds for the displayed
raw representation: $f_A-\kappa$ and $\varrho-\kappa$ have weight at
least two, while
\[
 f_{BB}-\Lambda-\kappa u^2,\qquad
 f_{BC}+2\Lambda+4\kappa k^2,\qquad
 f_{CC}-\Lambda-\kappa k^2
\]
each has weight at least six. Similarly $S_B-\mu,S_C+\mu$ have
weight at least four separately, and $(R_A,R_B,R_C)$ have weights
at least $(3,5,5)$. Thus the chosen raw allocation of each exact
difference already has the claimed filtration; no low-weight
Gram coefficients are hidden by cancellations between raw products.

Apply the bounded conductor just proved to the exact differences in
(RI16)--(RI17). The scalar correction Gram coefficients have weight
at least two; the $M_B$ normal corrections have weight at least two;
the $N$ normal corrections have weight at least three. They are all
infinitesimal. The coefficients divisible by $g$ lift boundedly by the
same argument and are infinitesimal independently of every other
opening. Restoring the finite unit denominators preserves this fact.

For strictness use the two pairings on the separated monomials
\[
 a_i^2=(\mathsf A,u\mathsf B,k\mathsf C,\mathsf D),\qquad
 Q_1=a_1a_2-a_3a_4,\quad Q_2=a_1a_3-a_2a_4.
\]
Their product relation gives the exact two-pairing identity (RB15),
with the scaled roots. Its three-dimensional difference block has
determinant $1/2$, and both pairing pivots are positive. The complete
residue on
$(\mathsf x,\mathsf y,\mathsf z,\mathsf h,Q_1,Q_2,N,M_B)$
therefore consists of this strict circuit with coefficient
$\kappa>0$, the $N$ pivot $K_m+1/2>0$, and
\[
 \begin{pmatrix}
 1/(2ht)&1/(2h)\\
 1/(2h)&A_*+C_*+D_*+1/2
 \end{pmatrix},
 \quad A_*=t/(2h),\ C_*=(h+t)/(2l),\ D_*=1/[2(t-l)].
 \tag{RI19}
\]
Its determinant is exactly
\[
 \frac{C_*+D_*+1/2}{2ht}>0.                         \tag{RI20}
\]
All pivots are positive bounded units. The strict separated Gram
lemma WA10 yields a finite tensor-SOS identity for the complete
remainder. Restore the two actual nonnegative retained packets and
all separated clearing factors. This proves the entire local contact
chart, including every normal direction and arbitrary relative rates.

\paragraph{The $\omega$-unit state chart uses the original $n$ packet.}
When $u,v$ are infinitesimal, $Z=1$, and $\omega$ is a positive unit,
retain $nD_uD_vX^2$ itself, together with $(m-K_m)N^2$.
The complete remainder is (RB11) with its two negative $n\gamma$
and $n\gamma^2$ contributions absent. This choice is legitimate
directly from $n\geq0$ and does not require a contact lift away from
the contact state scale.

Use the raw roots
\[
 (A_o,\mathsf x,\mathsf D,N,M_B),\qquad A_o=\omega v.
\]
The only new inverse is $1/\omega$, which is pure testing and
bounded. The relations $\mathsf B=A_o/\omega$ and
$\mathsf C=A_o/\omega-\mathsf x$ put the entire scalar and both
first normals in this frame, by the same full formulas (RI9),
(RI12) with the absent $n$ terms deleted. All coefficients are
bounded; the constant term is still $\alpha D_u=\varrho\mathsf D^2$.
Explicitly $\mathsf A=\varepsilon A_o$, so the first scalar term
becomes $D_u^2D_v(p^2/2+\mu_2)A_o^2$ and the first $N$-normal
term becomes $D_u(p-2\mu_1)A_o$. The printed $\varepsilon$
denominators therefore cancel before any coefficient estimate.
Its exact residue is the $A_o$ pivot $1/2$, the $\mathsf D$ pivot
$\kappa$, the $N$ pivot $K_m+1/2$, and the coupled block (RI19).
The first $N$ normal has zero residue and the first $M_B$ normal
has residue $\mathsf x/h$. Thus this entire state chart
also has a strict finite separated Gram.

\paragraph{All other independent tests and states.}
Every source bound used by the noncontact joins of the preceding
$b$-unit theorem is (RB4), which is valid here without a lower bound
on $b$. For bounded unit $u$ and infinitesimal $v$, retain the original
$n,\alpha$ events and use the exact full-state form (RB22), whose
determinant (RB23) is positive. For both angles infinitesimal, the
SD5 exact sign reduction disposes of the nonnegative-core branch
and literal zero axes by original-event squares. On the negative
branch choose the exact gauges $u,v>0$. If the normalized state has
$Z$ infinitesimal, BE4 gives the strict full-state frame
$(X,Y,uW,vW)$ using positive real caps of $m,A$.
Otherwise normalize $Z=1$; $\omega>0$ is bounded, and its two
possibilities are exactly the contact chart and the unit chart just
proved. No source/testing ratio split is used.

For finite remote $v_0\ne0$, use the actual centered-event identity
R5.5, the strict kernel determinant RT10 (printed as RB24), and a
sufficiently large positive real cap below unbounded $m$.
For unbounded $u$ and bounded unit $v$, use R6.5 specialized to
$n_0=0$ and the source residues (RB4). For unbounded $u$ and
infinitesimal $v$, use the exact BE $(\infty,0)$ full frame
$(p_uX+Z,W,Y,vX)$, $p_u=1/u$, with positive caps of $C,A$.
For unbounded $v$ and bounded $u$, R6.6 at $n_0=0$ is strict.
For both angles unbounded, R6.7 at $n_0=0$ has determinants
$D_*/2$ and $A_*(K_m+1/2)$ on its two nontrivial blocks.
These uses import the displayed homogeneous identities with the
source bounds (RB4), without importing a source theorem from another
stratum.
Literal reciprocal zeros are polynomial values of those charts.
All signs and complex states follow by the same exact gauges and
the real/imaginary split. This exhausts all independent tests.

\paragraph{Imported contact algebra.}
The scaled identities are SE5--SE7 in
Section~\ref{component:pa7e744614f49}. Their required rational-support
hypotheses are proved here in (RI8)--(RI15). Only these identities are
used, with the present source estimates; the native small-$e$ source
theorem and its particular residues are not imported. The monic
square-lift and gradient repair are CC4--CC6 in
Section~\ref{component:pfabc29e28ab3}.


\section{\texorpdfstring{R6 at a positive residual value of $a^2B$: an all-tests theorem}{R6 at a positive residual value of a2B: an all-tests theorem}}
\label{component:pb60eba754641}
% Scoped mathematical certificate; see local_identity_audit.closure.json for exact changes.

\label{pb60eba754641:sec:r6-b-unit}

\paragraph{Scope.}
All notation and all physical source constraints are the original weak-C2
ones. In particular $d=1$, $\lambda=\beta=0$, $n\geq0$, and the actual
finite stationary contact, every strict gap, the total-scale condition,
and the complement of earlier admissions are retained. On R6,
$t,h,l=1-h,j,r$ are positive bounded units, $B,x=aB,y=tB,g$ are
positive infinitesimals, and here $a$ is unbounded. Define
\[
 s=a^{-1},\quad w_s=ae,\quad b=a^2B,\quad e=1-r,
 \qquad B=s^2b,\quad x=sb,\quad y=ts^2b.
 \tag{RB1}
\]
The existing exact source identity gives
\[
 2r\Delta_*n/s^2=w_sQ_*-bK_*,\qquad
 Q_*\sim\Delta_*\sim1,\quad K_*\sim h+t.
 \tag{RB2}
\]
The gap $g=1-w_s+hsb$ gives $w_s\sim1$. Hence $b$ is bounded and
\[
 \rho=n/e^2\sim\frac{1-b_0(h_0+t_0)}2,\qquad
 0\leq b_0(h_0+t_0)\leq1.                         \tag{RB3}
\]
Here and below a subscript zero means the real residue; a displayed
residue formula with bounded source letters has the same meaning.
No testing variable occurs in these source estimates.

\begin{theorem}
On the entire original R6 source intersection, assume in addition that
$b=a^2B$ is a positive bounded unit. Then the complete original parent
is a finite sum of squares in the separated source--testing tensor ring
at every independent projective testing pair and every real or complex
state, with either angle sign and every literal projective endpoint.
All positive $s,g$ opening rates are independent. Literal $n=0$ is
included. The theorem makes no assertion about $b$ infinitesimal.
\end{theorem}

\paragraph{Source bounds needed by the proof.}
Put $R=1+hx$, $K=1+x$ and $T=1+x(1-e)$. The original rational source
formulas, after the pure-source substitutions (RB1), have only unit
denominators after $mg$, $n/s^2$ and $\alpha/(s^3b^2)$ are formed.
Their joint residues are
\[
 mg\sim\tfrac12,\quad
 C\sim\frac{h+t}{2l},\quad A\sim\frac t{2h},\quad
 D\sim\frac1{2(t-l)},\quad
 \frac\alpha{s^3b^2}\sim\frac{h+t}2.               \tag{RB4}
\]
Thus $m$ is unbounded, $C,A,D$ are positive units, and all four have
positive real lower bounds. There is no $s/g$ or $g/s$ assumption.
These are rational cancellations in the actual physical source;
no formal negative-beta source is imported.

\paragraph{Two contact-zero packets remove the two troublesome terms.}
Write the complete parent in the normalization $P=F/4$:
\[
\begin{split}
 P={}&nD_uD_vX^2+mM^2+C D_v(uW-Y)^2+A D_uD_vY^2\\
 &+(D+\tfrac12)J^2+(\alpha+\tfrac12)D_u(vY-Z)^2
       +\tfrac12D_v(X+uZ)^2-\tfrac12D_uD_vZ^2,\\
 M={}&vW-X+uvY-uZ,\qquad J=uW+uvX-Y-vZ,
 \quad D_u=1+u^2,\quad D_v=1+v^2.
\end{split}                                                     \tag{RB5}
\]
In the finite state chart $Z=1$, use the pure-testing coordinate
$\omega=(W+uY)/D_u$ and the exact normals
\[
\begin{gathered}
 p=\frac{w_s}{R}=1-\frac gR,\quad c=\frac{l}{tT},\quad
 q_A=D_uv\omega,\quad q_B=v(1-uX),\\
 \qquad N=X+u-pq_A,\quad M_B=uW-Y-cq_B.
\end{gathered}                                    \tag{RB6}
\]
Their full inverse is
\[
\begin{gathered}
 X=N+pD_uv\omega-u,\quad q_B=v(D_u-uN-puD_uv\omega),\\
 \quad W=\omega+\frac{u(M_B+cq_B)}{D_u},\quad
 Y=u\omega-\frac{M_B+cq_B}{D_u}.
\end{gathered}                                    \tag{RB7}
\]
The only testing denominator is the positive pure-testing unit $D_u$.
Set
\[
 \delta=\frac{R}{K\sqrt{bt}},\quad E=\frac{tK}{R},\quad
 k=\sqrt x,\quad K_c=\delta^2Ek^2=\frac{sR}{K},\quad
 \gamma=\frac r{K_c}=\frac{rK}{sR}.                \tag{RB8}
\]
All of $\delta,E,b$ are positive bounded units in this theorem.
The actual contact in these coordinates is
\[
 (u,\omega,v;N,M_B)=(k,\delta k,\delta k^2E;0,0).
 \tag{RB9}
\]
For direct verification, its original state coordinates with $Z=1$ are
\[
 u_c=\sqrt{sb},\quad v_c=s\sqrt{bt},\quad
 W_c=\sqrt{s/t},\quad X_c=-r\sqrt{sb},\quad
 Y_c=hs\sqrt{b/t},\quad
 \omega_c=\sqrt{s/t}\,R/K.
 \tag{RB9a}
\]
They give $u_cW_c-Y_c=ls\sqrt{b/t}=c\,v_c(1-u_cX_c)$
and $X_c+u_c=e\sqrt{sb}=pD_{u_c}v_c\omega_c$.
Also $\gamma v_c\omega_c=r\sqrt{sb}=-X_c$.
Thus both normals and the retained $n$ row vanish exactly, before
any residue calculation. The same contact substitution also gives
$\omega_c=\delta k$, $v_c=\delta k^2E$, and $K_c=sR/K$.
For any fixed positive real $K_m$, retain the two whole squares
\[
 (m-K_m)N^2,
 \qquad nD_uD_v(X+\gamma v\omega)^2.                \tag{RB10}
\]
The first coefficient is positive because $m$ is unbounded. The second
vanishes identically when $n=0$. Both rows vanish at (RB9), with every
first jet of their squares. No inverse of $n$ is used.

The complete remaining polynomial is exactly
\[
\begin{split}
 G={}&K_mN^2-2\frac{mg}{R}q_AN+\frac{mg^2}{R^2}q_A^2\\
 &-2n\gamma D_uD_vX(v\omega)-n\gamma^2D_uD_v(v\omega)^2\\
 &+CD_v(uW-Y)^2+AD_uD_vY^2+(D+\tfrac12)J^2\\
 &+(\alpha+\tfrac12)D_u(vY-1)^2
       +\tfrac12D_v(X+u)^2-\tfrac12D_uD_v.
\end{split}                                                     \tag{RB11}
\]
This is the entire parent minus precisely (RB10), including every
first normal and every hard quadratic term. The exact cancellation
$M=(g/R)q_A-N$ removes the whole $g^{-1}$ pole before an estimate.
Moreover $n\gamma=O(s)$ and $n\gamma^2$ is bounded. All coefficients
in (RB11), after (RB7), are therefore bounded. The only nonconstant
state-normal degree is one or two.

\paragraph{Raw support before the conductor.}
Define
\[
 A_0=\omega v,\quad B_0=\delta v,\quad C_0=\delta E u\omega,
 \quad D_0=\delta^2k^3E,\quad X_0=B_0-C_0.           \tag{RB12}
\]
At $N=M_B=0$ let $V_0=1-puv\omega$. Then
\[
 q_B=D_uvV_0,\quad uW=u\omega+u^2cvV_0,\quad
 Y=u\omega-cvV_0,\quad J=-(1-c)D_uvV_0,
 \quad X+u=pD_uA_0.                               \tag{RB13}
\]
Thus every term of the scalar of (RB11) lies, with bounded coefficients,
in $(A_0,B_0,C_0)^2+(D_0^2)$. The only constant is from $\alpha$,
and $\alpha/D_0^2$ is bounded with residue $b(h+t)/2$.
For the new negative $n$ contribution use exactly
\[
\begin{split}
 -2n\gamma D_uD_vX A_0-n\gamma^2D_uD_vA_0^2
 ={}&2n\gamma D_uD_vuA_0\\
 &-(2n\gamma pD_u+n\gamma^2)D_uD_vA_0^2
 -2n\gamma D_uD_vA_0N,
\end{split}                                                     \tag{RB14}
\]
where $uA_0=B_0C_0/(\delta^2E)$. Both complete first normals lie in
$(A_0,B_0,C_0)$ with bounded coefficients. Every remaining hard entry
is bounded. These are full rational support statements, not a scalar
contact-value test. Clearing $D_u^2$ makes them polynomial; its inverse
is a pure-testing square factor when restored.

\paragraph{The contact Gram and its exact leading calculation.}
Let
\[
 Y_0=A_0-kC_0,\quad Z_0=uB_0-kC_0,\quad H_0=kC_0-D_0.
\]
Use the separated nonnegative monomials $a_i^2=(A_0,uB_0,kC_0,D_0)$
and the two pairings $Q_1=a_1a_2-a_3a_4$,
$Q_2=a_1a_3-a_2a_4$. Their exact circuit is
\[
\begin{split}
 T_4&=A_0^2+(uB_0)^2+(kC_0)^2+D_0^2-4k^2B_0C_0\\
 &=\tfrac12\{(Y_0-Z_0)^2+H_0^2+Y_0^2+(Z_0+H_0)^2\}
                   +Q_1^2+Q_2^2.
\end{split}                                      \tag{RB15}
\]
Its matrix on $(Y_0,Z_0,H_0)$ is
$\left(\begin{smallmatrix}1&-1/2&0\\-1/2&1&1/2\\0&1/2&1\end{smallmatrix}\right)$,
which is positive definite with determinant $1/2$. All radicals
separate: $a_1$ is pure testing, $a_2,a_3$ are products of pure-source
and pure-testing radicals, and $a_4$ is pure source.

Here is enough explicit algebra to reproduce the leading collection.
At $g=0$, the canceled source expansions through the needed orders are
\[
\begin{aligned}
 n/s^2&=\tfrac12\{1-b(h+t)\}+O(s),&
 mg&=\tfrac12+\tfrac12(bh-1)s+O(s^2),\\
 C&=\frac{h+t}{2l}-\frac{t+1}{2l}s+O(s^2),&
 A&=\frac t{2h}-\frac t{2h}s+O(s^2),\\
 D&=\frac1{2(t-l)}-
       \frac{bt(h+t)+t-l}{2(t-l)^2}s+O(s^2),&
 \alpha&=\tfrac12 b^2(h+t)s^3+O(s^4).
\end{aligned}                                                     \tag{RB16}
\]
Here $O(s^j)$ means exact divisibility by $s^j$ in a rational ring with
unit denominator, not an asymptotic substitution in the proof identity.
Give $k,u,\omega$ weight one and $v$ weight two. Put $q=b(h+t)$.
Using $s=k^2/b$ and
$\delta=(1-lk^2+O(k^4))/\sqrt{bt}$, $\delta E=\sqrt{t/b}$,
the scalar $G_0=G|_{N=M_B=0}$ satisfies
\[
 G_0=\frac b{2h}(1+u^2)X_0^2+\frac q2T_4
       -\frac{k^2[1+bh(h+t)]}{2h}X_0^2+O_{\rm wt}(8)
       \quad\pmod g.                               \tag{RB17}
\]
The complete first-normal coefficients obey
\[
 L_N=O_{\rm wt}(5),\qquad
 L_{M_B}=\frac{X_0}{h\delta}+O_{\rm wt}(4)\quad\pmod g,
 \tag{RB18}
\]
and the hard block has residue
\[
 \operatorname{diag}\left(K_m+\tfrac12,\,
           A_*+C_*+D_*+\tfrac12\right),
 \quad A_*={t\over2h},\ C_*={h+t\over2l},\ D_*={1\over2(t-l)}.
 \tag{RB19}
\]

Apply the inherited monic contact--conductor identities CC1--CC6 to
the exact difference in (RB17), and to the exact first-normal
differences in (RB18). Their hypotheses are exactly the raw support in
(RB13)--(RB14), together with zero contact value and full first jet.
Both follow from the original stationary parent and (RB10); the model
$X_0^2,T_4$ also has these zeros. No contact difference is inverted.
The conductor maps preserve this grading: $X_0$ has weight two,
the five circuit roots have weight three. The scalar difference of
weight at least eight and normal differences of the displayed weights
therefore have strictly positive-weight Gram coefficients. Every such
coefficient is infinitesimal. Terms divisible by $g$ also lift with
bounded coefficients and remain infinitesimal: the contact (RB9),
$\delta$ and $E$ are independent of $g$, and the canceled coefficient
ring has unit denominators at $g=0$. Clearing and restoring its finite
unit denominators changes none of these conclusions.

Consequently the residue on the eight roots
$(X_0,Y_0,Z_0,H_0,Q_1,Q_2,N,\delta M_B)$ is the strict circuit with
coefficient $q/2>0$, the $N$ pivot $K_m+1/2>0$, and the coupled block
\[
 \begin{pmatrix}
 b/(2h)&bt/(2h)\\
 bt/(2h)&bt(A_*+C_*+D_*+1/2)
 \end{pmatrix}.                                    \tag{RB20}
\]
Its determinant is
\[
 \frac{b^2t}{2h}(C_*+D_*+1/2)>0.                    \tag{RB21}
\]
All pivots are positive bounded units. The separated strict-residue
Gram lemma therefore gives finite tensor SOS for the exact remainder.
Restore (RB10) and the pure-source/pure-testing clearing factors.
This proves the whole local contact parent, for every independent rate
and every normal direction.

\paragraph{All remaining finite state and angle charts.}
The following details make the join an all-tests result.
For bounded $u$ with nonzero residue and infinitesimal $v$, retain
$nD_uD_vX^2$ and $\alpha D_u(vY-Z)^2$, and the centered $m$ excess
$(m-K_m)N^2$ from (RB6), now homogeneous in $Z$. Put
$V_Z=vZ$ and $R_t=uW-Y$. After $X=N-uZ+pv(W+uY)$, the whole
remainder is a bounded quadratic on $(N,R_t,Y,V_Z)$. Its residue
splits off $K_m+1/2$ on $N$, with the other form
\[
 C_*R_t^2+A_*(1+u^2)Y^2
 +(D_*+1/2)\{R_t-(1+u^2)V_Z\}^2
 -(1+u^2)YV_Z-\tfrac12(1+u^2)V_Z^2.               \tag{RB22}
\]
Its first two leading principal minors are positive, and its determinant
is exactly
\[
 (1+u^2)^2A_*C_*(D_*+1/2)u^2>0.                    \tag{RB23}
\]
This direct four-root proof covers every state, including $Z=0$, with
no negative-beta or small-$\rho$ argument.

It remains to consider both angles infinitesimal. The exact testing-core
sign reduction from SD5 handles nonnegative core and literal axes by
original-event squares. On its negative-core branch one may use the
standard exact gauges $u,v>0$. Normalizing the original state by a
maximal coordinate, the case $Z$ infinitesimal has $W$ a unit and is
covered by BE4, with positive real caps of $m,A$ supplied by (RB4).
Retain the actual nonnegative $n,\alpha$ packets. That exact selected
full-parent frame is $(X,Y,uW,vW)$ and has strict residue; no coefficient
bound involving $b$ enters its proof.

If $Z$ is a unit normalize $Z=1$. The same core identity gives
$\omega>0$, and $\omega$ is bounded. If $\omega$ is infinitesimal,
use (RB17)--(RB21). If $\omega$ is a positive unit, use the exact
relations $B_0=(\delta/\omega)A_0$ and $C_0=B_0-X_0$ in the raw
support (RB13)--(RB14). The inverse is pure testing. Its complete
residue on $(A_0,X_0,D_0,N,\delta M_B)$ consists of the positive
$q/2$ pivots on $A_0,D_0$, the $K_m+1/2$ normal pivot, and precisely
(RB20). Direct substitution in (RB11) gives these residues; all other
coefficients contain $k,g,u,v$ after the cancellations already recorded.
Thus this whole state chart is strict as well.

\paragraph{Every other projective test.}
Use the exact centered-event identity R5.5 for the finite remote pairs
$v_0\ne0$. All source residues are (RB4) and $n_0=\alpha_0=0$.
On the kernel of the limiting centered $m$ row, the determinant is
\[
 \frac{(1+U)^2(1+V)^2(h+t)}{8hl(t-l)}
       \{V^2+t(h+t)U^2V+hUV^2\}>0.                 \tag{RB24}
\]
The restriction to $Z=0$ is positive by the original $A,C$ packets
and the nonnegative signed-core restriction, exactly as in RT10.
A sufficiently large positive real cap $K_m$ makes the full bounded
residue strict. Its choice depends on the real testing residues only;
unbounded $m$ permits every such real cap.

For unbounded $u$ and bounded unit $v$, R6.5 at $n_0=0$ is strict
because $C_*W^2$, $A_*Y^2$ and $(D_*+1/2)(W+vX)^2$ have common
kernel zero. For unbounded $u$ and infinitesimal $v$, the exact
BE $(\infty,0)$ frame $(pX+Z,W,Y,vX)$ retains every event except
positive caps of $C,A$, and has strict residue. For unbounded $v$
and bounded $u$, R6.6 at $n_0=0$ is strict; the centered correction
vanishes on its $m$-row kernel. For both angles unbounded, R6.7 at
$n_0=0$ has block determinants $D_*/2$ and
$A_*(K_m+1/2)$, both positive. These are exact homogeneous full-parent
identities. Their literal reciprocal zeros include the actual projective
endpoints; undoing a nonzero homogeneous factor uses a pure-testing
square only.

This exhausts all independently chosen angle and state charts.
Apply the real identities separately to the real and imaginary parts
of a complex state. Every retained positive packet and original
normalization is restored. The new admission is the entire stated
$b$-unit source subdomain, rather than a scalar graph or a single
local testing corner.


\paragraph{Imported exact identities.}
The BE4 four-root state identity and the unnumbered $(\infty,0)$
identity in Section~\ref{component:p69995d829bb2} supply the two
selected full-parent frames used above. The latter is the paragraph
beginning ``At $(\infty,0)$ retain all events except $c,A$''.
The remote kernel determinant and its $Z=0$ restriction are printed
in (RB24) and the paragraph following it. The centered-event identity
is R5.5 in Section~\ref{component:p8da9ea6ed9e5}. The homogeneous
kernel and root matrices R6.5--R6.7 are printed in
Section~\ref{component:pd5fa361b4623}; they are used at $n_0=0$
with the source residues (RB4). Finally SD5 and its following sign
Schur identity in Section~\ref{component:p0acdb66bba2f} give the
pure-testing inverse and the exact sign reduction. Only these
algebraic identities are used with the present source estimates;
no C1 physical-source or negative-beta perturbation theorem is imported.


\section{\texorpdfstring{The unbounded-$a$ complement of R7, including its physical zero face}{The unbounded-a complement of R7, including its physical zero face}}
\label{component:p405d96527133}
% Scoped mathematical certificate; see local_identity_audit.closure.json for exact changes.

\label{p405d96527133:sec:r7-unbounded-a}

\paragraph{Quantified source and conclusion.}
Keep the entire original weak-C2 source: $d=1$, $\lambda=\beta=0$,
the actual rational stationary coefficients, $n\ge0$, every strict gap,
the original finite stationary contact, total scale, and all prior
exclusions. In R7, $B,z,y$ are positive infinitesimals,
$a=z/B$, $t=y/B$ is a positive bounded unit, and $h,g$ are positive
infinitesimals. The original exclusion of simultaneously infinitesimal
$r$ and the source gap $j_{\rm gap}$ is retained.

\begin{theorem}
On the whole original R7 source intersection with $a$ unbounded,
the complete original parent is a finite separated source--testing
tensor sum of squares for every independent projective testing pair
and every real or complex state. Both angle signs, all literal zero
and infinite angles, literal $n=0$, and every relative opening rate
of the source infinitesimals are included. No inverse or radical of
a mixed source--testing expression is used.
\end{theorem}

The inherited bounded-$a$ R7 theorem and this theorem exhaust R7.
This statement does not admit R3, R4, or the global finite-lift theorem.

\paragraph{Actual source scale and the physical inequality.}
Use
\[
 \begin{gathered}
 s=1/a,\quad b=a^2B,\quad x=sb=z,\quad B=s^2b,\\
 w_s=ae=1+hx-g,\quad e=sw_s,\quad r=1-e,\quad l=1-h.
 \end{gathered}
 \tag{RU1}
\]
Set $K=1+x$, $R=1+hx$, $T=1+xr$ and
$j_{\rm gap}=T-l/t$. The exact original source identity is
\[
\begin{gathered}
 2r\Delta n=eQ-BK_{\rm phys},\qquad
 2r\Delta_*n/s^2=w_sQ-bK_*,\qquad \Delta_*=s\Delta,\\
 Q=1+B(1+t+ah+2ar+at)
    +B^2\{a(h+t)+a^2[h(r+1)+t(1+r-r^2)]\},\\
 K_{\rm phys}=(a+1)(h+t)(1+x),\qquad
 K_*=(1+s)(h+t)(1+x).
\end{gathered}                                                     \tag{RU2}
\]
Before bounding $b$, $Q\sim\Delta_*\sim1$ since all their nonconstant
terms contain $s$ or $x$. Here $w_s\sim1$ and $K_*\sim t$ is a
positive unit. The original inequality $eQ\ge BK_{\rm phys}$
therefore gives $b$ bounded and $n=O(s^2)$, including its literal
equality case. In particular $r\sim1$, so the original R7 exclusion
is respected. The source gap gives $t_0\ge1$; $j_{\rm gap}$ is a
unit exactly when $t_0>1$, and otherwise $t_0=1$.

The original rational formulas, after cancellation, give jointly
\[
 mg\sim\tfrac12,\quad C\sim\tfrac t2,\quad
 A_h:=Ah\sim\tfrac t2,\quad
 \beta_D:=D(1-c)=Dj_{\rm gap}/T\sim\frac1{2t},
 \quad \frac\alpha{s^3b^2}\sim\tfrac t2,
 \qquad c=\frac l{tT}.
 \tag{RU3}
\]
Thus $m,A$ are unbounded; $C$ is a positive unit. If $t_0>1$,
$D\sim1/[2(t-1)]$ is a positive unit. If $t_0=1$, $D$ is unbounded.
The canceled coefficients $mg,C,A_h,\beta_D,n\gamma,n\gamma^2$
below have unit denominators through both small-gap ends. No ratio
$h/g$, $j_{\rm gap}/h$, or comparison with a testing rate is used.

\paragraph{Four whole packets and a bounded full-parent remainder.}
Use the complete parent $P=F/4$:
\[
\begin{aligned}
 P={}&nD_uD_vX^2+mM^2+CD_v(uW-Y)^2+AD_uD_vY^2\\
 &+(D+1/2)J^2+(\alpha+1/2)D_u(vY-Z)^2
                   +D_v(X+uZ)^2/2-D_uD_vZ^2/2,\\
 M={}&vW-X+uvY-uZ,\qquad J=uW+uvX-Y-vZ.
\end{aligned}
\]
The exact coordinate inverse (RB6)--(RB7) is
\[
\begin{gathered}
 D_u=1+u^2,\quad D_v=1+v^2,\quad Z=1,\quad
 p=w_s/R,\quad q_A=D_uv\omega,\\
 N=X+u-pq_A,\qquad q_B=v(1-uX),\qquad M_B=uW-Y-cq_B,\\
 X=N+pD_uv\omega-u,\quad q_B=v(D_u-uN-puD_uv\omega),\\
 W=\omega+u(M_B+cq_B)/D_u,\qquad
 Y=u\omega-(M_B+cq_B)/D_u.
\end{gathered}                                                     \tag{RU4}
\]
The sole testing inverse is $D_u^{-1}$. Put
\[
 \begin{gathered}
 \gamma=\frac{rK}{sR},\quad \eta=Y-huW,\quad q=1-c=j_{\rm gap}/T,\\
 \mu_1=mg/R,\quad \mu_2=(mg)g/R^2,\quad n_1=n\gamma,\quad n_2=n\gamma^2.
 \end{gathered}
\]
Retain the four entire nonnegative squares
\[
 mN^2,\qquad nD_uD_v(X+\gamma v\omega)^2,
 \qquad AD_uD_v\eta^2,\qquad DM_B^2.
 \tag{RU5}
\]
All four rows vanish at the actual contact; (RB9a) verifies the first
two, and directly $Y_c=hu_cW_c$, $M_{B,c}=0$.
The $n$ packet vanishes identically on the literal $n=0$ face.

The full $A,D$ cancellations are the exact identities
\[
\begin{aligned}
 AD_uD_v(Y^2-\eta^2)&=2A_hD_uD_vuWY-A_hhD_uD_v(uW)^2,\\
 D(J^2-M_B^2)&=-2\beta_DM_Bq_B+\beta_Dq q_B^2,
 \qquad J=M_B-q q_B.
\end{aligned}                                                     \tag{RU6}
\]
Since $M=(g/R)q_A-N$, subtraction of precisely (RU5) leaves
\[
\begin{aligned}
 G={}&-2\mu_1q_AN+\mu_2q_A^2
       -2n_1D_uD_vX(v\omega)-n_2D_uD_v(v\omega)^2\\
    &+CD_v(uW-Y)^2+2A_hD_uD_vuWY-A_hhD_uD_v(uW)^2\\
    &-2\beta_DM_Bq_B+\beta_Dq q_B^2+J^2/2\\
    &+(\alpha+1/2)D_u(vY-1)^2+D_v(X+u)^2/2-D_uD_v/2.
\end{aligned}                                                     \tag{RU7}
\]
This is a whole-parent identity, including every scalar, normal and
hard coefficient. Equations (RU3), (RU6) bound all of them before
any contact restriction. In particular neither $h^{-1}$ nor
$j_{\rm gap}^{-1}$ survives.

\paragraph{One regular contact frame for the two $b$ alternatives.}
Use the pure-source parameters
\[
 \varepsilon=\frac{K\sqrt{bt}}R,\quad
 \jmath=\sqrt{s/t}\,\frac RK,\quad k=\varepsilon\jmath=\sqrt x,
 \quad E=tK/R.
 \tag{RU8}
\]
The symbol $\jmath$ is a contact scale, not $j_{\rm gap}$.
Here $\jmath$ is infinitesimal; $\varepsilon$ is a unit or
infinitesimal exactly as $b$ is. The exact inverse source chart is
\[
 x=\varepsilon^2\jmath^2,\qquad
 b=\frac{\varepsilon^2R^2}{K^2t},\qquad
 s=\frac{t\jmath^2K^2}{R^2}.
 \tag{RU9}
\]
Only source units are inverted. The contact is
$(u,\omega,v;N,M_B)=(k,\jmath,\varepsilon\jmath^2E;0,0)$.

Define
\[
\begin{gathered}
 \mathsf A=\varepsilon\omega v,\quad\mathsf B=v,
 \quad\mathsf C=Eu\omega,\quad\mathsf D=\varepsilon^2\jmath^3E,\\
 \mathsf x=\mathsf B-\mathsf C,\quad
 \mathsf y=\mathsf A-k\mathsf C,\quad
 \mathsf z=u\mathsf B-k\mathsf C,\quad
 \mathsf h=k\mathsf C-\mathsf D.
\end{gathered}                                                     \tag{RU10}
\]
The two potentially singular source combinations are
\[
 F_0=p^2/2-n_2-2n_1p+\mu_2,\qquad A_N=p-2\mu_1-2n_1.
\]
Their exact canceled quotients satisfy
\[
 F_0/b\sim t/2,\qquad A_N/(bs)\sim t.
 \tag{RU11}
\]
Both denominators are units through $s=b=h=g=0$. These assertions
follow from the original rational source, not from the inequality
$n\ge0$. In particular
\[
 F_0/\varepsilon^2=\frac{R^2}{K^2t}\frac{F_0}b,
 \qquad \alpha/\mathsf D^2=\frac1t\frac\alpha{s^3b^2},
 \qquad A_N/\varepsilon=\frac{\varepsilon R^2s}{K^2t}
                               \frac{A_N}{bs}
 \tag{RU12}
\]
are bounded without losing the zero-$n$ face.

\paragraph{Exact full scalar and both first normals.}
Put $H=\alpha+1/2$, $V_0=1-puv\omega$. At $N=M_B=0$,
$uW=\mathsf C/E+u^2c\mathsf BV_0$ and
$Y=\mathsf C/E-c\mathsf BV_0$. The complete scalar is
\[
 G_0=f_A\mathsf A^2+f_{BB}\mathsf B^2+
       f_{BC}\mathsf B\mathsf C+f_{CC}\mathsf C^2+
       \varrho\mathsf D^2,                                      \tag{RU13}
\]
with the bounded coefficients
\[
\begin{aligned}
 f_A={}&D_u^2D_vF_0/\varepsilon^2,\qquad
 \varrho=\alpha D_u/\mathsf D^2,\\
 f_{BB}={}&[Cc^2D_u^2D_v+(\beta_Dq+q^2/2)D_u^2+HD_uv^2c^2
             -A_hD_uD_v(2u^2+hu^4)c^2]V_0^2\\
          &+2HD_ucV_0-D_u/2-D_u^2\mu_2v^2\omega^2,\\
 f_{BC}={}&[2A_hD_uD_v(lu^2-1)cV_0-2HD_uv^2cV_0
                   -2HD_u+2n_1D_uD_v]/E,\\
 f_{CC}={}&[A_h(2-h)D_uD_v+HD_uv^2+D_uD_vn_2v^2]/E^2.
\end{aligned}                                                     \tag{RU14}
\]
For the full normal expansion
$G=G_0+L_NN+L_MM_B+G_{\rm hard}$, the exact coefficients are
\[
 L_N=R_A\mathsf A+R_B\mathsf B+R_C\mathsf C,
 \qquad L_M=S_B\mathsf B+S_C\mathsf C,                           \tag{RU15}
\]
where
\[
\begin{aligned}
 R_A={}&D_uA_N/\varepsilon,\\
 R_B={}&D_u(p-2n_1)v^2\omega
 -2uv[Cc^2D_uD_v+(\beta_Dq+q^2/2)D_u]V_0\\
 &+2A_hc^2uvD_vu^2(2+hu^2)V_0-2Hcuv(1+cv^2V_0),\\
 R_C={}&[2A_hcuvD_v(1-lu^2)+2Hcuv^3]/E,\\
 S_B={}&2[CcD_uD_v-(\beta_D+q/2)D_u
                      -A_hcD_vu^2(2+hu^2)]V_0
                         +2H(1+cv^2V_0),\\
 S_C={}&[2A_hD_v(lu^2-1)-2Hv^2]/E.
\end{aligned}                                                     \tag{RU16}
\]
At $u=v=0$ the complete hard matrix on $(N,M_B)$ is exactly
$\operatorname{diag}(1/2,C+1/2)$.

\paragraph{Bounded monic conductor, including the small-$\varepsilon$ end.}
The identities (RI14)--(RI15) apply to (RU10) without a source-shape
assumption. Their crucial inputs are checked here on the new complete
coefficients. After finite source unit denominators and $D_u^2$ are
cleared, every $\omega$ occurrence of (RU14) is through $uv\omega$
or $v^2\omega^2$, and all remaining testing powers are even in $u,v$.
The scalar raw support is exactly the five terms in (RU13).

For completeness, first replace $\mathsf A=\mathsf y+k\mathsf C$,
$\mathsf B=\mathsf x+\mathsf C$, and
$\mathsf D=k\mathsf C-\mathsf h$. The coefficient requiring a
first conductor is $S=k^2f_A+f_{BB}+f_{BC}+f_{CC}+k^2\varrho$.
Divide its contact-zero difference monically in order $u,\omega,v$.
At $u=k=\varepsilon\jmath$, the $\omega$ divided difference is in
$(\varepsilon v,v^2)$; write it as $\varepsilon U+vV$ with
$U,V\in(v)$. The last $v$ divided difference is in
$(\varepsilon,v)$ because, modulo $\varepsilon$, the remaining
$v$ dependence is even. The first $u$ divided difference restricts
to zero at $\varepsilon=u=v=0$.

These facts permit the exact bounded lift (RI15), which uses the
identity for $\mathsf Cv(\omega-\jmath)$ and never divides by
$\varepsilon$. Subtract the gradient multiple
$2k\varrho_c(-2k,1,1,-1)$. The first three contact gradients have
determinant $-\varepsilon^3\jmath^4E^2\ne0$; the CC5 adjugate
argument cancels a zero identity, and introduces no inverse of that
determinant. Each second-stage coefficient at
$\varepsilon=u=v=0$ is independent of $\omega$: the $\mathsf x$
one restricts to $2f_{BB}+f_{BC}$, and the other three to zero.
Hence every second $\omega$ divided difference is in
$(\varepsilon,v)$ and has the same bounded lift. The final repair is
\[
 k\mathsf C(\mathsf y+\mathsf z-\mathsf h-2k\mathsf x)
      =Q_1^2-\mathsf y\mathsf z.
 \tag{RU17}
\]
The pairings $Q_i$ are defined below. For each complete normal in
(RU15), one conductor suffices: after replacing $\mathsf A,\mathsf B$,
its $\mathsf C$ coefficient restricted at
$\varepsilon=u=v=0$ is independent of $\omega$, directly from
(RU16). This is precisely the normal criterion in the same proof.
When $b$ is a unit, $\varepsilon^{-1}$ is bounded and the ordinary
CC1--CC6 conductor is also available.

All contact value and first-jet hypotheses follow from the stationary
original parent and the four exact contact-zero packets (RU5).
The algorithm is a bounded, source-linear map on the stated raw
coefficient class and preserves the filtrations used next. In
particular it preserves the source ideal $(h,g)$; no separate
contact-zero assertion for its individual coefficient summands is
needed. The source substitution in each monic step uses the actual
$E$, which is a bounded unit through $h=0$.

\paragraph{The complete filtered leading model.}
Set
\[
 \Lambda=\frac1{t(1+s)},\qquad \mu=\frac1{1+s},\qquad
 T_4=\mathsf A^2+(u\mathsf B)^2+(k\mathsf C)^2+
                \mathsf D^2-4k^2\mathsf B\mathsf C.
\]
The model scalar is
\[
\begin{aligned}
 G_{\rm mod}
 &=\Lambda\mathsf x^2+T_4/2-(1/t+1/2)k^2\mathsf x^2
                      +(k^2-u^2)\mathsf C\mathsf x/t\\
 &=\left[\Lambda+\frac{u(u+k)}t-(1/t+1/2)k^2\right]\mathsf x^2
                  +T_4/2-\frac{u+k}t\mathsf x\mathsf z.
\end{aligned}                                                     \tag{RU18}
\]
The second equality is the exact conductor identity
$\mathsf C(u-k)=\mathsf z-u\mathsf x$. Use the first expression
for the raw coefficient filtration, since it retains the even-$u$
class. Its angular raw coefficients are
\[
\begin{aligned}
 F_{BB}&=\Lambda+u^2/2-(1/t+1/2)k^2,\\
 F_{BC}&=-2\Lambda-u^2/t+(3/t-1)k^2,\\
 F_{CC}&=\Lambda+u^2/t-2k^2/t.
\end{aligned}                                                     \tag{RU19}
\]
At $h=g=0$, exact source collection gives the following two
coefficientwise filtrations:
\[
\begin{array}{c|c|c}
 &b\text{ a unit}&b\text{ infinitesimal}\\ \hline
 \text{variable weights}&(k,u,\omega,v)=(1,1,1,2)
     &(\varepsilon,\jmath,u,\omega,v)=(1,1,2,1,3)\\
 \operatorname{wt}(\mathsf x)&2&3\\
 \operatorname{wt}(\mathsf y,\mathsf z,\mathsf h,Q_1,Q_2)&3&5\\
 f_A-1/2,\ \varrho-1/2&\ge2&\ge2\\
 f_{BB}-F_{BB},\ f_{BC}-F_{BC},\ f_{CC}-F_{CC}&\ge4&\ge6\\
 S_B-\mu,\ S_C+\mu&\ge2&\ge4\\
 (R_A,R_B,R_C)&\ge(2,3,3)&\ge(3,5,5).
\end{array}                                                       \tag{RU20}
\]
These statements concern the chosen raw coefficients separately;
polynomial cancellation between raw products is not used to infer
them. Consequently
\[
\begin{array}{c|c|c}
 &b\text{ a unit}&b\text{ infinitesimal}\\ \hline
 G_0-G_{\rm mod}&O_{\rm wt}(8)&O_{\rm wt}(12)\\
 L_M-\mu\mathsf x&O_{\rm wt}(4)&O_{\rm wt}(7)\\
 L_N&O_{\rm wt}(5)&O_{\rm wt}(8).
\end{array}                                                       \tag{RU21}
\]
Here each equality is modulo the source ideal $(h,g)$.

To reproduce the collection, the necessary $b$-infinitesimal source
jets at $h=g=0$ are
\[
\begin{aligned}
 C&=t/2-(t+1)s/2+(t+1)s^2/2+O_{\deg(s,b)}(3),\\
 A_h&=t(1-s+s^2)/2+O_{\deg(s,b)}(3),\\
 \beta_D&=[1-s+s^2-(t+1)sb]/(2t)+O_{\deg(s,b)}(3),\\
 n_1&=(s-s^2-tsb)/2+O_{\deg(s,b)}(3).
\end{aligned}                                                     \tag{RU22}
\]
Also $c=(1-k^2)/t+O_{\rm wt}(6)$,
$E=t(1+k^2)$ at $h=0$, and
$s=t\jmath^2+O_{\rm wt}(6)$. For the unit-$b$ filtration use
$s=k^2/b$ and the respective coefficients of $s$:
$-(t+1)/2,-t/2,-[1+(t+1)b]/(2t),(1-bt)/2$.
The source cancellations (RU11)--(RU12) supply the first and last
rows of (RU20). Scalar weights are even in both
charts. The $M_B$ weights are even in the first and odd in the
second, giving exactly the next weights in (RU21).

Apply the bounded conductor to the exact differences in (RU21).
Their coefficientwise raw filtration makes every correction Gram
coefficient have strictly positive weight. Terms in $(h,g)$ lift
boundedly and remain infinitesimal at arbitrary relative rates.
The whole hard block differs infinitesimally from
$\operatorname{diag}(1/2,(t+1)/2)$, since all its coefficients are
bounded and $u,v,h,g,s$ are infinitesimal.

On the negative-core contact chart take $u,v,\omega>0$. The separated
nonnegative monomials
\[
 a_i^2=(\mathsf A,u\mathsf B,k\mathsf C,\mathsf D),\qquad
 Q_1=a_1a_2-a_3a_4,\quad Q_2=a_1a_3-a_2a_4
\]
have the exact product relation yielding
\[
 T_4=\tfrac12\{(\mathsf y-\mathsf z)^2+\mathsf h^2+
                 \mathsf y^2+(\mathsf z+\mathsf h)^2\}
                  +Q_1^2+Q_2^2.                                \tag{RU23}
\]
Its three-dimensional difference matrix has determinant $1/2>0$.
Each radical separates into source and testing factors.
The full eight-root residue on
$(\mathsf x,\mathsf y,\mathsf z,\mathsf h,Q_1,Q_2,N,M_B)$
therefore consists of half this strict circuit, the $N$ pivot $1/2$,
and the complete coupled block
\[
 \begin{pmatrix}1/t&1/2\\1/2&(t+1)/2\end{pmatrix},\qquad
 \det=\frac{t+2}{4t}>0.                          \tag{RU24}
\]
The additional cross term in the second line of (RU18) has
infinitesimal coefficient. WA10 now gives a finite separated SOS
for the exact remainder. Restore all four actual squares (RU5).
This proves the whole local contact chart in both $b$ alternatives.

\paragraph{The $\omega$-unit state chart.}
If $u,v$ are infinitesimal, $Z=1$, and $\omega$ is a positive unit,
replace the retained centered $n$ packet by $nD_uD_vX^2$ itself.
The complete remainder is (RU7) with the two $n_1,n_2$ terms removed.
Use the frame $(A_o,\mathsf x,\mathsf D,N,M_B)$, $A_o=\omega v$.
The exact identities
$\mathsf A=\varepsilon A_o$, $\mathsf B=A_o/\omega$,
$\mathsf C=A_o/\omega-\mathsf x$ use only a pure-testing unit
inverse. In particular the first scalar coefficient becomes
$D_u^2D_v(p^2/2+\mu_2)A_o^2$, and the first $N$ coefficient
becomes $D_u(p-2\mu_1)A_o$; the printed $\varepsilon$ denominators
cancel completely. The full residue has the $A_o,\mathsf D,N$
pivots $1/2,1/2,1/2$ and the same block (RU24).
The complete normal residues are $0,\mathsf x$. Hence this chart
also has a strict finite separated Gram.

\paragraph{A remote cap lemma and complete bounded parents.}
Away from the contact chart there is no need for centered excesses.
Retain the original nonnegative $n,\alpha$ packets and the positive
excesses of $m,A$ above chosen positive real caps $m_c,A_c$.
If $D$ is unbounded, retain its original excess above a real cap
$D_c$ too; otherwise leave $D$ at its actual bounded value.
The remainder is the full original parent with these three
coefficients capped, and all its coefficients are bounded.

We use the elementary finite-dimensional fact that a real quadratic
form positive definite on the common kernel of finitely many real
linear rows becomes positive definite after sufficiently large
positive real multiples of the squared rows are added. Choose
coordinates on the common kernel and its orthogonal complement;
the assertion follows by a Schur complement. Unbounded actual
$m,A,D$, when used, exceed each chosen real cap. Thus every cap
is an exact nonnegative packet subtraction in the original parent,
and a strict residue gives WA10 for the entire bounded remainder.
The real caps depend only on real testing residues, never on a
source/testing rate comparison.

\paragraph{Finite remote $v$ with nonzero residue.}
Write $U=u^2,V=v^2$ for real testing residues, with $V>0$.
On the common kernel $Y=M=0$, $X=vW-uZ$, and the remaining
two-state form is exactly
\[
 CD_vU W^2+(D+1/2)(uD_vW-vD_uZ)^2
                  +VD_vW^2/2-D_uVZ^2/2.                         \tag{RU25}
\]
Its determinant is
\[
 VD_uD_v\{DV/2+U[CD-(D+1/2)/2]+C(D+1/2)U^2\}.
 \tag{RU26}
\]
For $t_0>1$, insert $C=t/2,D=1/[2(t-1)]$; it equals
$D_uD_v[V^2+t^2U^2V]/[4(t-1)]>0$.
Its $W^2$ coefficient is positive since $V>0$.
For $t_0=1$, $C=1/2$ and one may choose a real $D$ cap: the braces
are $D(V+U^2)/2+(U^2-U)/4$, positive for sufficiently large $D$.
The cap lemma then chooses $m_c,A_c$ for a strict full state matrix.

\paragraph{Finite unit $u$ with infinitesimal $v$, including zero.}
Use $N=-M=X+uZ-v(W+uY)$, $R_t=uW-Y$, and $V_Z=vZ$.
After the original $n,\alpha$ packets and allowed positive excesses
are retained, the entire remainder lies in the square of the four
roots $(N,R_t,Y,V_Z)$. Only the pure-testing unit $u$ is inverted
to write $W=(R_t+Y)/u$; no inverse of $v$ is needed. Its exact
residue is
\[
\begin{aligned}
 (m_c+1/2)N^2+CR_t^2+A_cD_uY^2
 +(D+1/2)(R_t-D_uV_Z)^2-D_uYV_Z-D_uV_Z^2/2.
\end{aligned}                                                     \tag{RU27}
\]
For $t_0>1$, its restriction to $Y=0$ has determinant
$t^2D_u(D_u-1)/[4(t-1)]>0$ on $(R_t,V_Z)$.
Choose $A_c$ sufficiently large. For $t_0=1$, the common kernel
$Y=0,R_t-D_uV_Z=0$ has value $D_uu^2V_Z^2/2>0$;
large real $D_c,A_c$ give strictness. The normal pivot is positive.
The full four-root support and the residue are checked directly;
literal $v=0$ is their polynomial value.

\paragraph{Both finite angles infinitesimal, and all state charts.}
The exact SD5 pure-core sign identity deals with the nonnegative-core
branch and literal zero axes by original-event squares. On the
negative branch use its standard gauges $u,v>0$ and normalize the
state by a maximal coordinate. If $Z$ is infinitesimal, $W$ is a
unit and the exact BE4 selected four-root identity applies to
$(X,Y,uW,vW)$, with positive caps of $m,A$ provided by (RU3).
All other original events remain as nonnegative packets.
If $Z$ is a unit normalize $Z=1$; the same sign identity gives
bounded positive $\omega$. Its infinitesimal and unit alternatives
are exactly (RU18)--(RU24) and the subsequent $\omega$-unit chart.

\paragraph{Every reciprocal chart and literal projective infinity.}
For unbounded $u$, put $p_u=1/u$ and use the entire homogeneous
polynomial $p_u^2P$. If bounded $v$ has nonzero residue, the full
capped four-state residue is
\[
 CD_vW^2+A_cD_vY^2+(D_c+1/2)(W+vX)^2
                          +(m_c+1/2)(vY-Z)^2\succ0.             \tag{RU28}
\]
Here $D_c$ denotes either a positive real cap or the actual positive
unit residue. The four square rows have common kernel zero because
$v$ is a unit. If instead $v$ is infinitesimal, including zero,
use the exact BE $(\infty,0)$ selected full frame
$(p_uX+Z,W,Y,vX)$ with positive caps of $C,A$. Its strict residue
is the sum of positive multiples of
$(p_uX+Z)^2,(W+vX)^2,W^2,Y^2$; the other original packets are
retained. This handles the zero-$n$ degeneration at that corner.

For unbounded $v$ and bounded $u$, put $q_v=1/v$. The full capped
residue of $q_v^2P$ is
\[
 m_c(W+uY)^2+C(uW-Y)^2+(A_c+1/2)D_uY^2
                         +D_c(uX-Z)^2+D_uX^2/2\succ0.           \tag{RU29}
\]
Its common square kernel is zero for every bounded real $u$,
including zero. Finally, if both angles are unbounded, use
$p_u^2q_v^2P$ and the full frame
$(X,W,Y,P_Z=p_uZ,Q_Z=q_vZ)$. The exact residue is
\[
 CW^2+A_cY^2+D_c(X-P_Z)^2+X^2/2
                           +(m_c+1/2)(Y-Q_Z)^2\succ0.           \tag{RU30}
\]
The two nontrivial block determinants are $D_c/2$ and
$A_c(m_c+1/2)$, both positive. The entire polynomial has this
five-root support, with no inverse of $p_u$ or $q_v$ after
homogenization. All three formulas are checked from the original
parent. Their literal reciprocal zeros are polynomial values;
restoring a nonzero homogeneous factor uses a pure-testing square.

The charts exhaust independent testing pairs and states. Apply each
real identity to real and imaginary parts of a complex state, restore
every retained actual packet and separated clearing factor, and use
the exact sign gauges for both angle signs. This completes the stated
all-tests theorem without a restriction on any opening ratio.



\section{\texorpdfstring{The remaining small-$r$ local chart of weak-C2 R8}{The remaining small-r local chart of weak-C2 R8}}
\label{component:p482fcad1a2e2}
% Scoped mathematical certificate; see local_identity_audit.closure.json for exact changes.

\label{p482fcad1a2e2:sec:r17-r8-small-r}

\begin{theorem}
Retain the original physical stationary weak-C2 source, all strict gaps,
nonnegative event coefficients, the original positive-infinitesimal
total-scale intersection, and every inherited residual exclusion.
Suppose the source is in R8 and $r$ is infinitesimal. Thus $t=y/B$
and $g=G/B$ are positive bounded units, while $B,h=H/B,j,r$ are
positive infinitesimals. Then the entire original four-state parent is
a finite tensor SOS for every independent pair of infinitesimal or zero
testing half-angles $u,v$, with either sign and all real or complex states.
No source opening is compared with another opening or with a testing rate.
\end{theorem}

This is the missing local component of the R8 join. Its remote tests
and its $r$-unit local chart are supplied separately by the R17
joint-$D$ argument. The proof below does not rely on an unproved
extension of the RC normalization across $r=0$.

\paragraph{Source entry and the three independent hard scales.}
Write $a=z/B$, $l=1-h$, $x=aB$, and
\[
 j=1+xr-l/t,\qquad g=1+a(hB+r-1).
\]
The R8 and small-$r$ hypotheses imply that $a$ is bounded with residue
$\rho\in[0,1)$: $g_0=1-\rho>0$. Also $t_0=1$.
The unchanged exact source formulas WE1, jointly at $B=h=r=0,t=1$,
give
\[
 (nr)_0=(Ah)_0=(Dj)_0=\frac1{2(1+\rho)},\qquad
 \left(\frac{m}{a^2}\right)_0=\frac1{2(1-\rho^2)},
 \quad C_0=\left(\frac\alpha{B^2}\right)_0
       =\frac{1-\rho}{2(1+\rho)}>0.                 \tag{RD1}
\]
These are joint regular reductions. After replacing $n,A,D$ by
$nr,Ah,Dj$, all source denominators are products of units
$t,g,1-h,1-hB,\Delta$ with $\Delta_0=1+\rho$.
Thus $n,A,D$ are all positive unbounded, with completely unrestricted
relative rates, while $nr,Ah,Dj,\mu=m/a^2,C,\alpha_*=\alpha/B^2$
are positive bounded units. In particular $n=0$ is impossible on
this chart; no inverse of a vanishing event coefficient is being used.

\paragraph{A general exact centered-$D$ gauge of the complete WA matrix.}
Use the unnormalized complete WA3--WA8 identity, with the WE4 source
rescaling of the second scalar root:
\[
 \mathcal R=(\xi,\eta,R_1,R_2',T)^T,
 \quad \xi=X+ruZ,\quad\eta=Y-huW,
\]
\[
 R_1=uW-t^{-1}vZ,\quad R_2'=avW-uZ,\quad T=uvW-BZ,
 \qquad F/4=\mathcal R^TG\mathcal R.                \tag{RD2}
\]
Here $G$ contains the complete hard block, both complete first-normal
rows, and the complete scalar block. No packet or normal is omitted.
All its entries are affine in the formal event coefficient $D$.

For $U=u^2$, $V=v^2$, define three five-component columns
\begin{align*}
 d&=(uv,-1,l,ruv,-arv)^T,\\
 s&=(0,0,(x-U)v,-(B-V/t)u,U-aV/t)^T,\\
 k&=(-2rtu,0,0,-r^2tu,ar^2t)^T.
\end{align*}
The physical root relation and the centered $D$ event are exactly
\[
 s^T\mathcal R=0,\qquad
 L_D:=d^T\mathcal R
   =uW+uvX-Y-vZ+jvZ.                               \tag{RD3}
\]
The following is a matrix identity, valid before any residue operation:
\[
 \partial_DG-dd^T-\frac{sk^T+ks^T}{2}=jH_D,        \tag{RD4}
\]
where $H_D$ is symmetric and its only nonzero entries are
\[
 (H_D)_{13}=tuv,\quad (H_D)_{15}=-tu,\quad
 (H_D)_{34}=rtuv,\quad (H_D)_{35}=-artv.           \tag{RD5}
\]
Indices in RD5 refer to the displayed order in RD2 and start at one.
This explicitly resolves the failure of the ungauged matrix
$\partial_DG-dd^T$ to be divisible by $j$.

Let $G^{(0)}=G|_{D=0}$ denote a formal matrix substitution, with every
other source coefficient kept at its actual value. This is not a
physical modified source and no physicality theorem is applied to it.
Because the term involving $s$ vanishes exactly on the physical roots,
RD4 gives the entire-parent identity
\[
 F/4=D L_D^2+\mathcal R^T\{G^{(0)}+DjH_D\}\mathcal R.
                                                               \tag{RD6}
\]
The first-normal and hard-block entries in RD4 are part of this identity,
not merely its scalar restriction. No inverse of $u$, $v$, $j$, or a
state coordinate was used in deriving RD3--RD6.

\paragraph{The bounded $D$-free matrix through the $a$ opening.}
Normalize only the first two roots in $G^{(0)}$ initially, by
$\widehat\xi=\sqrt n\,\xi$ and $\widehat\eta=\sqrt A\,\eta$.
The resulting matrix, in
$(\widehat\xi,\widehat\eta,R_1,R_2',T)$, is a finite polynomial in
$u,v$ with bounded pure-source coefficients. Here is the complete
coefficient check needed for this assertion.

The hard entries after $D=0$ are
\begin{align*}
 h_{xx}/n&=(1+U)(1+V)+[m+UV/2+(1+V)/2]/n,\\
 h_{xy}/\sqrt{nA}&=-uv(m+1/2)/\sqrt{nA},\\
 h_{yy}/A&=(1+U)(1+V)
       +[C(1+V)+1/2+mUV+(\alpha+1/2)(1+U)V]/A.
\end{align*}
Every added coefficient is bounded and has zero residue.
In the first normal, the complete WA6 parameters now have the forms
\[
 p_0=t(nr+r/2),\quad p_1=t(nr+r/2),\quad
 q_1/a=nr+Bp_1,
\]
\[
 f_x=-mh+l/2-p_0-q_1,
 \quad (-m+q_1U)/a=-a\mu+(nr+Bp_1)U.
\]
They are bounded before multiplication by $1/\sqrt n$.
In the complete second normal, all WA7 coefficients involve only
the bounded products $Ah,m,C,\alpha,r,h$, and
\[
 q_y/a=-k_z-(Ah)/t-Bk_{uv}
\]
is bounded. Multiplication by $1/\sqrt A$ therefore bounds that
entire normal row as well. No ratio $n/A$, $A/n$, $h/r$, or $r/h$
appears.

For the scalar block use the exact WE5 expressions with $D=0$.
All occurrences of $n$ are through $nr^2=(nr)r$, and all occurrences
of $A$ are through $Ah^2=(Ah)h$. The only possible new $a$ denominator
is
\[
 Q_*=(b_2-\alpha_*)/a,\quad
 b_2=2a^2\mu h+C(1-h)^2+(Ah)h+(B^2\alpha_*+1/2)h^2. \tag{RD7}
\]
The exact WE1 expressions for $b_2$ and $\alpha_*$ involve no $D$
and no $j$ denominator. At $a=0$ their difference vanishes identically
for arbitrary $B,h,r,t$. Its cleared numerator is consequently
divisible by $a$; the canceled denominator at $B=h=r=0,t=1$ is
$2(a-1)(a+1)$, a unit for $0\le\rho<1$. The divided numerator
vanishes at that joint corner for every $a$. Hence $Q_*$ is bounded
(in fact infinitesimal), independently of $a,h,r,j$ rates.
This proves the bound on
\[
 f_s/a=\tfrac12[Q_*-a\,nr^2-aBt(n+1/2)r^2]
\]
and on the whole remaining scalar block. In particular the old
assumption that $j$ was a unit is unnecessary after the exact $D$
channel has been removed by RD6.

\paragraph{Simultaneous normalization of all three hard scales.}
Use the final five roots
\[
 \mathcal N=(\widehat\xi,\widehat\eta,\widehat D,R_2',T)^T,
 \qquad\widehat D=\sqrt D\,L_D.                    \tag{RD8}
\]
Solving only the source-unit coefficient $l=1-h$ in RD3 gives
\[
 R_1=\frac1l\left\{
 \frac{\widehat D}{\sqrt D}
 +\frac{\widehat\eta}{\sqrt A}
 -\frac{uv\widehat\xi}{\sqrt n}
 -ruvR_2'+arvT\right\}.                            \tag{RD9}
\]
Every coefficient in RD9 is bounded and separated. In particular,
independent small $r,h,j$ cause no ratio of their square roots.
Let $P(u,v)$ be the explicit matrix mapping $\mathcal N$ to
$\mathcal R$ according to RD9 and
$\xi=\widehat\xi/\sqrt n$, $\eta=\widehat\eta/\sqrt A$.
Then exactly
\[
 F/4=\mathcal N^T\widehat G\mathcal N,
 \qquad
 \widehat G=E_{33}+P^T\{G^{(0)}+DjH_D\}P.          \tag{RD10}
\]
The preceding bounds and RD9 prove that every coefficient of this
finite testing-polynomial matrix is bounded pure source. The factor
$Dj$ is a unit, while every entry of $H_D$ has positive testing degree.
Thus its contribution to the residue at $u=v=0$ is zero.

The normalized $n$ and $A$ hard diagonals have residue one. Since
$1/\sqrt n$, $1/\sqrt A$, $1/\sqrt D$, $r,h$ are infinitesimal,
the entire $R_1$ row of $P(0,0)$ has zero residue. The remaining
source scalar diagonals have residues $\mu_0$ and $\alpha_{*,0}=C_0$.
All cross entries have zero residue. Consequently
\[
 \overline{\widehat G}
   =\operatorname{diag}\left(1,1,1,
              \frac1{2(1-\rho^2)},
              \frac{1-\rho}{2(1+\rho)}\right)\succ0. \tag{RD11}
\]
This is the full five-root Gram residue, not a Hessian or state slice.

The positive pure-source matrix $\widehat G(0,0)$ has bounded inverse.
Whiten it purely in the source field and apply the finite separated
coefficient construction WA10 to the finitely many nonconstant
testing monomials. Their values are infinitesimal for arbitrary
independent infinitesimal $u,v$. This produces a finite SOS in
$K_c\otimes_{\mathbb R}K_t$ without any mixed inverse or mixed radical.
The argument includes either zero test and both signs because all
identities are polynomial in $u,v$. Applying it to real and imaginary
state parts proves the complex-state statement.



\section{\texorpdfstring{R8: a joint divergent-$D$ certificate, a paired-$D$ thin lift, and the remote join}{R8: a joint divergent-D certificate, a paired-D thin lift, and the remote join}}
\label{component:pe07041e033da}
% Scoped mathematical certificate; see local_identity_audit.closure.json for exact changes.


\paragraph{Exact inputs in this refined appendix.}
The generic source arrays and RC1/RC6 balances are printed in the native
weak-C2 dictionary. The complete generic identity BA4--BA7, with explicit
hard quadratic BA6, supplies the formerly cited RC7--RC13 algebra; it is
an exact identity before the bounded-$a$ source conditions are imposed.
The remote identities formerly called GR2--GR7 are reproduced as BA11,
the strict residue following BA11, BA12, BA13 and BA14. The double
homogeneous identity is BA15. Thus the old WC/R9 proof bodies and old
GR/BS source-domain proofs are not imported. The exact physical scale
identity BS1--BS3 is BA9. This section verifies every further bound
needed on its own source domain.

\label{pe07041e033da:sec:r17-R8-joint-D}

\paragraph{Scope and conventions.}
All identities use $F/4$ and the unchanged physical C2 parent, with
$\lambda_{\rm event}=\beta=0$, $d=1$, and $n\ge0$.  Every original
stationary-contact, nonnegative-coefficient, strict-gap, total-scale,
and prior-admission-complement intersection is retained.  The two real
closed source/testing fields are independent.  The symbol $\lambda$
below is the positive circuit coefficient, not an event coefficient.

R8 means $B,z,y,h,j>0$ infinitesimal, $t=y/B$ and $g=G/B$ positive
bounded units, $h=H/B$, $a=z/B$, $r\in(0,1)$, and $e=1-r$.  Put
$x=z=aB$, $l=1-h$.  The exact gap is
\[
                 j=1+xr-l/t.                         \tag{JD1}
\]
In particular $t_0=1$.  This section proves:
\begin{enumerate}
\item every independent test and state on the whole R8 subdomain
      where $r$ is a positive bounded unit;
\item the entire remote testing complement of both angles
      infinitesimal on all of R8, including $r$ infinitesimal.
\end{enumerate}
The simultaneous $r,h,j$-infinitesimal local domain is proved by the
separate joint-$n,A,D$ lemma in
Section~\ref{component:p482fcad1a2e2}; it is not
inferred from the present $r$-unit proof.  Its union with the two
assertions above closes the whole ordered R8 shape.  No R1 or R2
admission is asserted here.

The exact inputs are the native source arrays and RC1/RC6 contact
balances in Section~\ref{component:pd709027c7805}, the generic full-parent
identity BA4--BA7, R10.10--R10.12, and TL1--TL10 with the complete
support calculations printed in that thin-layer proof. The retained
CA/CS/SE contact conductors and critical-wall identities SEO10--SEO20
are used with the bounds verified below. No physical C1 theorem is
invoked at a negative formal beta.

\paragraph{A single contact-zero row removes every naked $D$ pole.}
Use the homogeneous R15/R16 roots
\[
 E=tuW-vZ,\quad N_y=(Y-huW)/\sqrt h,\quad \xi=X+ruZ,
 \quad R_D=v(u^2-x)Z,\quad F_v=(v^2-y)Z.
\]
The other roots $Q,N,T,R,J$ and source coefficients are exactly those
of R10.9.  Define
\[
 L_D={l\over t}E-rR_D-\sqrt hN_y+uv\xi,
 \qquad M_D=\sqrt D\,L_D.                            \tag{JD2}
\]
The original event row satisfies the exact homogeneous identity
\[
 uW+uvX-Y-vZ=L_D-jvZ.                               \tag{JD3}
\]
Thus $L_D$ vanishes at the actual source contact, including its full
state normal.  This is a legal separated root at every testing axis.

Let $\mathscr P$ be the right side of R10.9 with the two exact
positive squares $N^2,T^2$ omitted, and let $\mathscr P_0$ denote
its formal coefficient specialization $D=0$.  The other source
letters are held fixed in this notation; it does not change the
physical source.  Direct collection of the complete $D$ channel gives
\[
 \boxed{\quad \mathscr P=D L_D^2+\mathscr P_0
                        -2Dj\,uF_v\xi.\quad}        \tag{JD4}
\]
Here is the full coefficient calculation.  The hard contribution is
$D(lE/t-\sqrt hN_y+uv\xi)^2$; the remaining $D$ contributions are
$Dr^2R_D^2-2Dr(lE/t-\sqrt hN_y+uv\xi)R_D$ and
$2D\{l/t-1-rx\}uF_v\xi$.  Their sum is exactly JD4 by JD1.
This accounts for every hard and first-normal term, rather than just
the angular graph or its Hessian.

Use the polynomial inverse
\[
 E={t\over l}\{D^{-1/2}M_D+rR_D+\sqrt hN_y-uv\xi\}. \tag{JD5}
\]
On R8, $t/l$ is a unit, $D^{-1/2}$ is infinitesimal, and $Dj$ is
bounded.  Every coefficient of the new remainder is bounded whenever
the remaining source arrays of R10.9 are bounded.  In particular no
ratio $h/j$, $x/j$, or testing angle divided by a source opening occurs.
The full divergent rank-one block has been normalized jointly.

\paragraph{The exact residue after the joint normalization.}
Put $\sigma=t\lambda x/\sqrt q$ and $\mu=m-a^2q$ as in R10.9.
At the R8 source/testing corner, JD4--JD5 have the complete residue
\[
 \begin{split}
 M_D^2+r_0^2(C_0+n_0)R_D^2+(Ah)_0N_y^2
  +(n_0+\tfrac12-\mu_0)\xi^2
  +(c_F)_0F_v^2+\lambda_0J^2+R^2+Q^2-2\sigma_0Q\xi.
 \end{split}                                                   \tag{JD6}
\]
This formula applies when $r,n$ are bounded and $r$ a unit.  Its
derivation can retain $n,C,Ah,\mu,c_F,\lambda,\sigma$ as free source
letters: specialize $u=v=\sqrt h=D^{-1/2}=x=\alpha=0$ and $t=1$
in JD4--JD5.  In particular the $E,R_D,N_y$ terms become
$r^2(C+n)R_D^2+AhN_y^2$, with no missing $D$ contribution.
The separate $N^2,T^2$ remain intact.

\paragraph{Bounded $a$, unit $r$, and unit $e$.}
Positivity of $g=1-ae+hx$ makes $a$ bounded whenever $e$ is a unit.
The unchanged source formulas, with $B=h=0$ and $t=1$ only after
clearing unit denominators, give
\[
 \begin{gathered}
 n_0={e_0\over2r_0(1+a_0)},\quad
 (m/a^2)_0={e_0\over2g_0(1+a_0)},\\
 C_0={1+a_0(2r_0-1)\over2(1+a_0)},\quad
 (Ah)_0=(Dj)_0={1+a_0r_0\over2(1+a_0)},\\
 \lambda_0={r_0(1+a_0r_0)\over2(1+a_0)},\quad
 (c_F)_0={(1+a_0r_0)g_0\over2(1+a_0)}.                
 \end{gathered}
\tag{JD7}
\]
All displayed pivots are positive units.  Indeed
$1+a_0(2r_0-1)=g_0+a_0r_0>0$.  The source cancellation defining
$\lambda=(m-a^2q)/(ta^2B)$ is regular also at $a=0$; no division
by an infinitesimal $a$ survives its canceled formula.
Here $q$ is a unit, $\mu=ta^2B\lambda$ is infinitesimal and
$\sigma$ is infinitesimal.  Therefore JD6 is strictly positive.
All omitted entries have zero residue by the finite polynomial
JD4--JD5 and bounded source arrays.  The bounded separated-Gram
lemma gives the complete local parent for every independent ratio
$u/v$, every state, and both literal axes.

\paragraph{Bounded $a$, infinitesimal $e$, including $n=0$.}
Now $r_0=1$, $C_0=(Ah)_0=1/2$, $n,m,\mu$ are infinitesimal,
$\lambda_0=(c_F)_0=1/2$, and
\[
                    (\alpha/B^2)_0=(1+a_0)/2.         \tag{JD8}
\]
The physical inequality OP2 supplies the crucial source lower bound,
without dividing by $n$:
\[
 {e\over B}\ge
 {(1+x)(h+t)(1+ar)\over P_n/B},\qquad P_n/B\sim1.
                                                               \tag{JD9}
\]
It remains valid when $e/B$ is unbounded and at literal $n=0$.
Since $q\ge\alpha+e^2/2$, the $Q,\xi$ Schur pivot satisfies
\[
 \begin{split}
 \operatorname{st}\left(\tfrac12-\mu
                 -{t^2\lambda^2a^2B^2\over q}\right)
 &\ge {1\over2}-{a_0^2\over2(1+a_0)(2+a_0)}\\
 &= {2+3a_0\over2(1+a_0)(2+a_0)}>0.                  
 \end{split}\tag{JD10}
\]
This is a bound entirely in the source field.  The entry $\sigma$
is bounded by JD8, so the whole block is bounded and strict.
The remaining entries of JD6 are positive units.  This proves all
local tests and all states on this source domain, including $n=0$.

\paragraph{Unbounded $a$: the thick layer.}
Put $s=1/a$, $w=ae$, $b=a^2B$.  The exact R10 source identity
still reads
\[
 n=s^2{wQ_*-bK_*\over2r\Delta_*},\qquad
 Q_*\sim1,\ K_*\sim t,\ \Delta_*\sim1.             \tag{JD11}
\]
Thus $b/w$ is bounded, $b$ is bounded, $g_0=1-w_0>0$, and
\[
\begin{gathered}
 (m/w)_0={1\over2(1-w_0)},\quad C_0=(Ah)_0=1/2,
 \quad (Dj)_0=1/2,\\ \lambda_0=1/2,
 \quad(c_F)_0=(1-w_0+b_0)/2>0.
\end{gathered}                       \tag{JD12}
\]
The source formulas for these quantities have no $j$ denominator;
only $D$ itself is unbounded, and JD4 has removed that pole.
For $\rho=n/e^2$ positive noninfinitesimal, retain the exact R10.10
identities and R10.11 Schur pivot:
\[
 \bar q=a^2q=w^2\{\rho r^2+H\}+a^2\alpha,
 \quad H=m+\tfrac12,\quad \chi=m/w,
 \quad \chi-H\in\operatorname{Inf},
\]
\[
 S=\tfrac12-\mu-{(t\lambda x)^2\over q}
 =H-{m^2\over\bar q}
 \ge {H\rho r^2\over\rho r^2+H}
       +{H^2-\chi^2\over\rho r^2+H}.                 \tag{JD13}
\]
The first term has positive residue and the second is infinitesimal,
also for unbounded $\rho$.  Moreover $x/\sqrt q=b/\sqrt{\bar q}$
is bounded by JD11 and $\bar q\ge Hw^2$.  JD6 consequently has
strict residue on $(M_D,R_D,N_y,\xi,F_v,J,R,Q)$, with the only
possible nonzero cross residue in its strictly positive $Q,\xi$
block.  Restore $N^2,T^2$.  This proves the full local parent, not
just its contact-state restriction.

\paragraph{Unbounded $a$: the thin layer and the exact paired-$D$ repair.}
Assume $0<\rho=n/e^2$ infinitesimal.  JD11 gives
$b/w\sim1/t$, so $\psi=w/b$ is a unit.  The formal C1 contact is
exactly OP4, with native parameters
\[
 x_1=x,\quad\eta_1=t,\quad\zeta_1=s,\quad r_1=h,
 \quad\theta_1=e,\quad t_1=w.
\]
TL2 still gives $\beta_1=-s_n/2$ infinitesimal and
$\kappa=(1+2n)/(1-s_n)\sim1$.  The native gaps satisfy
$M_1=g$ a unit and $J_1=j>0$.  The canceled CA/CS/SE remainders
have no inverse $J_1$; their divergent paired-$D$ mass is retained
in SD3.  The formal source has $C_1,A_1h$ positive units,
$D_1j$ a positive unit, positive $m_1,\alpha_1$, and positive
$\beta_1+1/2$.  These facts follow directly from TL3 and JD12.
Its bounded opening ratio $t_1/e_{1,\mathrm{scale}}^2=w/(tb)$
is a unit.  Thus the explicit CA/CS/SE strict base and their finite
conductors apply exactly as in the R16 thin proof; the reservation
$(2+4\beta_1)\nu-\nu$ is positive because $\beta_1$ is
infinitesimal.  This is an algebraic use of those formulas, not an
application of a physical C1 theorem with negative beta.

Here is the one change needed to the R16 perturbation.  In native
notation put
\[
 T_1=c_a=1+xr,\quad a_1=l/t,\quad
 b_D={J_1\over T_1}={j\over c_a},\quad
 h_1={a_1\over T_1}=1-b_D,
\]
\[
 q_B=v(Z-uX),\quad M_B=(uW-Y)-h_1q_B,
 \quad R_{D,\mathrm{event}}=M_B-b_Dq_B.               \tag{JD14}
\]
The old retained paired-$D$ packet in $\kappa F_1/4$ is
$\kappa(D_1+1/2)h_1M_B^2$.  In the actual parent retain instead
the positive contact-zero packet $(D_2+1/2)h_1M_B^2$.
For $\delta_D=(D_2+1/2)-\kappa(D_1+1/2)$ the entire resulting
difference is
\[
 \boxed{\quad
 \delta_D\{R_{D,\mathrm{event}}^2-h_1M_B^2\}
 =\eta_D\{M_B^2-2M_Bq_B+b_Dq_B^2\},\quad
 \eta_D=\delta_D b_D.\quad}                          \tag{JD15}
\]
By the exact TL3 coefficients, with $\Theta=n/B$,
\[
 \eta_D={\Theta\{S_cP_D/t-jP_\kappa\}
                    \over K(h+t)c_a},\qquad
       \eta_D/\rho=O(w).                            \tag{JD16}
\]
Every denominator is a unit.  In particular the previously unbounded
$\rho w/j$ correction is absent; no relative $w/j$ bound is needed.

Keep the R16 positive contact-zero $n$ row TL4, the actual centered
$A$ packet, and every other unchanged retained packet.  The exact
full stress is TL6 with its $\delta_D R_{D,\mathrm{event}}^2$
line replaced by the right side of JD15.  Both parents and all
retained packets have zero contact value and full first jet, so this
new entire stress has the same exact vanishing.

For completeness its new term satisfies every finite-support premise
of the inherited conductor.  The exact inverse state formulas give
\[
 q_B=v\{D_u-uN-puD_uv\mathfrak w\},\qquad p=t_1/R_1.
                                                               \tag{JD17}
\]
At zero normals $q_B=vD_uV_0$.  Hence the new scalar term is in
$(v)^2$, both complete first normals are in $(v)$, and the remaining
hard entries are bounded.  There are no terms above normal degree
two.  Every dependence on $\mathfrak w$ has a factor $p$.
In the CA chart JD16 divided by $\rho$ has bounded coefficients
on the same raw roots as TL8.  In the SE chart
\[
 \Theta/\rho=\varepsilon^2\tau^2QE/K
\]
gives the same extra $\varepsilon^2$ factor to every new coefficient.
Modulo $\varepsilon$, scalar coefficients are even in $u,v$ and
the normal coefficients are independent of $\mathfrak w$.
Thus the exact support, full-jet, and finite divisibility hypotheses
of CC3--CC6 and SE5--SE8 are unchanged.  The complete stress divided
by $\rho$ has a bounded Gram on the same eight roots.  Multiplication
by $\rho$ makes it an infinitesimal perturbation of the strict formal
C1 base.  Restoring the positive packets proves the actual C2 parent
on the entire thin contact chart, including both full state normals.

At literal $n=0$, OP4 is a physical C1 identification with
$k_1,\zeta_1,r_1$ infinitesimal, $\eta_1=t$ a unit and
$e_{1,\mathrm{scale}}^2=tb$ bounded.  The physical C1.2/C1.3
all-tests theorems allow arbitrary positive $J_1$ and literal beta
zero, so they cover this boundary exactly.

\paragraph{The unit-$W$ entry on the unbounded-$a$ thin layer.}
The actual physical C2 critical coefficient is
\[
 C_*={(h+2Ah)(1+2D)\over2(4AhD-h)}.
\]
Use the polynomial UW3 numerator/denominator with $b=wz_*$,
$z_*=b/w$ bounded.  Its numerator is divisible by $sw$ because it
vanishes identically at $s=0$ and $w=0$.  The denominator after
canceling the $m$ numerator has residue $16t_0=16$, also when
$j_0=0$.  Consequently
\[
                   (C_*-C)/m=s\times O(1).            \tag{JD18}
\]
There is no $j$ inverse in this estimate.  SEO12 permits both $A,D$
unbounded: its entire hard block remains strict and its two radial
masses tend to $2,2$.  With $W=1$, the exact testing identity
$vZ=Z(R_m+X+uZ)/(1+uY)$ absorbs the full negative critical correction
by JD18 and the retained original $m$ square.  All denominators are
pure testing units.  This proves every local normalized unit-$W$
state, including $Z=0$; negative-core state normalization leaves
exactly the $Z=1$, $u,v,\mathfrak w>0$ infinitesimal chart proved
above.  This is the same actual-C2 state entry as UW, with its
source bound checked at the new $t_0=1$ boundary.

\paragraph{Finite remote tests when $e$ is infinitesimal.}
For bounded $a$, $C_*-C$ is infinitesimal because $C_*,C\sim1/2$.
For unbounded $a$ this follows also from JD18 and bounded $m$.
The full critical-wall identities SEO17, SEO18, SEO20, in the
normalized SEO12 frame, are strict whenever at least one finite
angle has nonzero residue.  They explicitly allow $A,D$ unbounded.
The entire correction $-4(C_*-C)_+D_v(uW-Y)^2$ has infinitesimal
bounded entries in that frame.  Retain original $n,m,\alpha$ and
any positive excess.  This proves every finite remote test, with
all states, for the entire $e$-infinitesimal R8 domain.

\paragraph{Finite remote tests when $e$ is a unit, including small $r$.}
Here $a$ is bounded, $m,C$ bounded, $C$ a positive unit, $A,D$
unbounded, and $N_*=n+1/2$ is either a unit or unbounded.  All
normalizations below are pure source.  Let
$\widehat X=\sqrt{N_*}X$, $\widehat Y=\sqrt A Y$, and
$\widehat D=\sqrt D\,R_{D,\mathrm{event}}$.

If $u$ is a bounded unit and $v$ bounded, use the exact inverse
\[
 W={D^{-1/2}\widehat D+A^{-1/2}\widehat Y+vZ
                         -uvN_*^{-1/2}\widehat X\over u}.
                                                               \tag{JD19}
\]
All full coefficients are bounded.  The $\widehat Y,\widehat D$
residue pivots are $(1+U)(1+V),1$, with zero cross residue, where
$U=u_0^2>0,V=v_0^2\ge0$.  For unit $r$, the remaining two-state
matrix before $X$ normalization is the complete original parent at
$Y=R_{D,\mathrm{event}}=0$, namely
\[
 \begin{split}
 &(nD_uD_v+mD_v^2+CD_vUV+D_v/2)X^2\\
 &+2\{-mD_v(V-U)/u-CD_vuV+uD_v/2\}XZ\\
 &+\{m(V-U)^2/U+CD_vV+(U-V)/2\}Z^2.
 \end{split}                                                   \tag{JD20}
\]
Insert the residue source values JD7, with $e=1-r$, $g=1-ae$.
Its determinant has the exact positive expression
\[
 \begin{split}
 U\det={D_uD_v(1+ar)\over4r(1+a)^2g}
 \{e(U-aV)^2+r(g+ar)U^2V+g(g+ar)UV^2\}.             
 \end{split}\tag{JD21}
\]
It is positive for every $U>0,V\ge0$.  The $X$ diagonal is
positive.  Division by the positive unit $N_*$ proves strictness
after normalization.  When $r$ is infinitesimal, divide JD21 by
$N_*\sim n$ before taking residue.  The result is
\[
 U\det_{\rm normalized}
 ={D_uD_v\over2(1+a_0)(1-a_0)}
       \{(U-a_0V)^2+(1-a_0)^2UV^2\}>0,             \tag{JD22}
\]
with $0\le a_0<1$.  The normalized $X$ diagonal tends to $D_uD_v$.
Thus JD19 proves this remote chart for arbitrary relative $n,A,D$
rates, without a comparison between any of them.

It remains to take $u$ infinitesimal and $v$ a bounded unit.  If
$a$ is a unit, $m$ is a unit.  Solve instead for $Z$ by the pure
testing unit $v$.  The full residue is
\[
 m_0(v_0W-N_{*,0}^{-1/2}\widehat X)^2
       +D_v\widehat X^2+D_v\widehat Y^2+\widehat D^2,
                                                               \tag{JD23}
\]
where $N_*^{-1/2}$ has residue zero if $n$ is unbounded.  This is
strict, with the two-state determinant $m_0V D_v>0$.
If $a$ is infinitesimal, use the extra legal root $S=aW$ and put
$R=uW$.  The exact D event gives
$R=vZ-uvX+Y+D^{-1/2}\widehat D$.  Substitute this expression
in every occurrence except the $m$ event, where use
$mR_m^2=(m/a^2)(vS-aX+auvY-auZ)^2$.
All coefficients are bounded, also when $n$ is unbounded after
the displayed $X$ normalization.  Its entire residue is diagonal
on $(S,\widehat X,\widehat Y,\widehat D,Z)$ with pivots
\[
                    ((m/a^2)_0V,D_v,D_v,1,V^2/2)>0.   \tag{JD24}
\]
These identities retain every original event and the complete signed
core; introducing a dependent extra root is harmless.  The finite
separated lift around the real angle residues proves each chart.

\paragraph{Projective endpoints, signs, and the complete join.}
On all of R8, $C$ is a positive unit and $A,D$ are unbounded.
For infinity charts retain the whole original $m,n,\alpha$ packets,
clip $C,A,D$ at suitable positive real lower bounds, and retain their
positive excesses.  The exact homogeneous infinity identities used
in SEO19 and its accompanying BE2 routing require only these lower
bounds and $d=1$; they do not require $D$ bounded or $m$ bounded
below.  In particular the delicate $(u,1/v)=(0,0)$ corner has the
strict selected residue
\[
 2X^2+4c(uW-Y)^2+(4A_*+2)Y^2+4D_*Z^2
\]
on $(X,Y,uW,Z)$ after retaining $m$.  All leftover $W$ occurrences
carry $u$.  The remaining single and double infinities are exactly
the same homogeneous event identities.  Literal endpoint values are
taken in those original homogeneous formulas; no higher-degree
contact graph is substituted at infinity.

For unit $r$, source $a$ is bounded or unbounded.  The bounded case
has $e$ unit or infinitesimal and is covered locally by JD7 or
JD8--JD10.  The unbounded case forces $e$ infinitesimal and is
covered by the thick/positive-thin/literal-zero split above.  These
are exhaustive pure-source alternatives.  All finite remote tests
use JD18--JD24 or the critical wall, and every infinity uses the
homogeneous join.  Original testing sign gauges, literal-axis event
identities, state homogenization, and the real/imaginary split give
all signs and complex states.  All inverses and radicals are pure
source or pure testing.  No independent source/testing rate is ever
compared.

This proves the entire R8 $r$-unit source subdomain all tests and
the remote complement for every R8 source.  If $r$ is infinitesimal,
$e$ is a unit and $g>0$ forces bounded $a$ with $a_0<1$; precisely
these hypotheses, with every other original intersection unchanged,
are those of the joint-$n,A,D$ local lemma.  That lemma supplies a
strict full-parent Gram with residue $(1,1,1,(m/a^2)_0,C_0)$ on
all remaining local tests, including the axes.  Therefore the two
source alternatives $r$ unit or infinitesimal now exhaust R8 with
an all-tests certificate.  Literal $n=0$ was retained in the unit-$r$
proof; it is impossible in the small-$r$ part because $nr$ is a unit.
The whole ordered R8 shape is closed, with no extra source or testing
rate restriction.


\section{\texorpdfstring{R10: forced source scale, every remote test, and a weighted local admission}{R10: forced source scale, every remote test, and a weighted local admission}}
\label{component:p9ee4ba7326a7}
% Scoped mathematical certificate; see local_identity_audit.closure.json for exact changes.


\paragraph{Exact inputs in this refined appendix.}
The generic source arrays and RC1/RC6 balances are printed in the native
weak-C2 dictionary. The complete generic identity BA4--BA7, with explicit
hard quadratic BA6, supplies the formerly cited RC7--RC13 algebra; it is
an exact identity before the bounded-$a$ source conditions are imposed.
The remote identities formerly called GR2--GR7 are reproduced as BA11,
the strict residue following BA11, BA12, BA13 and BA14. The double
homogeneous identity is BA15. Thus the old WC/R9 proof bodies and old
GR/BS source-domain proofs are not imported. The exact physical scale
identity BS1--BS3 is BA9. This section verifies every further bound
needed on its own source domain.


\paragraph{Scope and result.}
Retain the original weak-C2 R10 physical stationary source, every strict
source gap, $n\ge0$, the positive-infinitesimal total-scale intersection,
and every inherited exclusion. Thus $t,g,j$ are positive bounded units,
$B,h,e$ are infinitesimal, $a$ is unbounded, and $x=aB$ is infinitesimal.
The new all-tests admission proved here is the source subdomain
\[
                         n/e^2\notin\operatorname{Inf}.
\tag{R10.1}
\]
Every nonlocal independent testing chart is also proved on all of R10.
The positive-$n$, infinitesimal-$n/e^2$ local complement is handled by the following thin-layer and state-entry sections; it is not covered by this weighted pivot.
The literal $n=0$ boundary is already an inherited C1 admission, by the
exact map specified below. These statements retain all independent rates.

\paragraph{The decisive exact source scale.}
Use the pure-source variables
\[
 s=1/a,\qquad w=ae,\qquad b=a^2B=x/s,
 \quad B=s^2b,\quad x=sb,\quad e=sw,\quad r=1-sw.
\]
Initially $b$ is not assumed bounded. The gap $g=1-w+hx$ makes $w$
bounded and $1-w$ a positive unit in residue. In the exact positive
polynomials $Q_s,K_s$ printed in BA9 put
\[
 Q_*=1+B(1+t+ah+2ar+at)
 +B^2\{a(h+t)+a^2[h(r+1)+t(1+r-r^2)]\},
\]
\[
 K_*=(1+s)(h+t)(1+x),\qquad \Delta_*=s\Delta.
\]
Before imposing a bound on $b$, $Q_*\sim1$ and $\Delta_*\sim1$:
write them using bounded $w$, infinitesimal $s,x,h$; their only
nonconstant monomials contain $s$ or $x$. BA9 gives exactly
\[
 \boxed{\quad 2r\Delta_*\,n/s^2=wQ_*-bK_* .\quad}
\tag{R10.2}
\]
Consequently
\[
 0<b/w\le Q_*/K_*,\qquad
 0\le n/s^2\le wQ_*/(2r\Delta_*).
\tag{R10.3}
\]
In particular $b=O(w)$ is bounded, $n=O(s^2w)$, $an$ is
infinitesimal, and $ck=\sqrt{b/t}$ is bounded for
$c=\sqrt{a/t}$, $k=\sqrt{x}$. This is an original-source bound,
including equality $n=0$, not a comparison with either testing angle.

After this bound every displayed source denominator is a unit in
$(s,b,w,h,t)$ except the explicit positive monomials already canceled.
Direct finite rational cancellation in the source formulas gives
\[
 \begin{gathered}
 m/w\sim\frac1{2(1-w)},\quad
 C\sim Ah\sim t/2,\quad D\sim\frac1{2(t-1)},\quad
 \alpha/(s^3b^2)\sim t/2,\\
 \lambda=\frac{m-a^2q}{ta^2B}\sim\frac12,\qquad
 c_F\sim\frac{t-w+tb}{2t^2}>0,\qquad
 q=nr^2+(m+\tfrac12)e^2+\alpha.
 \end{gathered}
\tag{R10.4}
\]
The last strictness follows from $t_0>1$ and $w_0<1$.
These are joint residue identities in bounded letters; they do not
set $w$ or $b$ to a positive unit. In particular they hold when either
is infinitesimal. The cancellation defining $\lambda$ leaves no $b$
or $s$ denominator; the new checker verifies this explicitly.
We also have $m$ bounded, $n,\alpha,A^{-1/2}$ infinitesimal, and
\[
 d_4\sim\frac{t}{2(t-1)},\quad b_4\sim0,\quad
 a_E\sim-\frac1{t-1},\quad x_3\sim-1.
\tag{R10.5}
\]
All first-normal arrays printed in
Section~\ref{component:pd709027c7805} are bounded:
the only newly unbounded letter $a$ occurs there through $an$ or
$aq$, both infinitesimal. The factors $a^2n$ in $c_F$ are bounded.

\paragraph{All remote tests on all of R10.}
The exact full-parent identities printed with BA11--BA15 and the
literal-axis formulas following BA10 apply algebraically. The
bounded-$a$ proof used $m$ infinitesimal; under R10.4 its entire
nonnegative event packet is retained. In BA11 this adds a nonnegative
rank-one form to the strict six-root residue displayed immediately
after that identity; all its coefficients are bounded. For the chart
$u$ infinitesimal, $v$ a bounded unit, keep $S=\sqrt m W$.
The complete $(S,X)$ residue is now
\[
 (v_0S-\sqrt{m_0}X)^2+\frac{1+v_0^2}{2}X^2,
 \qquad \det=\frac{v_0^2(1+v_0^2)}2>0.
\tag{R10.6}
\]
The $(uW,Z)$ block remains BA12 and the normalized $Y$ pivot remains
strict; the extra $uvY,uZ$ terms in the $m$ packet are infinitesimal.
For unbounded $v$ its reciprocal identity has the same diagonal
residue BA13 because the potential $\sqrt m X$ term carries $1/v$.
On the remaining remote charts retain $mR_m^2$ as a full positive
square and apply the unchanged strict residue to the remainder.
The original homogeneous identities include either literal infinity,
and the literal-axis identities following BA10 retain $n=0$ without a division. The original
sign gauges and real/imaginary split are unchanged. Hence every test
outside both angles infinitesimal is covered for the entire R10 source.

\paragraph{The exact local full-parent identity.}
Use the generic contact balances RC1, RC6 in
Section~\ref{component:pd709027c7805} and the complete identity
BA4--BA7. These are rational identities of the original source and
never require bounded $a$. Set $U=u^2,V=v^2$, $y=tB$,
\[
 E=tuW-vZ,\quad N_y=(Y-huW)/\sqrt h,\quad \xi=X+ruZ,
\]
\[
 F_v=(V-y)Z,\quad R_D=v(U-x)Z,\quad J=(uv-x/c)Z,
\]
\[
 Q=\sqrt q(avW-uZ-a\xi),\qquad
 N=\sqrt n\{r[(a/t)V-U]Z+u\xi\},
\]
\[
 R=\sqrt\lambda\,tk(cv-u)W,\qquad
 T=\sqrt{2\lambda/c}\,k\sqrt{uv}(tcW-Z),\quad
 \mu=m-a^2q=ta\lambda x.
\]
At positive nonzero angles these are exactly the state-linear roots
in BA4, with no testing-ratio inverse. The entire hard quadratic
$\mathcal H_0$ is the explicit expression BA6. It has no inverse
testing angle or $a$ denominator; after $\eta=\sqrt hN_y$, every
occurrence of $A$ is through $Ah$, $\sqrt hAh$, or $hAh$.
Its R10 local residue is
\[
 \operatorname{diag}\left(\frac1{2(t_0-1)},\frac{t_0}2,\frac12\right)
 \quad\hbox{on }(E,N_y,\xi).
\tag{R10.7}
\]
This permits bounded $m$; every occurrence of $m$ in the remaining
hard entries carries a positive testing degree or $h$.

The principal angular--$\xi$ coupling is resolved as a whole.
Writing $R_c=kv(tcW-Z)$, direct multiplication gives
\[
 -2t\lambda ck R_c\xi+2t\lambda c(R_D-uJ)\xi
 =-2t\lambda x Q\xi/\sqrt q-2\mu\xi^2.
\tag{R10.8}
\]
Thus neither $cu$ nor $cR_D$ is left as a Gram coefficient. The
complete generic identity BA7, with the source arrays just specified,
gives the entire original $F/4$ exactly as
\[
\begin{split}
 &Q^2+N^2+T^2+R^2+\lambda J^2+d_4R_D^2
 +(c_F+b_4x/t^2)F_v^2+(b_4/t^2)vR_DF_v\\
 &\quad+\mathcal H_0-\lambda xE^2-\mu\xi^2
                  -2(t\lambda x/\sqrt q)Q\xi\\
 &\quad+E\{(a_E+c_EV)R_D+(b_E+c_Ex)vF_v\}\\
 &\quad+\sqrt hN_y\{(a_Y+c_YV)R_D+(b_Y+c_Yx)vF_v\}\\
 &\quad+\xi\{x_5uvR_D+u(x_3+x_5x)F_v\}.
\end{split}
\tag{R10.9}
\]
All coefficients are the exact source arrays in
Section~\ref{component:pd709027c7805} and BA5.
This identity retains every scalar, first-normal and hard-normal term.

\paragraph{A strict weighted pivot when $n/e^2$ is noninfinitesimal.}
Put $\rho=n/e^2$, $H=m+1/2$, $\chi=m/w$, and
$\overline q=a^2q$. Then
\[
 \overline q=w^2(\rho r^2+H)+a^2\alpha,
 \qquad \chi-H\in\operatorname{Inf}.
\tag{R10.10}
\]
The $Q,\xi$ source block in R10.9 has Schur pivot
\[
 S=\frac12-\mu-\frac{(t\lambda x)^2}{q}
   =H-\frac{m^2}{\overline q}.
\tag{R10.11}
\]
The equality uses $a^2q+\mu=m$ exactly. Since $\alpha\ge0$,
\[
 S\ge H-\frac{\chi^2}{\rho r^2+H}
 =\frac{H\rho r^2}{\rho r^2+H}
    +\frac{H^2-\chi^2}{\rho r^2+H}.
\tag{R10.12}
\]
Here $H$ is a positive bounded unit. Under R10.1 the first summand
has positive residue and the second is infinitesimal, also when
$\rho$ is unbounded. Hence $S$ is a positive bounded unit.
Moreover $b/w$ is bounded and $\overline q\ge Hw^2$, so
$x/\sqrt q=b/\sqrt{\overline q}$ is bounded. All entries of this
strict source block are therefore bounded.

In the rest of R10.9 the only nonzero off-diagonal residue is
$a_E$ between $E$ and $R_D$, giving
\[
 \frac1{2(t_0-1)}
 \begin{pmatrix}1&-1\\-1&t_0\end{pmatrix}\succ0.
\]
The $N_y,F_v,J,R$ pivots have positive residues
$t_0/2,(t_0-w_0+t_0b_0)/(2t_0^2),1/2,1$.
Every other entry is infinitesimal by R10.4--R10.5 and the explicit
positive testing degrees in R10.9. Retain $N^2,T^2$ as exact positive
squares. The bounded-Gram lemma supplies finite tensor SOS for the
entire parent on R10.1. No source/testing rate comparison occurs.
All roots homogenize without $Z$ or testing-ratio inverses, so literal
axes and $Z=0$ are included directly. The original signs, complex-state
split and the remote theorem complete the all-tests admission.

\paragraph{Literal $n=0$ and the exact remaining complement.}
At $n=0$, OP4 is a physical C1 identification with
\[
 x_1=x,\quad\eta_1=t,\quad\zeta_1=s,\quad r_1=h,
 \quad\theta_1=e,\quad e_{1,\mathrm{scale}}=\sqrt{tb}.
\]
Thus $k_1,\zeta_1,r_1$ are infinitesimal, $\eta_1$ a unit, and
$e_{1,\mathrm{scale}}$ bounded. Its infinitesimal type is covered by
Section~\ref{component:p29aeda9b68a4};
its positive-unit type is covered by the neighboring
Section~\ref{component:pccaa793fd0bf}. Both retain $\beta=0$.
This is C1.2/C1.3 routing, not the unbounded-$\eta$ C1.1 theorem.
For $n>0$ OP5 gives formal $\beta_1<0$, so this physical C1 theorem
cannot be extended by continuity or imported into the neighborhood.
The complement left for the companion thin-layer proof is exactly
\[
       0<n/e^2\in\operatorname{Inf},\qquad u,v\in\operatorname{Inf},
\]
with every original physical and total-scale intersection retained.
In this complement R10.11 can have zero residue: at a unit $w$ and
$b_0=w_0/t_0$, its residue is zero. Thus R10.9 is a useful full-parent
identity there, but its stated strict-Gram argument does not prove that
complement. The companion thin-layer proof and direct state join complete this complement; the final join states the whole-R10 theorem.


\section{\texorpdfstring{The positive-$n$ thin R10 layer: an exact extra row and a bounded conductor perturbation}{The positive-n thin R10 layer: an exact extra row and a bounded conductor perturbation}}
\label{component:pef37475d831e}
% Scoped mathematical certificate; see local_identity_audit.closure.json for exact changes.


\paragraph{Domain.}
All polynomial identities below use the convention $F/4$, as in
R10.9; in particular the original $n$ packet is $nD_uD_vX^2$.
All parent and packet identities in this section use the common normalization $F/4$.
Keep every original R10 condition and assume
$\rho=n/e^2>0$ infinitesimal. Use $s,w,b,x,B$ from R10.2.
Dividing that exact identity by $w$ gives
\[
 Q_*-(b/w)K_*=2r\Delta_*\rho w,
 \qquad b/w\sim1/t.
\tag{TL1}
\]
Thus $\psi=w/b$ is a positive bounded unit. No testing rate occurs.
The formal native C1 map OP4 is
\[
 x_1=x,\quad\eta_1=t,\quad\zeta_1=s,\quad r_1=h,\quad
 \theta_1=e,\quad t_1=\theta_1/\zeta_1=w.
\]
Write $P_d=P_n/B$. OP3--OP5 give the exact slack
\[
 s_n=2r\Delta_*\rho w/P_d,\quad
 \beta_1=-s_n/2,\quad \pi=1-s_n=1+2\beta_1>0,\quad
 \kappa=(1+2n)/\pi.
\tag{TL2}
\]
Hence $s_n/w=O(\rho)$, $\beta_1$ is infinitesimal, and
$\kappa\sim1$. The formal C1 parent is used through its explicit
identities, not imported as a physical C1 source theorem.

\paragraph{Exact coefficient comparison.}
Let $F_1$ be the native C1 parent at this same contact, normalized by
$d_1=1$, and $F_2$ the actual C2 parent, with $d_2=1$.
For a coefficient $L$ put $\delta L=L_2-\kappa L_1$.
Set $K=1+x$, $c_a=1+xr$, $S_c=1+B(1-h)+xr$, $Q_c=1-hB$,
and $\Theta=n/B=\rho w\psi$. The following short polynomials are
bounded with the displayed unit denominators:
\[
 P_\kappa=c_a^2+sx(1+t)+sx^2(r+he+t),
\]
\[
 P_D=tx^2r^2+x\{t-1+e(1-h-2t)\}-1,\quad
 P_\alpha=1+s-2e+x\{r^2+s(r+he)\}.
\]
Exact cancellation of the original C2 and native C1 formulas gives
\[
\begin{array}{c|l}
 \delta m&0\\
 \delta C&-\Theta c_a/(1-h)\\
 \delta(Ah)&-\Theta t c_aS_c/[K(h+t)]\\
 \delta D&\Theta S_cP_D/[K(h+t)tj]\\
 \delta\alpha&-nxP_\alpha/(KQ_c)\\
 \delta d=1-\kappa&-2\Theta P_\kappa/[K(h+t)]\\
 \delta\beta+\delta d/2&-n.
\end{array}
\tag{TL3}
\]
In particular $\delta C,\delta D,\delta d,\delta(Ah)=O(\rho w)$,
while $\delta\alpha=O(nx)$. There is no surviving $1/h$ after the
large $A$ event is centered.

\paragraph{The formal C1 local identities still supply a strict base.}
The native $\eta_1$ is a unit, $k_1,\zeta_1,r_1$ are infinitesimal,
and $e_{1,\mathrm{scale}}^2=tb$ is bounded. Separate its positive-unit
and infinitesimal types. Use respectively the exact centered-$A$
CA5/CA8 construction (or its CS1 scaling) and the SE4/SE9--SE11
construction. This use concerns the printed formulas and conductor
lemmas; the physical-source theorem is not being invoked at $\beta_1<0$.
Here are all changes to their source premises:

The formulas and native stationary identities are rational identities,
valid at the formal OP4 contact. All native gaps remain positive; $M=g$
and $J=j$ are units. The other native event coefficients are positive,
with $C_1,D_1,A_1h$ positive units and $A_1$ unbounded. The opening
ratio $t_1/e_{1,\mathrm{scale}}^2=w/(tb)$ is a positive bounded unit
by TL1, which supplies the only scale bound previously obtained from
$\beta_1\ge0$. All four retained packets SD3 involve the positive
paired coefficient $\beta_1+1/2=\pi/2$, never a standalone
$\beta_1$ square. In SE4, the coefficient of the reserved normal packet
is $(2+4\beta_1)\nu$, so its subtraction of $\nu$ is still positive,
since $\beta_1$ is infinitesimal. The remaining source bounds, finite
coquotients and strict residue matrices are unchanged. A positive-unit
normalization to $C_1=2$, when used by those formulas, changes none of
these bounds. Consequently their exact local construction yields the
same strict bounded eight-root Gram and positive retained packets for
$\kappa F_1$ on its contact chart.

\paragraph{A new positive contact-zero $n$ row.}
Use the inherited pure-testing coordinate and normals
\[
 \mathfrak w=(W+uY)/D_u,\quad A_t=v\mathfrak w,\quad
 X=N+pD_uA_t-uZ,\quad p=t_1/R_1,
\]
with $D_u=1+u^2,D_v=1+v^2$. In the exact CA chart the contact is
\[
 (u,\mathfrak w,v)=(k,\delta k,\delta k^2E),\quad
 K_c=\delta^2Ek^2=\zeta_1K/Q.
\]
Here $Q,E$ denote the inherited CA source letters, not the angular
weight of R10. The SE chart has $k=\varepsilon j$,
$\delta=1/\varepsilon$, so $K_c=j^2E$.
Put $\gamma=r/K_c$ and retain the entire state-linear square
\[
 \boxed{\quad \mathcal N^2=nD_uD_v(X+\gamma A_t)^2.\quad}
\tag{TL4}
\]
At the actual contact $A_t=K_ck$ and $X=-rk$, so its value and
all first derivatives vanish exactly. The square root of $D_uD_v$ is
pure testing; every other radical or inverse is pure source or pure
testing. The following entire-parent identity includes its full normal:
\[
\begin{split}
 nD_uD_vX^2-\mathcal N^2
 ={}&2n\gamma D_uD_vuA_tZ
 -(2n\gamma pD_u+n\gamma^2)D_uD_vA_t^2\\
 &-2n\gamma D_uD_vA_tN.
\end{split}
\tag{TL5}
\]
There is no remaining hard quadratic from this subtraction. This row
is necessary: without it the scalar $nu^2$ lies outside the old raw
square ideal, even if $n/k^4$ is infinitesimal.

\paragraph{Exact bounded stress after the positive packets.}
Let $A_{*,1}$ be the inherited bounded centered-$A$ cap. In the actual
parent retain $(A_2-\kappa A_{*,1})D_uD_v(Y-huW)^2$, which is positive
because $A_2$ is unbounded and $\kappa A_{*,1}$ bounded. Comparing with
the corresponding retained packet of $\kappa F_1$ replaces the whole
$A$ difference by
\[
 \delta(Ah)D_uD_v(2uWY-hu^2W^2).
\]
After retaining TL4, the full difference of the bounded remainders is
exactly
\[
\begin{split}
 \mathscr E={}&\bigl[nD_uD_vX^2-\mathcal N^2\bigr]
   -nD_v(X+uZ)^2+\delta C D_v(uW-Y)^2\\
 &+(\delta D+\delta d/2)R_{D,\mathrm{event}}^2
   +\delta\alpha D_u(vY-Z)^2\\
 &+(\delta d/2)D_u\{(vY-Z)^2-D_vZ^2\}\\
 &+\delta(Ah)D_uD_v(2uWY-hu^2W^2),
\end{split}
\tag{TL6}
\]
where $R_{D,\mathrm{event}}=uW+uvX-Y-vZ$ is the original event
row. This follows from TL3; in particular the paired beta difference
is exactly $-n$. Every scalar, first-normal and hard-normal term is
present. The full expression and its first derivatives vanish at the
unchanged source contact, because both full parents and every retained
square do.

\paragraph{Finite support and bounds for the unchanged conductor.}
Work on the local chart $Z=1$, with $u,v,\mathfrak w$ positive
infinitesimal. At zero normals set
\[
 V_0=1-puv\mathfrak w,\quad q_B=D_uvV_0,\quad
 uW=u\mathfrak w+u^2h_1vV_0,\quad
 Y=u\mathfrak w-h_1vV_0,\quad h_1=a_1/T_1.
\tag{TL7}
\]
These are exact inherited state inverses. The rows $uW-Y$ and
$R_{D,\mathrm{event}}$ are in $(v)$ at zero normals;
$X+u=pD_uA_t$. The centered $A$ expression in TL6 is a quadratic
in $(v,u\mathfrak w)$ with bounded coefficients. The signed-core
line has no constant: $(vY-1)^2-D_v$ is explicitly
$(vY)^2-2vY-v^2$. The sole state-constant term is
$\delta\alpha D_u$. Finally TL5 is a combination of
$uA_t=(v)(u\mathfrak w)$ and $A_t^2$, with first normal in $(A_t)$.
Thus this is an explicit raw allocation, not an ideal-membership guess.

For the unit-$e_{1,\mathrm{scale}}$ chart use
\[
 (A,B_c,C_c,D_0)
 =(\mathfrak wv,\delta v,\delta E u\mathfrak w,\delta^2k^3E).
\]
Every scalar coefficient of $\mathscr E/\rho$ in
$(A,B_c,C_c)^2+(D_0^2)$ is bounded with unit denominators. Both complete
first normals belong to $(A,B_c,C_c)$ with bounded coefficients, and
the hard-normal entries are bounded. The decisive source coquotients
are
\[
 \Theta/\rho=\delta^2QEw^2/K,\quad
 n\gamma^2/\rho=w^2r^2Q^2/K^2,\quad
 nx/(\rho D_0^2)=Q^2w^2/K^2.
\tag{TL8}
\]
The remaining $n\gamma$ coquotient has an additional $\delta^2k^2E$.
All apparent negative $k$ powers have therefore canceled before the
conductor is used.

For the infinitesimal-$e_{1,\mathrm{scale}}$ chart use the exact SE
parameters $w=t_1=\varepsilon^2\tau$,
$\zeta_1=j^2QE/K$, $k=\varepsilon j$, and
\[
 (A,B_c,C_c,D_0)
 =(\varepsilon\mathfrak wv,v,Eu\mathfrak w,\varepsilon^2j^3E).
\]
The full raw scalar allocation of $\mathscr E/\rho$ is
\[
 f_AA^2+f_{BB}B_c^2+f_{BC}B_cC_c+f_{CC}C_c^2+f_0D_0^2.
\tag{TL9}
\]
All five coefficients are bounded; after division by $\rho$ each has a
factor $\varepsilon^2$. The $A^2$ coefficient has this factor by TL10;
all its other terms carry $p=\varepsilon^2\tau/R_1$, because
$\delta m=0$. The coefficients of $B_c^2,B_cC_c,C_c^2$ have it
from $\Theta/\rho$, and the constant coefficient from the final
identity in TL10. In particular
\[
 \Theta/\rho=\varepsilon^2\tau^2QE/K,
\]
\[
 \frac{n\gamma^2}{\rho\varepsilon^2}
 =\frac{\varepsilon^2\tau^2r^2Q^2}{K^2},\quad
 \frac{n\gamma}{\rho\varepsilon}
 =\frac{\varepsilon^3j^2E\tau^2rQ^2}{K^2},\quad
 \frac{nx}{\rho D_0^2}
 =\frac{\varepsilon^2\tau^2Q^2}{K^2}.
\tag{TL10}
\]
The first-normal coefficients are in $(A,B_c,C_c)$ with bounded
arrays, also divisible by $\varepsilon^2$ after division by $\rho$;
the new $A$ coefficient has the stronger $\varepsilon^3j^2$ factor
printed in TL10. The extra term $-nD_v(X+u)^2$ has scalar coefficient
$-np^2D_u^2D_v$ on $A_t^2$, first normal
$-2npD_uD_vA_tN$, and hard term $-nD_vN^2$; TL10 and
$p=\varepsilon^2\tau/R_1$ bound each normalized coefficient.

Exactly as in SE8, all dependence of the raw coefficients on
$\mathfrak w$ comes through $V_0$ and has a factor
$p=O(\varepsilon^2)$. The additional TL5 coefficients are independent
of $\mathfrak w$. Modulo $\varepsilon$, scalar raw coefficients are
even in $u,v$; the normal raw coefficients are independent of
$\mathfrak w$. Clearing powers of $D_u$ and the source unit
denominators preserves these properties. These are precisely the
finite divisibility premises of SE7: the first scalar conductor has
$T_{\mathfrak w}\in(\varepsilon)$, the second-stage remainders have
$T_{\mathfrak w}\in(\varepsilon,v)$, and each full normal has
$T_{\mathfrak w}\in(\varepsilon)$.

For either chart, TL6 has exact contact value and first-jet zero.
Apply the inherited monic contact lift to $\mathscr E/\rho$:
CC3--CC6 in the first chart and the division-free SE7 lift in the
second. The raw support restriction retains only $D_0^2$, not arbitrary
mixed $D_0$ products. The proofs use only finite monic division and
bounded unit denominators; extending their coefficient ring by the
bounded source letters in TL8--TL10 changes nothing. They give a
bounded Gram matrix on the same eight inherited roots for
$\mathscr E/\rho$. Hence the exact stress $\mathscr E$ has an
infinitesimal Gram matrix on that list. No inverse of $\rho$ occurs in
its final coefficients: division by $\rho$ here records the common
positive source factor before a linear finite lift.

Adding this infinitesimal matrix to the strict formal-C1 base Gram
preserves its positive residue. Restore all its retained packets,
the modified positive centered-$A$ packet, and TL4. This proves the
full actual C2 parent on the thin-layer contact chart, with both
complete normals, all hard couplings and every independent rate.
The direct actual-C2 outer/state argument and the all-R10 remote
argument are required for the all-tests join; no negative formal-beta
packet is used in either join.



\section{\texorpdfstring{R10: a direct critical-wall certificate for every unit-$W$ state}{R10: a direct critical-wall certificate for every unit-W state}}
\label{component:pf1aac1dff7c7}
% Scoped mathematical certificate; see local_identity_audit.closure.json for exact changes.

\paragraph{Statement.}
On the entire original R10 source, every independent test with $u,v$
infinitesimal and normalized $W$ a testing unit is tensor SOS. This
includes bounded $Z$ of every sign, infinitesimal $Z$, and $Z=0$.
The theorem uses the actual physical C2 coefficients. No continuation
of a negative-beta C1 source is involved.

Use the R10 source scale $s=1/a$, $w=ae$, $b=a^2B$, and put
$z=b/w$. The source constraint proves $z$ bounded. Write $A_h=Ah$.
The source estimates give $A_h,C,D$ positive units, $h$ infinitesimal,
$m/w$ a positive unit, and $t_0>1$, $w_0<1$.
Define the positive pure-source critical coefficient
\[
 C_*={ (h+2A_h)(1+2D)\over 2(4A_hD-h)}
     ={(1+2A)(1+2D)\over2(4AD-1)}.
 \tag{UW1}
\]
Its denominator after the first normalization is a positive unit.
We prove the opening-uniform estimate
\[
                 {C_*-C\over m}=s\,R,
                 \qquad R\text{ bounded}.             \tag{UW2}
\]

\paragraph{A short polynomial proof of the source estimate.}
All the following letters are polynomials in bounded source variables
$(s,w,z,h,t)$ after $b=wz$, with $x=sb$, $B=s^2b$, $r=1-sw$:
\[
 c_a=1+xr,\quad S_c=1+B(1-h)+sbr,\quad
 g=1-w+hx,\quad J=t c_a-(1-h),\quad \Delta_*=s\Delta,
\]
\[
 M_*=(1+x)(s+r)(h+t)-sw(1+tB)(1-h).
\]
Use the exact source numerators and denominators
\[
 \begin{array}{ll}
 N_C=c_aM_*,&D_C=2(1-h)\Delta_*,\\
 N_A=t c_a(s+r)S_c,&D_A=2\Delta_*,\\
 N_D=(1+x)(s+r)(1-txr)S_c,&D_D=2J\Delta_*,\\
 N_m=w(1+tB)c_aS_c,&D_m=2g\Delta_*.
 \end{array}
\]
Thus $C=N_C/D_C$, $A_h=N_A/D_A$, $D=N_D/D_D$,
and $m=N_m/D_m$. Clear the denominator in UW1:
\[
 \begin{split}
 P={}&(hD_A+2N_A)(D_D+2N_D)D_C\\
     &-2N_C(4N_AN_D-hD_AD_D),\\
 Q={}&2(4N_AN_D-hD_AD_D)D_C.
 \end{split}                                                  \tag{UW3}
\]
Then $C_*-C=P/Q$. At $s=0$ one has
\[
 C={t+h\over2(1-h)},\quad A_h=t/2,\quad
 D={1\over2(t-1+h)},
\]
and substitution gives $C_*=C$ identically. At $w=0$ (hence
$b=0$), the same three formulas hold for every $s$, because all
factors $1+s$ cancel. Therefore the polynomial $P$ is divisible by
both $s$ and $w$, hence $P=swP_1$. Since $N_m=wN_{m,1}$,
\[
 {C_*-C\over m}=s\,{P_1D_m\over QN_{m,1}}.
\]
The denominator $QN_{m,1}$ at $s=h=0$ is $16t>0$.
This proves UW2 by finite polynomial cancellation, without comparing
$w$ to $s$ or $h$. The new checker verifies exactly this divisibility
and denominator. Its displayed residue also gives
$R\sim t(1-w)(tz+2)$, though the proof only needs boundedness.

\paragraph{Complete parent and state allocation.}
Set $\epsilon=(C_*-C)_+$. By UW2, $\epsilon/m$ is
infinitesimal. If $C>C_*$, retain the nonnegative excess $C-C_*$.
Retain the entire original $n$ and $\alpha$ packets; beta is zero.
The remaining signed core and $C_*,A,D$ packets are exactly the
critical parent SEO11, with $c=C_*$ and $d=1$, because
\[
                  2C_*(4AD-1)=(1+2A)(1+2D).
\]
The exact critical-wall identities SEO12--SEO15 give a bounded strict
Gram on
\[
 B_t=X+uZ,\quad H_A=\sqrt A(Y-yV),\quad
 H_D=\sqrt{4D+2}\{uW-Y+uvB_t-(b_*+u^2)V\},
 \quad uV,\ vV,
\]
where $V=vZ$, $y=1/(2A)$ and
$b_*=(4AD-1)/(A(4D+2))$. Here $y,b_*$ are bounded;
$C_*,D$ are units and $A$ may be unbounded. The hard residue is
strict and the two radial masses have positive residues $2$ and
$2/t_0$. This is an application of those full polynomial identities,
not of their original C1 source theorem.

Normalize $W=1$ by a pure-testing factor. The entire $m$ event obeys
\[
 R_m=v(1+uY)-B_t,\qquad
 V={Z\over1+uY}(R_m+B_t).                         \tag{UW4}
\]
The denominator is a positive testing unit. From the displayed
$H_D$ root,
\[
 uW-Y=(b_*+u^2)V-uvB_t+H_D/\sqrt{4D+2}.
\]
Thus $\sqrt\epsilon(uW-Y)$ is a linear combination of the strict
critical roots and $\sqrt mR_m$ with infinitesimal coefficients.
The only potentially small denominator is accounted for by
$\sqrt{\epsilon/m}$, infinitesimal by UW2. Consequently the complete
negative correction $-4\epsilon(1+v^2)(uW-Y)^2$ has an
infinitesimal bounded Gram. The $m$ packet supplies a strict additional
pivot. Restore all retained original packets to obtain the entire
C2 parent, including every hard state direction.

\paragraph{State entry for the remaining local work.}
Normalize a real testing state by a maximal coordinate. On the
negative pure-testing-core branch at $u,v$ infinitesimal, $X,Y$
are infinitesimal. If $W$ is a unit the theorem applies. Otherwise
$Z$ is a unit; set $Z=1$, and then
$w_{\rm test}=(W+uY)/(1+u^2)$ is infinitesimal. The inherited exact
Schur sign identity gives the positive contact chart. Literal axes
have the original event-square exits; complex states split into real
and imaginary parts. Together with the R10 remote theorem, the only
remaining task is the contact chart $u,v,w_{\rm test}>0$
infinitesimal, $Z=1$, on any source not already admitted locally.


\section{\texorpdfstring{A direct all-tests join for the five-root A-unbounded subchart}{A direct all-tests join for the five-root A-unbounded subchart}}
\label{component:p0d4eedd6c7a4}
% Scoped mathematical certificate; see local_identity_audit.closure.json for exact changes.

\label{p0d4eedd6c7a4:sec:r9-weak-A-remote}

\paragraph{Revision 10 scope update.}
Section~\ref{p296e824cb97a:sec:r10-weak-boundary} extends this admission to $z/B$
positive infinitesimal with all other source hypotheses retained.
The unit-ratio restriction below records the original revision-9 theorem.

This section supplies a new direct join for exactly the source
hypotheses (WA1). It does not import the strong-II remote theorem
or any other weak-C2 remote admission. All source physicality,
strict-gap and contact hypotheses from the local theorem remain
in force. No membership in the residual complement is needed;
applications to the original residual problem retain that intersection
without asserting that it is nonempty.

\begin{theorem}
Under the source hypotheses of the five-root local theorem, including
the positive-unit assumptions on $z/B$ and $1-r$, the entire original
parent is tensor SOS for every independent projective testing angle
pair and every real or complex projective state. Literal zero and
infinite half-angles and all signed charts are included.
This is an all-tests source-subdomain admission; it does not assert
nonemptiness of its intersection with the previous residual complement.
\end{theorem}

\paragraph{The limiting source is explicit and all four finite weights are strict.}
Put
\[
 \rho=\operatorname{st}(z/B)>0,\qquad
 \tau=\operatorname{st}(y/B)>0,\qquad
 r_0=\operatorname{st}(r)\in(0,1).
\]
The original unit gaps give
$g_0=1-\rho(1-r_0)>0$ and
$j_0=1-1/\tau>0$, hence $\tau>1$.
Exact reduction of the unchanged source formulas gives
\[
 \begin{gathered}
 n_0=\frac{1-r_0}{2r_0(1+\rho)},\qquad
 m_0=\frac{\rho^2(1-r_0)}{2g_0(1+\rho)},\\
 C_0=\frac{\tau(1+\rho r_0)-\rho(1-r_0)}{2(1+\rho)},\qquad
 D_0=\frac{1+\rho r_0}{2(1+\rho)(\tau-1)}.             
 \end{gathered}
\tag{WR1}
\]
Every displayed coefficient is strictly positive. For $C_0$ the
numerator equals $\tau g_0+\rho(\tau-1)+\rho r_0>0$.
In particular $n$ is a positive bounded unit on this subchart;
the stronger fact is used at one remote endpoint. The formulas
are joint regular reductions under $B\to0$, $H/B\to0$, with
the other source ratios tending to their residues; they do not
compare the two infinitesimal rates. Also
$\alpha\in\operatorname{Inf}$ and $A^{-1/2}\in\operatorname{Inf}$.

Normalize the hard state by $\widehat Y=\sqrt A\,Y$ in each
remote chart. This is a pure-source state change, and no parent
term is removed. On bounded testing charts its matrix has bounded
pure-source polynomial coefficients; all couplings of
$\widehat Y$ to $W,X,Z$ have zero residue. Its hard diagonal
residue is the positive angle factor multiplying $AY^2$.
It remains to check the complete three-state residue on $Y=0$.

\paragraph{Finite angles outside the local corner.}
For finite real angle residues $u_0,v_0$, write $U=u_0^2$,
$V=v_0^2$, $D_u=1+U$, $D_v=1+V$. Let $\mathsf P_3$ be the
matrix of the following complete residue of $F/4$ on $Y=0$:
\[
 \begin{split}
 F_3={}&n_0D_uD_vX^2+m_0(v_0W-X-u_0Z)^2
                    +C_0D_vU W^2\\
 &+(D_0+\tfrac12)(u_0W+u_0v_0X-v_0Z)^2
       +\tfrac12D_v(X+u_0Z)^2-\tfrac12D_uVZ^2.
 \end{split}                                           \tag{WR2}
\]
No hard-state relaxation term remains in this residue, because
the normalized $\widehat Y$ cross row is infinitesimal. Define
the strictly positive real constants
\[
 a_*=\frac{D_0m_0(2n_0+1)}2,\quad
 c_* =\frac{C_0(2m_0+1)(2D_0+2n_0+1)}4,\quad
 d_* =\frac{C_0D_0(2n_0+1)}2.
\]
The exact residue determinant is
\[
 \det\mathsf P_3=D_uD_v
 \left\{a_*\left(V-\frac\tau\rho U\right)^2
                         +c_*U^2V+d_*UV^2\right\}.    \tag{WR3}
\]
The checker obtains this directly from the full three-state
polynomial and (WR1), not from a numerical table. It is strictly
positive whenever $U+V>0$, including either axis. The restriction
of (WR2) to $Z=0$ is positive definite on $(W,X)$: its positive
$n_0,m_0,C_0$ event squares already have no common nonzero kernel
when $U+V>0$, and the core restriction is nonnegative. The real
Schur criterion and (WR3) therefore give
$\mathsf P_3\succ0$. Together with the hard $\widehat Y$ pivot,
the entire four-state residue is strictly positive definite.

For independent bounded testing angles with these residues,
expand the normalized full matrix around the real pair
$(u_0,v_0)$. Its nonconstant monomials in $u-u_0,v-v_0$ are
pure-testing infinitesimals. The finite separated Gram construction
(WA10) applies verbatim. If $u_0=v_0=0$, use the local five-root
theorem instead. Thus every bounded testing pair is covered.

\paragraph{One infinite angle.}
For $u$ unbounded use $p=1/u$ and the full homogeneous parent
$p^2F(1/p,v)$; at literal $u=\infty$ use its polynomial value
at $p=0$. If $v$ is bounded, the complete three-state residue on
$Y=0$ is
\[
 C_0(1+v_0^2)W^2+n_0(1+v_0^2)X^2
       +(D_0+\tfrac12)(W+v_0X)^2+(m_0+\tfrac12)Z^2.
                                                               \tag{WR4}
\]
It is positive definite for every finite real $v_0$, including
$v_0=0$. The hard $\widehat Y$ pivot is $1+v_0^2>0$.
For $v$ unbounded and $u$ bounded, put $q=1/v$ and use the
polynomial parent $q^2F(u,1/q)$. Its complete three-state residue is
\[
 (m_0+C_0u_0^2)W^2+(n_0+\tfrac12)(1+u_0^2)X^2
                                      +D_0(u_0X-Z)^2, \tag{WR5}
\]
again positive definite, with hard pivot $1+u_0^2$.
Both identities are verified directly against the original parent.
The same bounded separated Gram lift applies in the infinitesimal
variables $p,v-v_0$ or $u-u_0,q$. This includes literal infinity
and literal zero of the other half-angle.

\paragraph{Both infinite angles: keep the exact null state.}
Put $p=1/u$, $q=1/v$, and work with the full polynomial parent
$\widehat F=p^2q^2F(1/p,1/q)$. Introduce the five roots
\[
 (W,X,\widehat Y,P,Q),\qquad
 \widehat Y=\sqrt A\,Y,\quad P=pZ,\quad Q=qZ.
\]
Write $Y=A^{-1/2}\widehat Y$ in the following exact quadratic
identity, interpreted as a polynomial in these five roots:
\[
 \begin{split}
 \widehat F/4={}&n(1+p^2)(1+q^2)X^2
   +m(pW-pqX+Y-Q)^2+C(1+q^2)(W-pY)^2\\
 &+(1+p^2)(1+q^2)\widehat Y^2
   +D(qW+X-pqY-P)^2+\alpha(1+p^2)(Y-Q)^2\\
 &+\tfrac12\left\{(qW+X-pqY-P)^2+(1+p^2)(Y-Q)^2
       +p^2(1+q^2)X^2\right.\\
 &\hspace{39mm}\left.+2(1+q^2)XP-(1+q^2)P^2\right\}.
 \end{split}                                           \tag{WR6}
\]
The checker verifies this as the complete original four-state
parent after $P=pZ$, $Q=qZ$, not only to leading order.
All matrix coefficients in (WR6) are bounded source polynomials
in $p,q$. Its residue, in the printed root order, is
\[
 \begin{pmatrix}
 C_0&0&0&0&0\\
 0&n_0+D_0+1/2&0&-D_0&0\\
 0&0&1&0&0\\
 0&-D_0&0&D_0&0\\
 0&0&0&0&m_0+1/2
 \end{pmatrix}\succ0,                                 \tag{WR7}
\]
because the nontrivial determinant is $D_0(n_0+1/2)>0$.
The separated lift (WA10) therefore gives a finite tensor SOS
for arbitrary independently infinitesimal $p,q$, of either sign
or zero. At $p=q=0$ both $P,Q$ vanish, respecting the original
exact null direction $W=X=Y=0$; no positive $Z$ pivot is invented.

\paragraph{Coverage and restoration.}
Every finite element of an ordered extension is either bounded or
unbounded in absolute value. A bounded testing pair is covered by
(WA1)--(WA10) if both residues vanish, and by (WR2)--(WR3)
otherwise. A single unbounded coordinate is covered by (WR4) or
(WR5); two unbounded coordinates are covered by (WR6)--(WR7).
Literal projective endpoints use the polynomial chart values with
$p=0$ or $q=0$. For nonzero reciprocal coordinates, restoring
$F$ multiplies the established identity by the pure-testing
square $p^{-2}$, $q^{-2}$, or $(pq)^{-2}$; at literal infinity
no division is performed. All four homogeneous state coordinates
were retained in every chart, and splitting real and imaginary
state parts gives the complex-state statement. Thus the join
exhausts all independent testing types without a source/testing
rate assumption or an imported strong-face theorem.

\begin{corollary}[The polynomial obstruction family is nevertheless admitted]
The entire two-parameter physical family (SO1),
\[
 B=\varepsilon,\quad z=\varepsilon/2,\quad y=2\varepsilon,
 \quad r=9/10,\quad H=\varepsilon\eta,
 \quad L=\varepsilon(1-\eta),
\]
with independent positive infinitesimals $\varepsilon,\eta$, is
all-tests tensor SOS for the original full parent.
\end{corollary}
\begin{proof}
It has $z/B=1/2$, $1-r=1/10$, $H/y=\eta/2$ infinitesimal,
and $y/B,L/B,r,G/B,j$ positive bounded units. Its global
physicality and all original source signs were proved in the
obstruction module. It therefore satisfies every hypothesis of
the present all-tests theorem.
\end{proof}
The strict $K_1<0$ and (JG13) failures and the nonexistence of a
global polynomial SOS of its scalar graph remain correct.
The local tensor roots use pure-testing radicals as in (WA10),
so these facts are compatible. The whole family is now a proved
admission and is removed from the residual complement. No claim
that any other residual intersection is nonempty is implied.


\section{\texorpdfstring{The small-ratio weak-C2 opening is also admitted for all tests}{The small-ratio weak-C2 opening is also admitted for all tests}}
\label{component:p296e824cb97a}
% Scoped mathematical certificate; see local_identity_audit.closure.json for exact changes.

\label{p296e824cb97a:sec:r10-weak-boundary}

\begin{theorem}
Retain the original physical C2 source, its exact stationary contact,
all strict source gaps, $n\ge0$, and the normalization $d=1$.
Suppose $B,z,y>0$ are infinitesimal,
$r,1-r,G/B,j,y/B,L/B$ are positive bounded units, and $H/y$ is
infinitesimal. If $a=z/B$ is positive infinitesimal, the entire
original parent is a finite tensor SOS at every independent
projective testing pair, for all real or complex states, including
signed charts and literal zero and infinite half-angles.
No comparison between source infinitesimal rates or between source
and testing rates is imposed. In particular $\delta_{\rm tot}>0$
is infinitesimal. This admission does not assert nonemptiness of
its intersection with the previous residual complement.
\end{theorem}

The original omitted event coefficients remain $\lambda=\beta=0$.
Here $B=H+L$, $h=H/B$, $t=y/B$, and the original contact is
\[
 (u_c,v_c;W_c,X_c,Y_c,Z_c)
 =\left(\sqrt z,\sqrt y;B/\sqrt{zy},-r\sqrt z,H/\sqrt y,1\right).
\]
No residual-complement premise is used. The proof extends the exact
six-contact full-parent Gram (WA3)--(WA8) and its finite separated
lift (WA10); it does not use the disproved global polynomial splitter.

\paragraph{Exact cancellation, with no $h/a$ restriction.}
Put $\Delta=\delta/B$, $M_*=M/B^2$, $\mu=m/a^2$ and
$\alpha_*=\alpha/B^2$. The unchanged source formulas give
\[
\begin{gathered}
 c_a=1+aBr,\quad S=1+B(1-h+ar),\\
 g=G/B=1+a(hB+r-1),\quad j=c_a-(1-h)/t,\\
 \begin{aligned}\Delta={}&a+(1+aB)+(1+aB)B+2aB(1+aB)r\\ &+a^3B^2r^2
       +2a^2Br+tB[1+a+aB+a^2Br(2-r)]\end{aligned},\\
 M_*=(1+aB)(1+ar)(h+t)-a(1+tB)(1-r)(1-h),\\
 n=\frac{(1+ar)c_aS-r\Delta}{2r\Delta},\qquad
 \mu=\frac{(1+tB)(1-r)c_aS}{2g\Delta},\\
 C=\frac{c_aM_*}{2(1-h)\Delta},\quad
 Ah=\frac{t c_a(1+ar)S}{2\Delta},\\
 D=\frac{(1+aB)(1+ar)(1-atBr)S}{2tj\Delta},\quad
 \alpha_* =\frac{(1+ar)M_*}{2(1-hB)\Delta}.       \end{gathered}\tag{WE1}
\]
All displayed denominators are units at the joint corner
$B=a=h=0$; $j>0$ a unit implies $t_0=\operatorname{st}(t)>1$.
The subscript $0$ denotes residue, and $0<r_0<1$.
Thus
\[
 n_0=\frac{1-r_0}{2r_0},\quad \mu_0=\frac{1-r_0}2,
 \quad C_0=\frac{t_0}2,\quad D_0=\frac1{2(t_0-1)},
 \quad (Ah)_0=(\alpha_*)_0=\frac{t_0}2.          \tag{WE2}
\]
These are positive; $m=a^2\mu$ is infinitesimal, $A$ is
unbounded, and $Ah^2,\sqrt A h,A^{-1/2}$ are infinitesimal.
All other retained coefficients are bounded. Consequently
$\delta_{\rm tot}=1/(1+n+m+C+A+D+\alpha)$ is positive
infinitesimal, as required in the original source normalization.

The only new potentially unbounded local coefficient is
$(b_2-\alpha_*)/a$, where
\[
 b_2=2a^2\mu h+C(1-h)^2+(Ah)h+(B^2\alpha_*+1/2)h^2.
\]
At $a=0$, with $d_*=1+B(1+t)$, the exact values are
\[
 C=\frac{h+t}{2(1-h)d_*},\quad
 Ah=\frac{t[1+B(1-h)]}{2d_*},\quad
 \alpha_* =\frac{h+t}{2(1-Bh)d_*}.
\]
Substitution gives $b_2=\alpha_*$ identically for arbitrary
$B,h,t$. Clearing the unit denominators in (WE1), the numerator
of $b_2-\alpha_*$ is therefore divisible by $a$ in the polynomial
ring. Its quotient has a unit denominator at $B=a=h=0$.
Hence
\[
 Q_*=(b_2-\alpha_*)/a\quad\hbox{is bounded}.       \tag{WE3}
\]
This is a joint regularity argument, so it permits arbitrary
$h/a$, including an unbounded ratio. The checker verifies the
exact vanishing, polynomial divisibility, and nonzero denominator
at the joint corner.

\paragraph{The complete local five-root Gram survives.}
Keep $\kappa=B/y$ and every coefficient and normal in
(WA3)--(WA8), but set $\ell=1/a$ and replace $R_2$ by
\[
 \widetilde R_2=aR_2=avW-uZ.
\]
The entire full-parent matrix is the exact congruence
\[
 \widetilde G=EGE,\qquad E=\operatorname{diag}(1,1,1,a^{-1},1),
 \qquad \widetilde{\mathcal R}=(\xi,q,R_1,\widetilde R_2,T).
                                                               \tag{WE4}
\]
Its newly rescaled coefficients are
\[
\begin{gathered}
 m/a^2=\mu,\qquad
 \beta_s/a^2=(Bd_4+\kappa c_4)/\kappa,\\
 J_s(U)/a=\frac{aBd_4-a_4\kappa^2+d_2+d_4U}{2\kappa},\qquad
 f_s/a=\tfrac12[Q_*-ac_4-aBd_4/\kappa],\\
 (-m+q_1U)/a=-a\mu+(nr+Bp_1)U,\\
 q_y/a=-k_z-\kappa p_y'-Bk_{uv}.                  \end{gathered}\tag{WE5}
\]
All are bounded, by (WE1)--(WE3). The other coefficients are
bounded as before; the complete second normal remains divided
by $\sqrt A$. Thus no first normal or hard-state term is omitted.
For independent infinitesimal or zero $u,v$ the complete residue is
\[
 \overline{\widetilde G}
 =\operatorname{diag}(n_0+1/2,1,C_0+D_0+1/2,\mu_0,t_0/2)
 \succ0.                                               \tag{WE6}
\]
In particular the previous $\xi,R_2$ cross entry becomes
$-a\mu$ and has zero residue. Applying the pure-source and
pure-testing finite lift (WA10) proves the full local statement.

\paragraph{Bounded tests with nonzero $u$ residue.}
Normalize $\widehat Y=\sqrt A Y$ without centering and use the
complete three-state residue (WR2), now with $m_0=0$ and (WE2).
For $U=u_0^2$, $V=v_0^2$, its exact determinant is
\[
 \det P_3=(1+U)(1+V)U[b_*U+c_*UV+d_*V^2],             \tag{WE7}
\]
where
\[
 b_*=\frac{n_0(2C_0+2D_0+1)}4,\quad
 c_*=\frac{C_0(2D_0+2n_0+1)}4,\quad
 d_*=\frac{C_0D_0(2n_0+1)}2.
\]
Every coefficient is positive. For $U>0$, the restriction to
$Z=0$ is positive definite from the positive $n_0,C_0$ packets.
The determinant therefore gives a positive definite complete
residue, with the additional hard $\widehat Y$ pivot
$(1+U)(1+V)$. The finite separated lift around $(u_0,v_0)$ applies.

\paragraph{The missing finite-$v$ and reciprocal-$v$ charts.}
For $u$ infinitesimal or zero and $v$ a bounded unit, use the five
roots
\[
 (S_*,R_*,X,\widehat Y,Z)=(aW,uW,X,\sqrt A Y,Z),\qquad
 Y=A^{-1/2}\widehat Y.
\]
Writing $D_u=1+u^2$, $D_v=1+v^2$, the following is the entire
original $F/4$, as a quadratic polynomial in those roots:
\[
\begin{split}
 nD_uD_vX^2&+\mu(vS_*-aX+auvY-auZ)^2
 +CD_v(R_*-Y)^2+D_uD_v\widehat Y^2\\
 &+D(R_*+uvX-Y-vZ)^2+\alpha D_u(vY-Z)^2\\
 &+\tfrac12\{(R_*+uvX-Y-vZ)^2+D_u(vY-Z)^2
               +D_v(X+uZ)^2-D_uD_vZ^2\}.             \end{split}\tag{WE8}
\]
The coefficients are polynomials in $u,v$. Their pure-source factors $a,\mu,n,C,D,\alpha,A^{-1/2}$ are bounded.
At $u_0=0$, $v_0\ne0$, the separate pivots on
$(S_*,X,\widehat Y)$ are
$\mu_0v_0^2,(n_0+1/2)(1+v_0^2),1+v_0^2$.
The remaining block on $(R_*,Z)$ is
\[
 \begin{pmatrix}
 C_0(1+v_0^2)+D_0+1/2&-(D_0+1/2)v_0\\
 -(D_0+1/2)v_0&D_0v_0^2
 \end{pmatrix},\qquad \det=C_0D_0v_0^4>0.           \tag{WE9}
\]
Here $C_0D_0-D_0/2-1/4=0$ exactly. Thus (WA10) applies
in $u,v-v_0$, including literal $u=0$; roots need not be
independent for a positive definite Gram to give the asserted SOS.

For $v$ unbounded or literally infinite, put $q_v=1/v$ and
use the polynomial homogeneous parent $q_v^2F(u,1/q_v)/4$.
In the same five roots it is exactly
\[
\begin{split}
 nD_u(1+q_v^2)X^2
 &+\mu(S_*-aq_vX+auY-auq_vZ)^2
 +C(1+q_v^2)(R_*-Y)^2+D_u(1+q_v^2)\widehat Y^2\\
 &+D(q_vR_*+uX-q_vY-Z)^2+\alpha D_u(Y-q_vZ)^2\\
 &+\tfrac12\{(q_vR_*+uX-q_vY-Z)^2+D_u(Y-q_vZ)^2\\
 &\hspace{20mm}+(1+q_v^2)(X+uZ)^2-D_u(1+q_v^2)Z^2\}.
                                                               \end{split}\tag{WE10}
\]
For independent infinitesimal or zero $u,q_v$ its full residue
is diagonal with pivots
\[
 (\mu_0,C_0,n_0+1/2,1,D_0)>0.                       \tag{WE11}
\]
Its coefficients are bounded for the same reason as (WE8).
The separated lift again applies, including $q_v=0$ literally.
The checker expands (WE8) and (WE10) back to all six original
event packets and the complete signed core, and checks both residues.

\paragraph{Remaining endpoints and exhaustive coverage.}
If $u$ is unbounded or infinite and $v$ bounded, (WR4) is still
positive definite with $m_0=0$: its positive coefficients are
$C_0,n_0,D_0+1/2,1/2$, and the hard pivot is $1+v_0^2$.
If $v$ is unbounded or infinite and $u$ is a bounded unit,
(WR5) is positive definite because $C_0u_0^2>0$,
$n_0+1/2>0$, and $D_0>0$, with hard pivot $1+u_0^2$.
If both are unbounded or infinite, use the full homogeneous
five-root identity (WR6). Its residue (WR7) is still strictly
positive: the $Q$ pivot becomes $1/2$, and the other nontrivial
block determinant remains $D_0(n_0+1/2)>0$.
All these identities have bounded coefficients also at $m_0=0$.

For two bounded tests, both zero residues use (WE4)--(WE6);
nonzero $u$ residue uses (WE7); zero $u$ and nonzero $v$ residue
uses (WE8)--(WE9). A single unbounded coordinate is covered by
the preceding paragraph, except for unbounded $v$ and zero $u$
residue, which is (WE10)--(WE11). Two unbounded coordinates use
(WR6)--(WR7). These cases exhaust every independent testing type.
Literal projective endpoints use the exact polynomial chart value.
Away from literal endpoints, restoration multiplies by a
pure-testing square $p^{-2}$, $q_v^{-2}$, or $(pq_v)^{-2}$.
No division is made at a literal endpoint. All four original
homogeneous state coordinates are retained; real and imaginary
parts give the complex-state result. The proof establishes the
full-parent all-tests admission on this source opening, without
a scalar-to-parent inference or a new physicality assumption.

\paragraph{The unit-ratio assumption can now be removed.}
Under the retained hypotheses, $H\in\Inf$ and $1-r$ is a
positive unit. Since $g=1-a(1-r-H)>0$, one has
$0<a<(1-r-H)^{-1}$, so $a$ is bounded. Every positive such
$a$ is either a unit, covered by (WA1), or infinitesimal, covered
here. Their union therefore gives the same all-tests theorem
without the additional $z/B$-unit assumption. This argument does
not cover a vanishing $1-r$ or any other omitted source gap.


\section{\texorpdfstring{One complete Gram for every bounded test on the small-r opening}{One complete Gram for every bounded test on the small-r opening}}
\label{component:p99372d145f17}
% Scoped mathematical certificate; see local_identity_audit.closure.json for exact changes.

\label{p99372d145f17:sec:r11-small-r-local}

\begin{theorem}
Retain the original physical C2 source, its exact stationary contact,
all strict source gaps, and $d=1$. Suppose $B,z,y,r>0$ are
infinitesimal; $a=z/B$ is positive bounded (a unit or infinitesimal);
$G/B$, $j$, $y/B$, $L/B$ are positive bounded units; and $H/y$
is infinitesimal.
Then the entire original four-state parent is a finite tensor SOS
for arbitrary independent bounded testing half-angles $u,v$, of either sign, and for every real or complex state.
There is no comparison of source or testing rates.
The total-scale parameter $\delta_{\rm tot}$ is positive infinitesimal.
This covers all bounded tests; the projective endpoint join is a
separate obligation. No residual-complement nonemptiness is claimed.
\end{theorem}

Here $B=H+L$, $h=H/B$, $t=y/B$, $\kappa=1/t$, $\ell=1/a$;
the omitted event coefficients remain $\lambda=\beta=0$.
The contact is the unchanged original contact
\[
 (\sqrt z,\sqrt y;B/\sqrt{zy},-r\sqrt z,H/\sqrt y,1).
\]
No residual membership premise is used. All source physicality,
contact and strict-gap hypotheses remain in force.

\paragraph{Both unbounded source coefficients can be retained.}
Use the exact normalized source formulas (WE1), whose denominators
are regular at $B=h=r=0$ after replacing $n$ by $N=nr$.
Explicitly, with $c_a=1+aBr$, $S=1+B(1-h+ar)$ and
$\Delta=\delta/B$,
\[
 N=\frac{(1+ar)c_aS-r\Delta}{2\Delta}.
\]
Let $\rho=\operatorname{st}(a)\in[0,1)$ and
$\tau=\operatorname{st}(t)>1$. Direct joint reduction gives
\[
\begin{gathered}
 N_0=\frac1{2(1+\rho)},\qquad
 (m/a^2)_0=\frac1{2(1-\rho^2)},\qquad
 C_0=\frac{\tau-\rho}{2(1+\rho)},\\
 D_0=\frac1{2(1+\rho)(\tau-1)},\qquad
 (Ah)_0=\frac\tau{2(1+\rho)},\qquad
 (\alpha/B^2)_0=\frac{\tau-\rho}{2(1+\rho)}=C_0.    \end{gathered}\tag{SL1}
\]
Each displayed residue is strictly positive. Thus $n=N/r$ and
$A=(Ah)/h$ are positive unbounded, with no comparison between
them. In contrast $nr$ is a unit and $nr^2=Nr$ is infinitesimal;
$\mu=m/a^2,C,D,Ah,\alpha/B^2$ are units, $m=a^2\mu$ is
bounded, and $\alpha$ is infinitesimal.
All normalized source denominators are units since
$G/B\to1-\rho>0$, $j\to1-1/\tau>0$, and
$\Delta\to1+\rho>0$. These are joint regular reductions,
not iterated limits. In particular
$\delta_{\rm tot}=1/(1+n+m+C+A+D+\alpha)>0$ is infinitesimal.

\paragraph{Exact congruence of the complete full-parent identity.}
Keep all source coefficients, the scalar block $\mathsf S$, and
both complete normal rows $\mathsf L_x,\mathsf L_y$ in
(WA3)--(WA8), with $\ell=1/a$. Normalize the first root and
the second scalar root by pure-source factors:
\[
 \widehat\xi=\sqrt n(X+ruZ),\quad
 q=\sqrt A(Y-huW),\quad
 R_1=uW-\kappa vZ,\quad \widetilde R_2=avW-uZ,\quad T=uvW-BZ.
                                                               \tag{SL2}
\]
Every root still vanishes at the exact contact. If
$E_s=\operatorname{diag}(1,a^{-1},1)$ and
$\widehat{\mathcal R}=(\widehat\xi,q,R_1,\widetilde R_2,T)^T$, then
the entire original parent has the exact identity
\[
 F/4=\widehat{\mathcal R}^{T}\widehat G\widehat{\mathcal R},
 \qquad
 \widehat G=
 \begin{pmatrix}
 h_{xx}/n&h_{xy}/\sqrt{nA}&\mathsf L_x E_s/\sqrt n\\
 h_{xy}/\sqrt{nA}&h_{yy}/A&\mathsf L_y E_s/\sqrt A\\
 E_s\mathsf L_x^T/\sqrt n&E_s\mathsf L_y^T/\sqrt A&E_s\mathsf S E_s
 \end{pmatrix}.                                        \tag{SL3}
\]
This is the pure-source congruence of (WA8) by
$\operatorname{diag}(n^{-1/2},1,1,a^{-1},1)$, so it retains
the exact signed core, all event packets, both full normals and
the entire hard block. It invokes the previously verified six
original contact identities; no scalar-to-parent implication is used.

Here is a complete coefficient-bound check. In the scalar block,
every occurrence of $n$ is $nr^2=Nr$. The exact rescaled formulas
(WE5) show that $\beta_s/a^2,J_s/a,j_1/a$ are bounded (in fact
infinitesimal). The only possible new pole is $f_s/a$.
The exact source identity $b_2-\alpha/B^2=0$ at $a=0$ holds
for arbitrary $B,h,r,t$ and contains no inverse $r$.
Its cleared numerator is divisible by $a$, and a canceled
denominator reduces to $2(\rho-1)(\rho+1)\ne0$ at
$B=h=r=0$, $0\le\rho<1$.
The quotient numerator also vanishes identically at $B=h=r=0$,
for every $a$ in this residue interval. Consequently
$Q_*=(b_2-\alpha/B^2)/a$ is infinitesimal uniformly,
including joint $r,a$ infinitesimal with arbitrary rates. Formula
$f_s/a=[Q_*-ac_4-aBd_4/\kappa]/2$ is therefore infinitesimal.
All other rescaled scalar coefficients are bounded, with
$m/a^2=\mu$ a unit. In the first normal, the only
potentially large terms occur through
\[
 p_0=\frac{D+nr+r/2}{\kappa},\quad
 p_1=\frac{nr+r(D+1/2)}{\kappa},\quad
 q_1=\frac{nr+Bp_1}{\ell}.
\]
They are bounded, as is
$f_x=-mh+(D+1/2)(1-h)-p_0-q_1$; therefore every coefficient
of $\mathsf L_x E_s/\sqrt n$ is infinitesimal, since its
middle entry before division by $\sqrt n$ is
$(-m+q_1U)/a=-a\mu+(nr+Bp_1)U$.
The second normal depends only on bounded $m,C,D,\alpha,r,h,Ah$
and unit $\kappa$; its rescaled middle coefficient is bounded
because $q_y/a=-k_z-\kappa p_y'-Bk_{uv}$. Thus all its
coefficients become infinitesimal after division by $\sqrt A$.
Finally, writing
$U=u^2,V=v^2$, the complete hard entries are
\[
\begin{split}
 h_{xx}/n&=(1+U)(1+V)
        +[m+(D+1/2)UV+(1+V)/2]/n,\\
 h_{xy}/\sqrt{nA}&=-uv(m+D+1/2)/\sqrt{nA},\\
 h_{yy}/A&=(1+U)(1+V)
 +[C(1+V)+D+1/2+mUV+(\alpha+1/2)(1+U)V]/A.
\end{split}                                            \tag{SL4}
\]
No ratio $n/A$ or $A/n$ occurs. Thus $\widehat G$ is a finite
polynomial matrix in $u,v$ with bounded pure-source coefficients,
uniformly for arbitrary independent source rates.

For arbitrary bounded testing $u,v$, put $u_0=\operatorname{st}(u)$,
$v_0=\operatorname{st}(v)$, and $D_0^{\rm test}=(1+u_0^2)(1+v_0^2)>0$.
The scalar source coefficients
$a_4,\beta_s/a^2,J_s/a,j_1/a,\gamma_s,f_s/a$ are all infinitesimal.
Indeed $nr^2=Nr\in\Inf$, $d_2,d_4,Ah^2\in\Inf$, and
$Q_*\in\Inf$ by the just-proved exact quotient identity.
The complete normalized normal rows have infinitesimal coefficients.
Consequently the entire Gram has the strict diagonal residue
\[
 \overline{\widehat G}
 =\operatorname{diag}(D_0^{\rm test},D_0^{\rm test},
                     C_0+D_0+1/2,\mu_0,C_0)\succ0.       \tag{SL5}
\]
This is stronger than a both-infinitesimal testing statement.
Evaluate the source polynomial Gram at the real pair $(u_0,v_0)$;
its pure-source inverse is bounded. Whiten there and expand in the
pure-testing infinitesimals $u-u_0,v-v_0$. The finite separated
lift (WA10) then applies verbatim, using only pure-source
and pure-testing radicals. All bounded tests of either sign and
zero are covered, independently of their rates relative to any
source scale. Splitting real and imaginary state parts proves the
complex-state conclusion. A separate join below handles every
unbounded or literal infinite testing angle.


\section{\texorpdfstring{All projective tests on the small-r weak-C2 source opening}{All projective tests on the small-r weak-C2 source opening}}
\label{component:p34609a6519a6}
% Scoped mathematical certificate; see local_identity_audit.closure.json for exact changes.

\label{p34609a6519a6:sec:r11-small-r-remote}
\begin{theorem}
Retain the original physical weak-C2 source, exact stationary contact,
all strict source gaps, $n\ge0$, and $d=1$. Suppose $B,z,y,r>0$
are infinitesimal, $G/B,j,y/B,L/B$ are positive bounded units,
and $H/y$ is infinitesimal. Then the entire original parent is
finite tensor SOS at every independent projective testing pair,
with all signs, literal endpoints, and real or complex states.
There is no comparison of independent source and testing rates.
\end{theorem}
Write $a=z/B$, $t=y/B$, $h=H/B$. Positivity of
$G/B=1-a(1-r-H)$ makes $a$ bounded; its unit gap gives
$0\le\rho=\operatorname{st}(a)<1$. Put
$\tau=\operatorname{st}(t)>1$. The preceding local lemma proves
the source estimates, the exact full-parent five-root identity,
and its strict residue for every bounded testing pair. In particular
\[
 \mu_0=\frac1{2(1-\rho^2)},\quad m_0=\rho^2\mu_0,\quad
 C_0=\frac{\tau-\rho}{2(1+\rho)},\quad
 D_0=\frac1{2(1+\rho)(\tau-1)}.                 \tag{SR1}
\]
Here $\mu=m/a^2$, while $n,A$ are unbounded, independently.
This section supplies the complete remote join; no remote theorem
requiring bounded $n$ is imported as a domain admission.

\paragraph{Two normalized hard state coordinates.}
Use $\widehat X=\sqrt n X$, $\widehat Y=\sqrt A Y$, with
$e_x=n^{-1/2}$ and $e_y=A^{-1/2}$ positive infinitesimal.
Substitute $X=e_x\widehat X$, $Y=e_y\widehat Y$ into the
original full parent. Its $nX^2$ and $AY^2$ packets become the
exact hard squares, and every remaining coefficient is bounded.
Every cross term with these two normalized states has zero residue.
At bounded test residues $u_0,v_0$, put $U=u_0^2$, $V=v_0^2$.
The entire residue of $F/4$ is
\[
 (1+U)(1+V)(\widehat X^2+\widehat Y^2)+\Psi(W,Z),
\]
where
\[
\begin{split}
 \Psi={}&m_0(v_0W-u_0Z)^2+C_0(1+V)U W^2
 +(D_0+\tfrac12)(u_0W-v_0Z)^2\\
 &+\tfrac12(1+V)UZ^2-\tfrac12(1+U)VZ^2.
\end{split}                                                   \tag{SR2}
\]
This retains the complete signed core. If $P_2$ is its matrix on
$(W,Z)$, direct expansion using (SR1) gives
\[
\begin{split}
 \det P_2={}&D_0\mu_0(\rho V-\tau U)^2\\
 &+\frac{C_0(2m_0+1)}2U^2V+C_0D_0UV^2.          
\end{split}                                                   \tag{SR3}
\]
All three displayed coefficients are strictly positive.
For $U>0$, this determinant is positive for every $V\ge0$:
if $V=0$ its first term is positive; if $V>0$ its second is.
The $W$ diagonal $m_0V+[C_0(1+V)+D_0+1/2]U$ is positive.
Thus the full four-state residue is strictly positive.
For $U=0,V>0$ the same holds when $\rho>0$; here the two
remaining pivots are $m_0V,D_0V$.
Expanding around the real residues and applying the finite separated
lift (WA10) proves every such bounded testing chart.
All bounded pairs are already covered by the preceding strengthened
lemma. The determinant gives a second direct check away from its
vanishing locus.

\paragraph{The joint small-a and small-r corner.}
Only $\rho=0$, $u$ infinitesimal or zero, and nonzero or infinite
$v$ remains outside those arguments. The complete identities
(WE8) and (WE10) apply algebraically at arbitrary $n,A$.
In those exact identities make the same pure-source substitutions
$X=e_x\widehat X$, $Y=e_y\widehat Y$, using the five roots
\[
 (S,R,\widehat X,\widehat Y,Z)=(aW,uW,\sqrt n X,\sqrt A Y,Z).
\]
All resulting coefficients are bounded polynomials in the testing
coordinates with pure-source coefficients. The factors $nX^2$
and $AY^2$ again become the exact normalized squares; no term is
removed. For $v_0\ne0$ finite, the separate pivots on
$(S,\widehat X,\widehat Y)$ are
\[
 \mu_0V,\qquad 1+V,\qquad 1+V.
\]
The remaining block on $(R,Z)$ is
\[
 \begin{pmatrix}
 C_0(1+V)+D_0+1/2&-(D_0+1/2)v_0\\
 -(D_0+1/2)v_0&D_0V
 \end{pmatrix},\qquad \det=C_0D_0V^2>0.          \tag{SR4}
\]
Indeed at $\rho=0$, $C_0=\tau/2$,
$D_0=1/[2(\tau-1)]$, and
$C_0D_0-D_0/2-1/4=0$.
For $v$ unbounded put $q=1/v$ and use the polynomial parent
$q^2F(u,1/q)$. At $u=q=0$ the full five-root residue of the
normalized (WE10) is diagonal with pivots
\[
 (\mu_0,C_0,1,1,D_0)>0.                           \tag{SR5}
\]
The same finite separated lift covers both charts, including
literal $u=0$ and literal $q=0$. The roots may be dependent;
a strict Gram identity remains a valid finite SOS on them.

\paragraph{Every other projective endpoint.}
For $u$ unbounded use $p=1/u$ and the exact polynomial
$p^2F(1/p,v)$. With $v$ bounded its normalized full residue is
\[
 [C_0(1+V)+D_0+1/2]W^2
 +(1+V)(\widehat X^2+\widehat Y^2)+(m_0+1/2)Z^2.
                                                               \tag{SR6}
\]
It is strictly positive for every $V\ge0$, whether $\rho$ vanishes
or not. For $v$ unbounded and $u$ bounded the residue is
\[
 (m_0+C_0U)W^2+(1+U)(\widehat X^2+\widehat Y^2)+D_0Z^2.
                                                               \tag{SR7}
\]
This is strictly positive if $\rho>0$ or $U>0$; the missing
case is precisely (SR5).
For two unbounded coordinates use the complete polynomial identity
(WR6), replace $X,Y$ by their normalized expressions, and keep
$P=pZ$, $Q=qZ$ as separate roots. The full residue on
$(W,\widehat X,\widehat Y,P,Q)$ is diagonal with pivots
\[
 (C_0,1,1,D_0,m_0+1/2)>0.                         \tag{SR8}
\]
All exact coefficients are bounded. The finite lift applies in
$p,q$, also at their literal zero values. At $p=q=0$ both
$P,Q$ vanish, preserving the original exact null state direction.

Every bounded pair has both residues zero, $U>0$, or $U=0,V>0$.
These are covered by the local lemma, (SR3), and (SR3) or (SR4).
Single infinities use (SR5)--(SR7); double infinity uses (SR8).
Restoration away from a literal endpoint multiplies by a pure-testing
square; no inverse is taken at the endpoint. Every state coordinate
and first normal was retained. Applying the real identity separately
to real and imaginary state parts gives the complex-state result.
No nonemptiness of a residual intersection is asserted.

\begin{corollary}[Two source-unit assumptions removed]
In the revision-9 weak-C2 source theorem, both additional requirements
that $r$ and $z/B$ be bounded units can be omitted, provided $1-r$
and the other original gaps and ratios remain positive bounded units.
\end{corollary}
\begin{proof}
The original $0<r<1$ is either a unit or infinitesimal.
For unit $r$, revisions 9 and 10 cover unit and infinitesimal $z/B$.
For infinitesimal $r$, the present theorem covers both possibilities.
Positivity of $G/B=1-(z/B)(1-r-H)$ and the retained unit $1-r$
exclude unbounded $z/B$. These source-only alternatives are exhaustive.
\end{proof}
The remaining $1-r$-infinitesimal and other source-gap openings
are not admitted by this corollary. The failed simple rescaling on
that opposite boundary is recorded separately.


\section{\texorpdfstring{The complete directed weak-C2 source and testing join}{The complete directed weak-C2 source and testing join}}
\label{component:pec6c617d3a3d}
% Scoped mathematical certificate; see local_identity_audit.closure.json for exact changes.

\label{pec6c617d3a3d:sec:weak-c2-complete-join-rederived}

\paragraph{Theorem.}
Let \(K_c,K_t\) be arbitrarily chosen, separately labelled
real closed extensions of \(\mathbb R\); no
algebraic-disjointness assumption is made.  In particular
they may be two labelled copies of the same extension.
Use an original physical stationary C2 source
with \(d>0\), \(\lambda_{\rm event}=\beta=0\), all event
coefficients nonnegative, and every original strict source gap.
Normalize \(d=1\).  Suppose the positive actual source chart
has two infinitesimal contact half-angles and
\[
 C\notin\mathrm{Inf}_c,\qquad
 \alpha\in\mathrm{Inf}_c,\qquad
 \delta_{\rm tot}
   =\frac1{1+n+m+C+A+D+\alpha}\in\mathrm{Inf}_{c,>0}.
 \tag{HJ1}
\]
Then the entire original parent is a finite sum of squares in
\(K_c\otimes_{\mathbb R}K_t\) at every independent projective
testing pair and every real or complex state.
There is no earlier-admission-complement hypothesis.
The theorem retains literal \(n=0\) and every allowed
independent source opening.
By the exact reflection HJ3 it also proves the original
directed weak-\(C\), strong-\(\alpha\) HH source.
The entry into the physical positive source chart is a
separate global source-entry theorem, not a conclusion
inferred from this join.

\paragraph{The full parent and the reflection.}
For finite testing half-angles write
\[
 D_u=1+u^2,\quad D_v=1+v^2,\quad
 M=vW-X+uvY-uZ,\quad J=uW+uvX-Y-vZ .
\]
The quarter-parent is exactly
\[
\begin{split}
 \mathcal P={}&nD_uD_vX^2+mM^2+CD_v(uW-Y)^2
       +AD_uD_vY^2+DJ^2+\alpha D_u(vY-Z)^2\\
 &+\tfrac12\{J^2+D_u(vY-Z)^2+D_v(X+uZ)^2-D_uD_vZ^2\}.
\end{split}
\tag{HJ2}
\]
Every proof imported below certifies this entire polynomial.
The omission-preserving reflection is
\[
\begin{split}
 &(u,v;W,X,Y,Z;n,m,C,A,D,\alpha)\\
 &\hspace{4mm}\longmapsto
 (1/v,1/u;Z,X,Y,W;n,D,\alpha,A,m,C).
\end{split}
\tag{HJ3}
\]
Under HJ3, \(\mathcal P=u^2v^2\mathcal P'\) exactly.
This is verified packet by packet and for the full signed
core.  Its degree-\((2,2)\) homogeneous form exchanges
the corresponding angle coordinates, so it includes
the literal projective endpoints.
It preserves \(n,A\), and interchanges \(m,D\) and
\(C,\alpha\).  A weak-\(C\) HH input therefore becomes
strong-\(C\), weak-\(\alpha\) II.
No transformation \(t\mapsto1/t\) of the subsequent
source shape is asserted.

\subsubsection*{Actual source equations and the weak face}

Normalize the positive actual contact by \(Z_c=1\), and put
\[
\begin{gathered}
 z=u_c^2,\quad y=v_c^2,\quad H=v_cY_c,\quad
 L=v_c(u_cW_c-Y_c),\quad B=H+L,\quad r=-X_c/u_c,\\
 c_a=1+zr,\quad S_c=c_a+L,\quad T_c=B+zr,\quad
 j=c_a-L/y,\quad G=B+z(H+r-1).
\end{gathered}
\tag{HJ4}
\]
The unchanged source conditions are
\[
 z,y>0,\quad 0<H,r<1,\quad L>0,\quad
 j,G,B-zyr>0,\quad n\geq0 .
\]
In particular
\[
 (u_c,v_c;W_c,X_c,Y_c,Z_c)
 =\left(\sqrt z,\sqrt y;\frac{B}{\sqrt{zy}},
                  -r\sqrt z,\frac H{\sqrt y},1\right).
\]
The exact source denominator and positive numerator are
\[
\begin{split}
 \delta={}&z+(1+z)B+(1+z)B^2+2z(1+z)Br
                       +z^3r^2+2z^2r\\
 &+y\{(1+z)B+z+z^2r(2-r)\},\\
 M_s={}&(1+z)T_c(H+y)-z(1+y)(1-r)L\\
     ={}&(1+z)T_cH+yG+yzj+zrL(1+y).
\end{split}
\tag{HJ5}
\]
The six original coefficients, before any size specialization, are
\[
\begin{aligned}
 n&=\frac{T_cc_aS_c-r\delta}{2r\delta},&
 m&=\frac{z^2(1+y)(1-r)c_aS_c}{2G\delta},\\
 C&=\frac{c_aM_s}{2L\delta},&
 A&=\frac{yc_aT_cS_c}{2H\delta},\\
 D&=\frac{(1+z)T_c(B-zyr)S_c}{2yj\delta},&
 \alpha&=\frac{T_cM_s}{2(1-H)\delta}.
\end{aligned}
\tag{HJ6}
\]
All inequalities and all contact identities below refer
to these actual coefficients.  In particular the
numerator of \(n\) is never replaced by a limiting
source whose sign has not been proved.

Since \(z,y\) are infinitesimal, the exact \(j\) gap gives
\(L<y(1+zr)\), hence \(L\) infinitesimal and \(B\)
bounded.  If \(B\) is infinitesimal, then \(H\) is
infinitesimal and \(\delta\geq B+z\geq T_c\).
The first expression for \(M_s\) gives
\[
 0<\alpha
 \leq\frac{T_c(1+z)(H+y)}{2(1-H)}
 \in\mathrm{Inf}_c.
 \tag{HJ7}
\]
If \(B\) is a unit, \(H=B-L\), \(\delta\), and \(M_s\)
are units by the positive expression in HJ5.
Then \(\alpha\) is noninfinitesimal and
\(CL=c_aM_s/(2\delta)\) is a unit.
Thus \(\alpha\) is weak exactly when \(B\) is
infinitesimal.  The input HJ1 therefore has
\[
 B,z,y\in\mathrm{Inf}_{c,>0}.
 \tag{HJ8}
\]
This conclusion is an actual source estimate.
It imposes no power-law relation or unprinted rate
comparison between \(B,z,y\).

\subsubsection*{Uniform source alternatives}

Put
\[
 a=z/B,\quad t=y/B,\quad h=H/B,\quad l=L/B=1-h,\quad
 e=1-r,\quad g=G/B.
 \tag{HJ9}
\]
The quantities \(h,l,r,e,g,j\) are positive and bounded;
\(a,t\) may be infinitesimal, units, or unbounded.
Their exact relations give
\[
\begin{gathered}
 g=1-ae+hz>0,\qquad j=1+zr-l/t>0,\\
 ae<1+hz,\qquad
 e\text{ unit}\Longrightarrow a\text{ bounded},\\
 a\text{ unbounded}\Longrightarrow e\text{ infinitesimal},
 \qquad a\text{ infinitesimal}\Longrightarrow g\sim1,\\
 g\text{ infinitesimal}\Longrightarrow ae\sim1.
\end{gathered}
\tag{HJ10}
\]

Here is the source estimate behind the finite partition.
From HJ5,
\(\delta/(B+z)\sim1\).
Moreover \(T_c\geq B\), \(T_c>z(1-H)\), and
\(T_c\leq B+z\), so \(T_c/\delta\) is a positive unit.
Consequently
\[
 r(2n+1)=\frac{T_cc_aS_c}{\delta}\asymp1,\qquad
 \frac{AH}{y}=\frac{c_aT_cS_c}{2\delta}\asymp1.
 \tag{HJ11}
\]
The \(m\) formula can be rewritten exactly as
\[
 m=\frac{(1+y)c_aS_c}{2}\frac z\delta
               \frac{1+zH/B-g}{g}.
 \tag{HJ12}
\]
For unit \(g\) this is bounded.  For infinitesimal \(g\),
\(ae\sim1\), \(z/\delta\) is a unit, and \(m\asymp1/g\).
Thus no independent \(a=z/B\) rate creates an omitted
intermediate \(m\) case.

The positive decomposition of \(M_s\), and its first
expression in HJ5, similarly show
\[
\begin{array}{c|c}
 \text{original coefficient is unbounded}
       &\text{equivalent source condition}\\ \hline
 n&r\text{ infinitesimal}\\
 m&g\text{ infinitesimal}\\
 A&H/y\text{ infinitesimal}\\
 C&L/(H+y)\text{ infinitesimal}\\
 D&y/B\text{ infinitesimal or }j\text{ infinitesimal}.
\end{array}
\tag{HJ13}
\]
For completeness, the \(C\) upper bound is a bounded
multiple of \((H+y)/L\).
If \(L/(H+y)\) is infinitesimal, the subtracted
term in the first formula for \(M_s\) is
infinitesimal relative to \(T_c(H+y)\);
this gives the reverse estimate.
For \(D\), if \(B/y\) is infinitesimal, then
\(j\sim1\) and the exact formula gives \(D=O(B/y)\).
Otherwise \(zyr/B\) is infinitesimal, and
\(D\asymp B/(yj)\).
The original gap \(B-zyr>0\) is thereby retained in
the region where it cannot be replaced by \(B\).

It follows that all six original coefficients are bounded
if and only if
\[
             t,h,l,j,r,g\text{ are units}.
 \tag{HJ14}
\]
Indeed \(t\) infinitesimal makes \(C,D\) unbounded;
\(t\) unbounded makes \(C,A\) unbounded.
For unit \(t\), HJ13 gives exactly the other five
unit conditions.  HJ1 rules out HJ14, since the
sum in the denominator of \(\delta_{\rm tot}\)
would then be bounded.  No positive-\(d\) source
theorem is needed to remove this empty intersection.

There is also the exact source equivalence
\[
 C\text{ infinitesimal}
 \quad\Longleftrightarrow\quad
 r,g,H/y,j\text{ infinitesimal}.
 \tag{HJ15}
\]
To verify the forward direction, the four positive
terms of \(M_s/(L\delta)\) first force \(H/L\)
infinitesimal.  Bounded \(C\) and HJ13 give
\(L\asymp y\), hence \(H/y\) infinitesimal and
\(B\asymp y\).  The same four summands give
\(G/\delta,zr/\delta,zj/\delta\) infinitesimal.
The factor \(z/\delta\) cannot be infinitesimal:
that would imply \(B/\delta\sim1\) and
\(G/\delta\sim1\).  Hence \(r,j\) are infinitesimal,
and \(G>0\) together with \(z/\delta\) unit gives
\(B/z\) unit.  It follows that \(g\) is infinitesimal.
Conversely \(r,g\) infinitesimal imply \(B/z\sim1\).
With \(H/y,j\) infinitesimal the gap gives
\(L/y\sim1\), hence \(B/y\sim1\).
Every positive summand of \(M_s/(L\delta)\) is then
infinitesimal, proving the converse.
On the directed input HJ1 this entire intersection
is empty.  In particular it is not an allowed
unproved continuation of the R7 residue pivots.

\subsubsection*{The ten rows and every remaining exit}

The symbol \({\rm I}\) denotes positive infinitesimal,
\({\rm U}\) a positive unit, and \({\rm O}\) positive
unbounded.  Every row below includes HJ4--HJ10 and
the actual coefficient formulas HJ6.
\[
\begin{array}{c|l}
 {\rm R1}&t:{\rm I}\\
 {\rm R2}&t:{\rm O}\\
 {\rm R3}&t,h:{\rm U},\ l:{\rm I}\\
 {\rm R4}&t,h,l:{\rm U},\ j:{\rm I}\\
 {\rm R5}&t,h,l,j:{\rm U},\ r:{\rm I}\\
 {\rm R6}&t,h,l,j,r:{\rm U},\ g:{\rm I}\\
 {\rm R7}&t:{\rm U},\ h,g:{\rm I},
                  \quad r,j\text{ not both }{\rm I}\\
 {\rm R8}&t,g:{\rm U},\ h,j:{\rm I}\\
 {\rm R9}&t,g,j:{\rm U},\ h,e,a:{\rm I}\\
 {\rm R10}&t,g,j:{\rm U},\ h,e:{\rm I},\ a:{\rm O}.
\end{array}
\tag{HJ16}
\]
The rows are disjoint.  Their complements have the
following explicit destinations, not a blanket
earlier-admission clause.

Split \(t\) first.  Its infinitesimal and unbounded
cases are R1 and R2.
For \(t\) unit split \(h\).
If \(h\) is a unit, test in order whether \(l,j,r,g\)
is infinitesimal.  These are exactly R3--R6.
If none is, HJ14 contradicts the total-scale
condition.
If \(h\) is infinitesimal, \(l\sim1\).
For infinitesimal \(g\), simultaneous infinitesimal
\(r,j\) contradicts HJ15 and HJ1; otherwise this
is R7.
For unit \(g\), infinitesimal \(j\) gives R8.

Only \(t,g,j\) units and \(h\) infinitesimal remain.
If \(e\) is a unit, HJ10 makes \(a\) bounded.
For unit \(r\), unit \(a\) gives WA and infinitesimal
\(a\) gives WE.  For infinitesimal \(r\), either
bounded type of \(a\) gives SL.
If \(e\) is infinitesimal, \(r\sim1\).
Every bounded \(a\), including the former WC and R9 cases,
is covered by the complete BA1--BA15 certificate. Unbounded \(a\)
gives R10. The WA/WE/SL source conditions are the unit \(t,g,j,l,e\),
infinitesimal \(B,h,z,y\), and the indicated \(r,a\) types;
these are precisely the conditions used in their explicit formulas.
No \(h/a\), \(e/B\), or source/testing comparison
has been inserted.  This proves exhaustion of
the actual source, without asserting that every
formal type tuple is physically realized.

\subsubsection*{Actual complete proof for each row}

The following are explicit algebraic certificates on the source
conditions just displayed. Their equations use the actual stationary
weights and their inequalities use only those source gaps, coefficient
signs and source types. A historical residual-complement premise is
therefore not needed to apply these formulas. Each full-parent local
identity and each finite or homogeneous remote join is retained below.

\paragraph{R1 and R2.}
R1N1--R1N24 in
Section~\ref{component:pbaab204a8ab1}
prove all original R1 sources.
The exact compact-source and full-normal derivation
permits arbitrary positive normalized state ratio;
its full six-root cap identities and projective
entry supply all outer tests.
R2A1--R2A4 in
Section~\ref{component:p913a09bcdcda}
prove all original R2 sources.
That theorem explicitly splits bounded \(a\),
unbounded \(a\) with bounded \(n/(tB)\), and
unbounded \(a,n/(tB)\), using respectively
RL/SH, TU, and UL.
Its R2O/SJ join supplies every test and state.
The formal C1 comparison in TU/UL is implemented
by exact coefficients and the complete actual
\(n\) packet; it never assumes a negative-beta
source is physical.

\paragraph{R3 and R4.}
Every positive \(e<1\) is either a unit or
infinitesimal.  The unit cases are the actual
R3E and R4E complete all-tests bodies on exactly their displayed
positive-unit-$e$ source conditions.  Every infinitesimal
\(e\) case is EN1--EN25 in
Section~\ref{component:p1de8b5ffab69}.
The latter keeps both complete state normals,
the endpoint \(H\) or \(1-H\) opening, and the
entire source partition before its remote join.
These alternatives exhaust each original row.

\paragraph{R5--R7.}
R5.1--R5.12 and its homogeneous endpoint paragraphs give the
complete R5 certificate.
R6 splits into bounded \(a\) and unbounded \(a\);
in the latter case physicality bounds \(b=a^2B\),
whose unit and infinitesimal cases are RB and RI.
Thus its bounded-$a$ body and the RB/RI bodies supply every
source of R6.
R7 splits into bounded and unbounded \(a\);
The bounded-$a$ body and RU respectively supply these two cases.
Its sole genuine simultaneous-\(r,j\) exclusion
has already been dispatched by HJ15.
These statements preserve the actual \(n=0\)
boundary wherever it is possible.

\paragraph{R8--R10.}
R8 is the union of the unit-\(r\) local JD proof,
the infinitesimal-\(r\) local RD proof, and JD's
remote join valid for both cases.
R9 and the former regular WC exit are both covered by the new
BA1--BA15 whole-parent proof, including literal $n=0$, every state
and every projective testing pair.
R10 is the actual three-component all-tests join:
its weighted full-parent proof, its positive-thin
contact proof, and its unit-\(W\) complement.
For unbounded $a$, the source alternatives are $n/e^2$
noninfinitesimal, $0<n/e^2$ infinitesimal, and literal $n=0$.
They are exhaustive and are proved respectively by the weighted
identity, the positive-thin comparison, and the exact boundary map.
The literal \(n=0\) import is exactly OP4 followed
by the native C1 small-\(e\) or centered-\(A\)
body.  It contains no whole-C1 or whole-weak-C2
routing premise.

\subsubsection*{State, sign, and projective restoration}

Each row proof certifies the full HJ2, either directly
on four states or after an exact inverse with both
state normals present.
All normalizations use pure-source, pure-testing,
or positive separated factors.
At a nonzero finite chart they are reversed on
the SOS roots.  Literal zero and infinite testing
angles use the actual homogeneous parent, not a
limiting division by a zero coordinate.
The exact state sign gauges preserve every event
packet and the core; the negative-core state
normalizations and their complementary
nonnegative-core branches are part of the named
all-tests bodies.
No arbitrary state ratio is bounded by a source
scale in this assembly.

Every bounded complete Gram used in those bodies
has strict positive real residue on its retained
root list.  The strict bounded-Gram lemma, including
its WA10 polynomial version, therefore
gives a finite tensor SOS; it does not take a
square root or inverse of a mixed tensor.
Restoring all original positive excess packets
preserves that conclusion.
The original measurement matrices are real, so a
complex-state value is the sum of its real and
imaginary state values.
The source partition HJ16 and its four explicit
regular exits now prove the theorem.

\paragraph{Dependency boundary.}
This is a proof of the complete directed weak-C2
leaf and its exact reflected HH orientation.
The positive-total-scale interior is excluded
algebraically by HJ1 and HJ14, and the two-weak
intersection is excluded algebraically by HJ1
and HJ15.  Hence neither global A5 nor a global
C2 union is a dependency.
Global source entry, other omission patterns,
and other source-scale regimes remain the
responsibility of their own stated theorems.
The accompanying manifest lists actual proof
bodies and prerequisite parent/conductor
identities with their precise uses; status
documents are not mathematical inputs.


\section{\texorpdfstring{A single complete local proof for the bounded-$a$, small-$e$ weak-C2 corner}{A single complete local proof for the bounded-a, small-e weak-C2 corner}}
\label{sec:new-bounded-a-weak-c2}

This section replaces the separate WC and R9 local constructions. Its
source assumptions are $B,h,e>0$ infinitesimal, $a>0$ bounded,
$r=1-e$, and $t$ a positive unit with $t_0>1$, together with the
original physical stationary source and $n\ge0$. No lower bound on
$a$ or $n$ is assumed. The proof covers the whole four-state parent
at every pair of independent infinitesimal testing angles, including
both axes and all signs. The complete homogeneous remote join below
covers its complement.

\paragraph{Exact source and parent.}
Put $x=aB$, $y=tB$, $l=1-h$, and
\[
 c_a=1+xr,\quad S_c=1+B(l+ar),\quad g=1-ae+hx,\quad
 j=c_a-l/t,
\]
\[
 \Delta=a+(1+x)+(1+x)B+2x(1+x)r+a^3B^2r^2+2a^2Br
             +tB[1+a+x+a^2Br(2-r)],
\]
\[
 M=(1+x)(1+ar)(h+t)-a(1+tB)el.
\]
The unchanged stationary weights are
\[
 \begin{aligned}
 n&=\frac{(1+ar)c_aS_c-r\Delta}{2r\Delta},&
 m&=\frac{a^2(1+tB)e c_aS_c}{2g\Delta},\\
 C&=\frac{c_aM}{2l\Delta},&
 Ah&=\frac{t c_a(1+ar)S_c}{2\Delta},\\
 D&=\frac{(1+x)(1+ar)(1-txr)S_c}{2tj\Delta},&
 \alpha&=\frac{B^2(1+ar)M}{2(1-hB)\Delta}.
 \end{aligned}                                                     \tag{BA1}
\]
Writing $D_u=1+u^2$, $D_v=1+v^2$ and
$J_D=uW+uvX-Y-vZ$, the quarter-parent is
\[
 \begin{split}
 P={}&m(vW-X+uvY-uZ)^2+nD_uD_vX^2+CD_v(uW-Y)^2
       +AD_uD_vY^2\\
    &+(D+\tfrac12)J_D^2+(\alpha+\tfrac12)D_u(vY-Z)^2
      +\tfrac12D_v(X+uZ)^2-\tfrac12D_uD_vZ^2.
 \end{split}                                                       \tag{BA2}
\]

Define
\[
 q=nr^2+(m+\tfrac12)e^2+\alpha,\qquad
 \mu=m-a^2q,\qquad \lambda=\frac{\mu}{ta^2B}.
\]
All apparent $a$ and $B$ poles in $\lambda$ cancel in BA1.
Indeed at $B=0$ one has
$m/a^2=nr^2+(m+1/2)e^2$ identically; the remaining divided
numerator is polynomial and its denominator has a unit residue
for every bounded $a\ge0$. Direct evaluation gives
\[
 \lambda\sim\tfrac12,\quad m,n,\alpha,\mu\in\operatorname{Inf},
 \quad C\sim Ah\sim t/2,\quad D\sim[2(t-1)]^{-1},
 \quad \alpha/B^2\sim t(1+a)/2.                           \tag{BA3}
\]
Here and below a residue formula with bounded letters on its right
means the same formula evaluated at their real residues.
These cancellations are identities of the printed rational functions,
so no relative source rate has been fixed.

\paragraph{A full-parent identity with bounded matrix coefficients.}
Set $U=u^2$, $V=v^2$, $k=\sqrt{x}$, and $c=\sqrt{a/t}$.
Use the state-linear forms
\[
 E=tuW-vZ,\quad \eta=Y-huW=\sqrt h\,N_y,\quad
 \xi=X+ruZ,\quad F_v=(V-y)Z,\quad R_D=v(U-x)Z,
\]
\[
 J=(uv-x/c)Z,\quad Q=\sqrt q(avW-uZ-a\xi),\quad
 N=\sqrt n\{r[(a/t)V-U]Z+u\xi\},
\]
\[
 R=\sqrt\lambda\,tk(cv-u)W,\qquad
 T=\sqrt{2\lambda/c}\,k\sqrt{uv}(tcW-Z).              \tag{BA4}
\]
The square root in $T$ is pure testing. Initially take $u,v>0$;
literal axes are addressed below.

For a compact coefficient table put
\[
 \begin{aligned}
 b_2&=2mh+Cl^2+(A+\alpha+\tfrac12)h^2,&
 b_4&=h^2(m+A+\alpha+\tfrac12),\\
 d_4&=(n+D+\tfrac12)r^2,&
 e_1&=mhe+r(D+\tfrac12)l+(\alpha+\tfrac12)h,\\
 c_F&=(b_2-a^2nr^2)/t^2,&
 a_E&=2Ah^2/t^2-2e_1/t,\\
 b_E&=2b_2/t^2,&c_E&=2b_4/t^2,\\
 a_Y&=2[Ah/t-me+r(D+\tfrac12)-\alpha-\tfrac12],\\
 b_Y&=2[-Cl+Ah+m+(\alpha+\tfrac12)h]/t,\\
 c_Y&=2[Ah+h(m+\alpha+\tfrac12)]/t,\\
 x_3&=2[(D+\tfrac12)(l/t-1)+e/2-nr(1+a/t)-mh/t],&
 x_5&=-2r(n+D+\tfrac12).
 \end{aligned}                                                     \tag{BA5}
\]
Let $Y_*=hE/t+\eta$ and define the entire hard quadratic explicitly by
\[
 \begin{split}
 H_0={}&CD_v(lE/t-\eta)^2+AD_uD_vY_*^2
 +(D+\tfrac12)(lE/t-\eta+uv\xi)^2\\
 &+(\alpha+\tfrac12)D_uVY_*^2
 +[n(1+V+UV)+D_v/2]\xi^2\\
 &+2m(VE/t-uv\xi)Y_*+mUVY_*^2.
 \end{split}                                                       \tag{BA6}
\]
In particular BA6 has no inverse testing angle and no $a$ denominator.
Upon replacing $\eta$ by $\sqrt hN_y$, every occurrence of $A$
is through $Ah$, $\sqrt hAh$, or $hAh$.

With $\sigma=t\lambda x/\sqrt q$, the exact identity is
\[
 \begin{split}
 P={}&Q^2+N^2+T^2+R^2+\lambda J^2+d_4R_D^2
 +(c_F+b_4x/t^2)F_v^2+(b_4/t^2)vR_DF_v\\
 &+H_0-\lambda xE^2-\mu\xi^2-2\sigma Q\xi\\
 &+E[(a_E+c_EV)R_D+(b_E+c_Ex)vF_v]\\
 &+\sqrt hN_y[(a_Y+c_YV)R_D+(b_Y+c_Yx)vF_v]\\
 &+\xi[x_5uvR_D+u(x_3+x_5x)F_v].
 \end{split}\tag{BA7}
\]
This is the generic RC/R10 identity, not a scalar restriction.
One pen-and-paper verification is to substitute BA4 into BA7 and use
\[
 R^2+T^2+\lambda J^2
 =\lambda\{t^2x[U+(a/t)V]W^2-4txuvWZ+(UV+xy)Z^2\}.
\]
The hard quadratic is then exactly BA6. The remaining six scalar
and first-normal coefficient identities follow by substituting BA1;
a common denominator leaves zero polynomial numerators. Equivalently
these six identities are the original value and all five first-contact
conditions. Both the full identity and those six numerator equalities
have also been independently expanded during verification.

All coefficients in BA5 are bounded. Their relevant residues are
\[
 c_F\sim(2t)^{-1},\quad d_4\sim t/[2(t-1)],\quad
 a_E\sim-(t-1)^{-1},\quad b_4\in\operatorname{Inf}.
\]
The entire hard residue of BA6 on $(E,N_y,\xi)$ is
\[
 \operatorname{diag}\bigl([2(t_0-1)]^{-1},t_0/2,1/2\bigr).
                                                               \tag{BA8}
\]

\paragraph{The only nontrivial source Schur bound.}
The exact physical inequality is
\[
 2r\Delta n=eQ_s-BK_s\ge0,
\]
\[
 \begin{split}
 Q_s={}&1+B(1+t+ah+2ar+at)\\
       &+B^2\{a(h+t)+a^2[h(r+1)+t(1+r-r^2)]\},\\
 K_s={}&(1+a)(h+t)(1+aB).
 \end{split}                                                       \tag{BA9}
\]
Every term of $Q_s,K_s$ is nonnegative, $Q_s\sim1$, and
$K_s\sim t(1+a)$. Thus $e/B\ge K_s/Q_s$. Since
$q\ge e^2/2+\alpha$, BA3 gives the opening-uniform bound
\[
 \operatorname{st}(\sigma^2)
 \le \frac{t_0a_0^2}
 {2(1+a_0)[1+t_0(1+a_0)]}.
\]
If $e/B$ is unbounded the left side is zero, so the same estimate
holds. No inverse of $n$ occurs. Consequently the source block on
$(Q,\xi)$ has a strictly positive Schur residue:
\[
 \operatorname{st}(1/2-\mu-\sigma^2)
 \ge\frac{1+a_0+t_0(1+2a_0)}
 {2(1+a_0)[1+t_0(1+a_0)]}>0.                         \tag{BA10}
\]
It also follows that $\sigma$ is bounded. This includes $a_0=0$
and literal $n=0$.

Retain $N^2,T^2$ whole. The residue of the remaining bounded Gram
on $(E,R_D,N_y,\xi,F_v,J,R,Q)$ has only two nontrivial blocks.
The $(E,R_D)$ block is
\[
 \frac1{2(t_0-1)}
 \begin{pmatrix}1&-1\\-1&t_0\end{pmatrix}\succ0,
\]
and the $(Q,\xi)$ block is strict by BA10. The other pivots are
$t_0/2$, $1/(2t_0)$, $1/2$, and $1$. Every other entry contains
an infinitesimal source factor or a positive testing degree, exactly
as displayed in BA5--BA7. The finite separated strict-Gram lemma
therefore proves $P$ to be a tensor sum of squares.

There are no state-coordinate inverses and no testing-ratio inverses
in the final roots. The literal axes have the complete nonnegative identities
\[
 \begin{split}
 P|_{u=0}={}&(n+\tfrac12)D_vX^2+m(vW-X)^2
 +(C+A+\tfrac12)D_vY^2\\
 &+D(Y+vZ)^2+\alpha(vY-Z)^2,
 \end{split}
\]
\[
 P|_{v=0}=nD_uX^2+(m+\tfrac12)(X+uZ)^2
 +(C+D+\tfrac12)(uW-Y)^2+AD_uY^2+\alpha D_uZ^2.
\]
The sign gauges
$(u,Y,Z)\mapsto(-u,-Y,-Z)$ and
$(v,X,Z)\mapsto(-v,-X,-Z)$ preserve BA2 and restore all signs.
Applying the real identity to the real and imaginary state parts gives
the complex-state statement. The proof is uniform for all independent
source and testing rates and includes every homogeneous state.

\paragraph{The finite separated lift used above.}
For completeness, let $G=G_0+\sum_{i=1}^L G_i t_i$ be a finite
symmetric matrix, with bounded pure-source coefficients, with
$G_0$ positive definite and bounded inverse, and with each $t_i$
a pure-testing infinitesimal. Whiten $G_0$ purely in the source
field and split the identity matrix into $L$ equal parts. Diagonalize
each resulting pure-source symmetric matrix. Each scalar coefficient
is $L^{-1}+\ell t_i$ with $\ell$ bounded. Choose a positive real
$R>L|\ell|$ and use
\[
 L^{-1}+\ell t_i=\frac1{2L}
 \{(1+L\ell/R)(1+Rt_i)+(1-L\ell/R)(1-Rt_i)\}.
\]
All four factors are positive in their respective real closed fields.
Their square roots give finite separated squares. There are finitely
many eigenvalues, so one real $R$ works for all. The case $L=0$
is immediate. Expanding at real testing residues proves the same
statement around any finite test pair. This is the only strict-Gram
principle needed in this section.

\paragraph{The complete remote join.}
The following direct identities close every remaining projective
testing pair under BA1--BA3. Put $\widehat Y=\sqrt A Y$, so
$Y=A^{-1/2}\widehat Y$. Write $\tau=t_0>1$,
$c_0=\tau/2$ and $d_0=[2(\tau-1)]^{-1}$. The source weights
$n,m,\alpha$ are infinitesimal, while $C\sim c_0$, $D\sim d_0$,
and $A^{-1/2}$ is infinitesimal.

First suppose $u$ is a bounded unit or unbounded and $v$ is
infinitesimal. Put $p=1/u$, $D_p=1+p^2$ and use
\[
 N_\infty=\sqrt n X,\quad V_X=vX,\quad K=pX+Z,\quad
 R_\infty=\sqrt\alpha pX,\quad W,\quad\widehat Y.
\]
Let $J_\infty=W+D_pV_X-pY-pvK$. The entire homogeneous parent is
\[
 \begin{split}
 p^2P(1/p,v)={}&D_pD_vN_\infty^2+m(pvW+vY-K)^2
 +CD_v(W-pY)^2+D_pD_v\widehat Y^2\\
 &+DJ_\infty^2+D_p[R_\infty+\sqrt\alpha(vY-K)]^2\\
 &+\tfrac12\{J_\infty^2+(1-p^2v^2)K^2\\
 &\hspace{9mm}+D_p[v^2Y^2-2vYK+2pYV_X+2pvKV_X-p^2V_X^2]\}.
 \end{split}                                                       \tag{BA11}
\]
It follows directly by substituting $Z=K-pX$ in BA2. Its bounded
six-root residue, with $d_p=1+p_0^2$, is
\[
 d_p(N_\infty^2+R_\infty^2+\widehat Y^2)+K^2/2
 +c_0W^2+(d_0+\tfrac12)(W+d_pV_X)^2-d_pp_0^2V_X^2/2.
\]
The $W,V_X$ determinant is $d_pc_0(d_0+1/2)>0$, using
$c_0(d_0+1/2)=(c_0+d_0+1/2)/2$. All other pivots are positive.
No inverse of $n$ or $\alpha$ occurs, and either root may vanish.
This covers literal $u=\infty$ by $p=0$ as well.

Next suppose $u$ is infinitesimal and $v$ is a bounded unit.
Use the five roots
\[
 (S,R_u,X,\widehat Y,Z)=(\sqrt mW,uW,X,\sqrt A Y,Z).
\]
The exact full polynomial is obtained from BA2 by the two replacements
\[
 m(vW-X+uvY-uZ)^2=(vS-\sqrt mX+\sqrt m\,uvY-\sqrt m\,uZ)^2,
 \qquad uW=R_u
\]
in all other terms. This operation uses no inverse of $m$ or $u$.
Its separate residue pivots on $S,X,\widehat Y$ are
$V,(1+V)/2,1+V$, where $V=v_0^2>0$. The remaining $R_u,Z$ block is
\[
 \begin{pmatrix}
 c_0(1+V)+d_0+1/2&-(d_0+1/2)v_0\\
 -(d_0+1/2)v_0&d_0V
 \end{pmatrix},\qquad \det=c_0d_0V^2>0.             \tag{BA12}
\]
Here $c_0d_0-d_0/2-1/4=0$. If instead $v$ is unbounded, use
$q_v=1/v$ and the same five roots. Multiplying the same exact
polynomial by $q_v^2$ gives, at $u=q_v=0$, the diagonal residue
\[
 (S,R_u,X,\widehat Y,Z):\quad (1,c_0,1/2,1,d_0)>0.   \tag{BA13}
\]
This homogeneous substitution is polynomial and also applies at
literal $q_v=0$.

If both finite testing residues are nonzero, the hard
$\widehat Y$ pivot is $(1+U)(1+V)$. The complete remaining
three-state residue on $(W,X,Z)$ is
\[
 c_0(1+V)UW^2+(d_0+\tfrac12)(u_0W+u_0v_0X-v_0Z)^2
 +\tfrac12(1+V)(X+u_0Z)^2-\tfrac12(1+U)VZ^2,
\]
where $U=u_0^2,V=v_0^2$. Its determinant is
\[
 (1+U)(1+V)UV\left\{\frac{c_0(2d_0+1)}4U
                              +\frac{c_0d_0}2V\right\}>0.        \tag{BA14}
\]
The $Z=0$ restriction is strictly positive, so the full matrix is
positive definite. Every coefficient is bounded after the hard-state
normalization.

The one-infinity cases not already covered by BA11 or BA13 are
strict directly. At $u=\infty$ and a bounded unit $v$, the residue is
\[
 c_0(1+V)W^2+(d_0+\tfrac12)(W+v_0X)^2+Z^2/2
 +(1+V)\widehat Y^2\succ0.
\]
At $v=\infty$ and a bounded unit $u$, it is
\[
 c_0UW^2+(1+U)X^2/2+d_0(u_0X-Z)^2
 +(1+U)\widehat Y^2\succ0.
\]
These are the degree-two homogeneous values of BA2 in the respective
angles; nearby reciprocal coordinates are pure-testing infinitesimals.

For two infinite angles put $p=1/u$, $q_v=1/v$ and introduce
$P_\infty=pZ$, $Q_\infty=q_vZ$. The entire exact polynomial is
\[
 \begin{split}
 p^2q_v^2P(1/p,1/q_v)={}&n(1+p^2)(1+q_v^2)X^2
 +m(pW-pq_vX+Y-Q_\infty)^2\\
 &+C(1+q_v^2)(W-pY)^2
 +(1+p^2)(1+q_v^2)\widehat Y^2\\
 &+D(q_vW+X-pq_vY-P_\infty)^2
 +\alpha(1+p^2)(Y-Q_\infty)^2\\
 &+\tfrac12\{(q_vW+X-pq_vY-P_\infty)^2
 +(1+p^2)(Y-Q_\infty)^2\\
 &\hspace{9mm}+(1+q_v^2)[p^2X^2+2XP_\infty-P_\infty^2]\}.
 \end{split}                                                       \tag{BA15}
\]
Its residue on $(W,X,\widehat Y,P_\infty,Q_\infty)$ is
\[
 c_0W^2+(d_0+\tfrac12)X^2-2d_0XP_\infty
 +d_0P_\infty^2+\widehat Y^2+Q_\infty^2/2\succ0.
\]
The nontrivial determinant is $d_0/2>0$. At $p=q_v=0$ the two
extra roots vanish, respecting the original null state rather than
inventing a positive $Z$ pivot.

The finite separated lift applies to every strict residue just listed.
Away from literal endpoints one restores the parent by multiplying by
the appropriate pure-testing square; at a literal endpoint one uses
the homogeneous polynomial value itself. Local BA7, BA11--BA15,
the one-infinity residues, sign gauges and real/imaginary splitting
exhaust all independent projective tests and all real or complex
states. This proves the all-tests admission of the whole bounded-$a$,
small-$e$ source domain.

\paragraph{Dependencies and safe deletions.}
BA1 is the common native weak-C2 source formula. The local identity
BA7 follows from the displayed coefficient calculation and the six
RC1/RC6 contact balances printed in
Section~\ref{component:pd709027c7805}; the hard quadratic is explicitly
BA6. The remote formulas BA11--BA15 are direct homogeneous rewritings
of BA2. Thus this proof imports no physical C1 theorem, contact
conductor, or omitted WC/R9/GR/BS body. It also does not import the
R10 source estimates: those estimates concern a different source
range and use the same generic algebra.

For the qualitative theorem this proof replaces the full WC and R9
local-and-join bodies and the older bounded-ratio R9 alternative.
Retain the native source dictionary and generic R10 identity for the
unbounded-$a$ and R8 branches, together with their own source bounds.

